% Revised chapter source for inclusion in the main manuscript.
% Construction-first prequantum chapter, repaired against Chapters 9--17.
% Companion bibliography: Brave_New_Prequantum_Field_Theory.bib
% Requires the manuscript preamble, in particular amsmath/amssymb and tikz-cd.

\chapter{Brave-New Prequantum Field Theory}
\label{ch:prequantum}

\section{Purpose, continuity, and scope}
\label{sec:prequantum-purpose-continuity-scope}

The preceding chapters have reached the point at which brave-new variational
geometry, its formal infinitesimal symmetries, its differential density sector,
and its critical reduction theory have all been constructed.  Chapters~9--10
supplied the Thom-stable differential coefficient theory, the BNV deformation and
cycle sectors, the intrinsic relative de Rham calculus, coherently closed
infinitesimal cochains, differential exceptional integration, and differential
transgression from actual coherent coefficient traces.  Chapter~12 supplied
differential Lagrangians, the higher field-moduli action, relative Lepage
realizations, and the universal coherently closed presymplectic transgression.
Chapters~13--16 supplied formal Lie groupoids, spectral differentiation,
partition-Lie algebroids, Atiyah--Spencer operations, tangent representations,
and native Lie integration.  Finally, Chapter~\ref{ch:symmetry-noether} connected
the variational and symmetry branches by constructing the brave-new First and
Second Noether theorems, critical formalization and reduction, the mixed
symmetry--vertical--horizontal calculus, Hamiltonian realization, residual
presymplectic transgression, complete presymplectic descent, and reduced currents.

The purpose of the present chapter is to construct prequantum geometry from those
objects without inserting a line bundle, gerbe, integral lift, connection,
orientation, Cartan transport on an arbitrary coefficient, Hamiltonian-liftability
condition, equivariant-prequantizability condition, isotropy-exactness condition,
or reduction-compatibility condition as input.  Whenever one of those structures
is not already canonically present, its entire derived realization object is
formed.  The stable closed differential coefficient is formed by a pullback in
the parameterized stable category; the primary prequantization object is then the
relative infinite-loop-space fibre over the specified coherently closed class.
Ordinary
line bundles with connection, higher gerbes, and shifted prequantizations enter
only through the comparison-functor realization constructed below; a point of
that realization supplies the actual bridge and all required curvature and mixed
compatibilities.  This organization is compatible with the
Kostant--Souriau and higher-prequantum pictures
\cite{Kostant1970,FiorenzaRogersSchreiber2013,Schreiber2016} while retaining the
native spectral distinctions established earlier.

The field-theoretic emphasis is equally important.  Differential actions,
boundary prequantum objects, and higher-codimension descendants are assembled from
one differential coefficient by exceptional transgression.  Coefficient-map,
coherent de Rham-stage-map, and differential fundamental-class-candidate objects
are represented by pointed mapping realizations; their refinements of prescribed
motivic or stagewise data are represented by derived lift fibres.  Their
composition is inherited from
exceptional Beck--Chevalley and trace pasting.  Thus the gluing law belongs to the
prequantum differential geometry rather than being imposed later as a
quantization rule.  The classical Chern--Simons and higher Wess--Zumino mechanisms
remain comparison benchmarks \cite{Freed1995,FiorenzaSatiSchreiber2013}.

The chapter constructs the complete prequantum target described below and ends
at the provenance boundary proved in
Proposition~\ref{prop:prequantum-output-type-boundary}.  That proposition assigns
every native construction a point of a tagged derivation family whose signature has no
quantization constructor.  It does not claim that the underlying categorical
types of prequantum and quantum outputs are disjoint.  On the
shifted-geometric-quantization comparison branch the additional input and changed
target are exhibited explicitly rather than attributed to an unnamed prequantum
compatibility field.

Throughout, fix
\[
x:X\longrightarrow S,
\qquad
Q:=Q_{X/S},
\]
with the notation of Chapters~10--12, a Cartan field object $(A,\rho_A)$ with
descent object
\[
a:\overline A\longrightarrow Q,
\]
a stable variational coefficient $M\in E_Q$, and a degree-$n$ Lagrangian
\[
L:\mathbf1_Q\longrightarrow\operatorname{Sec}^0(A;M)[n].
\]
Retain the relative Lepage objects $\operatorname{Lep}^q(L)$, the universal
presymplectic transgressions $\omega_L^{(q)}$, the one-stage critical object
$C_L:=\operatorname{Crit}(L)$, and a formal symmetry groupoid
\[
\mathcal G_\bullet\in\operatorname{FormLieGpd}(\overline A/Q).
\]
Write
\[
\mathcal K_{\mathcal G,L,\bullet}\rightrightarrows C_L
\]
for the canonically formalized critical symmetry groupoid of
Chapter~\ref{ch:symmetry-noether} and
\[
q_{\mathcal G,L}:C_L\twoheadrightarrow C_L//\mathcal G
\]
for its effective quotient.

The logical progression is
\[
\text{big differential coefficients}
\longrightarrow
\text{stable curvature diagrams}
\longrightarrow
\text{prequantization},
\]
\[
\text{Cartan transport and mixed prisms}
\longrightarrow
\text{quantomorphism/Hamiltonian groupoids}
\longrightarrow
\text{integrated KS geometry},
\]
\[
\text{partition-Lie differentiation}
\longrightarrow
\text{native and differentiated Noether lifting}
\longrightarrow
\text{critical descent and reduction}.
\]
The Thom-stable differential, intrinsic stable, nonlinear formal-groupoid,
differentiated partition-Lie, and represented comparison branches meet only
through the morphisms and realization objects constructed below.

\subsection{Chapter-level target}
\label{subsec:prequantum-chapter-level-target}

The constructions below provide:
\begin{enumerate}
\item the affine-point extension $T\mapsto\mathcal R_T^{\partial}$ of the
      Chapter~9 Thom-stable differential categories, represented recovery,
      arbitrary pullback, effective descent, and big differential mapping spectra;
\item stable curvature-arrow and curvature-diagram realization objects, the closed
      differential class coefficient, universal prequantization, its connecting
      obstruction, flat torsor, and higher automorphisms;
\item the internal differential Lagrangian refinement object
      $\operatorname{PreQLag}$ and the separate differential Lepage-refinement
      realization $\operatorname{PreQLep}$, whose fibres measure lifting of actual
      relative Lepage realizations;
\item the complete closed presymplectic coefficient, its closed universal section,
      the variational curvature map, and the canonical map
      $\operatorname{PreQLep}\to\operatorname{PreQ}_{\omega,L}^{q}$;
\item orientation-free differential action, pointed fundamental-class-candidate
      and coefficient-trace-candidate realization objects, their lift fibres over
      prescribed motivic shadows, the explicit Chapter~10 realized stage/cochain
      diagram, and coherent-trace extension fibres on that diagram;
\item the universal commutative-square realization $\operatorname{SqReal}$,
      the typed transgression-operator realization
      $\operatorname{TransCurvReal}_{\partial}$, coherent composition, and
      Fubini transgression;
\item differential lift realizations for Chapter~8 correspondence $2$-cells and the
      resulting genuinely decorated traced-correspondence $(\infty,2)$-category;
\item the functor-level BNV--Stokes inverse locus, its small detection
      presentation, and the coefficientwise discrepancy, kept separate from
      infinitesimal Hamiltonianity;
\item Cartan-transport realization of an arbitrary curvature span, raw and formal
      quantomorphism and curvature-symmetry groupoids, and the exact raw/formal
      flat-gauge kernels;
\item the vertical coefficient-origin realization $\operatorname{VertCurvReal}$ and
      the class-specific origin realization $\operatorname{VertClassReal}$, replacing
      any implicit relation-level presentation of a generic complete closed coefficient
      or of its selected closed class;
\item the mixed rectangular relation, the ordered Alexander--Whitney morphism, the
      complete closed first prism identity, the two remaining normalized prism
      identities, the long-edge multiplicativity obstruction, and the complete
      Hamiltonian multiplicativity realization whose universal point is a genuine
      raw Hamiltonian groupoid;
\item the raw Hamiltonian-prequantum pullback, its effective raw liftable image,
      the raw KS torsor extension, its formalization, and the distinct inverse loci
      $\operatorname{KSEq}^{\mathrm{raw}}$ and
      $\operatorname{KSEq}^{\mathrm{for}}$, together with the raw-recovery fibre
      over a selected formal inverse;
\item exact integrated relative isotropy, with flat gauge identified as the formal
      relative kernel and no separate integrated isotropy-exactness condition;
\item separate differentiation of prequantum, liftable-Hamiltonian, quantomorphism,
      and actual-flat groupoids; the intrinsic algebroid KS fibre; the comparison
      from actual flat gauge; its source-module discrepancy and inverse locus; and
      the prequantum-to-quantomorphism comparison, cofiber, and inverse locus;
\item the source-module connecting KS class and the complete Chapter~15
      morphism-corolla--latching lift tower, with anchors recovered automatically
      from the actual algebroid projection;
\item the stable flat Cartan coefficient descended before applying
      $\Omega^\infty$, its comparison with kernel conjugation, its independently
      differentiated tangent representation, observable controller, labelled
      operations, coefficient connection, curvature filler, Bianchi coherence,
      and higher flatness hierarchy;
\item the Chapter~17 complete Noether transgression $U^{(q)}$, the separate
      realization $\operatorname{HamCompReal}$ lifting the Chapter~17 Hamiltonian
      identity into the generic transgression fibre, and the complete native
      Hamiltonian-Noether matching tower;
\item coherent factorization of that Hamiltonian section through the formal
      liftable image, the native Prequantum Noether lift realization, its
      differentiated lift and higher operation towers;
\item the universal differential variational refinement
      $(L,\widehat L,\eta_{\mathrm{Lag}},\widehat\ell)$, its complete
      symmetry-totalization realization, and its compatibility with the
      independent Prequantum Noether lift;
\item universal critical restriction over the full prequantization parameter,
      coefficient-level descent of the complete curvature span, the internally
      parameterized $\operatorname{PreQDesc}$ object, and the reduced prequantum
      section;
\item the fixed-pair and universal theorems that prequantization commutes with the
      constructed critical reduction;
\item descent of critical formal isotropy as the Chapter~13 formal adjoint
      group and the reduced Prequantum Noether action
      $\operatorname{Ad}^{\mathrm{for}}\to\operatorname{Aut}_{\mathrm{PreQ}}$,
      with individual automorphisms obtained by evaluation at selected reduced
      symmetry parameters;
\item the separately represented nondegeneracy inverse locus, the complete
      discrepancy/anomaly system, and typed comparison realizations between
      discrepancy objects;
\item the indexed comparison realization
      $\operatorname{CmpPreQReal}_{\mathcal B,b}$, whose actual exchange-occurrence
      poset $\mathcal K_b$ indexes the bridge of coefficient systems together
      with the dependent-product, internal-Hom, complete-totalization, Cartan,
      relation, and mixed inverse loci used by the branch;
\item arbitrary native pullback and effective descent for the intrinsic stable
      branch, the actual Chapter~10 realized de Rham exchange morphisms together
      with their inverse loci $\operatorname{DREq}$, finite Nisnevich descent, and
      separately scoped spectral-\'etale transport of the differentiated branch.
\end{enumerate}

\section{Audit of retained distinctions}
\label{sec:prequantum-audit-retained-distinctions}

The constructions below meet coefficient theories which were deliberately kept
separate in the preceding chapters.  The following distinctions remain in force.

First, a Thom-stable differential coefficient is not itself a coherently closed
infinitesimal coefficient.  Arbitrary represented base change gives the actual
Chapter~10 exchange morphisms, while invertibility is separately represented by
$\operatorname{DREq}$ outside the spectral-\'etale theorem domain.

Second, a morphism between stable coefficients is not silently attached to the
name of a coefficient.  Curvature and underlying-class arrows are points of their
mapping realizations when not already constructed.

Third, a differential Lagrangian refinement and a differential refinement of a
chosen relative Lepage realization are different lifting problems.  The former is
$\operatorname{PreQLag}$; the latter is the fibre of
$\operatorname{PreQLep}\to\operatorname{PreQLag}\times_Q\operatorname{Lep}^q(L)$.
Only the canonical zero-comparison Lepage point lifts functorially without an
additional fibre problem.

Fourth, a generic complete closed coefficient does not carry relation-level
vertical presentation data by notation, and a selected closed section need not lift
to that origin.  Hamiltonian contraction is defined only after points of
$\operatorname{VertCurvReal}$ and $\operatorname{VertClassReal}$ have constructed
the coefficient origin and the class origin separately.
The variational coefficient has a canonical point of that realization.

Fifth, a curvature morphism and a point of its prequantization fibre are distinct.
Prequantization is a lift of a specified complete closed class through the stable
curvature fibre sequence, and the fibre may be empty.

Sixth, the unrestricted coefficient-map, stage-diagram-map,
fundamental-class-candidate, and underlying-operator mapping realizations are
canonically pointed by zero.  The genuine existence problems are their fibres
over prescribed motivic shadows, selected component traces, or fixed underlying
operators.  Curvature compatibility for fixed operators is the path object between
the two composites in the resulting square.

Seventh, a geometric correspondence $2$-cell and a differential lift of its
natural transformation are different objects.  The latter is represented by
$\operatorname{TwoCellReal}_{\partial}$ and can be empty.

Eighth, the BNV cycle sector and the intrinsic infinitesimal de Rham sector remain
separate.  BNV interval integration controls the differential-Umkehr/Stokes
comparison; Hamiltonian primitives come from the mixed infinitesimal relation.

Ninth, the ordered staircase family is the ordered Alexander--Whitney morphism.
The primary first prism identity lives in the complete coherently closed
coefficient; its unary normalized formula is only its associated-graded shadow.

Tenth, curvature preservation, Hamiltonianity, raw arrowwise prequantum
liftability, formal liftability, and a coherent section of the KS extension are
separate.  In particular
$\operatorname{KSRawLiftReal}$ is the fibre which measures recovery of raw
inverse data from a selected point of $\operatorname{KSEq}^{\mathrm{for}}$.

Eleventh, actual flat gauge and the intrinsic algebraic KS fibre obtained after
partition-Lie differentiation are distinct and are connected only by
$\gamma^{\mathrm{KS,diff}}$.  Connecting morphisms and cofibres are formed after
forgetting to the stable source-module category, not in the anchored slice.

Twelfth, the stable flat Cartan representation is descended at the spectrum level.
Conjugation of $G_{\mathrm{flat}}$ is its infinite-loop comparison, not a device
for promoting an arbitrary group action to a stable action.

Thirteenth, Chapter~17's current $J^{(q)}$ and its complete transgression
$U^{(q)}$ are different.  The generic Hamiltonian primitive is $U^{(q)}$ after a
point of $\operatorname{HamCompReal}$ lifts the Chapter~17 Hamiltonian path to the
transgression fibre.  The current is the horizontal-tail shadow.

Fourteenth, symmetry is formed before the underlying Lagrangian is frozen.  The
universal differential symmetry realization stabilizes
$(L,\widehat L,\eta_{\mathrm{Lag}},\widehat\ell)$ as one refinement point and
recovers the fixed-$L$ problem only by pullback.

Fifteenth, effective descent of a selected prequantum pair and descent of the
whole prequantization moduli problem are different theorems.  The reduced Noether
output is a formal-adjoint-group-parameterized action on the reduced prequantum
automorphism object; the reduced current is its Hamiltonian primitive shadow.

Sixteenth, prequantization and nondegeneracy are separately constructed.  The
latter is the represented inverse locus of the exact $(1,1)$ contraction, and
any comparison with another discrepancy is represented by
$\operatorname{DiscCmpReal}$.

Seventeenth, external recognition is not a conditional guardrail.  The indexed
comparison frame and every right-adjoint or complete-limit exchange used by a
branch are represented inside $\operatorname{CmpPreQReal}_{\mathcal B,b}$.
The exchange-occurrence poset $\mathcal K_b$ is extracted from the mathematical branch
diagram and gives the actual index of those inverse loci.

Eighteenth, on the Safronov $n=1$ comparison branch, the twisted global-sections
category is universally a presentable $\operatorname{QCoh}(B)$-module.  Its
invertibility as such a module is a separate recognition problem represented by
$\operatorname{Gerb}^{1}_{\mathrm{Pic}}(B)$; the Picard-valued target is used
only after that inverse locus has supplied coherent inverse data.

\begin{convention}[Non-conflation]
\label{conv:prequantum-non-conflation}
In this chapter:
\begin{enumerate}
\item a ``closed differential class'' is a point of
      $\mathscr D\times_{\mathscr U}\mathscr C$;
\item a ``prequantization'' is a point of its homotopy fibre over a specified
      complete closed class;
\item ``flat ambiguity'' is $\operatorname{fib}(\kappa)$;
\item a ``differential Lepage refinement'' is a point of
      $\operatorname{PreQLep}$ over a selected underlying Lepage point;
\item a ``vertical coefficient origin'' is a point of
      $\operatorname{VertCurvReal}$ and a ``vertical class origin'' is a point of
      $\operatorname{VertClassReal}$; neither is an attribute of arbitrary
      $\mathscr C$ or of an arbitrary closed section;
\item ``curvature symmetry'' means preservation of the complete closed class;
\item ``Hamiltonian symmetry'' means a point of the mixed primitive fibre whose
      complete closed prism supplies curvature preservation;
\item ``raw liftable Hamiltonian image'' means the effective image inside
      $\operatorname{HamRaw}$; its formalization is
      $\operatorname{Ham}^{\mathrm{lift}}$;
\item ``Prequantum Noether primitive'' means the complete transgression $U^{(q)}$;
      $J^{(q)}$ is its horizontal-tail current shadow;
\item ``prequantum reduction'' means effective descent of the complete closed
      differential coefficient--section pair;
\item ``reduced Prequantum Noether action'' means the descended homomorphism
      $\operatorname{Ad}^{\mathrm{for}}(\mathcal K)\to
      \operatorname{Aut}^{\partial}_{\mathrm{PreQ}}$; an individual reduced
      automorphism is its evaluation at a selected reduced symmetry parameter;
\item ``comparison'' means a point of the indexed comparison frame together
      with the explicit inverse loci of every exchange map actually used by the
      branch;
\item ``anomaly'' is reserved for a represented image of an explicitly constructed
      discrepancy.
\end{enumerate}
\end{convention}

\begin{convention}[Theorem domains]
\label{conv:prequantum-theorem-domains}
The big differential category, parameterized stable curvature constructions,
vertical-origin and mixed relation realizations, raw/formal stabilizers, and
critical descent live in the native big spectral Nisnevich $\infty$-topos.
Differential exceptional integration is used only in the represented
Chapter~9--10 theorem range.  Partition-Lie differentiation and tangent
representation theory use the locally coherent spectral coefficient sector of
Chapters~14--15.  Ordinary tangent, line-bundle, gerbe, shifted-symplectic, or
dg/$L_\infty$ formulas occur only on their independently proved comparison
domains.  These are properties of objects to which the theorems apply; they are
not fields retained on a prequantization, formal symmetry, or Noether lift.
\end{convention}

\begin{construction}[Walking-equivalence and inverse-locus realizations]
\label{constr:prequantum-walking-equivalence-inverse-locus}
Let $\mathcal C$ be an internal $\infty$-category over a parameter object $T$,
let $\mathbb I=[1]$ be the walking arrow, and let $\mathbb E[1]$ be the
indiscrete groupoid on two objects.  For two $T$-points $A,B$ of $\mathcal C$
define
\[
\operatorname{Mor}_{\mathcal C}(A,B)
:=
\operatorname{hofib}_{(A,B)}
\left(\mathcal C^{\mathbb I}\longrightarrow
      \mathcal C\times_T\mathcal C\right)
\]
and
\[
\boxed{
\operatorname{Eqv}_{\mathcal C}(A,B)
:=
\operatorname{hofib}_{(A,B)}
\left(\mathcal C^{\mathbb E[1]}\longrightarrow
      \mathcal C\times_T\mathcal C\right).
}
\]
Restriction along $\mathbb I\to\mathbb E[1]$ gives
$\operatorname{Eqv}_{\mathcal C}(A,B)\to
\operatorname{Mor}_{\mathcal C}(A,B)$ by remembering the forward arrow.  If
$f:P\to\operatorname{Mor}_{\mathcal C}(A,B)$ is a family of arrows, define its
inverse locus by the actual pullback
\[
\boxed{
\operatorname{Inv}_{P}(f)
:=
P\times_{\operatorname{Mor}_{\mathcal C}(A,B)}
\operatorname{Eqv}_{\mathcal C}(A,B).
}
\]
When $P=T$ it is abbreviated $\operatorname{Inv}(f)$.  Its points are a point of
$P$, a reverse arrow, the two inverse homotopies, and all coherences encoded by
$\mathbb E[1]$.  Consequently it is inhabited precisely over the locus on which
$f$ is an equivalence.  Cotensoring and pullback are limits, so this construction
commutes with every base change and limit used below.
\end{construction}

\section{The big Thom-stable differential coefficient stack}
\label{sec:prequantum-big-differential-curvature}

The Chapter~9 categories $\mathcal R_U^0$ are defined on represented spectral
bases in the differential theorem range and satisfy finite spectral Nisnevich
descent.  Prequantum constructions must also vary over nonrepresented parameter
objects, including formal quotients and derived realization spaces.  We extend the
coefficient category itself by affine-point right Kan descent.

\subsection{Big differential coefficients}
\label{subsec:prequantum-big-differential-coefficients}

Let
\[
\mathfrak R^{0}_{\mathrm{aff}}:
(\operatorname{AffSp})^{\mathrm{op}}
\longrightarrow
\operatorname{CAlg}(\operatorname{Pr}^{L}_{\mathrm{st}}),
\qquad
h_U\longmapsto\mathcal R_U^0,
\]
be the Chapter~9 Thom-stable differential coefficient system, with transition
functors the constructed exact strong symmetric monoidal pullbacks.  Let
\[
h:\operatorname{AffSp}\longrightarrow\mathsf{XSp}
\]
be the affine Yoneda functor into the native big spectral Nisnevich
$\infty$-topos.  For $T\in\mathsf{XSp}$ write
$\mathcal C^{\mathrm{aff}}_{/T}$ for a fixed small skeleton of its affine-point
category.

\begin{definition}[Big Thom-stable differential coefficient category]
\label{def:prequantum-big-differential-coefficient-category}
Define the big differential coefficient category to be the affine right Kan
extension
\[
\boxed{
\mathcal R_T^{\partial}
:=
\operatorname{Ran}_{h^{\mathrm{op}}}
(\mathfrak R^{0}_{\mathrm{aff}})(T)
\simeq
\lim_{(h_U\to T)\in
      (\mathcal C^{\mathrm{aff}}_{/T})^{\mathrm{op}}}
\mathcal R_U^0.
}
\]
The limit is formed in
$\operatorname{CAlg}(\operatorname{Pr}^{L}_{\mathrm{st}})$.  Thus an object
$\widehat E\in\mathcal R_T^{\partial}$ is a Cartesian affine-point family
\[
(h_U\xrightarrow{t}T)\longmapsto\widehat E_t\in\mathcal R_U^0
\]
with coherent equivalences
\[
g^{*,0}\widehat E_t\simeq\widehat E_{t'}
\]
for every affine-point morphism
$(h_V\xrightarrow{t'}T)\xrightarrow{g}(h_U\xrightarrow{t}T)$.
\end{definition}

The right-Kan formulation is part of the construction.  In particular no
additional stack axiom for $T\mapsto\mathcal R_T^{\partial}$ is imposed; its
pullback, represented recovery, and descent properties are consequences of the
affine coefficient system and the universal property below.

\begin{lemma}[Affine Nisnevich basis recovery]
\label{lem:prequantum-affine-nisnevich-basis-recovery}
Let $\mathcal A$ be the category of affine spectral schemes and let
$\mathcal C$ be the category of represented qcqs spectral schemes in the
Chapter~9 differential range.  For every complete $\infty$-category
$\mathcal D$, restriction along $j:\mathcal A\hookrightarrow\mathcal C$
induces an equivalence
\[
\boxed{
j^*:
\operatorname{Shv}_{\mathrm{Nis}}(\mathcal C;\mathcal D)
\xrightarrow{\ \sim\ }
\operatorname{Shv}_{\mathrm{Nis}}(\mathcal A;\mathcal D),
}
\]
whose inverse is right Kan extension.
\end{lemma}

\begin{proof}
Every object $Z\in\mathcal C$ has a finite affine Zariski cover and hence a
finite affine Nisnevich cover.  If $U\to Z$ is affine and
$\{U_i\to Z\}$ is such a cover, every pullback $U\times_ZU_i$ is qcqs and
has a finite affine Zariski cover.  Thus $\mathcal A$ is a Nisnevich basis of
$\mathcal C$: every covering sieve is generated after pullback by covering
sieves whose domains are affine.

The dense-basis comparison lemma for $\infty$-sites now identifies the two
space-valued sheaf topoi.  The same statement holds for a complete target
$\mathcal D$: apply the space-valued comparison after postcomposition with
$\operatorname{Map}_{\mathcal D}(d,-)$ for every $d\in\mathcal D$.  These
functors preserve limits and jointly detect equivalences.  Under this
comparison, the inverse to restriction is necessarily the pointwise right Kan
extension
\[
(\operatorname{Ran}_jG)(Z)
\simeq
\lim_{(U\to Z)\in(\mathcal A_{/Z})^{\mathrm{op}}}G(U).
\]
This proves the assertion on objects, maps, and all higher mapping spaces.  It
is the dense-basis consequence of the universal property of presheaves and
sheaf localization; see \emph{Higher Topos Theory},
Theorem~5.1.5.6 and Sections~6.2.2--6.2.3.
\end{proof}

\begin{proposition}[Pullback and represented recovery]
\label{prop:prequantum-big-differential-recovery-pullback}
The assignment $T\mapsto\mathcal R_T^{\partial}$ is contravariantly functorial.
Every $f:T'\to T$ induces an exact strong symmetric monoidal pullback
\[
f^{*,\partial}:\mathcal R_T^{\partial}
\longrightarrow
\mathcal R_{T'}^{\partial}.
\]
If $T=h_U$ is affine represented, evaluation at the identity point gives
\[
\boxed{
\mathcal R_{h_U}^{\partial}\simeq\mathcal R_U^0.
}
\]
If $Z$ is a represented qcqs spectral scheme in the Chapter~9 differential range,
restriction gives a canonical equivalence
\[
\boxed{
\mathcal R_Z^{\partial}\simeq\mathcal R_Z^0.
}
\]
\end{proposition}

\begin{proof}
Composition with $f$ sends an affine point of $T'$ to one of $T$ and therefore
induces the restriction functor on the defining limit.  Exact and symmetric
monoidal structures are computed componentwise.

For $h_U$, the identity affine point is terminal in
$\mathcal C^{\mathrm{aff}}_{/h_U}$ and hence initial after taking opposites, so the
limit evaluates to $\mathcal R_U^0$.

For represented qcqs $Z$, the Chapter~9 assignment
$Z\mapsto\mathcal R_Z^0$ is a
$\operatorname{CAlg}(\operatorname{Pr}^{L}_{\mathrm{st}})$-valued finite
Nisnevich stack.  Lemma~\ref{lem:prequantum-affine-nisnevich-basis-recovery},
applied to this complete target, identifies it with the right Kan extension of
its affine restriction.  The latter right Kan extension is exactly the defining
limit for $\mathcal R_Z^{\partial}$.  This gives
$\mathcal R_Z^{\partial}\simeq\mathcal R_Z^0$ as an equivalence of stable
presentable symmetric monoidal categories.  Since the basis lemma is an
equivalence of $\infty$-categories rather than only of maximal subgroupoids,
the same calculation identifies all mapping spectra and proves full
faithfulness.
\end{proof}

\begin{theorem}[Effective descent of big differential coefficients]
\label{thm:prequantum-effective-descent-big-differential}
Let $u:T'\twoheadrightarrow T$ be an effective epimorphism in the native big
spectral Nisnevich $\infty$-topos, with \v{C}ech nerve $T'_\bullet$.  Restriction
induces a canonical exact symmetric monoidal equivalence
\[
\boxed{
\mathcal R_T^{\partial}
\simeq
\operatorname{Tot}_{[p]\in\Delta}
\mathcal R_{T'_p}^{\partial}.
}
\]
\end{theorem}

\begin{proof}
The affine coefficient stack extends, by
Lemma~\ref{lem:prequantum-affine-nisnevich-basis-recovery} and the universal
property of the Yoneda embedding, to the unique limit-preserving functor
\[
\mathfrak R^{\partial}:
\mathsf{XSp}^{\mathrm{op}}
\longrightarrow
\operatorname{CAlg}(\operatorname{Pr}^{L}_{\mathrm{st}})
\]
whose value is the affine-point right Kan formula of
Definition~\ref{def:prequantum-big-differential-coefficient-category}.  Before
Nisnevich localization this is the limit-preserving extension from the opposite
presheaf category.  Since the affine functor satisfies Nisnevich descent, the
extension inverts local equivalences and factors through
$\mathsf{XSp}^{\mathrm{op}}$; the factored functor remains limit preserving.

An effective epimorphism is the geometric realization of its augmented
\v{C}ech nerve.  Hence
\[
\mathcal R_T^{\partial}
=\mathfrak R^{\partial}(T)
\simeq
\mathfrak R^{\partial}(|T'_\bullet|)
\simeq
\lim_{[p]\in\Delta}\mathfrak R^{\partial}(T'_p).
\]
This is the displayed totalization.  The calculation occurs in
$\operatorname{CAlg}(\operatorname{Pr}^{L}_{\mathrm{st}})$, so it retains the
exact and symmetric-monoidal structures, maps, mapping spectra, and all higher
descent coherences.  The \v{C}ech-resolution characterization used here is
\emph{Higher Topos Theory}, Corollary~6.2.3.5.
\end{proof}

\subsection{Big differential mapping spectra and section spectra}
\label{subsec:prequantum-big-differential-mapping-spectra}

\begin{construction}[Big differential mapping spectrum]
\label{constr:prequantum-big-differential-mapping-spectrum}
For $\widehat E,\widehat F\in\mathcal R_T^{\partial}$ define on an affine point
$t:h_U\to T$
\[
\operatorname{map}^{\partial}_T(\widehat E,\widehat F)(U,t)
:=
\operatorname{map}_{\mathcal R_U^0}(\widehat E_t,\widehat F_t).
\]
Finite Nisnevich descent of $\mathcal R_U^0$ and the fact that mapping spectra in
a limit of stable $\infty$-categories are limits make this a spectral Nisnevich
sheaf on the affine-point site.  Hence it defines
\[
\boxed{
\operatorname{map}^{\partial}_T(\widehat E,\widehat F)\in E_T.
}
\]
Write $\mathbf1_T^{\partial}$ for the tensor unit of
$\mathcal R_T^{\partial}$ and define the differential section spectrum
\[
\boxed{
\operatorname{Sec}^{\partial}_T(\widehat E)
:=
\operatorname{map}^{\partial}_T
(\mathbf1_T^{\partial},\widehat E)\in E_T.
}
\]
\end{construction}

\begin{proposition}[Mapping-spectrum base change and global recovery]
\label{prop:prequantum-mapping-spectrum-base-change}
For every $f:T'\to T$ there is a canonical equivalence
\[
\boxed{
f^*\operatorname{map}^{\partial}_T(\widehat E,\widehat F)
\simeq
\operatorname{map}^{\partial}_{T'}
(f^{*,\partial}\widehat E,f^{*,\partial}\widehat F).
}
\]
Moreover
\[
\boxed{
\Gamma_T^{\mathrm{st}}
\operatorname{map}^{\partial}_T(\widehat E,\widehat F)
\simeq
\operatorname{map}_{\mathcal R_T^{\partial}}(\widehat E,\widehat F).
}
\]
If $T=Z$ is represented in the Chapter~9 differential range, this recovers
\[
\Gamma_Z^{\mathrm{st}}
\operatorname{map}^{\partial}_Z(\widehat E,\widehat F)
\simeq
\operatorname{map}_{\mathcal R_Z^0}(\widehat E,\widehat F).
\]
\end{proposition}

\begin{proof}
The base-change equivalence is pointwise on affine tests.  Both global mapping
spectra are the limit, over the affine-point category, of the same affine mapping
spectra.  Represented recovery follows from
Proposition~\ref{prop:prequantum-big-differential-recovery-pullback}.
\end{proof}

\section{Stable curvature diagrams and closed differential classes}
\label{sec:prequantum-stable-curvature-diagrams}

The complete curvature problem is formed in the parameterized stable category.
This avoids two false identifications: arbitrary pullback exchange for the
Thom-stable realization of the de Rham package is not promoted to Cartesianity,
and the field-theoretic presymplectic target is not replaced by a scalar de Rham
coefficient.  It also prevents a curvature map or an underlying-class map from
being hidden inside the name of a coefficient.

\begin{construction}[Curvature-arrow realization]
\label{constr:prequantum-curvature-realization-object}
For $\mathscr D,\mathscr U\in E_T$ define
\[
\boxed{
\operatorname{CurvReal}_T(\mathscr D,\mathscr U)
:=
\Omega_T^\infty
\underline{\operatorname{Hom}}_{E_T}(\mathscr D,\mathscr U).
}
\]
Over this parameter object, evaluation constructs the universal morphism
\[
\kappa_{\mathrm{univ}}:
\mathscr D_{\mathrm{univ}}\longrightarrow\mathscr U_{\mathrm{univ}}.
\]
\end{construction}

\begin{construction}[Stable curvature-diagram realization]
\label{constr:prequantum-curvature-diagram-realization}
For $\mathscr D,\mathscr U,\mathscr C\in E_T$ define
\[
\boxed{
\operatorname{CurvDiagReal}_T(\mathscr D,\mathscr U,\mathscr C)
:=
\operatorname{CurvReal}_T(\mathscr D,\mathscr U)
\times_T
\operatorname{CurvReal}_T(\mathscr C,\mathscr U).
}
\]
Its universal point is a span
\[
\mathscr D_{\mathrm{univ}}
\xrightarrow{\kappa_{\mathrm{univ}}}
\mathscr U_{\mathrm{univ}}
\xleftarrow{u_{\mathrm{univ}}}
\mathscr C_{\mathrm{univ}}.
\]
If one of the two arrows has already been constructed, all later constructions are
pulled back along the corresponding section of this realization and that parameter
is suppressed.
\end{construction}

\begin{definition}[Stable prequantum curvature diagram]
\label{def:prequantum-stable-curvature-diagram}
A stable prequantum curvature diagram over $T$ is an actual point of
$\operatorname{CurvDiagReal}_T(\mathscr D,\mathscr U,\mathscr C)$, written
\[
\mathscr D\xrightarrow{\kappa}\mathscr U\xleftarrow{u}\mathscr C.
\]
The object $\mathscr D$ is the differential section coefficient,
$\mathscr C$ is the complete coherently closed coefficient, and $\mathscr U$ is
the common underlying coefficient in which their comparison is formed.
\end{definition}

The special case $\mathscr U=\mathscr C$ and $u=\mathrm{id}$ recovers the usual
picture of a curvature map landing directly in closed classes.  The more general
span is needed for variational transgression: the complete closed coefficient has
a constructed underlying-transgression map to the unclosed transgression
coefficient.

\begin{definition}[Closed differential class coefficient]
\label{def:prequantum-closed-differential-class-coefficient}
For a curvature diagram define
\[
\boxed{
\mathscr D^{\mathrm{cl}}_{\kappa,u}
:=
\mathscr D\times_{\mathscr U}\mathscr C
\in E_T.
}
\]
Write
\[
\operatorname{curv}_{\kappa,u}:
\mathscr D^{\mathrm{cl}}_{\kappa,u}
\longrightarrow
\mathscr C
\]
for the second projection.
\end{definition}

A point of $\Omega_T^\infty\mathscr D^{\mathrm{cl}}_{\kappa,u}[r]$ is therefore a
differential class $\widehat\alpha\in\Omega_T^\infty\mathscr D[r]$, a coherently
closed class $\alpha^{\mathrm{cl}}\in\Omega_T^\infty\mathscr C[r]$, and a
specified homotopy
\[
\kappa\widehat\alpha\simeq u\alpha^{\mathrm{cl}}
\]
in $\mathscr U[r]$.  Closedness is part of the point and is not reduced to the
unary equation in an associated-graded cochain complex.

\begin{proposition}[Pullback and descent of curvature diagrams]
\label{prop:prequantum-curvature-diagram-pullback-descent}
Curvature-arrow and curvature-diagram realizations, closed differential class
coefficients, and their displayed structure maps commute with arbitrary native
pullback.  They satisfy effective descent whenever the underlying coefficients do.
\end{proposition}

\begin{proof}
Chapter~10 closed base change gives
\[
f^*\underline{\operatorname{Hom}}(A,B)
\simeq
\underline{\operatorname{Hom}}(f^*A,f^*B)
\]
coherently in $f,A,B$.  Relative zeroth space and finite fibre products therefore
transport the realization objects.  The closed differential class coefficient is
a finite stable limit, and effective descent in $E$ is computed by the same limit
operations.
\end{proof}

\section{Parameterized prequantization, obstruction, and ambiguity}
\label{sec:prequantum-parameterized-prequantization-realization}

Fix a stable curvature diagram
\[
\mathscr D\xrightarrow{\kappa}\mathscr U\xleftarrow{u}\mathscr C
\]
over $T$.  All definitions remain parameterized in both displayed arrows by pulling back
from $\operatorname{CurvDiagReal}_T(\mathscr D,\mathscr U,\mathscr C)$ whenever
they have not already been constructed.

\begin{definition}[Differential and closed class moduli]
\label{def:prequantum-differential-and-closed-class-moduli}
For $r\in\mathbb Z$ put
\[
\operatorname{DiffClCls}_T^r(\kappa,u)
:=
\Omega_T^\infty
\mathscr D^{\mathrm{cl}}_{\kappa,u}[r],
\qquad
\operatorname{ClCls}_T^r(\mathscr C)
:=
\Omega_T^\infty\mathscr C[r].
\]
The projection defines
\[
\operatorname{curv}_{\kappa,u}:
\operatorname{DiffClCls}_T^r(\kappa,u)
\longrightarrow
\operatorname{ClCls}_T^r(\mathscr C).
\]
\end{definition}

\begin{definition}[Prequantization of a closed class]
\label{def:prequantum-prequantization-of-closed-class}
Let $P\to T$ and let
\[
\omega:P\longrightarrow\operatorname{ClCls}_T^r(\mathscr C)
\]
be a coherently closed class.  Define
\[
\boxed{
\operatorname{PreQ}_{\kappa,u}(\omega)
:=
P\times_{\operatorname{ClCls}_T^r(\mathscr C)}
\operatorname{DiffClCls}_T^r(\kappa,u).
}
\]
A point consists of a differential class, a closed curvature representative, and
specified identifications with $\omega$.
\end{definition}

\begin{definition}[Flat ambiguity coefficient]
\label{def:prequantum-flat-differential-ambiguity}
Define
\[
\boxed{
\mathscr F_\kappa:=\operatorname{fib}
(\mathscr D\xrightarrow{\kappa}\mathscr U)
\in E_T.
}
\]
\end{definition}

\begin{proposition}[Pulled-back fibre sequence]
\label{prop:prequantum-pulled-back-fibre-sequence}
There is a canonical stable fibre sequence
\[
\boxed{
\mathscr F_\kappa
\longrightarrow
\mathscr D^{\mathrm{cl}}_{\kappa,u}
\xrightarrow{\operatorname{curv}_{\kappa,u}}
\mathscr C
\xrightarrow{\partial_{\kappa,u}}
\mathscr F_\kappa[1].
}
\]
\end{proposition}

\begin{proof}
The defining square for
$\mathscr D^{\mathrm{cl}}_{\kappa,u}=\mathscr D\times_{\mathscr U}\mathscr C$
is a pullback.  Pulling the fibre sequence
$\mathscr F_\kappa\to\mathscr D\to\mathscr U$ back along
$u:\mathscr C\to\mathscr U$ gives the displayed fibre sequence.  Stability
identifies its fibre and cofiber connecting morphisms.
\end{proof}

\begin{construction}[Prequantization obstruction]
\label{constr:prequantum-obstruction-section}
Define
\[
\boxed{
\mathfrak o_{\mathrm{PreQ}}(\omega)
:=
\partial_{\kappa,u}[r]\,\omega:
P\longrightarrow
\Omega_T^\infty\mathscr F_\kappa[r+1].
}
\]
\end{construction}

\begin{theorem}[Obstruction recognition]
\label{thm:prequantum-obstruction-recognition}
There is a canonical equivalence
\[
\boxed{
\operatorname{PreQ}_{\kappa,u}(\omega)
\simeq
Z_P\bigl(\mathfrak o_{\mathrm{PreQ}}(\omega)\bigr).
}
\]
Thus the prequantization space is inhabited precisely when the constructed
connecting section admits a nullhomotopy; a point of the zero locus retains that
nullhomotopy and all of its higher compatibility data.
\end{theorem}

\begin{proof}
For a fibre sequence $F\to E\to C\xrightarrow{\partial}F[1]$ in a stable
$\infty$-category, a lift of a point of $C$ to $E$ is equivalent to a specified
nullhomotopy of its connecting image.  Apply this universal property to
Proposition~\ref{prop:prequantum-pulled-back-fibre-sequence}, shift by $r$, and
apply relative zeroth space.
\end{proof}

\begin{theorem}[Flat torsor and higher automorphisms]
\label{thm:prequantum-flat-torsor}
If a fibre $\operatorname{PreQ}_{\kappa,u}(\omega)_x$ is inhabited, it carries a
canonical simply transitive action of
\[
\boxed{
\Omega^\infty\mathscr F_\kappa[r]|_x.
}
\]
After choosing a prequantization $\widehat\omega$, the $j$th higher automorphism
space is canonically
\[
\boxed{
\operatorname{Aut}^{(j)}(\widehat\omega)
\simeq
\Omega^\infty\mathscr F_\kappa[r-j]|_x,
\qquad j\ge1.
}
\]
In particular ordinary automorphisms are governed by
$\Omega^\infty\mathscr F_\kappa[r-1]$.
\end{theorem}

\begin{proof}
Work internally over the selected base point $x$.  Put
\[
E_x:=
\operatorname{PreQ}_{\kappa,u}(\omega)_x,
\qquad
G_x:=\Omega^\infty\mathscr F_\kappa[r]|_x.
\]
Addition in the infinite-loop objects restricts to an action
$a:G_x\times E_x\to E_x$.  The action is principal because
\[
(a,\operatorname{pr}_2):
G_x\times E_x\longrightarrow E_x\times E_x
\]
is an equivalence.  Its inverse sends $(e_1,e_2)$ to
$(e_1-e_2,e_2)$; the difference lies in $G_x$ since both $e_1$ and $e_2$
map to the same closed class.  Thus the torsor is defined without choosing a
point.  A selected prequantization $\widehat\omega$ gives the translation
equivalence $G_x\simeq E_x$.  Iterated based looping then gives
\[
\Omega^j_{\widehat\omega}E_x
\simeq
\Omega^jG_x
\simeq
\Omega^\infty\mathscr F_\kappa[r-j]|_x,
\]
which is the stated higher-automorphism formula.
\end{proof}

\section{Universal prequantization over closed-class moduli}
\label{sec:prequantum-universal-prequantization}

\begin{construction}[Universal prequantization family]
\label{constr:prequantum-universal-prequantization-family}
Let
\[
\mathscr C_T^r:=\operatorname{ClCls}_T^r(\mathscr C)
\]
with its tautological point
\[
\omega_{\mathrm{univ}}:\mathscr C_T^r
\longrightarrow
\operatorname{ClCls}_T^r(\mathscr C).
\]
Define
\[
\boxed{
\operatorname{PreQMod}_{\kappa,u}^r
:=
\operatorname{PreQ}_{\kappa,u}(\omega_{\mathrm{univ}})
\longrightarrow
\mathscr C_T^r.
}
\]
Equivalently it is the curvature map
$\operatorname{DiffClCls}_T^r(\kappa,u)\to\operatorname{ClCls}_T^r(\mathscr C)$,
regarded as a family over its target.
\end{construction}

\begin{theorem}[Individual prequantizations are pullbacks]
\label{thm:prequantum-individual-as-pullback}
For every $\omega:P\to\mathscr C_T^r$,
\[
\boxed{
\omega^*\operatorname{PreQMod}_{\kappa,u}^r
\simeq
\operatorname{PreQ}_{\kappa,u}(\omega).
}
\]
Paths and higher homotopies of closed classes pull the universal family to the
corresponding paths and higher homotopies of prequantization problems.
\end{theorem}

\begin{proof}
This is fibre-pasting for the defining homotopy pullback.
\end{proof}

\section{Differential Lagrangian refinement and presymplectic prequantization}
\label{sec:prequantum-lagrangians}

The field-theoretic prequantum coefficient is obtained from the actual
Chapter~12 stable differential density sector.  No tensor product between a
Thom-stable differential object and an unrelated parameterized stable coefficient
is introduced.  There are two successive refinement problems and they are kept
separate: first the Lagrangian itself is lifted through a coefficient morphism,
and then a chosen relative Lepage realization is lifted through the induced map of
Lepage moduli.  The second problem is essential because an arbitrary point of
$\operatorname{Lep}^q(L)$ is not determined by a lift of $L$ alone.

\begin{definition}[Differential Lagrangian curvature realization]
\label{def:prequantum-differential-lagrangian-curvature-realization}
The Chapter~12 construction is functorial on the maximal subgroupoid and gives
\[
M^\partial_{x,-}:(\mathcal R_S^0)^\simeq\longrightarrow(E_Q)^\simeq.
\]
Define the differential-density curvature realization by the Grothendieck
construction
\[
\boxed{
\operatorname{DiffDenCurvReal}_x(M)
:=
\int_{D\in(\mathcal R_S^0)^\simeq}
\Omega_Q^\infty
\underline{\operatorname{Hom}}_{E_Q}(M^\partial_{x,D},M).
}
\]
Over it there are a universal differential coefficient
$\widehat M_{\mathrm{univ}}=M^\partial_{x,D_{\mathrm{univ}}}$ and a universal
coefficient morphism
\[
\kappa_{M,\mathrm{univ}}:\widehat M_{\mathrm{univ}}\longrightarrow M.
\]
For a fixed $D\in\mathcal R_S^0$, its fibre is
\[
\boxed{
\operatorname{CurvLag}_Q(M^\partial_{x,D},M)
:=
\Omega_Q^\infty
\underline{\operatorname{Hom}}_{E_Q}(M^\partial_{x,D},M).
}
\]
\end{definition}

\begin{proposition}[Universal property of differential-density curvature]
\label{prop:prequantum-differential-density-curvature-universal-property}
For every test $T\to Q$, maps
$T\to\operatorname{DiffDenCurvReal}_x(M)$ classify a family
$D_T\in(\mathcal R_S^0)^\simeq$ together with a coefficient morphism
$M^\partial_{x,D_T}|_T\to M|_T$.  The realization commutes with native base
change and satisfies effective descent.
\end{proposition}

\begin{proof}
The first assertion is the universal property of the Grothendieck construction
applied to the displayed functor of mapping objects.  The Chapter~12 coefficient
construction and the internal Hom commute with their constructed pullbacks.
Affine-testwise descent reconstructs $D_T$, its coefficient, and the morphism;
hence it reconstructs the Grothendieck construction.
\end{proof}

Thus a relation between a differential density coefficient and the already fixed
variational coefficient is not assumed, nor is a differential coefficient
supplied under a provenance label.  The differential origin $D$, its constructed
coefficient, and the entire derived mapping object are retained universally.

\begin{definition}[Differential Lagrangian refinement object]
\label{constr:prequantum-differential-lagrangian-refinement}
After pullback along a point
$(D,\kappa_M)$ of $\operatorname{DiffDenCurvReal}_x(M)$, write
$\widehat M:=M^\partial_{x,D}$ and put
\[
\operatorname{LagCls}_Q^n(\widehat M)
:=
\Omega_Q^\infty
\underline{\operatorname{Hom}}_{E_Q}
\bigl(\mathbf1_Q,\operatorname{Sec}^0(A;\widehat M)[n]\bigr)
\]
and similarly for $M$.  Functoriality in the coefficient gives
\[
\kappa_{\mathrm{Lag}}:
\operatorname{LagCls}_Q^n(\widehat M)
\longrightarrow
\operatorname{LagCls}_Q^n(M).
\]
Regard the fixed Lagrangian as the section
$L:Q\to\operatorname{LagCls}_Q^n(M)$ and define the internal homotopy pullback
\[
\boxed{
P_{\mathrm{Lag}}
=
\operatorname{PreQLag}_{\kappa_M}(L)
:=
Q\times_{\operatorname{LagCls}_Q^n(M)}
\operatorname{LagCls}_Q^n(\widehat M).
}
\]
Over $P_{\mathrm{Lag}}$ there are a tautological differential Lagrangian
$\widehat L_{\mathrm{univ}}$ and a tautological specified homotopy
\[
\eta_{\mathrm{Lag}}:
\kappa_{\mathrm{Lag}}(\widehat L_{\mathrm{univ}})
\simeq
L.
\]
\end{definition}

The possible nonexistence of a differential Lagrangian refinement is therefore
represented by an actual internal fibre.  No ``differentially refinable
Lagrangian'' subcategory is introduced.

\subsection{Differential Lepage refinement}
\label{subsec:prequantum-differential-lepage-refinement}

Chapter~12 constructs relative Lepage realizations functorially in coefficient
morphisms and in the Lagrangian.  A lift of the Lagrangian therefore induces a
map of Lepage moduli, but it does not choose a lift of an arbitrary point of the
underlying Lepage moduli.  We represent that second lifting problem.

Base-change the Cartan object, coefficients, and variational calculus along
$P_{\mathrm{Lag}}\to Q$.  Apply the Chapter~12 relative Lepage construction to
$\widehat L_{\mathrm{univ}}$ and define
\[
\boxed{
\operatorname{PreQLep}_{\kappa_M}^{q}(L)
:=
\operatorname{Lep}^{q}
\bigl(\widehat L_{\mathrm{univ}}\bigr)
\longrightarrow
P_{\mathrm{Lag}}.
}
\]
Functoriality in $\kappa_M$ gives
\[
\operatorname{Lep}^{q}(\widehat L_{\mathrm{univ}})
\longrightarrow
\operatorname{Lep}^{q}
\bigl(\kappa_M\widehat L_{\mathrm{univ}}\bigr).
\]
Transport along $\eta_{\mathrm{Lag}}$ and the Chapter~12 universal-pullback
identification for Lepage families gives a canonical comparison
\[
\boxed{
c_{\mathrm{Lep}}:
\operatorname{PreQLep}_{\kappa_M}^{q}(L)
\longrightarrow
P_{\mathrm{Lag}}\times_Q\operatorname{Lep}^{q}(L).
}
\]

\begin{proposition}[Meaning of differential Lepage refinement]
\label{prop:prequantum-differential-lepage-refinement-meaning}
Let $T\to P_{\mathrm{Lag}}$ classify
$(\widehat L,\eta_{\mathrm{Lag}})$ and let
$\ell:T\to\operatorname{Lep}^{q}(L)$ be an underlying relative Lepage
realization.  The homotopy fibre of $c_{\mathrm{Lep}}$ over
$(\widehat L,\eta_{\mathrm{Lag}},\ell)$ is the derived space of relative Lepage
realizations $\widehat\ell$ of $\widehat L$ whose image under $\kappa_M$ is
identified with $\ell$, including all higher compatibility data of the Chapter~12
homotopy-pullback definitions.  The fibre may be empty.
\end{proposition}

\begin{proof}
A point of $\operatorname{Lep}^{q}(\widehat L)$ is, by the Chapter~12
representing property, a coherently closed full lift of the first Euler source
together with a stage-$q$ horizontal realization of its difference from the
canonical full Euler lift.  Applying $\kappa_M$ carries both pieces and their
higher homotopies to the corresponding data for $\kappa_M\widehat L$.  Transport
through $\eta_{\mathrm{Lag}}$ identifies the target with the pulled-back Lepage
object of $L$.  Taking the homotopy fibre over $\ell$ therefore imposes exactly
the stated comparison and nothing else.
\end{proof}

The distinguished zero-comparison point of every Chapter~12 Lepage object is
functorial in coefficient morphisms.  Consequently every point of
$P_{\mathrm{Lag}}$ has a canonical lift of that distinguished Lepage point.  A
lift of an arbitrary Lepage realization, however, is measured by the fibre above
and is not inferred from the Lagrangian refinement alone.

\subsection{Complete closed presymplectic coefficient}
\label{subsec:prequantum-complete-closed-presymplectic-coefficient}

The universal presymplectic transgression of Chapter~12 already carries the full
coherent vertical closedness hierarchy.  We package that hierarchy as the closed
coefficient needed for prequantization.

Let
\[
C_{\mathrm{var}}^{r,s}(A;M)
\]
be the ambient vertical--horizontal bicosimplicial stable coefficient of
Chapter~12, but leave the vertical degree $r$ unnormalized.  For every vertical
degree put
\[
S_L^{r}(M)
:=
\operatorname{Tot}_H C_{\mathrm{var}}^{r,\bullet}(A;M),
\qquad
S_L^{r,<q}(M)
:=
\operatorname{Tot}_{H,\le q-1}
C_{\mathrm{var}}^{r,\bullet}(A;M),
\]
and define the stage-$q$ horizontal tail
\[
F_H^qS_L^{r}(M)
:=
\operatorname{fib}
\left(
S_L^{r}(M)
\longrightarrow
S_L^{r,<q}(M)
\right).
\]
Now define the vertical cosimplicial stage-$q$ transgression object degreewise by
\[
\boxed{
\mathbb T_L^{r,q}(M)
:=
\operatorname{fib}
\left(
F_H^qS_L^{r}(M)
\longrightarrow
S_L^{r}(M)
\right).
}
\]
Naturality in the unnormalized vertical simplicial operators makes
$r\mapsto\mathbb T_L^{r,q}(M)$ a vertical cosimplicial stable object.  Vertical
normalization is a finite matching fibre and therefore commutes with the
horizontal totalizations, horizontal tail, and transgression fibre.  Hence
\[
N_V^r\mathbb T_L^{\bullet,q}(M)
\simeq
\operatorname{Trans}^{r,q}_H(A;M),
\]
in particular
\[
N_V^2\mathbb T_L^{\bullet,q}(M)
\simeq
\operatorname{Trans}^{2,q}_H(A;M).
\]
Form the complete vertical filtration
\[
F_V^r\operatorname{Tot}_V\mathbb T_L^{\bullet,q}(M)
:=
\operatorname{fib}
\left(
\operatorname{Tot}_V\mathbb T_L^{\bullet,q}(M)
\longrightarrow
\operatorname{Tot}_{V,\le r-1}\mathbb T_L^{\bullet,q}(M)
\right).
\]

\begin{definition}[Complete closed presymplectic coefficient]
\label{def:prequantum-complete-closed-presymplectic-coefficient}
Define
\[
\boxed{
\operatorname{ClTrans}^{2,q}_H(A;M;n)
:=
F_V^2\operatorname{Tot}_V
\mathbb T_L^{\bullet,q}(M)[n+2].
}
\]
Stable Dold--Kan gives the underlying-transgression morphism
\[
\boxed{
u_{L,q}:
\operatorname{ClTrans}^{2,q}_H(A;M;n)
\longrightarrow
\operatorname{Trans}^{2,q}_H(A;M)[n].
}
\]
\end{definition}

\begin{proposition}[Closed lift of the universal presymplectic transgression]
\label{prop:prequantum-closed-lift-presymplectic-transgression}
The complete vertical closedness of $\omega_L^{(q)}$ canonically determines a
section
\[
\boxed{
\omega_L^{(q),\mathrm{cl}}:
\mathbf1_{\operatorname{Lep}^q(L)}
\longrightarrow
\operatorname{ClTrans}^{2,q}_H(A;M;n)
}
\]
whose image under $u_{L,q}$ is $\omega_L^{(q)}$.
\end{proposition}

\begin{proof}
Chapter~12 constructs $\omega_L^{(q)}$ from a vertical coherent cochain and proves
that its first vertical boundary is nullhomotopic with the entire higher
compatibility hierarchy inherited from the complete vertical relation.  The
complete-filtration/coherent-cochain equivalence of Chapter~10 identifies exactly
such data with a point of
$F_V^2\operatorname{Tot}_V\mathbb T_L^{\bullet,q}(M)[n+2]$ whose associated
normalized degree-two section is $\omega_L^{(q)}$.  This gives the displayed lift
and its full coherence.
\end{proof}

\subsection{Variational curvature morphism and presymplectic prequantization}
\label{subsec:prequantum-variational-curvature-morphism}

A point $\kappa_M:\widehat M\to M$ of
$\operatorname{CurvLag}_Q(\widehat M,M)$ acts functorially on every coefficient
construction of the variational chapter.  In particular it gives
\[
\kappa_{\mathrm{tr}}^{(q)}:
\operatorname{Trans}^{2,q}_H(A;\widehat M)[n]
\longrightarrow
\operatorname{Trans}^{2,q}_H(A;M)[n].
\]
Put
\[
P_\omega^q:=\operatorname{Lep}^q(L)
\]
and, after pullback to $P_\omega^q$, write
\[
\mathscr D_{L,q}
:=
\operatorname{Trans}^{2,q}_H(A;\widehat M)[n],
\qquad
\mathscr U_{L,q}
:=
\operatorname{Trans}^{2,q}_H(A;M)[n],
\]
\[
\mathscr C_{L,q}
:=
\operatorname{ClTrans}^{2,q}_H(A;M;n).
\]
Then
\[
\mathscr D_{L,q}
\xrightarrow{\kappa_{\mathrm{tr}}^{(q)}}
\mathscr U_{L,q}
\xleftarrow{u_{L,q}}
\mathscr C_{L,q}
\]
is a stable prequantum curvature diagram over $P_\omega^q$.

\begin{theorem}[Curvature commutes with differential Lepage transgression]
\label{thm:prequantum-curvature-variational-transgression}
Over $\operatorname{PreQLep}_{\kappa_M}^{q}(L)$, let
$\widehat\ell_{\mathrm{univ}}$ be the tautological differential Lepage realization
and let $\ell_{\mathrm{univ}}$ be its underlying image in
$\operatorname{Lep}^{q}(L)$.  The Chapter~12 presymplectic construction gives
\[
\omega_{\widehat L}^{(q)}:
\mathbf1\longrightarrow
\operatorname{Trans}^{2,q}_H(A;\widehat M)[n]
\]
and there is a canonical coherent homotopy
\[
\boxed{
\kappa_{\mathrm{tr}}^{(q)}
\bigl(\omega_{\widehat L}^{(q)}\bigr)
\simeq
u_{L,q}
\bigl(\omega_L^{(q),\mathrm{cl}}\bigr)
=
\omega_L^{(q)}
}
\]
after pullback along $\ell_{\mathrm{univ}}$.  Consequently the pair
$\bigl(\omega_{\widehat L}^{(q)},\omega_L^{(q),\mathrm{cl}}\bigr)$ determines a
canonical point of the closed differential pullback
\[
\mathscr D_{L,q}\times_{\mathscr U_{L,q}}\mathscr C_{L,q}.
\]
\end{theorem}

\begin{proof}
Every stage of the Chapter~12 variational construction is functorial in
coefficient morphisms.  Apply that functoriality to the actual differential
Lepage realization $\widehat\ell_{\mathrm{univ}}$.  Its image is identified by
the defining point of $\operatorname{PreQLep}$ with the pulled-back underlying
Lepage realization, so naturality of vertical variation and horizontal
transgression gives
$\kappa_{\mathrm{tr}}^{(q)}\omega_{\widehat L}^{(q)}\simeq\omega_L^{(q)}$.
Proposition~\ref{prop:prequantum-closed-lift-presymplectic-transgression}
identifies the latter with the underlying image of the complete closed lift.  The
homotopy-pullback universal property gives the final closed differential section.
No preservation of an unconstructed sequential inverse limit by Thom-stable
realization is used.
\end{proof}

\begin{definition}[Presymplectic prequantization family]
\label{def:prequantum-presymplectic-prequantization-family}
Define
\[
\boxed{
\operatorname{PreQ}_{\omega,L}^{q}
:=
\operatorname{PreQ}_{\kappa_{\mathrm{tr}}^{(q)},u_{L,q}}
\bigl(\omega_L^{(q),\mathrm{cl}}\bigr)
\longrightarrow
P_\omega^q.
}
\]
The preceding theorem induces a canonical map
\[
\boxed{
\operatorname{PreQLep}_{\kappa_M}^{q}(L)
\longrightarrow
\operatorname{PreQ}_{\omega,L}^{q}
}
\]
over its underlying map to $P_\omega^q$.
\end{definition}

Thus a differential Lagrangian refinement together with an actual differential
Lepage refinement constructs a prequantization of the corresponding underlying
presymplectic transgression.  The canonical zero-comparison Lepage point gives a
canonical special case for every differential Lagrangian refinement.  For an
arbitrary point of $\operatorname{Lep}^{q}(L)$, the complete lift space is
exactly the fibre represented by
$\operatorname{PreQLep}_{\kappa_M}^{q}(L)$, so existence and ambiguity are
already part of the constructed moduli problem.

\section{The prequantum action and fundamental-class realization}
\label{sec:prequantum-action-functional}

Work on the represented smooth separated finite-presentation differential action
domain of Chapter~12, and specialize $\widehat M$ to the Chapter~12 differential
density coefficient $M^\partial_{x,D}$.  Let $D\in\mathcal R_S^0$ be the
differential coefficient used there and let $M^\partial_{x,D}\in E_Q$ and
$\mathcal A_{x,D}^{\partial}\in E_S$ be the differential density and action-value
coefficients already constructed.  We retain the Chapter~12 notation
$x_{e!}\dashv x_e^!$ for the differential exceptional adjunction.

\begin{construction}[Orientation-free prequantum action]
\label{constr:prequantum-action}
A differential Lagrangian $\widehat L$ determines the Chapter~12 morphism of field
moduli
\[
\boxed{
\widehat S_{\widehat L}^{\mathrm{or}}:
\operatorname{Fld}_S(A)
\longrightarrow
\Omega_S^\infty\mathcal A_{x,D}^{\partial}[n].
}
\]
On a global field $\phi$, its value is represented by the differential exceptional
morphism
\[
\boxed{
x_{e!}\mathbf1_X^{\partial}
\longrightarrow
D[n].
}
\]
This is the orientation-free prequantum action.
\end{construction}

\begin{definition}[Differential fundamental-class candidate realization]
\label{def:prequantum-fundamental-class-realization}
Define the parameterized realization
\[
\boxed{
\operatorname{FundReal}_{\partial}(x)
:=
\Omega_S^\infty
\operatorname{map}^{\partial}_S
\bigl(\mathbf1_S^{\partial},x_{e!}\mathbf1_X^{\partial}\bigr).
}
\]
After pullback to this realization, evaluation gives the tautological differential
fundamental-class candidate
\[
[X]^\partial_{\mathrm{univ}}:
\mathbf1_S^{\partial}
\longrightarrow
x_{e!}\mathbf1_X^{\partial}.
\]
\end{definition}

The zero morphism gives a canonical basepoint of
$\operatorname{FundReal}_{\partial}(x)$.  Thus the unrestricted mapping
realization is never an existence obstruction and a point is not, by itself, an
orientation.

\begin{construction}[Differential lift of a prescribed motivic fundamental class]
\label{constr:prequantum-fundamental-class-shadow-lift}
Define
\[
\operatorname{FundReal}_{\mathrm{mot}}(x)\in\mathsf{XSp}_{/S}
\]
by the affine-point rule
\[
\operatorname{FundReal}_{\mathrm{mot}}(x)(T\to S)
:=
\operatorname{Map}_{\operatorname{SH}(T)}
\left(\mathbf1_T,x_{T,!}\mathbf1_{X_T}\right),
\]
followed by spectral Nisnevich descent.  Here
$x_T:X\times_ST\to T$.  Stable affine-point presentation and motivic descent
make this a represented object of the native slice.
Applying homotopification to a differential candidate and using the constructed
exceptional exchange
\[
H_Sx_{e!}\mathbf1_X^\partial\simeq x_!H_X\mathbf1_X^\partial
\simeq x_!\mathbf1_X
\]
gives a canonical shadow map
\[
\operatorname{sh}_{\mathrm{Fund},x}:
\operatorname{FundReal}_{\partial}(x)
\longrightarrow
\operatorname{FundReal}_{\mathrm{mot}}(x).
\]
For a constructed or selected motivic class
$\eta^{\mathrm{mot}}:P\to\operatorname{FundReal}_{\mathrm{mot}}(x)$ define
\[
\boxed{
\operatorname{FundLiftReal}_{\partial}(x;\eta^{\mathrm{mot}})
:=
P\times_{\operatorname{FundReal}_{\mathrm{mot}}(x)}
\operatorname{FundReal}_{\partial}(x).
}
\]
A point is a differential candidate together with a specified identification of
its motivic shadow with $\eta^{\mathrm{mot}}$.  This lift realization may be
empty.
\end{construction}

\begin{construction}[Orientation inverse locus]
\label{constr:prequantum-orientation-inverse-locus}
Fix an internal stable $\infty$-category
\[
\mathcal C\longrightarrow S
\]
and objects $A,B\in\mathcal C$ over $S$, equivalently sections
$A,B:S\to\mathcal C$.  Regard the internal morphism object and its
subobject of equivalences as objects of the native slice:
\[
\mathcal H_{A,B}
:=
\Omega_S^\infty
\underline{\operatorname{Hom}}_{\mathcal C}(A,B)
\in\mathsf{XSp}_{/S},
\qquad
\operatorname{Eqv}_{\mathcal C}(A,B)
\longrightarrow
\mathcal H_{A,B}.
\]
Let
\[
\mathcal F_x:=\operatorname{FundReal}_{\partial}(x)
\in\mathsf{XSp}_{/S}.
\]
Instead of requiring a branch to carry an assignment by name, define its entire
assignment realization
\[
\boxed{
\operatorname{DualAssignReal}_x(A,B)
:=
\underline{\operatorname{Map}}_S(\mathcal F_x,\mathcal H_{A,B})
\in\mathsf{XSp}_{/S}.
}
\]
The internal evaluation morphism
\[
\operatorname{ev}:
\mathcal F_x\times_S\operatorname{DualAssignReal}_x(A,B)
\longrightarrow\mathcal H_{A,B}
\]
gives the universal pair $(\eta,\theta_\eta)$.  Define
\[
\boxed{
\operatorname{OrientReal}_{\mathrm{univ}}(A,B)
:=
\bigl(\mathcal F_x\times_S\operatorname{DualAssignReal}_x(A,B)\bigr)
\mathop{\times}_{\mathcal H_{A,B},\,\operatorname{ev}}
\operatorname{Eqv}_{\mathcal C}(A,B).
}
\]
Here the walking-equivalence fibre and its forward-arrow map are those of
Construction~\ref{constr:prequantum-walking-equivalence-inverse-locus} in the
ambient internal stable category $\mathcal C$.

An assignment $\theta:\mathcal F_x\to\mathcal H_{A,B}$ over $S$ constructed by
an exceptional-duality formula has an exponential transpose
\[
\bar\theta:S\longrightarrow\operatorname{DualAssignReal}_x(A,B).
\]
Pullback of the universal orientation locus along the graph of $\bar\theta$
defines
\[
\begin{aligned}
\operatorname{OrientReal}(\theta)
&:=
\mathcal F_x
\mathop{\times}_{\mathcal F_x\times_S
\operatorname{DualAssignReal}_x(A,B)}
\operatorname{OrientReal}_{\mathrm{univ}}(A,B)\\
&\simeq
\mathcal F_x\mathop{\times}_{\mathcal H_{A,B},\theta}
\operatorname{Eqv}_{\mathcal C}(A,B).
\end{aligned}
\]
Thus the chapter neither assumes that such an assignment exists nor suppresses
the space of possible assignments.  When an earlier duality theorem supplies a
specific formula, its constructed section selects the corresponding fibre.

For a prescribed motivic class define the base change
\[
\boxed{
\operatorname{OrientReal}(\theta;\eta^{\mathrm{mot}})
:=
\operatorname{FundLiftReal}_{\partial}(x;\eta^{\mathrm{mot}})
\times_{\operatorname{FundReal}_{\partial}(x)}
\operatorname{OrientReal}(\theta).
}
\]
A point is exactly a differential lift of the prescribed motivic class for
which the induced duality morphism is equipped with coherent inverse data.
The universal version additionally retains the assignment itself.  Thus neither
the assignment nor orientation is an attribute hidden in the name of the
mapping realization.
\end{construction}

\begin{proposition}[Universal property of the orientation realization]
\label{prop:prequantum-orientation-realization-universal-property}
For $p:P\to S$, a map
$P\to\operatorname{OrientReal}_{\mathrm{univ}}(A,B)$ over $S$ is exactly:
\begin{enumerate}
\item a fundamental-class candidate $\eta:P\to\mathcal F_x$ over $S$;
\item an assignment
      $\theta_P:\mathcal F_x\times_SP\to\mathcal H_{A,B}$ over $S$;
\item coherent inverse data proving that the evaluated arrow
      $\theta_P(\eta)$ is an equivalence in $\mathcal C_P$.
\end{enumerate}
The fibre over $(\eta,\theta_P)$ is empty if the evaluated arrow is not an
equivalence and contractible if it is.  This realization commutes with native
base change and satisfies effective descent.
\end{proposition}

\begin{proof}
The exponential adjunction in $\mathsf{XSp}_{/S}$ identifies a $P$-point of
$\operatorname{DualAssignReal}_x(A,B)$ with a morphism
$\mathcal F_x\times_SP\to\mathcal H_{A,B}$.  Pairing that point with
$\eta:P\to\mathcal F_x$ and composing with evaluation produces
$\theta_P(\eta):P\to\mathcal H_{A,B}$.  The defining pullback lifts this map
through
$\operatorname{Eqv}_{\mathcal C}(A,B)\to\mathcal H_{A,B}$ and hence supplies
precisely an inverse and both triangle homotopies.  The space of such inverse
data is contractible exactly when the arrow is an equivalence and is otherwise
empty.  Pullback between slices preserves internal exponentials and finite
limits; these operations are computed in the native $\infty$-topos and satisfy
effective descent.  This proves the asserted base-change and descent
statements.
\end{proof}

\begin{construction}[Scalarized prequantum action]
\label{constr:prequantum-scalarized-action}
Over $\operatorname{FundReal}_{\partial}(x)$ define
\[
\boxed{
\widehat S_{\widehat L}:
\operatorname{Fld}_S(A)
\longrightarrow
\Omega_S^\infty
\operatorname{map}^{\partial}_S
(\mathbf1_S^{\partial},D)[n]
}
\]
by precomposition of the orientation-free action with
$[X]^\partial_{\mathrm{univ}}$.  All source and target objects are understood after
structural pullback to the realization base.
\end{construction}

\begin{theorem}[Action as top-dimensional differential transgression]
\label{thm:prequantum-action-top-dimensional-transgression}
The orientation-free action is the Chapter~12 differential exceptional Umkehr of
the differential Lagrangian.  Over
$\operatorname{FundReal}_{\partial}(x)$ its scalarization is the corresponding
top-dimensional differential transgression.  Coefficient morphisms and every
finite cubical variation commute with both forms of the action through the natural
equivalences constructed in Chapter~12.  Any curvature comparison is obtained only
after pullback to the corresponding curvature-diagram realization.
\end{theorem}

\begin{proof}
The orientation-free statement is the differential exceptional adjunction of
Chapter~12.  Scalarization is precomposition with the universal point of
$\operatorname{FundReal}_{\partial}(x)$.  Naturality of the exceptional adjunction
in coefficient morphisms and the Chapter~12 finite-cubical theorem give the stated
compatibilities.  No map
$\mathbf1_S^\partial\to x_{e!}\mathbf1_X^\partial$ is used before its realization
object is formed.
\end{proof}

\section{Differential transgression from realized traces and fundamental classes}
\label{sec:prequantum-differential-transgression}

Consider a smooth correspondence in the Chapter~10 theorem domain
\[
\begin{tikzcd}[column sep=huge]
& Z \arrow[dl,"e"'] \arrow[dr,"\pi"] & \\
X && B,
\end{tikzcd}
\]
with $\pi$ smooth, separated, and of finite presentation.

\begin{definition}[Differential coefficient-trace candidate realization]
\label{def:prequantum-coefficient-trace-realization}
For $\widehat E_X\in\mathcal R_X^0$ and
$\widehat E_B\in\mathcal R_B^0$ define
\[
\boxed{
\operatorname{TrReal}_{\partial}
(e,\pi;\widehat E_X,\widehat E_B)
:=
\Omega_B^\infty
\Pi_\pi
\operatorname{map}^{\partial}_Z
\left(e^{*,0}\widehat E_X,\pi_e^!\widehat E_B\right).
}
\]
For a test $T\to B$, a point is exactly a differential coefficient-map candidate
on the base-changed correspondence, by stable dependent-product adjunction and the
global mapping-spectrum theorem.  The zero map is a canonical section, so this
unrestricted mapping realization is never empty.
\end{definition}

\begin{construction}[Motivic trace shadow and differential trace lift]
\label{constr:prequantum-coefficient-trace-shadow-lift}
Define
\[
\operatorname{TrReal}_{\mathrm{mot}}
(e,\pi;H_X\widehat E_X,H_B\widehat E_B)
\in\mathsf{XSp}_{/B}
\]
by declaring its $T\to B$ points to be
\[
\operatorname{Map}_{\operatorname{SH}(Z_T)}
\left(e_T^*H_{X_T}\widehat E_{X_T},
      \pi_T^!H_T\widehat E_T\right).
\]
The motivic exceptional base-change theorem and Nisnevich descent make these
spaces a Cartesian sheaf over $B$.
Adjunction identifies a differential candidate with a map
\[
\pi_{e!}e^{*,0}\widehat E_X\longrightarrow\widehat E_B.
\]
Apply $H_B$, use the constructed exceptional comparison
\[
H_B\pi_{e!}e^{*,0}\widehat E_X
\simeq
\pi_!e^*H_X\widehat E_X,
\]
and adjoint back.  This constructs
\[
\operatorname{sh}_{\mathrm{Tr}}:
\operatorname{TrReal}_{\partial}
(e,\pi;\widehat E_X,\widehat E_B)
\longrightarrow
\operatorname{TrReal}_{\mathrm{mot}}
(e,\pi;H_X\widehat E_X,H_B\widehat E_B).
\]
For a constructed or selected motivic trace
\[
\gamma^{\mathrm{mot}}:
P\longrightarrow\operatorname{TrReal}_{\mathrm{mot}}
(e,\pi;H_X\widehat E_X,H_B\widehat E_B)
\]
define
\[
\boxed{
\operatorname{TrLiftReal}_{\partial}(\gamma^{\mathrm{mot}})
:=
P\times_{\operatorname{TrReal}_{\mathrm{mot}}}
\operatorname{TrReal}_{\partial}.
}
\]
A point is a differential trace candidate with a specified identification of its
motivic shadow with $\gamma^{\mathrm{mot}}$.  This is the trace existence problem
and may be empty.
\end{construction}

\begin{definition}[Correspondence fundamental-class candidate realization]
\label{def:prequantum-correspondence-fundamental-class-realization}
Define
\[
\boxed{
\operatorname{FundReal}_{\partial}(\pi)
:=
\Omega_B^\infty
\operatorname{map}^{\partial}_B
\left(\mathbf1_B^{\partial},\pi_{e!}\mathbf1_Z^{\partial}\right).
}
\]
A point is a differential fundamental-class candidate for the right leg.  Over
the realization there is the tautological section
\[
[Z/B]^\partial_{\mathrm{univ}}:
\mathbf1_B^\partial\longrightarrow\pi_{e!}\mathbf1_Z^\partial.
\]
The zero section canonically points this realization.  Applying
Construction~\ref{constr:prequantum-fundamental-class-shadow-lift} with
$x=\pi$ gives
\[
\operatorname{sh}_{\mathrm{Fund},\pi}:
\operatorname{FundReal}_{\partial}(\pi)
\longrightarrow
\operatorname{FundReal}_{\mathrm{mot}}(\pi)
\]
and the genuine lift fibre
$\operatorname{FundLiftReal}_{\partial}(\pi;\eta^{\mathrm{mot}})$ over any
prescribed motivic fundamental class.
\end{definition}

\begin{construction}[Orientation-free differential transgression]
\label{constr:prequantum-orientation-free-differential-transgression}
Let
\[
\widehat\alpha:
\mathbf1_X^{\partial}\longrightarrow\widehat E_X[r].
\]
Over $\operatorname{TrReal}_{\partial}$, let $\widehat\gamma$ denote the universal
coefficient-map candidate.  Define
\[
\boxed{
\tau^{\partial,\mathrm{or}}_{(e,\pi),\widehat\gamma}
(\widehat\alpha):
\pi_{e!}\mathbf1_Z^{\partial}
\longrightarrow
\widehat E_B[r]
}
\]
by
\[
\pi_{e!}\mathbf1_Z^{\partial}
\xrightarrow{\pi_{e!} e^{*,0}\widehat\alpha}
\pi_{e!}e^{*,0}\widehat E_X[r]
\xrightarrow{\pi_{e!}\widehat\gamma}
\pi_{e!}\pi_e^!\widehat E_B[r]
\xrightarrow{\operatorname{tr}^{\partial}_\pi}
\widehat E_B[r].
\]
\end{construction}

\begin{construction}[Scalar differential prequantum transgression]
\label{constr:prequantum-differential-transgression}
Over
$\operatorname{TrReal}_{\partial}(e,\pi;\widehat E_X,\widehat E_B)
\times_B\operatorname{FundReal}_{\partial}(\pi)$ define
\[
\boxed{
\tau^{\partial}_{(e,\pi)}(\widehat\alpha):
\mathbf1_B^{\partial}
\longrightarrow
\widehat E_B[r]
}
\]
as
\[
\mathbf1_B^{\partial}
\xrightarrow{[Z/B]^\partial_{\mathrm{univ}}}
\pi_{e!}\mathbf1_Z^{\partial}
\xrightarrow{\tau^{\partial,\mathrm{or}}_{(e,\pi),\widehat\gamma}
(\widehat\alpha)}
\widehat E_B[r].
\]
\end{construction}

\subsection{Coherent de Rham-stage traces and curvature compatibility}
\label{subsec:prequantum-curvature-compatible-trace-realization}

A single coefficient trace is sufficient for differential transgression.  It is
not sufficient to conclude compatibility with normalized or coherently closed de
Rham cochains.  Chapter~10 proves that conclusion from a morphism of the actual
realized stage/cochain diagram.  We therefore use that constructed diagram
itself rather than naming an unspecified ``coherent trace package.''

\begin{construction}[The Chapter~10 realized stage diagram]
\label{constr:prequantum-realized-dr-stage-diagram}
Let $\mathcal I_F:=\mathbb N^{\mathrm{op}}$ with vertices $F_q$ and generating
arrows $F_{q+1}\to F_q$, and let
$\mathcal I_\Omega:=\Xi_{\ge0}^{\mathrm{op}}$ be the pointed coherent-cochain
index of Chapter~10.  Let $\mathcal W_{\mathrm{cof}}$ be the free small stable
$\infty$-category on a cofiber sequence $a\to b\to c\to a[1]$, and let
$\mathcal T_q$ be its $q$th copy.  Thus $\mathcal T_q$ contains the walking
pointed cofiber square with vertices $a_q,b_q,c_q,0$, arrows
\[
a_q\longrightarrow b_q,\qquad a_q\longrightarrow0,\qquad
b_q\longrightarrow c_q,\qquad0\longrightarrow c_q,
\]
its commutative-square cell, and the induced connecting arrow
$\delta_q:c_q\to a_q[1]$.

Form $\mathcal I_{\mathrm{DR}}$ as the following homotopy colimit of free small
stable $\infty$-categories.  Glue the copies $\mathcal T_q$ to the free stable
envelope of $\mathcal I_F$ and to that of $\mathcal I_\Omega$ by
\[
a_q=F_{q+1},\qquad b_q=F_q,\qquad c_q=\Omega_q[-q],
\]
and impose, by a homotopy pushout along the walking parallel pair, the relation
that the shifted coherent differential is the adjacent-boundary composite
\[
d_q[-q]
=
\left(
\Omega_q[-q]
\xrightarrow{\delta_q}F_{q+1}[1]
\xrightarrow{t_{q+1}[1]}\Omega_{q+1}[-q]
\right),
\]
where $t_{q+1}:F_{q+1}\to\Omega_{q+1}[-q-1]$ is the next cofiber quotient.
Attach the nullhomotopy of $d_{q+1}d_q$ and every higher cochain cell along the
corresponding cells of $\Xi_{\ge0}^{\mathrm{op}}$.  This completely specifies a
small stable indexing category: it is a countable homotopy colimit of small free
stable categories and walking-cell attachments.  Exact functors out of it are
precisely an adjacent-stage cofiber system together with the compatible coherent
normalized-cochain object.  Theorem~10.8.1 supplies the
complete-filtration/coherent-cochain cell system,
Propositions~10.15.3--10.15.4 construct its Thom-stable exact realization, and
Theorem~10.17.2 constructs morphisms of these exact diagrams.

For every represented differential base $Z$ in the Chapter~10 theorem range,
write
\[
\boxed{
\mathbf{DRStage}^{\partial}_Z:
\mathcal I_{\mathrm{DR}}\longrightarrow\mathcal R_Z^0
}
\]
for the exact functor whose $F_q$-value is the realized stage
$\widehat F^q_{\mathrm{dR},Z}$ and whose $\Omega_q$-value is the realized normalized
cochain $\widehat\Omega^q_{\inf,Z}$.  This is the diagram already constructed in
Chapter~10; no vertex, arrow, or coherence is selected in the present chapter.
\end{construction}

For exact diagrams $\mathbf E,\mathbf F:\mathcal I_{\mathrm{DR}}\to\mathcal R_Z^0$
define their big diagram-mapping coefficient by the twisted-arrow end
\[
\boxed{
\operatorname{map}^{\partial}_{Z,\mathrm{DR}}
(\mathbf E,\mathbf F)
:=
\lim_{(i\to j)\in\operatorname{Tw}(\mathcal I_{\mathrm{DR}})}
\operatorname{map}^{\partial}_Z(E_i,F_j).
}
\]
Its stable global sections are the mapping spectrum in
$\operatorname{Fun}(\mathcal I_{\mathrm{DR}},\mathcal R_Z^0)$.

\begin{definition}[Coherent de Rham-stage trace realization]
\label{def:prequantum-coherent-stage-trace-realization}
Define
\[
\boxed{
\operatorname{TrCohReal}_{\partial}(e,\pi)
:=
\Omega_B^\infty\Pi_\pi
\operatorname{map}^{\partial}_{Z,\mathrm{DR}}
\left(
 e^{*,0}\mathbf{DRStage}^{\partial}_X,
 \pi_e^!\mathbf{DRStage}^{\partial}_B
\right).
}
\]
A point is exactly a morphism of the already constructed Chapter~10 realized
stage/cochain diagrams.  Hence compatibility with every stage transition,
adjacent exact diagram, coherent differential, exact-cochain classifier, and
higher cochain cell is a consequence of being such a morphism rather than an
independent field of trace data.
The zero diagram morphism canonically points this realization.
\end{definition}

\begin{definition}[Coherent extension of a selected stage trace]
\label{def:prequantum-coherent-stage-trace-lift-realization}
Evaluation at the vertex $F_q$ constructs
\[
\operatorname{ev}_q:
\operatorname{TrCohReal}_{\partial}(e,\pi)
\longrightarrow
\operatorname{TrReal}_{\partial}
\left(e,\pi;\widehat F^q_{\mathrm{dR},X},
                 \widehat F^q_{\mathrm{dR},B}\right).
\]
Write this target as $\operatorname{TrReal}_{\partial,q}(e,\pi)$.
For a selected stage trace
\[
\gamma_q:P\longrightarrow
\operatorname{TrReal}_{\partial}
\left(e,\pi;\widehat F^q_{\mathrm{dR},X},
                 \widehat F^q_{\mathrm{dR},B}\right)
\]
define
\[
\boxed{
\operatorname{TrCohLiftReal}_{\partial}(e,\pi;\gamma_q)
:=
P\times_{\operatorname{TrReal}_{\partial,q}(e,\pi)}
\operatorname{TrCohReal}_{\partial}(e,\pi).
}
\]
A point is a morphism of the complete realized stage/cochain diagrams whose
$q$th component is identified with the prescribed trace.  This fibre may be
empty; it is the coherent-stage trace discrepancy, rather than the pointed
unrestricted diagram-mapping realization.
\end{definition}

\begin{construction}[Parameterized commutative-square realization]
\label{constr:prequantum-square-realization}
Let $a:A\to A'$ and $b:B\to B'$ be morphisms in $E_T$.  Define the stable
square-mapping coefficient
\[
\operatorname{SqCoeff}_T(a,b)
:=
\underline{\operatorname{Hom}}(A,B)
\times_{\underline{\operatorname{Hom}}(A,B')}
\underline{\operatorname{Hom}}(A',B')
\in E_T
\]
and its parameterized realization
\[
\boxed{
\operatorname{SqReal}_T(a,b)
:=
\Omega_T^\infty\operatorname{SqCoeff}_T(a,b).
}
\]
The first map to the middle internal Hom is postcomposition by $b$ and the
second is precomposition by $a$.  A point is a pair
\[
f:A\to B,
\qquad
f':A'\to B',
\]
together with a specified coherent homotopy
\[
bf\simeq f'a.
\]
The construction is a finite limit of internal mapping objects, so it commutes
with arbitrary pullback and effective descent.
\end{construction}

\subsection{Curvature compatibility at the transgression-operator level}
\label{subsec:prequantum-transgression-operator-curvature}

The differential exceptional adjunction constructs the differential trace only
inside the Thom-stable differential categories.  We therefore compare the
\emph{actual transgression operator} with an independently constructed stable
underlying/closed transgression; no unconstructed exceptional pullback is
introduced in the parameterized stable category.

For the selected differential coefficients put
\[
\Gamma_X^{\partial}(\widehat E_X)
:=
\operatorname{map}_{\mathcal R_X^0}
(\mathbf1_X^{\partial},\widehat E_X),
\qquad
\Gamma_B^{\partial}(\widehat E_B)
:=
\operatorname{map}_{\mathcal R_B^0}
(\mathbf1_B^{\partial},\widehat E_B).
\]
For brevity abbreviate these spectra by $\Gamma_X^{\partial}$ and
$\Gamma_B^{\partial}$ throughout this subsection.  Write
\[
\operatorname{TrSp}_{\partial}(e,\pi)
:=
\operatorname{Map}_B
\bigl(B,\operatorname{TrReal}_{\partial}(e,\pi)\bigr),
\qquad
\operatorname{FundSp}_{\partial}(\pi)
:=
\operatorname{Map}_B
\bigl(B,\operatorname{FundReal}_{\partial}(\pi)\bigr)
\]
for the spaces of global trace and fundamental-class candidates.  For prescribed
motivic shadows $\gamma^{\mathrm{mot}}$ and $\eta^{\mathrm{mot}}$ define instead
the actual parameter object
\[
\boxed{
P_{\partial}^{\mathrm{tr,mot}}
:=
\operatorname{Map}_B
\bigl(B,\operatorname{TrLiftReal}_{\partial}(\gamma^{\mathrm{mot}})\bigr)
\times
\operatorname{Map}_B
\bigl(B,\operatorname{FundLiftReal}_{\partial}
(\pi;\eta^{\mathrm{mot}})\bigr).
}
\]
The motivically prescribed branch is the base change of the constructions below
along $P_{\partial}^{\mathrm{tr,mot}}\to P_{\partial}^{\mathrm{tr}}$.

\begin{construction}[Differential transgression operator]
\label{constr:prequantum-differential-transgression-operator}
Over
\[
P_{\partial}^{\mathrm{tr}}
:=
\operatorname{TrSp}_{\partial}(e,\pi)
\times
\operatorname{FundSp}_{\partial}(\pi)
\]
let $\widehat\gamma$ and $\eta$ denote the universal trace and
fundamental-class candidates.
They construct a canonical map of spectra
\[
\boxed{
T^{\partial}_{\widehat\gamma,\eta}:
\Gamma_X^{\partial}(\widehat E_X)
\longrightarrow
\Gamma_B^{\partial}(\widehat E_B)
}
\]
by the composite
\[
\begin{aligned}
\operatorname{map}_{\mathcal R_X^0}
(\mathbf1_X^{\partial},\widehat E_X)
&\longrightarrow
\operatorname{map}_{\mathcal R_Z^0}
(\mathbf1_Z^{\partial},e^{*,0}\widehat E_X)
\\
&\xrightarrow{\widehat\gamma_*}
\operatorname{map}_{\mathcal R_Z^0}
(\mathbf1_Z^{\partial},\pi_e^!\widehat E_B)
\\
&\simeq
\operatorname{map}_{\mathcal R_B^0}
(\pi_{e!}\mathbf1_Z^{\partial},\widehat E_B)
\\
&\xrightarrow{\eta^*}
\operatorname{map}_{\mathcal R_B^0}
(\mathbf1_B^{\partial},\widehat E_B).
\end{aligned}
\]
The construction is functorial in paths and higher homotopies of
$\widehat\gamma$ and $\eta$.
\end{construction}

Let $\Gamma_X^U$ and $\Gamma_B^U$ be the stable underlying/closed section spectra
used by the comparison problem.  Define the section-curvature realization
\[
\boxed{
\operatorname{SecCurvReal}(\Gamma_X^{\partial},\Gamma_X^U;
                            \Gamma_B^{\partial},\Gamma_B^U)
:=
\Omega^\infty\operatorname{Hom}_{\operatorname{Sp}}
(\Gamma_X^{\partial},\Gamma_X^U)
\times
\Omega^\infty\operatorname{Hom}_{\operatorname{Sp}}
(\Gamma_B^{\partial},\Gamma_B^U).
}
\]
Over it write
\[
\kappa_X^{\Gamma}:\Gamma_X^{\partial}\to\Gamma_X^U,
\qquad
\kappa_B^{\Gamma}:\Gamma_B^{\partial}\to\Gamma_B^U
\]
for the universal interfaces.  If either interface has already been constructed,
pull back along its section and suppress the corresponding factor.

\begin{definition}[Underlying transgression realization]
\label{def:prequantum-underlying-transgression-realization}
Define
\[
\boxed{
\operatorname{UTrReal}(e,\pi;\Gamma_X^U,\Gamma_B^U)
:=
\Omega^\infty
\operatorname{Hom}_{\operatorname{Sp}}
(\Gamma_X^U,\Gamma_B^U).
}
\]
Over it write
\[
T^U_{\mathrm{univ}}:\Gamma_X^U\longrightarrow\Gamma_B^U
\]
for the universal stable transgression.  If a geometric, closed, or represented
underlying transgression has already been constructed, all later objects are
pulled back along its section and this factor is suppressed.
The zero map canonically points the unrestricted realization; only a lift or
compatibility fibre over a fixed underlying operator can be empty.
\end{definition}

\begin{definition}[Curvature-compatible transgression realization]
\label{def:prequantum-curvature-compatible-transgression-realization}
Put
\[
S_{\mathrm{curv}}
:=
\operatorname{SecCurvReal}
(\Gamma_X^{\partial},\Gamma_X^U;
 \Gamma_B^{\partial},\Gamma_B^U).
\]
In the slice over $S_{\mathrm{curv}}$ define
\[
M_{\partial}
:=
S_{\mathrm{curv}}
\times\operatorname{Map}
(\Gamma_X^{\partial},\Gamma_B^{\partial}),
\qquad
M_U
:=
S_{\mathrm{curv}}
\times\operatorname{Map}(\Gamma_X^U,\Gamma_B^U).
\]
The universal interfaces are maps of spectrum objects over
$S_{\mathrm{curv}}$, so
Construction~\ref{constr:prequantum-square-realization} gives an object
\[
\operatorname{SqReal}_{S_{\mathrm{curv}}}
(\kappa_X^\Gamma,\kappa_B^\Gamma)
\longrightarrow M_{\partial}\times_{S_{\mathrm{curv}}}M_U.
\]
Let
\[
P_{\mathrm{tr}}
:=
S_{\mathrm{curv}}
\times P_{\partial}^{\mathrm{tr}}
\times
\operatorname{UTrReal}(e,\pi;\Gamma_X^U,\Gamma_B^U).
\]
The universal pair $(T^{\partial}_{\mathrm{univ}},T^U_{\mathrm{univ}})$
defines a morphism over $S_{\mathrm{curv}}$
\[
P_{\mathrm{tr}}
\longrightarrow M_{\partial}\times_{S_{\mathrm{curv}}}M_U.
\]
Define the following pullback in the single slice
$\mathsf{XSp}_{/S_{\mathrm{curv}}}$:
\[
\boxed{
\operatorname{TransCurvReal}_{\partial}(e,\pi)
:=
P_{\mathrm{tr}}
\times_{M_{\partial}\times_{S_{\mathrm{curv}}}M_U}
\operatorname{SqReal}_{S_{\mathrm{curv}}}
(\kappa_X^{\Gamma},\kappa_B^{\Gamma}),
}
\]
where the square realization is specialized fibrewise to spectra.  Thus a point
over fixed interfaces is
\[
(\widehat\gamma,\eta,T^U,H_{\mathrm{curv}})
\]
with
\[
\boxed{
H_{\mathrm{curv}}:
\kappa_B^{\Gamma}T^{\partial}_{\widehat\gamma,\eta}
\simeq
T^U\kappa_X^{\Gamma}.
}
\]
The unrestricted object is canonically inhabited by taking zero trace,
fundamental-class candidate, interfaces, and underlying operator.  Its pullback
over independently fixed data need not be inhabited.
\end{definition}

Define the coherent trace-section space and its evaluation map at the coefficient
vertex used in the transgression by
\[
\operatorname{TrCohSp}_{\partial}(e,\pi)
:=
\operatorname{Map}_B
\bigl(B,\operatorname{TrCohReal}_{\partial}(e,\pi)\bigr),
\qquad
\operatorname{ev}_q:
\operatorname{TrCohSp}_{\partial}(e,\pi)
\longrightarrow\operatorname{TrSp}_{\partial}(e,\pi).
\]

\begin{definition}[Coherent-stage curvature-compatible transgression]
\label{def:prequantum-coherent-stage-curvature-transgression}
Define the explicit base change
\[
\boxed{
\operatorname{TransCurvCohReal}_{\partial}(e,\pi)
:=
\operatorname{TransCurvReal}_{\partial}(e,\pi)
\times_{\operatorname{TrSp}_{\partial}(e,\pi)}
\operatorname{TrCohSp}_{\partial}(e,\pi).
}
\]
A point contains an actual morphism of the complete realized Chapter~10
stage/cochain diagrams whose selected coefficient component is the trace used by
the curvature-compatible transgression.
\end{definition}

\begin{definition}[Fixed transgression-square lift realization]
\label{def:prequantum-fixed-transgression-square-lift-realization}
After fixing maps
\[
T^\partial:\Gamma_X^\partial\to\Gamma_B^\partial,\qquad
T^U:\Gamma_X^U\to\Gamma_B^U,\qquad
\kappa_X^\Gamma:\Gamma_X^\partial\to\Gamma_X^U,\qquad
\kappa_B^\Gamma:\Gamma_B^\partial\to\Gamma_B^U,
\]
define
\[
\boxed{
\operatorname{TransSqLiftReal}
(T^\partial,T^U;\kappa_X^\Gamma,\kappa_B^\Gamma)
:=
\operatorname{Path}
\left(\kappa_B^\Gamma T^\partial,
      T^U\kappa_X^\Gamma\right).
}
\]
Equivalently, this is the fibre of
$\operatorname{SqReal}_{\operatorname{Sp}}
(\kappa_X^\Gamma,\kappa_B^\Gamma)$ over the fixed pair
$(T^\partial,T^U)$.  It may be empty, and it is the exact curvature-compatibility
discrepancy for the selected operators.
\end{definition}

Put $I:=\mathcal I_{\mathrm{DR}}$.  For $Z=X,B$, let
\[
\Theta_Z:\operatorname{Fun}(I,\operatorname{Sp})\longrightarrow
\operatorname{Sp}
\]
be the actual evaluation, terminal-output, or complete-totalization functor used
to form $\Gamma_Z^U$ on the branch under consideration.  Define the complete
realization of an underlying stage system with the required output by
\[
\boxed{
\operatorname{StageSys}_Z^U(\Gamma_Z^U)
:=
\int_{\mathbf G\in\operatorname{Fun}(I,\operatorname{Sp})^\simeq}
\operatorname{Eqv}_{\operatorname{Sp}}
(\Theta_Z\mathbf G,\Gamma_Z^U).
}
\]
This space is regarded as a constant object of the native slice.  Over it there
are a universal stage diagram $\mathbf\Gamma_Z^U$ and a universal output
equivalence $\epsilon_Z:\Theta_Z\mathbf\Gamma_Z^U\simeq\Gamma_Z^U$.

Let
\[
B_{\mathrm{stage}}
:=
\operatorname{TransCurvCohReal}_{\partial}(e,\pi)
\times\operatorname{StageSys}_X^U(\Gamma_X^U)
\times\operatorname{StageSys}_B^U(\Gamma_B^U).
\]
The Chapter~10 realization supplies over this base the differential stage
diagrams $\mathbf\Gamma_X^\partial,\mathbf\Gamma_B^\partial$, their output
identifications, the differential stage map $\mathbf T^\partial$, and the
terminal maps $T^\partial,T^U,\kappa_X^\Gamma,\kappa_B^\Gamma$ together with
$H_{\mathrm{curv}}$.

Output under $\Theta_X,\Theta_B$ and conjugation by $\epsilon_X,\epsilon_B$
define restriction maps from stagewise mapping spaces to terminal mapping
spaces.  Form their relative fibres over the universal terminal maps:
\[
\boxed{
\mathcal U
:=
B_{\mathrm{stage}}
\times_{\operatorname{Map}_{\operatorname{Sp}}
(\Gamma_X^U,\Gamma_B^U)}
\operatorname{Map}_{\operatorname{Fun}(I,\operatorname{Sp})}
(\mathbf\Gamma_X^U,\mathbf\Gamma_B^U),
}
\]
\[
\boxed{
\mathcal K_X
:=
B_{\mathrm{stage}}
\times_{\operatorname{Map}_{\operatorname{Sp}}
(\Gamma_X^\partial,\Gamma_X^U)}
\operatorname{Map}_{\operatorname{Fun}(I,\operatorname{Sp})}
(\mathbf\Gamma_X^\partial,\mathbf\Gamma_X^U),
}
\]
and define $\mathcal K_B$ by replacing $X$ with $B$.  Their universal points are
denoted $\mathbf T^U,\mathbf\kappa_X,\mathbf\kappa_B$.

\begin{definition}[Complete stagewise transgression-square realization]
\label{def:prequantum-complete-stagewise-transgression-square}
Put
\[
P_{\mathrm{stage}}
:=
\mathcal U\times_{B_{\mathrm{stage}}}
\mathcal K_X\times_{B_{\mathrm{stage}}}\mathcal K_B.
\]
Restriction under $\Theta_X,\Theta_B$ gives a morphism of path objects
\[
\operatorname{Path}_{\operatorname{Fun}(I,\operatorname{Sp})}
(\mathbf\kappa_B\mathbf T^\partial,
 \mathbf T^U\mathbf\kappa_X)
\longrightarrow
\operatorname{Path}_{\operatorname{Sp}}
(\kappa_B^\Gamma T^\partial,T^U\kappa_X^\Gamma).
\]
Define the relative pullback
\[
\boxed{
\operatorname{TransStageCurvReal}_{\partial}(e,\pi)
:=
P_{\mathrm{stage}}
\times_{
\operatorname{Path}_{\operatorname{Sp}}
(\kappa_B^\Gamma T^\partial,T^U\kappa_X^\Gamma)}
\operatorname{Path}_{\operatorname{Fun}(I,\operatorname{Sp})}
(\mathbf\kappa_B\mathbf T^\partial,
 \mathbf T^U\mathbf\kappa_X),
}
\]
where the first map is the universal path $H_{\mathrm{curv}}$.
A point is precisely an underlying stage system with its output comparison, a
stagewise lift of the underlying transgression, stagewise lifts of both
interfaces, and a natural curvature-square homotopy restricting to
$H_{\mathrm{curv}}$.
\end{definition}

\begin{proposition}[Universal property and locality of the stage realization]
\label{prop:prequantum-stage-realization-universal-property}
The preceding object represents complete stagewise curvature-compatible
transgression.  It commutes with native base change and satisfies effective
descent.
\end{proposition}

\begin{proof}
For a test object $T$, mapping $T$ into $\operatorname{StageSys}_Z^U$ selects a
diagram and an output equivalence.  Mapping into $\mathcal U,\mathcal K_X$, and
$\mathcal K_B$ selects lifts of the three universal terminal maps.  Mapping into
the final pullback selects a natural path whose output is the prescribed
$H_{\mathrm{curv}}$.  Fibre-pasting gives the asserted representing property.
All constituents are functor-category mapping objects, equivalence objects,
path objects, Grothendieck constructions, and limits.  Their affine-testwise
forms commute with pullback, and effective descent reconstructs their universal
points and all comparison paths.  When $\Theta_Z$ is a complete totalization,
its output equivalence is part of the represented stage-system point rather than
an unrecorded preservation assertion.
\end{proof}

\begin{theorem}[Curvature compatibility of differential transgression]
\label{thm:prequantum-transgression-curvature-compatibility}
Over $\operatorname{TransCurvReal}_{\partial}(e,\pi)$, every differential class
$\widehat\alpha\in\Omega^\infty\Gamma_X^{\partial}[r]$ satisfies the canonical
coherent identity
\[
\boxed{
\kappa_B^{\Gamma}
\bigl(T^{\partial}_{\widehat\gamma,\eta}(\widehat\alpha)\bigr)
\simeq
T^U\bigl(\kappa_X^{\Gamma}(\widehat\alpha)\bigr).
}
\]
After pullback to
$\operatorname{TransCurvCohReal}_{\partial}(e,\pi)$, the
differential operator $T^{\partial}$ is assembled stagewise from a morphism of
the actual Chapter~10 realized diagrams and therefore preserves normalized
cochains, coherently closed cochains, the coherent differential, exact-cochain
classifiers, and every higher cell already present in that differential diagram.
Over $\operatorname{TransStageCurvReal}_{\partial}(e,\pi)$ the curvature
identity is a natural homotopy of the complete stage diagrams.
\end{theorem}

\begin{proof}
The first statement is evaluation of the universal square in
$\operatorname{TransCurvReal}_{\partial}$ on $\widehat\alpha$.  A point of
$\operatorname{TransCurvCohReal}_{\partial}(e,\pi)$ contains an actual natural transformation
between the Chapter~10 realized differential diagrams.  Exact functoriality
carries the adjacent stage triangles and their connecting morphisms to the
corresponding target triangles, so every normalized and coherently closed
structure is transported with its specified higher coherences.  The defining
homotopy fibre of $\operatorname{TransStageCurvReal}_{\partial}$ retains an
actual natural transformation of the universally realized underlying stage diagrams, both
stagewise interfaces, and the path whose terminal evaluation is
$H_{\mathrm{curv}}$.  This proves the final assertion.  No stable
exceptional $\pi^!$ on the underlying parameterized coefficient theory is used.
\end{proof}

\section{Composition of traces, fundamental classes, and Fubini}
\label{sec:prequantum-iterated-transgression-flags}

The composite trace and fundamental class are constructed from the successive
ones; they are not separately chosen compatibility fields.

Consider composable correspondences
\[
X\xleftarrow{e}Z\xrightarrow{\pi}B,
\qquad
B\xleftarrow{e'}W\xrightarrow{\rho}C,
\]
and let $P:=Z\times_BW$.  Write
$\widetilde e:P\to X$ and $\widetilde\rho:P\to C$ for the composite legs.

\begin{construction}[Composite differential trace]
\label{constr:prequantum-composite-differential-trace}
For trace points
\[
\widehat\gamma:e^{*,0}\widehat E_X\to\pi_e^!\widehat E_B,
\qquad
\widehat\delta:e'^{*,0}\widehat E_B\to\rho_e^!\widehat E_C,
\]
define
\[
\boxed{
\widehat\delta\star\widehat\gamma:
\widetilde e^{*,0}\widehat E_X
\longrightarrow
\widetilde\rho_e^!\widehat E_C
}
\]
by the exceptional Beck--Chevalley composite
\[
\operatorname{pr}_Z^{*,0}e^{*,0}\widehat E_X
\xrightarrow{\operatorname{pr}_Z^{*,0}\widehat\gamma}
\operatorname{pr}_Z^{*,0}\pi_e^!\widehat E_B
\simeq
(\operatorname{pr}_W)_e^!e'^{*,0}\widehat E_B
\xrightarrow{(\operatorname{pr}_W)_e^!\widehat\delta}
(\operatorname{pr}_W)_e^!\rho_e^!\widehat E_C
\simeq
\widetilde\rho_e^!\widehat E_C.
\]
The same formula applied in
$\operatorname{Fun}(\mathcal I_{\mathrm{DR}},\mathcal R^0)$ composes coherent
de Rham-stage traces.
\end{construction}

\begin{construction}[Composite differential fundamental class]
\label{constr:prequantum-composite-fundamental-class}
Given
\[
[Z/B]^\partial:
\mathbf1_B^\partial\to\pi_{e!}\mathbf1_Z^\partial,
\qquad
[W/C]^\partial:
\mathbf1_C^\partial\to\rho_{e!}\mathbf1_W^\partial,
\]
pull $[Z/B]^\partial$ to $W$ and use smooth Beck--Chevalley to obtain
\[
\mathbf1_W^\partial
\longrightarrow
(\pi_W)_{e!}\mathbf1_P^\partial.
\]
Define
\[
\boxed{
[P/C]^\partial:
\mathbf1_C^\partial
\longrightarrow
\widetilde\rho_{e!}\mathbf1_P^\partial
}
\]
by
\[
\mathbf1_C^\partial
\to
\rho_{e!}\mathbf1_W^\partial
\to
\rho_{e!}(\pi_W)_{e!}\mathbf1_P^\partial
\simeq
\widetilde\rho_{e!}\mathbf1_P^\partial.
\]
\end{construction}

\begin{construction}[Composition of underlying transgressions]
\label{constr:prequantum-underlying-transgression-composition}
For the two correspondences, composition of maps of spectra constructs
\[
\boxed{
\operatorname{comp}^U:
\operatorname{UTrReal}(e,\pi;\Gamma_X^U,\Gamma_B^U)
\times
\operatorname{UTrReal}(e',\rho;\Gamma_B^U,\Gamma_C^U)
\longrightarrow
\operatorname{UTrReal}
(\widetilde e,\widetilde\rho;\Gamma_X^U,\Gamma_C^U),
}
\]
\[
(T_1^U,T_2^U)\longmapsto T_2^U\circ T_1^U.
\]
This map is associative and unital with the canonical coherences of
$\operatorname{Sp}$.  If an independently constructed operator $T_{21}^U$ for
the composite is selected, define its exact comparison fibre
\[
\boxed{
\operatorname{UTrCompLiftReal}
(T_{21}^U;T_2^U,T_1^U)
:=
\operatorname{Path}(T_{21}^U,T_2^UT_1^U).
}
\]
The same construction in
$\operatorname{Fun}(\mathcal I_{\mathrm{DR}},\operatorname{Sp})$ composes
stagewise underlying transgressions and maps to the displayed terminal
composition under evaluation.
\end{construction}

\begin{theorem}[Coherent composition]
\label{thm:prequantum-coherent-trace-fundamental-composition}
The preceding compositions are associative and unital through the canonical
exceptional pasting coherences.  They are natural in paths and higher homotopies
of trace and fundamental-class points.  Coherent de Rham-stage traces are
preserved under composition.  The underlying stable transgressions compose by
Construction~\ref{constr:prequantum-underlying-transgression-composition}; the
two operator-level curvature squares paste and hence preserve
$\operatorname{TransCurvReal}_{\partial}$.  On the complete stagewise
realization their natural homotopies paste in
$\operatorname{Fun}(\mathcal I_{\mathrm{DR}},\operatorname{Sp})$.
\end{theorem}

\begin{proof}
Chapter~8 constructs the exceptional correspondence pseudofunctor and proves that
its finite-string Beck--Chevalley and exceptional-composition equivalences are
compatible with every face, degeneracy, and subdivision.  Apply those coherences
to the trace composites.  The fundamental-class composite uses the same
Beck--Chevalley and exceptional-composition transformations.  Functor-category
composition gives the coherent-stage statement.  Construction~\ref{constr:prequantum-underlying-transgression-composition}
constructs the underlying composite rather than assuming it.  Pasting the two
commutative squares in spectra gives the
$\operatorname{TransCurvReal}_{\partial}$ statement; the same calculation in
the stage-diagram functor category gives the final assertion.  If an independently
constructed composite operator is fixed, append the selected path in
$\operatorname{UTrCompLiftReal}$.
\end{proof}

\begin{theorem}[Fubini theorem for differential prequantum transgression]
\label{thm:prequantum-fubini-transgression}
For any finite flag of smooth correspondences equipped successively with points of
the trace and fundamental-class realization objects, the recursively constructed
composite trace and fundamental class satisfy
\[
\boxed{
\tau^\partial_{p_k\odot\cdots\odot p_1}
\simeq
\tau^\partial_{p_k}\circ\cdots\circ\tau^\partial_{p_1}.
}
\]
Here $p_k\odot\cdots\odot p_1$ denotes the correspondence obtained by first
applying $p_1$, then $p_2$, and so on; $\odot$ is used to distinguish
correspondence composition from composition of the resulting operators.
The equivalence is associative and unital.  Over
$\operatorname{TransStageCurvReal}_{\partial}$ it is compatible with the full
realized closed/underlying transgression diagrams and is coherent under
refinement of the flag.
\end{theorem}

\begin{proof}
Write $F_i$ for the differential exceptional functor associated to $p_i$,
\[
\bar\gamma_i:F_i(\widehat E_{i-1})\longrightarrow\widehat E_i
\]
for the adjoint trace, and
\[
\eta_i:\mathbf1_i\longrightarrow F_i(\mathbf1_{i-1})
\]
for its differential fundamental class.  The corresponding operator is
\[
T_i(\alpha)
=
\bar\gamma_i\circ F_i(\alpha)\circ\eta_i.
\]
For two successive correspondences, exceptional Beck--Chevalley and trace
pasting give
\[
F_{21}\simeq F_2F_1,
\qquad
\bar\gamma_{21}
=\bar\gamma_2\circ F_2(\bar\gamma_1),
\qquad
\eta_{21}=F_2(\eta_1)\circ\eta_2.
\]
Consequently
\[
\begin{aligned}
T_{21}(\alpha)
&=
\bar\gamma_2\circ F_2(\bar\gamma_1)
\circ F_2F_1(\alpha)\circ F_2(\eta_1)\circ\eta_2\\
&=
\bar\gamma_2\circ
F_2\!\left(\bar\gamma_1\circ F_1(\alpha)\circ\eta_1\right)
\circ\eta_2\\
&=T_2\bigl(T_1(\alpha)\bigr).
\end{aligned}
\]
Induction gives the displayed order for a finite flag.  The identity,
associativity, face, degeneracy, and subdivision coherences are those of the
finite-string Chapter~8 exceptional correspondence pasting, so the induction is
coherent rather than merely an equality in the homotopy category.  Applying the
same calculation in the Chapter~10 realized diagram functor category proves the
coherent-stage statement.  Pasting the operator-level curvature squares and
using Theorem~\ref{thm:prequantum-transgression-curvature-compatibility} proves
the final assertion.
\end{proof}

\section{The traced-correspondence prequantum category}
\label{sec:prequantum-traced-correspondence-category}

A decorated $(\infty,2)$-category requires differential lifts not only of
correspondence $1$-morphisms but also of the geometric $2$-cells of the Chapter~8
correspondence category.  We construct those lifts rather than silently reusing
the motivic $2$-cells.

Let $c_i:X\xleftarrow{e_i}Z_i\xrightarrow{\pi_i}B$ $(i=0,1)$ be two
correspondences with smooth right legs and let
\[
\alpha:c_0\Longrightarrow c_1
\]
be a Chapter~8 $2$-cell.  Put
\[
F_i^\partial:=\pi_{i,e!}e_i^{*,0}:
\mathcal R_X^0\longrightarrow\mathcal R_B^0.
\]
Let $F_i^{\mathrm{mot}}$ be their motivic shadows and let
$\alpha^{\mathrm{mot}}:F_0^{\mathrm{mot}}\Rightarrow F_1^{\mathrm{mot}}$ be the
Chapter~8 transformation.

\begin{definition}[Differential lift realization of a correspondence $2$-cell]
\label{def:prequantum-two-cell-lift-realization}
Using the constructed equivalences
$H_BF_i^\partial\simeq F_i^{\mathrm{mot}}H_X$, define
\[
\boxed{
\operatorname{TwoCellReal}_{\partial}(\alpha)
:=
\operatorname{hofib}_{\alpha^{\mathrm{mot}}H_X}
\left(
\operatorname{Map}_{\operatorname{Fun}(\mathcal R_X^0,\mathcal R_B^0)}
(F_0^\partial,F_1^\partial)
\longrightarrow
\operatorname{Map}_{\operatorname{Fun}(\mathcal R_X^0,\operatorname{SH}(B))}
(H_BF_0^\partial,H_BF_1^\partial)
\right).
}
\]
A point is an actual natural transformation
$\widetilde\alpha^\partial:F_0^\partial\Rightarrow F_1^\partial$ together with a
specified identification of its motivic shadow with the prescribed Chapter~8
$2$-cell.  The space may be empty.
\end{definition}

For a correspondence $c$ write
\[
\operatorname{Dec}(c)
:=
\operatorname{TrReal}_{\partial}(c)
\times_B
\operatorname{FundReal}_{\partial}(c)
\]
for its internal decoration object and
\[
\operatorname{DecSp}(c):=\operatorname{Map}_B(B,\operatorname{Dec}(c))
\]
for its space of actual decorations.  Let a decoration
$d_i=(\widehat\gamma_i,\eta_i)$ on $c_i$ be fixed.  Adjunction turns the trace into
\[
\tau_i:
F_i^{\partial}(\widehat E_X)
\longrightarrow
\widehat E_B,
\]
while the fundamental class is
\[
\eta_i:
\mathbf1_B^{\partial}
\longrightarrow
F_i^{\partial}(\mathbf1_X^{\partial}).
\]
Here ``decoration'' means an actually selected pair of candidate maps.  For
prescribed motivic decorations define the actual lift object
\[
\operatorname{DecLift}(c)
:=
\operatorname{TrLiftReal}_{\partial}(\gamma_c^{\mathrm{mot}})
\times_B
\operatorname{FundLiftReal}_{\partial}(\pi_c;\eta_c^{\mathrm{mot}}).
\]
These two pieces have opposite variance under a natural transformation
$F_0^{\partial}\Rightarrow F_1^{\partial}$ and therefore are not combined into a
single covariant decoration-transport map.

\begin{definition}[Decorated differential $2$-cell realization]
\label{def:prequantum-decorated-two-cell-realization}
Over $\operatorname{TwoCellReal}_{\partial}(\alpha)$ let
$\widetilde\alpha^{\partial}_{\mathrm{univ}}$ be the universal differential lift.
Define
\[
\boxed{
\begin{aligned}
\operatorname{TwoCellDecReal}_{\partial}(\alpha;d_0,d_1)
:={}&
\operatorname{TwoCellReal}_{\partial}(\alpha)
\\
&\times
\operatorname{Path}
\left(
\tau_0,
\tau_1\circ
\widetilde\alpha^{\partial}_{\mathrm{univ},\widehat E_X}
\right)
\\
&\times
\operatorname{Path}
\left(
\widetilde\alpha^{\partial}_{\mathrm{univ},\mathbf1_X}
\circ\eta_0,
\eta_1
\right).
\end{aligned}
}
\]
Thus a point is a differential natural transformation together with a trace
compatibility
\[
\tau_0\simeq
\tau_1\widetilde\alpha^{\partial}_{\widehat E_X}
\]
and a fundamental-class compatibility
\[
\widetilde\alpha^{\partial}_{\mathbf1_X}\eta_0\simeq\eta_1.
\]
The coherent-stage, curvature, and prescribed-motivic refinements of this path
space are constructed below by a right fibration over the decorated hom
category.  No single map
$\operatorname{DecSp}(c_0)\to\operatorname{DecSp}(c_1)$ is used.
\end{definition}

\begin{construction}[The differential-lift hom $\infty$-category]
\label{constr:prequantum-differential-lift-hom-category}
Fix represented differential coefficient pairs
$(X,\widehat E_X)$ and $(B,\widehat E_B)$ in the Chapter~10 transgression
range.  Put
\[
\mathcal D_{X,B}^{\partial}
:=
\operatorname{Fun}^{L}(\mathcal R_X^0,\mathcal R_B^0),
\qquad
\mathcal D_{X,B}^{\mathrm{sh}}
:=
\operatorname{Fun}(\mathcal R_X^0,\operatorname{SH}(B)).
\]
Postcomposition with homotopification gives
\[
H_B\circ-:\mathcal D_{X,B}^{\partial}\longrightarrow
\mathcal D_{X,B}^{\mathrm{sh}},
\]
while the Chapter~8 exceptional correspondence functor and precomposition with
$H_X$ give
\[
\Phi_{X,B}^{\mathrm{mot}}:
\operatorname{Corr}_E(X,B)\longrightarrow
\mathcal D_{X,B}^{\mathrm{sh}},
\qquad
c\longmapsto F_c^{\mathrm{mot}}H_X.
\]
Form the pullback $\infty$-category
\[
\mathcal L_{X,B}^{\partial}
:=
\operatorname{Corr}_E(X,B)
\times_{\mathcal D_{X,B}^{\mathrm{sh}}}
\mathcal D_{X,B}^{\partial}
\]
and take its full subcategory on the distinguished objects
\[
(c,F_c^{\partial},
  H_BF_c^{\partial}\simeq F_c^{\mathrm{mot}}H_X)
\]
constructed above.  Thus the mapping space between two such objects is
\[
\operatorname{Map}_{\operatorname{Corr}_E}(c_0,c_1)
\times_{
 \operatorname{Map}_{\mathcal D^{\mathrm{sh}}}
 (H_BF_0^{\partial},H_BF_1^{\partial})}
\operatorname{Map}_{\mathcal D^{\partial}}
(F_0^{\partial},F_1^{\partial}).
\]
Its fibre over a fixed geometric $2$-cell $\alpha$ is exactly
$\operatorname{TwoCellReal}_{\partial}(\alpha)$.
\end{construction}

\begin{construction}[Decorated hom $\infty$-category]
\label{constr:prequantum-decorated-hom-category}
Evaluation at the two selected source coefficients defines
\[
\operatorname{ev}_{\widehat E_X},\operatorname{ev}_{\mathbf1_X}:
\mathcal L_{X,B}^{\partial}\longrightarrow\mathcal R_B^0.
\]
Use the source projection from the slice and the target projection from the
coslice to form
\[
\mathcal T_{X,B}
:=
\mathcal L_{X,B}^{\partial}
\times_{\mathcal R_B^0}
(\mathcal R_B^0)_{/\widehat E_B},
\]
\[
\mathcal F_{X,B}
:=
\mathcal L_{X,B}^{\partial}
\times_{\mathcal R_B^0}
(\mathcal R_B^0)_{\mathbf1_B/},
\]
where the maps from $\mathcal L_{X,B}^{\partial}$ are respectively
$\operatorname{ev}_{\widehat E_X}$ and $\operatorname{ev}_{\mathbf1_X}$, and
the maps from the slice and coslice are respectively source and target.
Define
\[
\boxed{
\operatorname{Corr}^{\partial}_{\mathrm{tr,fund}}
((X,\widehat E_X),(B,\widehat E_B))
:=
\mathcal T_{X,B}
\times_{\mathcal L_{X,B}^{\partial}}
\mathcal F_{X,B}.
}
\]
A morphism in the slice component is a differential lift
$\widetilde\alpha:F_0^{\partial}\to F_1^{\partial}$ together with
\[
\tau_0\simeq\tau_1\widetilde\alpha_{\widehat E_X},
\]
whereas a morphism in the coslice component carries
\[
\widetilde\alpha_{\mathbf1_X}\eta_0\simeq\eta_1.
\]
Consequently the mapping-space fibre over
$\alpha:c_0\Rightarrow c_1$ is precisely
$\operatorname{TwoCellDecReal}_{\partial}(\alpha;d_0,d_1)$.
\end{construction}

\begin{construction}[Motivic and complete-stage refinement fibrations]
\label{constr:prequantum-refined-decorated-hom-category}
Write
\[
\mathcal C_{X,B}
:=
\operatorname{Corr}^{\partial}_{\mathrm{tr,fund}}
((X,\widehat E_X),(B,\widehat E_B)).
\]
For an object $(c,d)$ of $\mathcal C_{X,B}$ define
\[
\operatorname{DecLiftSp}(c)
:=
\operatorname{Map}_B(B,\operatorname{DecLift}(c)),
\]
\[
R_{\mathrm{mot}}(c,d)
:=
\operatorname{fib}_{d}
\bigl(\operatorname{DecLiftSp}(c)\to\operatorname{DecSp}(c)\bigr),
\]
and
\[
R_{\mathrm{stage}}(c,d)
:=
\operatorname{fib}_{d}
\bigl(
\operatorname{TransStageCurvReal}_{\partial}(c)
\longrightarrow\operatorname{DecSp}(c)
\bigr).
\]
The motivic factor is included only on a branch with prescribed motivic trace
and fundamental-class points.

Put
\[
\operatorname{Sq}_I
:=
\operatorname{Fun}
\bigl([1]\times[1],\operatorname{Fun}(I,\operatorname{Sp})\bigr),
\qquad
\operatorname{Sq}_{\mathrm{out}}
:=
\operatorname{Fun}([1]\times[1],\operatorname{Sp}).
\]
A point $r\in R_{\mathrm{stage}}(c,d)$ determines a square
$S_r\in\operatorname{Sq}_I$ with horizontal maps
$\mathbf T^\partial,\mathbf T^U$, vertical maps
$\mathbf\kappa_X,\mathbf\kappa_B$, and square cell the selected stagewise
curvature homotopy.  The output functors $\Theta_X,\Theta_B$ determine
$\Theta_\square:\operatorname{Sq}_I\to\operatorname{Sq}_{\mathrm{out}}$.

Let $\alpha:(c_0,d_0)\to(c_1,d_1)$ be a raw decorated morphism.  Its
differential $2$-cell and two decoration paths determine a fixed-boundary
morphism
\[
\overline\alpha:
\Theta_\square S_{r_0}\longrightarrow\Theta_\square S_{r_1}
\]
in $\operatorname{Sq}_{\mathrm{out}}$.  Define the complete stage refinement
of $\alpha$ by
\[
\boxed{
\operatorname{TwoCellStageReal}(\alpha;r_0,r_1)
:=
\operatorname{fib}_{\overline\alpha}
\left(
\underline{\operatorname{Map}}^{\mathrm{fix}}_{\operatorname{Sq}_I}
(S_{r_0},S_{r_1})
\longrightarrow
\underline{\operatorname{Map}}^{\mathrm{fix}}_{\operatorname{Sq}_{\mathrm{out}}}
(\Theta_\square S_{r_0},\Theta_\square S_{r_1})
\right).
}
\]
Here ``fix'' means that the four vertex maps are the maps already supplied by
$\alpha$; the fibre selects the square morphism and all of its coherence cells.
It may be empty.  On the prescribed-motivic branch, the same fixed-boundary
construction in the motivic trace/fundamental-class diagram is the explicit
fibre
\[
\boxed{
\operatorname{TwoCellMotReal}(\alpha;m_0,m_1)
:=
\operatorname{fib}_{\overline\alpha^{\mathrm{mot}}}
\left(
\underline{\operatorname{Map}}^{\mathrm{fix}}
_{\operatorname{DecLift}}(m_0,m_1)
\longrightarrow
\underline{\operatorname{Map}}^{\mathrm{fix}}
_{\operatorname{Dec}}(d_0,d_1)
\right),
}
\]
where $\overline\alpha^{\mathrm{mot}}$ is the motivic trace/fundamental-class
morphism induced by the two decoration paths of $\alpha$.

Define the objects of $\mathcal C_{X,B}^{\mathrm{coh,curv}}$ to be triples
$((c,d),r)$ and their mapping spaces by
\[
\boxed{
\operatorname{Map}_{\mathcal C_{X,B}^{\mathrm{coh,curv}}}
(((c_0,d_0),r_0),((c_1,d_1),r_1))
:=
\int_{\alpha\in
\operatorname{Map}_{\mathcal C_{X,B}}
((c_0,d_0),(c_1,d_1))}
\operatorname{TwoCellStageReal}(\alpha;r_0,r_1).
}
\]
Objects of $\mathcal C_{X,B}^{\mathrm{mot,coh,curv}}$ additionally carry
$m\in R_{\mathrm{mot}}(c,d)$, and the integrand is the product of the stage and
motivic two-cell realizations over the same $\alpha$.

Composition is induced by composition in $\operatorname{Sq}_I$ and by pasting
the terminal comparison paths.  Since restriction to
$\operatorname{Sq}_{\mathrm{out}}$ is a functor, the composite remains in the
prescribed fibre.  Associativity, units, and all higher coherences are inherited
from the two functor categories.  Consequently the displayed mapping spaces
define actual $\infty$-categories.  No contravariant action of a raw decorated
$2$-cell on a complete stage refinement is asserted or required.

Constructions~\ref{constr:prequantum-composite-differential-trace},
\ref{constr:prequantum-composite-fundamental-class}, and
\ref{constr:prequantum-underlying-transgression-composition}, together with
stagewise square pasting, construct
\[
R_{\mathrm{stage}}(c_1,d_1)\times R_{\mathrm{stage}}(c_2,d_2)
\longrightarrow
R_{\mathrm{stage}}(c_2\odot c_1,d_2\star d_1),
\]
and the analogous motivic map.  Exceptional finite-string pasting and
associativity in $\operatorname{Sp}$ prove the associative and unital
coherences.  The refined hom categories therefore form
$\operatorname{Cat}_{\infty}$-enriched categories.
\end{construction}

\begin{definition}[Traced differential correspondence $(\infty,2)$-category]
\label{def:prequantum-traced-differential-correspondence-category}
Objects are represented differential coefficient pairs
$(X,\widehat E_X)$.  The hom $\infty$-categories are those of
Construction~\ref{constr:prequantum-decorated-hom-category}.  Horizontal
composition is induced by derived fibre product of correspondences, exceptional
Beck--Chevalley pasting, and
Constructions~\ref{constr:prequantum-composite-differential-trace} and
\ref{constr:prequantum-composite-fundamental-class}.  The complete
Segal nerve of this $\operatorname{Cat}_{\infty}$-enriched category is denoted
\[
\operatorname{Corr}^{\partial}_{\mathrm{tr,fund}}.
\]
\end{definition}

\begin{theorem}[Functorial prequantum transgression]
\label{thm:prequantum-functorial-transgression}
The preceding data satisfy the enriched pseudofunctor conditions of Chapter~8
and therefore define an $(\infty,2)$-category
$\operatorname{Corr}^{\partial}_{\mathrm{tr,fund}}$.  Differential class spaces
define a functor on it whose value on a $1$-morphism is scalar differential
transgression and whose value on a $2$-morphism is the selected differential
natural transformation.  On the coherent-stage and curvature-compatible
hom categories
$\mathcal C_{X,B}^{\mathrm{coh,curv}}$ and
$\mathcal C_{X,B}^{\mathrm{mot,coh,curv}}$ this functor maps to the
corresponding closed-class transgression functor and the two are related by the
complete stagewise curvature natural transformation.
\end{theorem}

\begin{proof}
The pullback $\mathcal L_{X,B}^{\partial}$ is an actual $\infty$-category, not
only a specification of objects and mapping spaces.  Its composition is the
composition in the two functor categories and in
$\operatorname{Corr}_E(X,B)$, restricted by the pullback universal property.
The slice and coslice are also $\infty$-categories, and their composition pastes
the two decoration paths with the required opposite variances.  Explicitly,
\[
\tau_0\simeq\tau_1\widetilde\alpha
      \simeq\tau_2\widetilde\beta\widetilde\alpha,
\qquad
\widetilde\beta\widetilde\alpha\eta_0
\simeq\widetilde\beta\eta_1
\simeq\eta_2.
\]

For horizontal composition,
Theorem~\ref{thm:prequantum-coherent-trace-fundamental-composition} supplies the
functor on the decorated hom categories.  The finite-string exceptional
Beck--Chevalley equivalences and
Theorem~\ref{thm:prequantum-fubini-transgression} supply its associativity,
unit, face, degeneracy, and subdivision coherences.  The Chapter~8 enriched
extension criterion therefore makes these hom categories into a
$\operatorname{Cat}_{\infty}$-enriched category.  Its complete Segal nerve is an
$(\infty,2)$-category.  Evaluation on differential classes is
Construction~\ref{constr:prequantum-differential-transgression}, and curvature
compatibility is
Theorem~\ref{thm:prequantum-transgression-curvature-compatibility}.
The final paragraph of
Construction~\ref{constr:prequantum-refined-decorated-hom-category} supplies
horizontal composition on the refined hom categories; its mapping-space formula
supplies the required refined $2$-cells.  Hence their complete Segal nerves give
the asserted coherent-stage and curvature-compatible $(\infty,2)$-categories.
\end{proof}

Thus the correspondence-level compositional part of extended prequantum field
theory is an actual decorated $(\infty,2)$-construction.  No differential lift of
a geometric $2$-cell is hidden inside the notation of the correspondence
category.

\section{Differential gluing and the BNV--Stokes comparison}
\label{sec:prequantum-gluing-stokes}

\begin{theorem}[Primary differential gluing]
\label{thm:prequantum-primary-differential-gluing}
For a decomposition into composable traced correspondences, the differential
prequantum transgression of the composite is canonically the composite of the two
transgressions.  After pullback to
$\operatorname{TransStageCurvReal}_{\partial}$, its curvature shadow is
the corresponding closed/underlying gluing comparison.  Its motivic shadow is the
Chapter~9--10 homotopification comparison on the theorem domain where that shadow
is defined.
\end{theorem}

\begin{proof}
The differential statement is the two-stage case of
Theorem~\ref{thm:prequantum-fubini-transgression}.  The closed statement is
Theorem~\ref{thm:prequantum-transgression-curvature-compatibility} applied to the
composite coherent trace diagram.  Homotopification is the already constructed
Chapter~9--10 exceptional comparison; applying it gives precisely the stated
motivic shadow on its theorem domain.
\end{proof}

\begin{construction}[Prequantum BNV--Stokes discrepancy]
\label{constr:prequantum-bnv-stokes-discrepancy}
For a smooth separated finite-presentation morphism $p$, retain the Chapter~9
exact functor-valued obstruction
\[
\operatorname{Obs}_{\mathrm{BNV}}(p)
=
\operatorname{cofib}(\omega_p).
\]
For a differential coefficient $\widehat E$ used in a prequantum transgression,
define the evaluated discrepancy
\[
\boxed{
C^{\mathrm{Stokes}}_{p,\widehat E}
:=
\operatorname{Obs}_{\mathrm{BNV}}(p)(\widehat E).
}
\]
The full subcategory on coefficients for which this evaluation vanishes is the
objectwise equivalence locus of $\kappa_p$ by Chapter~9,
Theorem~9.32.1.  A small complete detection family is constructed below.
\end{construction}

\begin{definition}[BNV--Stokes inverse locus]
\label{def:prequantum-bnv-stokes-inverse-locus}
Let
\[
\kappa_p:p_{e!}\longrightarrow\widehat p_{e!}
\]
be the Chapter~9 canonical natural transformation from the geometric differential
exceptional lift to its BNV correction.  Define
\[
\boxed{
\operatorname{StokesEq}_{p}
:=
\operatorname{Inv}(\kappa_p)
}
\]
in the natural-transformation mapping object.  A point is coherent inverse data for
the entire functor-level BNV comparison, not merely for its value on one
coefficient.
\end{definition}

\begin{theorem}[Small detection of the functor-level BNV comparison]
\label{thm:prequantum-bnv-small-detection}
There is a regular cardinal $\lambda$ such that $\mathcal R_X^0$ is
$\lambda$-accessible and both
\[
p_{e!},\widehat p_{e!}:\mathcal R_X^0\longrightarrow\mathcal R_Y^0
\]
preserve $\lambda$-filtered colimits.  For a small skeleton
$(\mathcal R_X^0)^\lambda$ of the $\lambda$-compact objects define
\[
\boxed{
\operatorname{StokesTestEq}_p^\lambda
:=
\prod_{G\in(\mathcal R_X^0)^\lambda}
\operatorname{Inv}\bigl(\kappa_p(G)\bigr).
}
\]
Evaluation induces a canonical equivalence
\[
\boxed{
\operatorname{StokesEq}_p
\simeq
\operatorname{StokesTestEq}_p^\lambda.
}
\]
Thus the functor-level comparison is detected by an actual small family, although
it is not detected by one arbitrary coefficient.
\end{theorem}

\begin{proof}
The geometric exceptional lift is accessible.  The Chapter~9 correction is
constructed from accessible localization and coreflection functors, the geometric
lift, and finite pullbacks in presentable stable categories, hence is accessible.
Choose a regular cardinal $\lambda$ large enough that the source category is
$\lambda$-accessible and both functors preserve $\lambda$-filtered colimits.

An equivalence of natural transformations is objectwise, so a point of
$\operatorname{StokesEq}_p$ gives a point of the displayed product.  Conversely,
write an arbitrary $E\in\mathcal R_X^0$ as a $\lambda$-filtered colimit
$E\simeq\operatorname*{colim}_iG_i$ of $\lambda$-compact objects.  Naturality and
$\lambda$-filtered-colimit preservation identify
\[
\kappa_p(E)
\simeq
\operatorname*{colim}_i\kappa_p(G_i).
\]
If the test inverse data are present, every arrow in the colimit is an
equivalence, hence so is $\kappa_p(E)$.  Therefore $\kappa_p$ is objectwise an
equivalence and admits its contractibly unique inverse as a natural
transformation.  The two constructions are inverse.  Enlarging $\lambda$ gives
the same inverse locus because both sides identify with
$\operatorname{Inv}(\kappa_p)$.
\end{proof}

\begin{theorem}[BNV--Stokes compatibility]
\label{thm:prequantum-bnv-stokes-compatibility}
Over $\operatorname{StokesEq}_{p}$, differential prequantum transgression preserves
the universal BNV interval primitive and hence the normalized endpoint-difference
factorization.  The resulting endpoint identity is the brave-new Stokes homotopy
for every selected prequantum coefficient in the pulled-back family.  In
particular every evaluated discrepancy $C^{\mathrm{Stokes}}_{p,\widehat E}$
vanishes there.
\end{theorem}

\begin{proof}
Chapter~9 proves that $\kappa_p$ is an equivalence exactly when the BNV cycle
comparison is an equivalence.  On that inverse locus the original geometric
exceptional lift agrees with its canonical BNV correction.  The Chapter~9--10
BNV--Umkehr naturality theorem then preserves the universal interval integral and
its endpoint-difference factorization.  Evaluation on $\widehat E$ gives the
vanishing of the displayed coefficientwise cofiber.  The exact converse is the
small-family criterion of
Theorem~\ref{thm:prequantum-bnv-small-detection}.
\end{proof}

\section{Formal automorphisms of a selected prequantization}
\label{sec:prequantum-formal-automorphisms}

Let $y:Y\to T$ be native and let
\[
\mathcal H_\bullet\in\operatorname{FormLieGpd}(Y/T),
\qquad
q_{\mathcal H}:Y\twoheadrightarrow Q_{\mathcal H}
\]
be its effective formal quotient.  Fix a stable curvature diagram on $Y$
\[
\mathscr D\xrightarrow{\kappa}\mathscr U\xleftarrow{u}\mathscr C
\]
and a selected degree-$r$ prequantization
\[
\widehat\omega:
\mathbf1_Y\longrightarrow
\mathscr D^{\mathrm{cl}}_{\kappa,u}[r].
\]
Write
\[
\omega:=\operatorname{curv}_{\kappa,u}(\widehat\omega):
\mathbf1_Y\longrightarrow\mathscr C[r].
\]

An arbitrary coefficient on $Y$ does not carry canonical transport merely because
its two vertex maps have the same composite to $T$.  We therefore construct the
transport realization before defining stabilizers.

\subsection{Cartan transport of a curvature diagram}
\label{subsec:prequantum-cartan-transport-curvature-diagram}

Let
\[
J_{\mathcal H}:=q_{\mathcal H}^*\Pi_{q_{\mathcal H}}
\]
be the stable jet comonad.  For a test $c:P\to Q_{\mathcal H}$ put
$Y_P:=P\times_{Q_{\mathcal H}}Y$ and let $J_{\mathcal H,P}$ be the base-changed
jet comonad.

\begin{definition}[Cartan-transport realization of a curvature diagram]
\label{def:prequantum-cartan-curvature-realization}
Let $\mathscr J$ be the walking span.  Define
\[
\boxed{
\operatorname{CartCurvReal}_{\mathcal H}(\kappa,u)
\longrightarrow Q_{\mathcal H}
}
\]
by the space-valued formula
\[
\boxed{
\operatorname{CartCurvReal}_{\mathcal H}(\kappa,u)(P)
:=
\operatorname{Fun}\!\left(
 \mathscr J,
 \operatorname{Coalg}_{J_{\mathcal H,P}}(E_{Y_P})
\right)^{\simeq}
\times_{\operatorname{Fun}(\mathscr J,E_{Y_P})^{\simeq}}
\left\{
\mathscr D_P\xrightarrow{\kappa_P}\mathscr U_P
\xleftarrow{u_P}\mathscr C_P
\right\}.
}
\]
Thus a point consists of Cartan coactions on all three coefficients for which
both displayed arrows are coalgebra morphisms.  Passing to maximal subgroupoids
before taking the fibre is essential: the categorical fibre of the forgetful
functor is not itself the space-valued realization used below.
\end{definition}

\begin{proposition}[Descent and base change of Cartan curvature transport]
\label{prop:prequantum-cartan-curvature-descent-base-change}
The preceding presheaf is a spectral Nisnevich sheaf over $Q_{\mathcal H}$ and is
compatible with arbitrary raw spectral base change.  Over it, finite limits of the
curvature diagram, in particular
$\mathscr D^{\mathrm{cl}}_{\kappa,u}$ and $\mathscr F_\kappa$, inherit canonical
Cartan transport.
\end{proposition}

\begin{proof}
Chapter~17 constructs the corresponding Cartan-lift realization for one arbitrary
stable coefficient and proves effective descent and arbitrary raw base change.
Apply that argument in the finite diagram category
$\operatorname{Fun}(\mathscr J,E_{Y_P})$, where $\mathscr J$ is the walking span.
Eilenberg--Moore categories for the compatible comonad diagram are computed by the
same limits, and the forgetful functor creates limits.  Hence the realization is a
sheaf and base-change compatible.  The inherited transport on pullbacks and fibres
is then created by the forgetful functor.
\end{proof}

All constructions through the integrated KS section are henceforth understood
after pullback to $\operatorname{CartCurvReal}_{\mathcal H}(\kappa,u)$.  If the
curvature diagram has already descended from $Q_{\mathcal H}$, this realization
has its canonical point and the parameter is suppressed.

\subsection{Raw differential stabilizer}
\label{subsec:prequantum-raw-differential-stabilizer}

For $p\ge0$ write
\[
v_i:\mathcal H_p\longrightarrow Y,
\qquad 0\le i\le p,
\]
for the ordered vertex maps.  Chapter~13 formal-groupoid transport identifies a
jet-comonad coalgebra with coherent descent transport; hence the universal Cartan
point supplies equivalences
\[
\theta_{ij}:v_i^*\mathscr D^{\mathrm{cl}}_{\kappa,u}
\xrightarrow{\sim}
v_j^*\mathscr D^{\mathrm{cl}}_{\kappa,u}
\]
with their complete \v{C}ech cocycle coherence.  Let $\mathbb E[p]$ be the
indiscrete groupoid on the ordered set $\{0,\ldots,p\}$.

\begin{construction}[Coherent equivalence simplex]
\label{constr:prequantum-coherent-equivalence-simplex}
Put
\[
\mathcal D_p
:=
\Omega_{\mathcal H_p}^\infty
(v_0^*\mathscr D^{\mathrm{cl}}_{\kappa,u}[r])
\]
and transport every vertex section to the zeroth coefficient:
\[
\widehat\omega_i^{(0)}
:=
\theta_{i0}(v_i^*\widehat\omega).
\]
Using the internal cotensor by $\mathbb E[p]$, define
\[
\boxed{
\operatorname{Eq}^{\partial}_p(\widehat\omega)
:=
\operatorname{hofib}_{(\widehat\omega_0^{(0)},\ldots,
                         \widehat\omega_p^{(0)})}
\left(
\mathcal D_p^{\mathbb E[p]}
\longrightarrow
\mathcal D_p^{p+1}
\right).
}
\]
A point is a fully coherent system of equivalences among the transported vertex
restrictions of the selected closed differential class.
\end{construction}

\begin{theorem}[Raw differential quantomorphism groupoid]
\label{thm:prequantum-raw-differential-stabilizer-groupoid}
The objects
\[
\operatorname{QuantRaw}_{\widehat\omega,p}(\mathcal H)
:=
\operatorname{Eq}^{\partial}_p(\widehat\omega)
\]
form an internal groupoid
\[
\boxed{
\operatorname{QuantRaw}_{\widehat\omega,\bullet}(\mathcal H)
\rightrightarrows Y
}
\]
and the evident projection is a fixed-object groupoid morphism to
$\mathcal H_\bullet$.
\end{theorem}

\begin{proof}
Restriction along a monotone map $[p']\to[p]$ transports the Cartan descent
coherences and induces a functor $\mathbb E[p']\to\mathbb E[p]$, so the displayed
objects are simplicial over $\mathcal H_\bullet$.  The indiscrete groupoid
$\mathbb E[p]$ is the iterated pushout of $p$ copies of $\mathbb E[1]$ along their
endpoints.  Cotensor converts this finite colimit into the corresponding limit.
After fixing the transported vertex values, the Segal map is therefore the
iterated fibre product of the degree-one equivalence objects.  Units and inverses
are identity and inverse equivalences in the infinite-loop section object; all
higher coherence comes from the cotensor groupoids and the Cartan descent
coherence.
\end{proof}

\begin{definition}[Formal differential quantomorphism groupoid]
\label{def:prequantum-formal-quantomorphism-groupoid}
Define
\[
\boxed{
\operatorname{Quant}_{\widehat\omega,\bullet}(\mathcal H)
:=
\widehat{\operatorname{QuantRaw}_{\widehat\omega,\bullet}(\mathcal H)}^{\,\mathrm{dR}/T}.
}
\]
\end{definition}

\begin{theorem}[Canonical quantomorphism anchor]
\label{thm:prequantum-canonical-quantomorphism-anchor}
There is a canonical fixed-object formal-groupoid morphism
\[
\boxed{
q_{\widehat\omega}:
\operatorname{Quant}_{\widehat\omega,\bullet}(\mathcal H)
\longrightarrow
\mathcal H_\bullet.
}
\]
It is natural in morphisms of Cartan-transported curvature diagrams, selected
prequantizations, fixed-object formal-groupoid maps, and arbitrary native change of
spectral base.
\end{theorem}

\begin{proof}
Write $G_\bullet:=
\operatorname{QuantRaw}_{\widehat\omega,\bullet}(\mathcal H)$ and let
$q_{\mathrm{raw}}:G_\bullet\to\mathcal H_\bullet$ be the raw projection.
Chapter~13, Theorem~13.3.4 supplies the coreflection counit
\[
\varepsilon_G:\widehat G_\bullet^{\,\mathrm{dR}/T}
\longrightarrow G_\bullet.
\]
Define
\[
q_{\widehat\omega}:=q_{\mathrm{raw}}\varepsilon_G.
\]
Equivalently, formalize $q_{\mathrm{raw}}$ and use the canonical equivalence
$\widehat{\mathcal H}^{\,\mathrm{dR}/T}\simeq\mathcal H$, since the target is
already formal.  Naturality of the coreflection counit identifies these two
descriptions and proves naturality in every listed variable.  Notice that the
arrow points from the formalized source to the raw source; no oppositely directed
factorization is used.
\end{proof}

\subsection{Curvature symmetries}
\label{subsec:prequantum-curvature-symmetries}

Apply the transported stabilizer construction to
$\omega:\mathbf1_Y\to\mathscr C[r]$ itself.

\begin{definition}[Raw and formal curvature-symmetry groupoids]
\label{def:prequantum-curvature-symmetry-groupoid}
Let
\[
\operatorname{CurvSymRaw}_{\omega,\bullet}(\mathcal H)
\rightrightarrows Y
\]
be the raw endpoint-fixed Cartan stabilizer of the complete closed class $\omega$,
constructed by the same indiscrete-groupoid cotensors as
$\operatorname{QuantRaw}_{\widehat\omega,\bullet}(\mathcal H)$.  Define
\[
\boxed{
\operatorname{CurvSym}_{\omega,\bullet}(\mathcal H)
:=
\widehat{\operatorname{CurvSymRaw}_{\omega,\bullet}(\mathcal H)}^{\,\mathrm{dR}/T}.
}
\]
The projection
$\mathscr D^{\mathrm{cl}}_{\kappa,u}\to\mathscr C$ is a Cartan morphism and hence
gives compatible raw and formal groupoid morphisms
\[
\boxed{
\operatorname{QuantRaw}_{\widehat\omega,\bullet}(\mathcal H)
\longrightarrow
\operatorname{CurvSymRaw}_{\omega,\bullet}(\mathcal H),
}
\]
\[
\boxed{
\operatorname{Quant}_{\widehat\omega,\bullet}(\mathcal H)
\longrightarrow
\operatorname{CurvSym}_{\omega,\bullet}(\mathcal H).
}
\]
\end{definition}

\section{Flat gauge transformations and the exact curvature kernel}
\label{sec:prequantum-flat-gauge-kernel}

Retain $\mathscr F_\kappa=\operatorname{fib}(\kappa)$.  Its stable Cartan transport
is created by the finite-limit statement of
Proposition~\ref{prop:prequantum-cartan-curvature-descent-base-change}.

\begin{definition}[Raw and formal flat gauge groupoids]
\label{def:prequantum-flat-gauge-groupoid}
Put
\[
G_{\mathrm{flat}}(\widehat\omega)
:=
\Omega_Y^\infty\mathscr F_\kappa[r-1].
\]
This is a fibrewise grouplike $E_\infty$-object over $Y$.  Let
$B_YG_{\mathrm{flat}}(\widehat\omega)$ be its one-object bar groupoid and define
\[
\boxed{
\operatorname{FlatGau}_{\widehat\omega,\bullet}
:=
\widehat{B_YG_{\mathrm{flat}}(\widehat\omega)}^{\,\mathrm{dR}/T}.
}
\]
\end{definition}

\begin{theorem}[Flat gauge is the raw and formal curvature kernel]
\label{thm:prequantum-flat-gauge-curvature-kernel}
There are canonical equivalences
\[
\boxed{
B_YG_{\mathrm{flat}}(\widehat\omega)
\simeq
\operatorname{hofib}_{\mathrm{id}_\omega}
\left(
\operatorname{QuantRaw}_{\widehat\omega,\bullet}(\mathcal H)
\longrightarrow
\operatorname{CurvSymRaw}_{\omega,\bullet}(\mathcal H)
\right),
}
\]
and, after unitwise formalization,
\[
\boxed{
\operatorname{FlatGau}_{\widehat\omega,\bullet}
\simeq
\operatorname{hofib}_{\mathrm{id}_\omega}
\left(
\operatorname{Quant}_{\widehat\omega,\bullet}(\mathcal H)
\longrightarrow
\operatorname{CurvSym}_{\omega,\bullet}(\mathcal H)
\right).
}
\]
\end{theorem}

\begin{proof}
Before formalization, take the homotopy fibre over the identity curvature
symmetry.  In degree one this is the based loop of the fibre of
\[
\Omega^\infty
\mathscr D^{\mathrm{cl}}_{\kappa,u}[r]
\longrightarrow
\Omega^\infty\mathscr C[r]
\]
at the selected lift.  The stable fibre sequence of prequantization identifies
that loop object with $\Omega^\infty\mathscr F_\kappa[r-1]$.  The higher degrees
are its bar nerve because an endpoint-fixed equivalence simplex in the fibre is a
coherent string of such loops.  Unitwise formalization is a right adjoint and
therefore preserves the finite pullback defining the formal kernel.
\end{proof}

\section{The mixed Hamiltonian calculus}
\label{sec:prequantum-hamiltonian-realization}

Hamiltonianity requires a relation-level origin for the complete closed
coefficient.  Such an origin is canonical for the variational closed coefficient
constructed above, but it is not part of the name of an arbitrary
$\mathscr C\in E_Y$.  We therefore represent it before defining generic
contraction.

Fix $Y\to T$ and put
\[
V_\bullet:=R_{Y/T,\bullet},
\qquad
G_Y:=Q_{Y/T}.
\]
Write
\[
a_s:V_s\longrightarrow T
\]
for the structural map and
\[
\pi_s:V_s\longrightarrow Y
\]
for the zeroth-vertex projection.  The distinction is essential: the complete
cochain object is formed over the common structural base $T$, while local
Hamiltonian contractions are obtained only afterwards by adjunction to the
relation and symmetry simplices.  This is the same variance convention used by
the intrinsic de Rham and variational cochains of Chapters~10 and~12.

\begin{construction}[Relation-level coefficient category and structural cochains]
\label{constr:prequantum-relation-level-coefficient-category}
The simplicial relation and contravariant parameterized-stable pullback define a
cosimplicial diagram of stable coefficient categories
\[
[s]\longmapsto E_{V_s}.
\]
Define
\[
\operatorname{RelCoeff}_V(Y/T)
:=
\lim_{[s]\in\Delta}E_{V_s}.
\]
An object $\mathcal V_\bullet$ is therefore a family
$\mathcal V_s\in E_{V_s}$ with the coherent pullback identifications forced by
all faces, degeneracies, and higher simplices of $V_\bullet$.

For such a family define the structural cochains
\[
\boxed{
C_V^s(\mathcal V)
:=
\Pi_{a_s}\mathcal V_s\in E_T.
}
\]
If $r_\alpha:V_m\to V_n$ is a simplicial operator, then
$a_m=a_nr_\alpha$.  The adjunction unit for
$r_\alpha^*\dashv\Pi_{r_\alpha}$ and the Cartesian comparison
$r_\alpha^*\mathcal V_n\simeq\mathcal V_m$ therefore give a canonical
restriction map
\[
\Pi_{a_n}\mathcal V_n
\longrightarrow
\Pi_{a_m}\mathcal V_m.
\]
Functoriality of adjunction units and the Cartesian coherences make
$C_V^\bullet(\mathcal V)$ an actual cosimplicial object of the single stable
category $E_T$.

Define
\[
\operatorname{Und}_V^2(\mathcal V)
:=
N_V^2C_V^\bullet(\mathcal V),
\]
\[
\operatorname{Cl}_V^2(\mathcal V)
:=
F_V^2\operatorname{Tot}_VC_V^\bullet(\mathcal V)[2]
\]
in $E_T$.  The complete-filtration/coherent-cochain equivalence gives the
canonical underlying map
\[
 u_V:
\operatorname{Cl}_V^2(\mathcal V)
\longrightarrow
\operatorname{Und}_V^2(\mathcal V).
\]
\end{construction}

\begin{definition}[Vertical-origin realization of a closed curvature coefficient]
\label{def:prequantum-vertical-origin-realization}
Let
\[
y:Y\longrightarrow T
\]
be the structural morphism and let $u:\mathscr C\to\mathscr U$ be the selected
underlying-class morphism in $E_Y$.  Define
\[
\boxed{
\operatorname{VertCurvReal}_V(\mathscr C\xrightarrow{u}\mathscr U)
}
\]
to be the derived realization of quadruples
\[
(\mathcal V_\bullet,\chi_C,\chi_U,H)
\]
where
\begin{enumerate}
\item $\mathcal V_\bullet\in\operatorname{RelCoeff}_V(Y/T)$;
\item
\[
\chi_C:\mathscr C\xrightarrow{\sim}
y^*\operatorname{Cl}_V^2(\mathcal V)
\]
is an equivalence in $E_Y$;
\item
\[
\chi_U:y^*\operatorname{Und}_V^2(\mathcal V)\longrightarrow\mathscr U
\]
is a morphism;
\item $H$ is a specified coherent homotopy
\[
\chi_U\,y^*u_V\,\chi_C\simeq u.
\]
\end{enumerate}
The realization is formed by unstraightening the maximal subgroupoid of
$\operatorname{RelCoeff}_V(Y/T)$, adjoining the two mapping objects above, and
taking the indicated walking-equivalence and path fibres.  Thus no
``relation-level presentation datum'' is attached to $\mathscr C$ by name: the
entire space of origins and comparison homotopies is constructed.
\end{definition}

A vertical coefficient origin still does not assert that the selected closed
class itself comes from the complete relation cochains.  That is a second,
class-specific lifting problem.

\begin{definition}[Vertical closed-class origin realization]
\label{def:prequantum-vertical-class-origin-realization}
Let
\[
\omega:\mathbf1_Y\longrightarrow\mathscr C[r]
\]
be the selected complete closed class.  Over
$\operatorname{VertCurvReal}_V$ the inverse equivalence $\chi_C^{-1}$ and
pullback along $y$ induce a map of section spaces
\[
\Omega^\infty\operatorname{map}_{E_T}
 (\mathbf1_T,\operatorname{Cl}_V^2(\mathcal V)[r])
\longrightarrow
\Omega^\infty\operatorname{map}_{E_Y}(\mathbf1_Y,\mathscr C[r]).
\]
Define
\[
\boxed{
\operatorname{VertClassReal}_V(\omega)
}
\]
to be its homotopy fibre over $\omega$.  A point is therefore a complete
relation-level closed class
\[
\omega_V^{\mathrm{cl}}:
\mathbf1_T\longrightarrow\operatorname{Cl}_V^2(\mathcal V)[r]
\]
together with a specified equivalence
\[
\chi_C^{-1}\bigl(y^*\omega_V^{\mathrm{cl}}\bigr)\simeq\omega.
\]
Its normalized degree-two adjoint is denoted
\[
\bar\omega:
\mathbf1_{V_2}\longrightarrow\mathcal V_2[r].
\]
\end{definition}

\begin{proposition}[Descent and the variational canonical point]
\label{prop:prequantum-vertical-origin-descent-canonical-point}
The coefficient and class origin realizations commute with arbitrary native
pullback and satisfy effective descent.  On the variational branch, take the
vertical relation and coefficient family to be the actual Chapter~12 vertical
relation and the vertical cosimplicial transgression family
$\mathbb T_L^{\bullet,q}(M)[n]$, after the structural base changes already used
to form $\operatorname{ClTrans}^{2,q}_H(A;M;n)$.  Then the normalization,
complete-filtration, and underlying-transgression identifications construct
canonical points of both $\operatorname{VertCurvReal}_V$ and
$\operatorname{VertClassReal}_V(\omega_L^{(q),\mathrm{cl}})$.
\end{proposition}

\begin{proof}
The relation objects and their structural maps commute with pullback.
Parameterized spectra satisfy effective descent; right Beck--Chevalley transports
the structural dependent products; normalization is a finite matching fibre and
complete totalization is a limit.  Mapping objects, walking-equivalence fibres,
and the class-lifting homotopy fibre therefore descend as well.

On the variational branch, Chapter~12 forms every double cochain by dependent
product along the structural map to the common variational base and proves that
the resulting coefficient is bicosimplicial before normalization.  The object
$\mathbb T_L^{\bullet,q}(M)$ is obtained from that actual diagram by horizontal
totalization, a horizontal tail, and a stable fibre, all functorially in the
vertical simplicial degree.  Its complete vertical stage is exactly
$\operatorname{ClTrans}^{2,q}_H(A;M;n)$ and its normalized degree two is the
underlying presymplectic transgression.  The closed section
$\omega_L^{(q),\mathrm{cl}}$ is already a point of that complete stage.  These
identifications give the two canonical realization points without choosing a
new presentation.
\end{proof}

The class-origin realization is already an object over the coefficient-origin
realization.  Put
\[
\boxed{
P_{\mathrm{mix}}
:=
\operatorname{VertClassReal}_V(\omega).
}
\]
All Hamiltonian constructions below are formed in the slice over
$P_{\mathrm{mix}}$.  Its universal point includes
$\mathcal V_\bullet$, the coefficient comparisons, the complete relation-level
class $\omega_V^{\mathrm{cl}}$, and its comparison with $\omega$.  Thus the
coefficient and class origins are both present without duplicating the
coefficient-origin factor or introducing a second mixed base.

\subsection{Rectangular mixed relation and local mixed coefficients}
\label{subsec:prequantum-rectangular-mixed-relation}

For $p,s\ge0$ define
\[
\operatorname{Rect}_V(p,s)
:=
Y^{\times_{G_Y}(p+1)(s+1)},
\]
with factors indexed by
$(i,j)\in\{0,\ldots,p\}\times\{0,\ldots,s\}$.  Projection to every horizontal row
gives
\[
\operatorname{Rect}_V(p,s)
\longrightarrow
(V_p)^{\times_{G_Y}(s+1)}.
\]
The canonical formal anchor
$\phi_{\mathcal H,p}:\mathcal H_p\to V_p$ gives the lifted row product.

\begin{definition}[Mixed symmetry--vertical relation]
\label{def:prequantum-mixed-symmetry-vertical-relation}
Define
\[
\boxed{
M^{\mathcal H,V}_{p,s}
:=
\operatorname{Rect}_V(p,s)
\times_{(V_p)^{\times_{G_Y}(s+1)}}
(\mathcal H_p)^{\times_{G_Y}(s+1)}.
}
\]
\end{definition}

\begin{proposition}[Bisimplicial mixed geometry]
\label{prop:prequantum-bisimplicial-mixed-geometry}
The objects $M^{\mathcal H,V}_{p,s}$ are canonically bisimplicial, with axes
\[
M^{\mathcal H,V}_{p,0}\simeq\mathcal H_p,
\qquad
M^{\mathcal H,V}_{0,s}\simeq V_s.
\]
\end{proposition}

\begin{proof}
Deletion and repetition of the $i$-coordinate act simultaneously on every row and
on every lifted $\mathcal H_p$-row; deletion and repetition of the $j$-coordinate
act on the rows.  They commute because they act on different indices.  The anchor
is simplicial, so the defining pullback is preserved.  The axis formulas are the
one-row and one-column cases.
\end{proof}

Projection to the zeroth lifted row and the zeroth column give canonical maps
\[
\lambda_{p,s}:M^{\mathcal H,V}_{p,s}\longrightarrow\mathcal H_p,
\qquad
\chi_{p,s}:M^{\mathcal H,V}_{p,s}\longrightarrow V_s.
\]
Define the local mixed coefficient
\[
\boxed{
\underline C^{p,s}_{\mathcal H,V}
:=
\Pi_{\lambda_{p,s}}\chi_{p,s}^*\mathcal V_s
\in E_{\mathcal H_p}.
}
\]
The vertical simplicial maps and the relation-level Cartesian comparisons of
$\mathcal V_\bullet$, transported by right Beck--Chevalley, make
$s\mapsto\underline C^{p,s}_{\mathcal H,V}$ a vertical cosimplicial object in
$E_{\mathcal H_p}$.  Put
\[
\underline\Omega^{p,s}_{\mathcal H,V,\mathrm{loc}}
:=N_V^s\underline C^{p,\bullet}_{\mathcal H,V}.
\]
For brevity, an underlined symbol without the suffix $\mathrm{loc}$ below always
means this vertical-only local normalization.  It never denotes the local
adjoint of a horizontally normalized global term.
Stable Dold--Kan supplies its coherent vertical differential
\[
d_V^{\mathrm{mix}}:
\underline\Omega^{p,s}_{\mathcal H,V,\mathrm{loc}}
\longrightarrow
\underline\Omega^{p,s+1}_{\mathcal H,V,\mathrm{loc}}.
\]
We do not form a horizontal cochain object in the varying categories
$E_{\mathcal H_p}$.  The horizontal coherent differential and its interchange
with $d_V^{\mathrm{mix}}$ are constructed only after the following structural
dependent product places all bidegrees in the single stable category $E_T$.

For every $p,s$ there is an ordered staircase
\[
\tau_{p,s}:M^{\mathcal H,V}_{p,s}\longrightarrow V_{p+s}
\]
retaining the vertices
\[
(0,0),(1,0),\ldots,(p,0),(p,1),\ldots,(p,s).
\]
The coefficient-pullback comparisons of the universal relation origin make
pullback along $\tau_{p,s}$ a morphism of the corresponding local stable
coefficients.

\subsection{Global mixed coefficients and the complete ordered Alexander--Whitney morphism}
\label{subsec:prequantum-ordered-contraction-cartan-homotopy}

The local mixed objects live over the varying bases $\mathcal H_p$.  The
coherent Alexander--Whitney argument must instead be carried out in one stable
category.  Write
\[
h_p:\mathcal H_p\longrightarrow T
\]
for the structural map.  Since
\[
h_p\lambda_{p,s}=a_s\chi_{p,s},
\]
define
\[
\boxed{
C_{\mathcal H,V}^{p,s}
:=
\Pi_{h_p}\underline C_{\mathcal H,V}^{p,s}
\simeq
\Pi_{a_s\chi_{p,s}}\chi_{p,s}^*\mathcal V_s
\in E_T.
}
\]
Every horizontal and vertical simplicial operator of
$M^{\mathcal H,V}_{\bullet,\bullet}$ is a morphism over $T$.  Applying the
adjunction units and the Cartesian comparisons of $\mathcal V_\bullet$ therefore
makes
\[
C_{\mathcal H,V}^{\bullet,\bullet}
\in
\operatorname{Fun}(\Delta\times\Delta,E_T)
\]
an actual bicosimplicial stable object.  No source-changing vertex transport is
needed to define this global object: all faces have the same structural target
$T$.  Put
\[
\Omega_{\mathcal H,V}^{p,s}
:=
N_{\mathcal H}^{p}N_V^{s}
C_{\mathcal H,V}^{\bullet,\bullet}.
\]
Call these the global horizontally--vertically normalized coefficients.  Their
local adjoints carry canonical maps
\[
\operatorname{ad}_{p,s}:
h_p^*\Omega_{\mathcal H,V}^{p,s}
\longrightarrow
\underline\Omega^{p,s}_{\mathcal H,V,\mathrm{loc}}
\]
obtained from the inclusion of the horizontal normalization matching fibre and
the counit of $h_p^*\dashv\Pi_{h_p}$.  These maps are not declared to be
equivalences.  In particular, a global $(2,0)$ component and the local
vertical-only $(2,0)$ coefficient are connected by $\operatorname{ad}_{2,0}$,
not identified by notation.

Retain the structurally defined vertical cosimplicial object
\[
C_V^s(\mathcal V)=\Pi_{a_s}\mathcal V_s\in E_T,
\qquad
\Omega_V^s:=N_V^sC_V^\bullet(\mathcal V).
\]
Define the total coherent double-cochain object by
\[
\operatorname{Tot}_{\oplus}(\Omega_{\mathcal H,V})^n
:=
\bigoplus_{p+s=n}\Omega_{\mathcal H,V}^{p,s},
\]
with
\[
d_{\mathrm{tot}}=d_{\mathcal H}+d_V^{\mathrm{mix}},
\qquad
 d_V^{\mathrm{mix}}|_{(p,s)}:=(-1)^p d_V.
\]
The sum is finite in every degree.

\begin{theorem}[Coherent ordered Alexander--Whitney morphism]
\label{thm:prequantum-ordered-aw-mixed-morphism}
Before normalization, the ordered staircase $\tau_{p,s}$ and the universal
relation-origin comparison construct
\[
\begin{aligned}
C_V^{p+s}(\mathcal V)
&=
\Pi_{a_{p+s}}\mathcal V_{p+s}
\\
&\longrightarrow
\Pi_{h_p}\Pi_{\lambda_{p,s}}
\tau_{p,s}^*\mathcal V_{p+s}
\\
&\xrightarrow{\sim}
\Pi_{h_p}\Pi_{\lambda_{p,s}}
\chi_{p,s}^*\mathcal V_s
=
C_{\mathcal H,V}^{p,s}.
\end{aligned}
\]
The first morphism is restriction of dependent sections along $\tau_{p,s}$;
it is well typed because
$a_{p+s}\tau_{p,s}=h_p\lambda_{p,s}$.  The second is the Cartesian comparison
between the two pullbacks of the universal relation-level coefficient.  The
component is natural for the horizontal and vertical operators which preserve
the two consecutive blocks.  At the unique break face, the two adjacent
components are related by the break identity used below; individual components
are not claimed to be natural under that face.

A horizontal codegeneracy repeats two vertices in the initial horizontal block
of the staircase and a vertical codegeneracy repeats two vertices in its final
vertical block.  Hence the resulting $(p+s)$-simplex of $V_{p+s}$ is degenerate.
The matching-fibre universal property therefore gives canonical normalized maps
\[
\operatorname{AW}_{p,s}:
\Omega_V^{p+s}
\longrightarrow
\Omega_{\mathcal H,V}^{p,s}.
\]
They assemble to a morphism
\[
\boxed{
\operatorname{AW}^{\mathrm{coh}}:
\Omega_V^\bullet
\longrightarrow
\operatorname{Tot}_{\oplus}
(\Omega_{\mathcal H,V}^{\bullet,\bullet})
}
\]
in the $\infty$-category of coherent nonnegative cochains in $E_T$.
Its homotopy-category shadow satisfies
\[
\boxed{
(d_{\mathcal H}+d_V^{\mathrm{mix}})
\operatorname{AW}^{\mathrm{coh}}
=
\operatorname{AW}^{\mathrm{coh}}d_V.
}
\]
The morphism itself retains the specified nullhomotopies of consecutive
differentials, horizontal--vertical interchange cells, and all higher
compatibilities.
\end{theorem}

\begin{proof}
The matching-fibre factorization was established above.  Put $N=p+s$ and let
\[
f_{p,s}:[p]\hookrightarrow[N],\quad f_{p,s}(i)=i,
\qquad
b_{p,s}:[s]\hookrightarrow[N],\quad b_{p,s}(j)=p+j
\]
be the consecutive-block injections.  Denote the corresponding unnormalized
staircase component by $F_{p,s}$.  The ordinal identities of Chapter~10,
Proposition~10.10.3 give, away from the break,
\[
d^iF_{p,s}=F_{p+1,s}d^i\quad(0\le i\le p),
\qquad
d^{p+j}F_{p,s}=F_{p,s+1}d^j\quad(1\le j\le s+1),
\]
with the coefficient transports inserted by the Cartesian
Beck--Chevalley comparisons.  At the break they give the actual identity
\[
\boxed{
F_{p+1,s}(d^{p+1}x)
=
F_{p,s+1}(d^0x).
}
\]
In the alternating-coface expansion these two occurrences have signs
$(-1)^{p+1}$ and $(-1)^p$ and cancel.  All remaining terms are exactly the
horizontal differential and the rephased vertical differential
$(-1)^p d_V$.  Therefore
\[
(d_{\mathcal H}+d_V^{\mathrm{mix}})\operatorname{AW}
=\operatorname{AW}d_V
\]
before passage to normalized homotopy classes.

This calculation is performed on the actual bicosimplicial restriction maps
before passing to the homotopy category.  Stable normalization is the total
fibre of the finite codegeneracy matching cubes and stable Dold--Kan is natural
for exact functors.  Consequently the alternating-coface calculation is the
shadow of a morphism in the coherent-cochain $\infty$-category, with the higher
cells supplied by the simplicial identities and the Beck--Chevalley pasting
coherences of the structural dependent products.
\end{proof}

\begin{construction}[Complete ordered Alexander--Whitney morphism]
\label{constr:prequantum-complete-aw-morphism}
Let
\[
F_V^\bullet\operatorname{Tot}_VC_V^\bullet
\]
be the complete vertical totalization filtration in $E_T$.  Let
\[
F_{\mathrm{tot}}^\bullet(\mathcal H,V)
\]
be the complete nonnegative filtration corresponding, under the Chapter~10
complete-filtration/coherent-cochain equivalence, to
$\operatorname{Tot}_{\oplus}(\Omega_{\mathcal H,V})$.
Transporting $\operatorname{AW}^{\mathrm{coh}}$ across that equivalence gives
\[
\boxed{
\operatorname{AW}^{\mathrm{comp}}:
F_V^\bullet\operatorname{Tot}_VC_V^\bullet
\longrightarrow
F_{\mathrm{tot}}^\bullet(\mathcal H,V).
}
\]
This is the required linear complete-filtration morphism.  The totalized de
Rham objects simultaneously retain the commutative-algebra multiplication already
constructed in Chapter~10, so no second multiplicative construction is needed for
the prism argument.
\end{construction}

Over $\operatorname{VertClassReal}_V(\omega)$ retain the universal complete
relation-level lift $\omega_V^{\mathrm{cl}}$ and its normalized degree-two
adjoint
\[
\bar\omega:
\mathbf1_{V_2}\longrightarrow\mathcal V_2[r].
\]
Applying $\operatorname{AW}^{\mathrm{comp}}$ to
$\omega_V^{\mathrm{cl}}$ and then adjointing its symmetry-degree components along
$h_p^*\dashv\Pi_{h_p}$ produces the local mixed sections used below.

\begin{construction}[Ordered contraction and degree-two correction]
\label{constr:prequantum-ordered-hamiltonian-contraction}
The total-degree-two components of $\operatorname{AW}^{\mathrm{coh}}$ applied to
the normalized degree-two component of $\omega_V^{\mathrm{cl}}$ have bidegrees
$(0,2)$, $(1,1)$, and $(2,0)$.  Their local adjoints are denoted
\[
\boxed{
\iota_{\mathcal H}\omega:
\mathbf1_{\mathcal H_1}
\longrightarrow
\underline\Omega^{1,1}_{\mathcal H,V}[r]
}
\]
and
\[
\boxed{
K_{\mathcal H}(\omega):
\mathbf1_{\mathcal H_2}
\longrightarrow
\underline\Omega^{2,0}_{\mathcal H,V}[r].
}
\]
Explicitly they are the $(1,1)$ and $(2,0)$ staircase pullbacks; in particular
\[
K_{\mathcal H}(\omega)
=\phi_{\mathcal H,2}^*\bar\omega
\]
because $\tau_{2,0}=\phi_{\mathcal H,2}$.
\end{construction}

For symmetry degree one define
\[
\underline{\operatorname{Cl}\Omega}^{1,2}_{\mathcal H,V}
:=
F_V^2\operatorname{Tot}_V
\underline C^{1,\bullet}_{\mathcal H,V}[2]
\]
and let
\[
d^{\mathrm{mix}}_{V,\mathrm{cl}}:
\underline\Omega^{1,1}_{\mathcal H,V}
\longrightarrow
\underline{\operatorname{Cl}\Omega}^{1,2}_{\mathcal H,V}
\]
be the exact-cochain classifier induced by the adjacent complete-filtration
stages.  Cartan transport gives the complete source--target difference
\[
\Delta^{\mathrm{cl}}_{\mathcal H}\omega:
\mathbf1_{\mathcal H_1}
\longrightarrow
\underline{\operatorname{Cl}\Omega}^{1,2}_{\mathcal H,V}[r].
\]

\begin{theorem}[Complete closed prism and normalized components]
\label{thm:prequantum-ordered-mixed-prism-identities}
The $(1,\ge2)$ component of
$\operatorname{AW}^{\mathrm{comp}}$ gives the canonical coherent equivalence
\[
\boxed{
\Delta^{\mathrm{cl}}_{\mathcal H}\omega
\simeq
-\,d^{\mathrm{mix}}_{V,\mathrm{cl}}
\bigl(\iota_{\mathcal H}\omega\bigr).
}
\]
Its associated coherent total cocycle has normalized total-degree-three
components
\[
\boxed{
d_{\mathcal H}\bar\omega
\simeq
-\,d_V^{\mathrm{mix}}
(\iota_{\mathcal H}\omega),
}
\]
\[
\boxed{
d_{\mathcal H}(\iota_{\mathcal H}\omega)
\simeq
-\,d_V^{\mathrm{mix}}K_{\mathcal H}(\omega),
}
\]
and
\[
\boxed{
d_{\mathcal H}K_{\mathcal H}(\omega)\simeq0.
}
\]
All three identities are components of the single coherent total-cocycle
condition and retain its higher compatibility cells.
\end{theorem}

\begin{proof}
The complete closed class is a point of the second stage of the complete vertical
filtration.  Apply $\operatorname{AW}^{\mathrm{comp}}$.  The symmetry-degree-one
piece of the resulting filtered morphism identifies the complete source--target
difference with the exact-cochain classifier of the $(1,1)$ component, giving
the first displayed equivalence.

Under the complete-filtration/coherent-cochain equivalence, the image is a
coherently closed total degree-two cochain.  Splitting its total differential by
bidegree gives the $(1,2)$, $(2,1)$, and $(3,0)$ equations above.  The local
formulas are obtained by adjunction from the global normalized coefficients.
Because the source is a coherent cochain object, the nullhomotopy of the total
differential and all higher compatibility cells are carried through the same
construction.
\end{proof}

\begin{theorem}[Brave-new Cartan homotopy formula]
\label{thm:prequantum-cartan-homotopy-formula}
The unary normalized shadow of the complete prism is
\[
\boxed{
d_{\mathcal H}\bar\omega
\simeq
-\,d_V^{\mathrm{mix}}
\bigl(\iota_{\mathcal H}\omega\bigr).
}
\]
The complete closed identity, not this unary shadow, is the input to the
curvature-symmetry construction below.
\end{theorem}

\subsection{Hamiltonian primitive arrows and the multiplicativity realization}
\label{subsec:prequantum-hamiltonian-primitive-arrows}

Put
\[
\mathscr J_{\mathcal H}
:=
\underline\Omega^{1,0}_{\mathcal H,V}[r]
\in E_{\mathcal H_1}.
\]
Its mixed vertical boundary maps to
$\underline\Omega^{1,1}_{\mathcal H,V}[r]$.

\begin{definition}[Raw Hamiltonian arrow object]
\label{def:prequantum-raw-hamiltonian-arrow-object}
Define $\operatorname{HamRaw}_{\omega,1}(\mathcal H)$ over $\mathcal H_1$ by
\[
\boxed{
\operatorname{HamRaw}_{\omega,1}(\mathcal H)
:=
\operatorname{hofib}_{\iota_{\mathcal H}\omega}
\left(
\Omega^\infty\mathscr J_{\mathcal H}
\xrightarrow{d_V^{\mathrm{mix}}}
\Omega^\infty\underline\Omega^{1,1}_{\mathcal H,V}[r]
\right).
}
\]
A point over $h:y_0\to y_1$ is a coherent primitive $J_h$ with
\[
d_V^{\mathrm{mix}}J_h\simeq\iota_h\omega.
\]
The primitive equation is a genuine lifting problem and its fibre may be empty.
\end{definition}

For $p\ge1$ put
\[
A_p:=\underline\Omega^{p,0}_{\mathcal H,V,\mathrm{loc}}[r],
\qquad
B_p:=\underline\Omega^{p,1}_{\mathcal H,V,\mathrm{loc}}[r].
\]
The degree-zero matching fibre is the normalized coefficient itself.  Hence the
axis comparison gives a canonical equivalence
\[
\vartheta_{02}:A_2\xrightarrow{\sim}e_{02}^*A_1.
\]
This equivalence comes from the degree-zero normalization axis; it is not an
inverse to a mixed horizontal coface.  Cartan transport supplies
\[
\operatorname{tr}_{01}^A:e_{12}^*A_1\xrightarrow{\sim}e_{02}^*A_1.
\]
In degree one there are instead canonical horizontal coface maps
\[
\beta_{ij}:e_{ij}^*B_1\longrightarrow B_2.
\]
In particular, no inverse to $\beta_{02}$ is part of the native coefficient
package.

Let
\[
P_1:=\operatorname{HamRaw}_{\omega,1}(\mathcal H),
\qquad
C_2:=e_{01}^*P_1\times_{\mathcal H_2}e_{12}^*P_1.
\]
Over $C_2$ define the candidate long-edge primitive
\[
\mu_K
:=
J_{01}+\operatorname{tr}_{01}^A(J_{12})
+\vartheta_{02}(K_{012})
\in\Omega^\infty e_{02}^*A_1
\]
and its Hamiltonian-equation defect
\[
\epsilon_2
:=
d_V^{\mathrm{mix}}\mu_K-e_{02}^*(\iota_{\mathcal H}\omega)
\in\Omega^\infty e_{02}^*B_1.
\]

\begin{lemma}[Mixed image of the Hamiltonian composition defect]
\label{lem:prequantum-hamiltonian-composition-correction}
There is a canonical coherent nullhomotopy
\[
\boxed{
\beta_{02}(\epsilon_2)\simeq0.
}
\]
Moreover the degree-zero component satisfies, in its normalized
$p=3$ coefficient,
\[
\boxed{
d_{\mathcal H}K_{\mathcal H}(\omega)\simeq0.
}
\]
The first statement does not, without an additional hypothesis, imply
$\epsilon_2\simeq0$.
\end{lemma}

\begin{proof}
Restrict the $(2,1)$ normalized component of
Theorem~\ref{thm:prequantum-ordered-mixed-prism-identities} to
$\mathcal H_2$.  Every term then lies in $B_2$.  Naturality of the vertical
boundary and the Cartan exchange gives
\[
\beta_{02}(d_V^{\mathrm{mix}}\mu_K)
\simeq
\beta_{01}(\iota_{01}\omega)
+\beta_{12}(\iota_{12}\omega)
+d_V^{\mathrm{mix}}K_{012}.
\]
The $(2,1)$ prism identity identifies the right-hand side with
$\beta_{02}(\iota_{02}\omega)$, and hence nullhomotopes
$\beta_{02}(\epsilon_2)$.  The $(3,0)$ component gives the second assertion.
All nullhomotopies and their compatibility cells are inherited from the single
coherent total-cocycle condition.  Since $\beta_{02}$ has not been proved an
equivalence or a monomorphism, no cancellation of $\beta_{02}$ is made.
\end{proof}

The preceding nullhomotopy lifts the defect canonically to
\[
\bar\epsilon_2:
\mathbf1_{C_2}
\longrightarrow
\operatorname{fib}(\beta_{02})|_{C_2}.
\]
Let
\[
\pi_2:C_2\longrightarrow P_{\mathrm{mix}}
\]
be the structural projection.

\begin{construction}[Binary Hamiltonian multiplicativity realization]
\label{constr:prequantum-composition-hamiltonian-primitives}
Let
\[
Z_2:=Z_{C_2}(\bar\epsilon_2)\longrightarrow C_2
\]
be the zero locus, equivalently the nullhomotopy fibre, of the lifted defect.
Define
\[
\boxed{
\operatorname{HamMultReal}_{\omega,2}(\mathcal H)
:=
\Pi_{\pi_2}Z_2.
}
\]
For every test object $T\to P_{\mathrm{mix}}$, dependent-product adjunction
gives
\[
\operatorname{Map}_{/P_{\mathrm{mix}}}
\bigl(T,\operatorname{HamMultReal}_{\omega,2}(\mathcal H)\bigr)
\simeq
\Gamma(C_{2,T},Z_{2,T}).
\]
Thus a point is one nullhomotopy of $\bar\epsilon_2$ varying coherently over the
entire family of composable pairs.  Such a nullhomotopy identifies
$d_V^{\mathrm{mix}}\mu_K$ with $e_{02}^*\iota_{\mathcal H}\omega$ compatibly
with the already fixed nullhomotopy after applying $\beta_{02}$.  Hence $\mu_K$
is an actual point of the long-edge Hamiltonian primitive fibre and defines a
binary multiplication
\[
m_2:C_2\longrightarrow e_{02}^*P_1
\]
whose underlying primitive is $\mu_K$.
\end{construction}

For $m\ge2$ put
\[
S_m:=P_1\times_Y\cdots\times_YP_1
\]
with $m$ factors, and define the prescribed degree objects
\[
\boxed{
N_0:=Y,
\qquad
N_1:=P_1,
\qquad
N_m:=\mathcal H_m
\times_{\mathcal H_1\times_Y\cdots\times_Y\mathcal H_1}S_m.
}
\]
The second projection is the canonical spine map
\[
\sigma_m:N_m\longrightarrow S_m.
\]
It is an equivalence because $\mathcal H_\bullet$ is a groupoid object and its
Segal map in degree $m$ is an equivalence.  In degree two, $N_2=C_2$.

\begin{construction}[Internal mapping object of fixed-object simplicial objects]
\label{constr:prequantum-internal-fixed-groupoid-mapping-object}
For simplicial objects $A_\bullet,B_\bullet$ in a slice of the native
$\infty$-topos define
\[
\underline{\operatorname{Nat}}_\Delta(A,B)
:=
\lim_{(i\to j)\in\operatorname{Tw}(\Delta^{\mathrm{op}})}
\underline{\operatorname{Hom}}(A_i,B_j).
\]
If $A_\bullet,B_\bullet\to G_\bullet$ have common degree-zero object $Y$, set
\[
\boxed{
\underline{\operatorname{Map}}^{\mathrm{fix}}_{/G}(A,B)
:=
\operatorname{hofib}_{(p_A,\mathrm{id}_Y)}
\left(
\underline{\operatorname{Nat}}_\Delta(A,B)
\longrightarrow
\underline{\operatorname{Nat}}_\Delta(A,G)
\times
\underline{\operatorname{Hom}}(Y,Y)
\right).
}
\]
The first component is postcomposition with $B\to G$ and the second is
evaluation in degree zero.  Local Cartesian closure makes the internal Homs
represent maps after arbitrary test-object base change, while the twisted-arrow
end imposes all naturality cells.  Replacing $\Delta$ by
$\Delta_{\le m}$ gives the truncated mapping object, and
\[
\underline{\operatorname{Map}}^{\mathrm{fix}}_{/G}(A,B)
\simeq
\lim_m
\underline{\operatorname{Map}}^{\mathrm{fix}}_{/G,\le m}(A,B).
\]
When $A$, $B$, and $G$ are internal groupoids, this is the internal mapping
object of fixed-object groupoid morphisms, since groupoid objects form a full
limit-closed subcategory of simplicial objects.
\end{construction}

\begin{construction}[Fixed-object Reedy coherence tower]
\label{constr:prequantum-fixed-object-reedy-coherence-tower}
Work first over
$P_{\mathrm{Cat},\le1}:=
\operatorname{HamMultReal}_{\omega,2}(\mathcal H)$ and retain the reflexive
graph $N_1\rightrightarrows N_0$ whose unit is the zero primitive on identity
arrows.  Suppose that the prescribed objects $N_0,\ldots,N_{m-1}$ have acquired
an $(m-1)$-truncated simplicial structure over
$\mathcal H_{\le m-1}$.  Form the relative matching object
\[
M_m^{/\mathcal H}N
:=
M_mN\times_{M_m\mathcal H}\mathcal H_m.
\]
The latching object $L_mN$ has its canonical map to $\mathcal H_m$, and the
truncated structure gives
\[
\lambda_m:L_mN\longrightarrow M_m^{/\mathcal H}N
\]
in the slice over $\mathcal H_m$.  Extensions over $\mathcal H_\bullet$ with the
already prescribed degree object $N_m\to\mathcal H_m$ are represented by the
internal factorization object
\[
\boxed{
\operatorname{Fact}_m(N_m)
:=
\operatorname{hofib}_{\lambda_m}
\left(
\underline{\operatorname{Hom}}_{/\mathcal H_m}(L_mN,N_m)
\times
\underline{\operatorname{Hom}}_{/\mathcal H_m}
  (N_m,M_m^{/\mathcal H}N)
\xrightarrow{\ \circ\ }
\underline{\operatorname{Hom}}_{/\mathcal H_m}
  (L_mN,M_m^{/\mathcal H}N)
\right).
}
\]
Take the further homotopy fibre imposing that the spine map induced by the
factorization is the canonical $\sigma_m$; call it
$\operatorname{Fact}_m^{\sigma}(N_m)$.  In degree two take one additional fibre
requiring the inner face $d_1:N_2\to N_1$ to be the multiplication $m_2$.
Define $P_{\mathrm{Cat},\le m}$ to be the resulting total realization over
$P_{\mathrm{Cat},\le m-1}$ and put
\[
\boxed{
P_{\mathrm{Cat}}
:=
\lim_{m\ge2}P_{\mathrm{Cat},\le m}.
}
\]
The degree-two stage supplies the two unit homotopies.  The degree-three stage
supplies associativity and its boundary compatibilities, and the later stages
supply all higher simplicial identities.  Here $\Delta^{\mathrm{op}}$ is used
with its Reedy structure; it is not a direct Reedy category.  The direct tower is
the skeletal degree filtration $m\in\mathbb N$.
\end{construction}

By the Reedy extension universal property, a $T$-point of $P_{\mathrm{Cat}}$ is
exactly a category-object structure on the fixed reflexive graph with prescribed
degree objects $N_m$, binary multiplication $m_2$, and Segal maps $\sigma_m$.
Let $\mathbb E[1]$ be the walking equivalence of
Construction~\ref{constr:prequantum-walking-equivalence-inverse-locus}.  The
internal functor object and restriction along
$\Delta^1\to\mathbb E[1]$ give
\[
q_{\mathrm{eq}}:
E(N):=\underline{\operatorname{Fun}}(\mathbb E[1],N)
\longrightarrow
\underline{\operatorname{Fun}}(\Delta^1,N)\simeq N_1.
\]
For $m\ge1$, let $\operatorname{Ed}(m)$ be the finite set of nondegenerate
edges of $\Delta^m$ and define
\[
N_m^\sim
:=
N_m
\times_{\prod_{e\in\operatorname{Ed}(m)}N_1}
\prod_{e\in\operatorname{Ed}(m)}E(N),
\qquad
N_0^\sim:=N_0.
\]
Simplicial operators carry equivalence edges to equivalence edges, so these
finite pullbacks assemble into a simplicial object
$N_\bullet^\sim\to N_\bullet$.  Every edge of $N^\sim$ carries coherent inverse
data and hence $N^\sim$ is a groupoid object.  Conversely, a map from a groupoid
object to $N$ sends every edge through $E(N)$, and the resulting factorization
through $N^\sim$ is contractibly unique.  Thus $N^\sim$ is the finite-limit
groupoid core.  Its counit is an equivalence in degree zero, and the Segal maps
show that it is an equivalence in every degree exactly when it is an equivalence
in degree one.  Define the groupoid locus by
\[
\boxed{
P_{\mathrm{Gpd}}
:=
\operatorname{Inv}
\bigl(q_{\mathrm{eq}}:N_1^\sim\longrightarrow N_1\bigr).
}
\]
Over this locus the universal $N_\bullet$ is an internal groupoid; the inverse
data form a contractible space whenever they exist.

The complete first prism gives a fixed-object morphism of reflexive graphs
\[
\gamma_{\le1}:
N_{\le1}
\longrightarrow
\operatorname{CurvSymRaw}_{\omega,\le1}(\mathcal H).
\]
Over $P_{\mathrm{Gpd}}$ define its complete extension realization by
\[
\boxed{
P_{\mathrm{curv}}
:=
\operatorname{hofib}_{\gamma_{\le1}}
\left(
\underline{\operatorname{Map}}^{\mathrm{fix}}_{/\mathcal H}
\bigl(N,\operatorname{CurvSymRaw}_{\omega}(\mathcal H)\bigr)
\longrightarrow
\underline{\operatorname{Map}}^{\mathrm{fix}}_{/\mathcal H_{\le1}}
\bigl(N_{\le1},
\operatorname{CurvSymRaw}_{\omega,\le1}(\mathcal H)\bigr)
\right).
}
\]

\begin{definition}[Hamiltonian multiplicativity realization]
\label{def:prequantum-complete-raw-hamiltonian-nerve}
Define
\[
\boxed{
\operatorname{HamMultReal}_{\omega}(\mathcal H)
:=P_{\mathrm{curv}}.
}
\]
It is an object of the native slice over $P_{\mathrm{mix}}$.  Its test points
are precisely the coherent binary nullification, fixed-object Segal groupoid
structure, and curvature-symmetry extension described above.
\end{definition}

\begin{theorem}[Hamiltonian groupoids over the multiplicativity realization]
\label{thm:prequantum-raw-formal-hamiltonian-groupoids}
Over $\operatorname{HamMultReal}_{\omega}(\mathcal H)$ there is a universal
internal groupoid
\[
\operatorname{HamRaw}_{\omega,\bullet}(\mathcal H)
\rightrightarrows Y
\]
over $\mathcal H_\bullet$.  Its arrow object is $P_1$, its binary
multiplication has underlying primitive
\[
J_{01}+\operatorname{tr}_{01}^A(J_{12})+\vartheta_{02}(K_{012}),
\]
and its Segal maps are equivalences.  Its unitwise formalization
\[
\boxed{
\operatorname{Ham}_{\omega,\bullet}(\mathcal H)
:=
\widehat{\operatorname{HamRaw}_{\omega,\bullet}(\mathcal H)}^{\,\mathrm{dR}/T}.
}
\]
is a formal fixed-object groupoid.
\end{theorem}

\begin{proof}
The universal section of $\Pi_{\pi_2}Z_2$ constructs $m_2$.  The universal
Reedy factorizations construct all faces, degeneracies, and higher simplicial
coherences on the prescribed objects $N_m$; the fibres over $\sigma_m$ make the
Segal maps equivalences.  Pullback to $P_{\mathrm{Gpd}}$ identifies the category
object with its groupoid core.  Finally the universal point of
$P_{\mathrm{curv}}$ extends $\gamma_{\le1}$ to the complete fixed-object
groupoid morphism.  Chapter~13 unitwise formalization is a right adjoint and
therefore preserves the finite limits defining the groupoid object.
\end{proof}

\begin{proposition}[Base change and descent of Hamiltonian multiplicativity]
\label{prop:prequantum-hamiltonian-multiplicativity-base-change-descent}
The objects $\operatorname{HamMultReal}_{\omega,2}(\mathcal H)$ and
$\operatorname{HamMultReal}_{\omega}(\mathcal H)$ commute with arbitrary native
base change and satisfy effective descent.
\end{proposition}

\begin{proof}
Pullback between slices of the native $\infty$-topos has both adjoints and hence
preserves small limits and colimits.  Right Beck--Chevalley identifies the base
change of $\Pi_{\pi_2}$ with the dependent product of the base-changed zero
object, and local Cartesian closure identifies the base change of every internal
Hom.  Latching and matching objects, the factorization fibres, the skeletal
limit, the finite-limit groupoid core, the inverse locus, the twisted-arrow end,
and the curvature-extension fibre therefore all commute with base change.

For an effective epimorphism, the coefficient origins, mixed diagrams, and
binary defect descend by the preceding constructions.  Applying the same
base-change calculation at every Reedy stage gives a descent diagram of the
stage realizations.  Effective descent in the native topos reconstructs each
stage, and taking the inverse limit reconstructs $P_{\mathrm{Cat}}$,
$P_{\mathrm{Gpd}}$, and $P_{\mathrm{curv}}$.  This proves both assertions.
\end{proof}

For comparison, let
\[
\beta_p^\bullet:
e_{0p}^*\underline C^{1,\bullet}_{\mathcal H,V,\mathrm{loc}}
\longrightarrow
\underline C^{p,\bullet}_{\mathcal H,V,\mathrm{loc}}
\qquad (p\ge2)
\]
be the long-edge horizontal comparison of coherent vertical coefficient
diagrams.  Let
\[
h_p:\mathcal H_p\longrightarrow P_{\mathrm{mix}}
\]
be the structural map and define, in the slice over $\mathcal H_p$,
\[
L_p
:=
\operatorname{Inv}_{\operatorname{Fun}(\Delta,E_{\mathcal H_p})}
(\beta_p^\bullet)
\longrightarrow\mathcal H_p.
\]
Let $\Delta_{\mathrm{edge}}$ be the category whose objects are $[p]$ for
$p\ge2$ and whose arrows are endpoint-preserving monotone maps.  Naturality of
the long-edge comparison identifies the pullbacks of the $L_p$ along the
simplicial maps of $\mathcal H_\bullet$.  Right Beck--Chevalley therefore makes
\[
E_p:=\Pi_{h_p}L_p
\in(\mathsf{XSp})_{/P_{\mathrm{mix}}}
\]
an actual diagram
$E_{\mathrm{edge}}:\Delta_{\mathrm{edge}}\to
(\mathsf{XSp})_{/P_{\mathrm{mix}}}$.
Define the sufficient exchange locus in this single slice by
\[
\boxed{
\operatorname{HamEdgeEq}_{\omega}(\mathcal H)
:=
\lim_{\Delta_{\mathrm{edge}}}E_{\mathrm{edge}}.
}
\]
Equivalently, this is the internal object of Cartesian sections of the
unstraightening of the $L_p$ over the endpoint-preserving simplex category.  A
test point is a family of inverses of all $\beta_p^\bullet$, varying over
$\mathcal H_p$, together with every simplicial compatibility square.

\begin{proposition}[Long-edge equivalences kill the obstruction]
There is a canonical morphism
\[
\operatorname{HamEdgeEq}_{\omega}(\mathcal H)
\longrightarrow
\operatorname{HamMultReal}_{\omega}(\mathcal H).
\]
\end{proposition}

\begin{proof}
On the source every $\beta_p^\bullet$ is invertible as a natural transformation,
with inverse natural transformation $(\beta_p^\bullet)^{-1}$ and its complete
coherent inverse cells.  Transport the ordered Alexander--Whitney total cocycle
through these inverses.  In degree two, transport of the $(2,1)$ identity
nullhomotopes $\epsilon_2$ itself and hence canonically gives the section of
$Z_2$ defining the binary stage.  Normalization transports the degeneracy cells
to the two unit fillers.

The transported $(3,0)$ identity $d_{\mathcal H}K\simeq0$ identifies the two
parenthesizations of the multiplication.  In additive primitive notation, the
inverse equation forces
\[
J_{h^{-1}}
=
-\operatorname{tr}_{h^{-1}}J_h
-\operatorname{tr}_{h^{-1}}K(h^{-1},h).
\]
The horizontal cocycle equation for $K$, evaluated on
$(h,h^{-1},h)$ and $(h^{-1},h,h^{-1})$, gives the two inverse identities; its
degeneracy components give their unit compatibilities.  In degree $m$, the
transported total-cocycle boundary is exactly the map
$L_mN\to M_mN$, while the transported degree-$m$ component factors it through
$N_m$.  These factorizations are compatible with the canonical spine maps, so
they define a point of every stage $P_{\mathrm{Cat},\le m}$.  The inverse
identities place the resulting category object on $P_{\mathrm{Gpd}}$.

Finally the complete first-prism components transport in the same way.  Their
degree-$m$ boundaries are the naturality matching maps for a morphism to
$\operatorname{CurvSymRaw}_{\omega}$, and their complete cells provide the
fillers.  They therefore give a point of $P_{\mathrm{curv}}$.  The construction
is natural in the coherent inverses $(\beta_p^\bullet)^{-1}$ and hence defines
the displayed canonical morphism.
\end{proof}

\begin{theorem}[Hamiltonian multiplicativity constructs complete curvature symmetry]
\label{thm:prequantum-hamiltonian-implies-curvature-symmetry}
Over $\operatorname{HamMultReal}_{\omega}(\mathcal H)$ there are canonical raw
and formal groupoid morphisms
\[
\boxed{
\operatorname{HamRaw}_{\omega,\bullet}(\mathcal H)
\longrightarrow
\operatorname{CurvSymRaw}_{\omega,\bullet}(\mathcal H),
}
\]
\[
\boxed{
\operatorname{Ham}_{\omega,\bullet}(\mathcal H)
\longrightarrow
\operatorname{CurvSym}_{\omega,\bullet}(\mathcal H).
}
\]
The degree-one graph map follows from the Hamiltonian primitive equation; its
extension through all compositions is part of the multiplicativity realization.
\end{theorem}

\begin{proof}
The raw morphism is the universal point of the fixed-object mapping fibre
$P_{\mathrm{curv}}$.  Its restriction to arrows is the complete first-prism
path obtained by applying the exact-cochain classifier to
$d_V^{\mathrm{mix}}J\simeq\iota\omega$.  The Reedy mapping fibre supplies its
compatibility with every face, degeneracy, multiplication, inverse, and higher
simplex.  Unitwise formalization gives the second morphism.
\end{proof}

All constructions in the remainder of the chapter that use a Hamiltonian
groupoid are pulled back to
$\operatorname{HamMultReal}_{\omega}(\mathcal H)$; this base is suppressed from
the notation.  On $\operatorname{HamEdgeEq}_{\omega}(\mathcal H)$ the preceding
proposition supplies the required canonical point.

\section{Hamiltonian prequantum realization and the integrated KS extension}
\label{sec:prequantum-integrated-kostant-souriau}

The raw and formal KS extensions are constructed at different levels and are not
identified by a single liftability condition.

\begin{definition}[Raw and formal Hamiltonian-prequantum groupoids]
\label{def:prequantum-hamiltonian-prequantum-lift-realization}
Define first the raw pullback
\[
\boxed{
\operatorname{HamPreQRaw}_{\widehat\omega,\bullet}(\mathcal H)
:=
\operatorname{HamRaw}_{\omega,\bullet}(\mathcal H)
\times_{\operatorname{CurvSymRaw}_{\omega,\bullet}(\mathcal H)}
\operatorname{QuantRaw}_{\widehat\omega,\bullet}(\mathcal H).
}
\]
Then define its unitwise formalization
\[
\boxed{
\operatorname{HamPreQ}_{\widehat\omega,\bullet}(\mathcal H)
:=
\widehat{\operatorname{HamPreQRaw}_{\widehat\omega,\bullet}(\mathcal H)}^{\,\mathrm{dR}/T}.
}
\]
Because formalization is a right adjoint, there is a canonical equivalence
\[
\operatorname{HamPreQ}_{\widehat\omega}
\simeq
\operatorname{Ham}_{\omega}
\times_{\operatorname{CurvSym}_{\omega}}
\operatorname{Quant}_{\widehat\omega}.
\]
\end{definition}

\begin{theorem}[Effective image of a fixed-object groupoid morphism]
\label{thm:prequantum-effective-image-subgroupoid}
Let
\[
F_\bullet:G_\bullet\longrightarrow H_\bullet
\]
be a fixed-object morphism of internal groupoids over $Y$.  Factor its degree-one
map by effective image
\[
G_1\twoheadrightarrow I_1\hookrightarrow H_1.
\]
Then $I_1\rightrightarrows Y$ carries a contractibly unique internal groupoid
structure for which $G_\bullet\to I_\bullet$ is degreewise effective and
$I_\bullet\hookrightarrow H_\bullet$ is a fixed-object groupoid monomorphism.
The construction is functorial under fixed-object groupoid morphisms and
arbitrary base change.
\end{theorem}

\begin{proof}
The source and target maps factor through $I_1$ because they agree with those of
$H_1$ after the effective cover $G_1\twoheadrightarrow I_1$.  The unit and
inversion factor for the same reason.  Effective epimorphisms are stable under
pullback, so
\[
G_1\times_YG_1\twoheadrightarrow I_1\times_YI_1
\]
is effective.  On this cover the composite in $H$ satisfies
\[
m_H(F_1\times F_1)=F_1m_G
\]
and hence lands in the monomorphism $I_1\hookrightarrow H_1$.  Factorization
through a monomorphism is local for effective epimorphisms, giving the
composition
\[
m_I:I_1\times_YI_1\longrightarrow I_1
\]
through a contractible space.  Associativity, unit, and inverse identities hold
after the effective cover and therefore descend uniquely.  The same local
factorization argument is natural under base change and fixed-object maps.
\end{proof}

\begin{construction}[Raw liftable Hamiltonian image and formalization]
\label{constr:prequantum-liftable-hamiltonian-image}
Factor the degree-one raw projection by effective image
\[
\operatorname{HamPreQRaw}_{\widehat\omega,1}
\twoheadrightarrow
I^{\mathrm{Ham,raw}}_{\widehat\omega,1}
\hookrightarrow
\operatorname{HamRaw}_{\omega,1}.
\]
Theorem~\ref{thm:prequantum-effective-image-subgroupoid} constructs a contractibly unique
internal subgroupoid
\[
I^{\mathrm{Ham,raw}}_{\widehat\omega,\bullet}
\hookrightarrow
\operatorname{HamRaw}_{\omega,\bullet}
\]
and a degreewise effective morphism
\[
\operatorname{HamPreQRaw}_{\widehat\omega,\bullet}
\twoheadrightarrow
I^{\mathrm{Ham,raw}}_{\widehat\omega,\bullet}.
\]
Define
\[
\boxed{
\operatorname{Ham}^{\mathrm{lift}}_{\widehat\omega,\bullet}
:=
\widehat{I^{\mathrm{Ham,raw}}_{\widehat\omega,\bullet}}^{\,\mathrm{dR}/T}.
}
\]
The raw image inclusion formalizes to a canonical morphism
\[
c_{\mathrm{KS}}^{\mathrm{for}}:
\operatorname{Ham}^{\mathrm{lift}}_{\widehat\omega,\bullet}
\longrightarrow
\operatorname{Ham}_{\omega,\bullet}.
\]
When no confusion can arise we abbreviate
$I^{\mathrm{Ham}}:=I^{\mathrm{Ham,raw}}$.
\end{construction}

\begin{definition}[Raw and formal Hamiltonian liftability inverse loci]
\label{def:prequantum-hamiltonian-liftability-inverse-locus}
Define separately
\[
\boxed{
\operatorname{KSEq}^{\mathrm{raw}}_{\widehat\omega}(\mathcal H)
:=
\operatorname{Inv}
\left(
I^{\mathrm{Ham,raw}}_{\widehat\omega}
\hookrightarrow
\operatorname{HamRaw}_{\omega}
\right),
}
\]
and
\[
\boxed{
\operatorname{KSEq}^{\mathrm{for}}_{\widehat\omega}(\mathcal H)
:=
\operatorname{Inv}(c_{\mathrm{KS}}^{\mathrm{for}}).
}
\]
Formalization carries an equivalence to an equivalence and therefore induces a
canonical map
\[
\operatorname{KSEq}^{\mathrm{raw}}_{\widehat\omega}
\longrightarrow
\operatorname{KSEq}^{\mathrm{for}}_{\widehat\omega}.
\]
For later brevity write
$\operatorname{KSEq}_{\widehat\omega}:=
\operatorname{KSEq}^{\mathrm{for}}_{\widehat\omega}$ unless the raw locus is
explicitly named.
\end{definition}

\begin{definition}[Recovery of raw liftability from a formal inverse]
\label{def:prequantum-raw-liftability-recovery-fibre}
For a formal inverse
\[
\xi_{\mathrm{for}}:
P\longrightarrow\operatorname{KSEq}^{\mathrm{for}}_{\widehat\omega}
\]
define
\[
\boxed{
\operatorname{KSRawLiftReal}_{\widehat\omega}(\xi_{\mathrm{for}})
:=
P\times_{\operatorname{KSEq}^{\mathrm{for}}_{\widehat\omega}}
\operatorname{KSEq}^{\mathrm{raw}}_{\widehat\omega}.
}
\]
A point is raw inverse data whose formalization is identified with the prescribed
formal inverse.  This fibre is the complete converse problem and may be empty;
no raw information is recovered merely by naming the formal inverse.
\end{definition}

\begin{theorem}[Raw and formal integrated brave-new Kostant--Souriau extension]
\label{thm:prequantum-integrated-kostant-souriau-extension}
There is an unconditional raw effective extension
\[
\boxed{
B_YG_{\mathrm{flat}}(\widehat\omega)
\longrightarrow
\operatorname{HamPreQRaw}_{\widehat\omega,\bullet}
\twoheadrightarrow
I^{\mathrm{Ham,raw}}_{\widehat\omega,\bullet},
}
\]
whose inhabited arrow fibres are the full left and right torsors of flat
prequantum modifications.  Formalization gives an unconditional formal
homotopy-kernel sequence
\[
\boxed{
\operatorname{FlatGau}_{\widehat\omega,\bullet}
\longrightarrow
\operatorname{HamPreQ}_{\widehat\omega,\bullet}
\longrightarrow
\operatorname{Ham}^{\mathrm{lift}}_{\widehat\omega,\bullet}.
}
\]
Over $\operatorname{KSEq}^{\mathrm{raw}}_{\widehat\omega}$ the raw target is the
full raw Hamiltonian groupoid.  Over
$\operatorname{KSEq}^{\mathrm{for}}_{\widehat\omega}$ the formal target is the
full formal Hamiltonian groupoid.  The second assertion does not imply the first.
\end{theorem}

\begin{proof}
The arrow kernel of the raw projection to its effective image consists of
quantomorphisms with identity curvature symmetry and zero Hamiltonian primitive,
hence is $B_YG_{\mathrm{flat}}$ by the raw curvature-kernel theorem.  Chapter~13's
kernel-bitorsor theorem identifies every inhabited image fibre with the
corresponding left and right kernel torsors.  Formalization preserves the finite
pullback defining the kernel and sends the raw effective morphism to the formal
KS morphism, giving the formal homotopy-kernel sequence.

On the raw inverse locus the image monomorphism is an equivalence by definition,
so the raw target is $\operatorname{HamRaw}$.  On the formal inverse locus
$c_{\mathrm{KS}}^{\mathrm{for}}$ is an equivalence, so the formal target is
$\operatorname{Ham}$.  Unitwise formalization is a localization/coreflection and
can erase raw nonlinear information; therefore no raw equivalence is deduced from
the formal one.
\end{proof}

This is the native integrated KS construction.  Raw arrowwise liftability,
formal liftability, and a coherent section of the KS extension remain three
separate derived problems.

\section{Relative prequantum isotropy and the integrated kernel}
\label{sec:prequantum-isotropy-discrepancy}

The whole isotropy group of $\operatorname{HamPreQ}$ need not map trivially to
Hamiltonian isotropy.  The correct integrated comparison is therefore relative
isotropy, and at the integrated level it is already exactly the flat kernel.

\begin{theorem}[Relative isotropy is the integrated flat kernel]
\label{prop:prequantum-universal-isotropy-exact-extension}
Let
\[
I^{\mathrm{for}}_{\mathrm{PreQ}},
\qquad
I^{\mathrm{for}}_{\mathrm{Ham,lift}}
\]
be the actual formal isotropy group objects of the two groupoids in the formal KS
sequence, and let
\[
I^{\mathrm{for}}_{\mathrm{Flat}}
:=
I^{\mathrm{for}}
\bigl(\operatorname{FlatGau}_{\widehat\omega,\bullet}\bigr)
\]
be the isotropy group object of the one-object flat gauge groupoid.  Then there is
a canonical equivalence of group objects over $Y$
\[
\boxed{
I^{\mathrm{for}}_{\mathrm{Flat}}
\simeq
I^{\mathrm{for}}_{\mathrm{PreQ}}
\times_{I^{\mathrm{for}}_{\mathrm{Ham,lift}}}Y,
}
\]
where $Y\to I^{\mathrm{for}}_{\mathrm{Ham,lift}}$ is the identity-isotropy section.
Equivalently,
\[
\operatorname{FlatGau}_{\widehat\omega,\bullet}
\simeq
B_Y\left(
I^{\mathrm{for}}_{\mathrm{PreQ}}
\times_{I^{\mathrm{for}}_{\mathrm{Ham,lift}}}Y
\right).
\]
Thus no additional integrated isotropy-exactness condition, discrepancy, or inverse
locus is required.
\end{theorem}

\begin{proof}
For any fixed-object groupoid morphism $F:G_\bullet\to H_\bullet$, its arrow kernel
over the identity is
\[
G_1\times_{H_1}Y.
\]
An arrow in this pullback has identical source and target because its $H$-image is
the identity, so it factors through $I_G$.  Conversely an isotropy arrow whose
$I_H$-image is the identity is an arrow-kernel point.  Fibre-pasting therefore
gives an equivalence of group objects
\[
G_1\times_{H_1}Y
\simeq
I_G\times_{I_H}Y.
\]
Apply this to the formal KS morphism.  Its arrow-kernel group object is the
isotropy of the one-object groupoid
$\operatorname{FlatGau}_{\widehat\omega,\bullet}$ by the formal homotopy-kernel
theorem.  Delooping the resulting group-object equivalence over $Y$ gives the
second displayed equivalence.  Every step is a finite-limit construction, so the
group laws and higher coherence are preserved.
\end{proof}

The genuinely nontrivial comparison appears only after partition-Lie
differentiation, because no finite-limit preservation theorem for native
differentiation is inserted.  That discrepancy is constructed next.

\section{Differentiation of the integrated prequantum extension}
\label{sec:prequantum-differentiation-quantomorphism}

Enter the locally coherent spectral coefficient domain of Chapters~14--15.  We
differentiate the native formal groupoids separately; no finite-limit preservation
of the full native groupoid differentiation functor is used.

\begin{definition}[Differentiated prequantum controllers]
\label{def:prequantum-quantomorphism-controller}
Define
\[
\boxed{
\mathbb A_{\widehat\omega}^{\mathrm{PreQ}}
:=
\operatorname{Diff}^{\pi,E_\infty}
\bigl(\operatorname{HamPreQ}_{\widehat\omega,\bullet}(\mathcal H)\bigr),
}
\]
\[
\boxed{
\mathbb A_{\widehat\omega}^{\mathrm{Ham,\,lift}}
:=
\operatorname{Diff}^{\pi,E_\infty}
\bigl(\operatorname{Ham}^{\mathrm{lift}}_{\widehat\omega,\bullet}(\mathcal H)\bigr),
}
\]
and independently
\[
\boxed{
\mathbb A_{\widehat\omega}^{\mathrm{Flat}}
:=
\operatorname{Diff}^{\pi,E_\infty}
\bigl(\operatorname{FlatGau}_{\widehat\omega,\bullet}\bigr).
}
\]
The integrated KS projection differentiates to an actual spectral partition-Lie
algebroid morphism
\[
p_{\mathrm{KS}}:
\mathbb A_{\widehat\omega}^{\mathrm{PreQ}}
\longrightarrow
\mathbb A_{\widehat\omega}^{\mathrm{Ham,\,lift}}.
\]
\end{definition}

Define also
\[
\mathbb A_{\widehat\omega}^{\mathrm{Quant}}
:=
\operatorname{Diff}^{\pi,E_\infty}
\bigl(\operatorname{Quant}_{\widehat\omega,\bullet}(\mathcal H)\bigr).
\]
Projection from the raw Hamiltonian-prequantum pullback and formalization give
\[
q_{\mathrm{PQ}}:
\operatorname{HamPreQ}_{\widehat\omega}
\longrightarrow
\operatorname{Quant}_{\widehat\omega}.
\]
Its fibre over a quantomorphism $\widehat h$ is canonically
\[
\boxed{
\operatorname{fib}_{\widehat h}(q_{\mathrm{PQ}})
\simeq
\operatorname{Ham}_{\omega}
\times_{\operatorname{CurvSym}_{\omega}}
\{\operatorname{curv}(\widehat h)\},
}
\]
the complete space of Hamiltonian primitives compatible with the curvature
symmetry of $\widehat h$.  Differentiation constructs the actual comparison
\[
\boxed{
a_{\mathrm{PQ}}:
\mathbb A_{\widehat\omega}^{\mathrm{PreQ}}
\longrightarrow
\mathbb A_{\widehat\omega}^{\mathrm{Quant}}.
}
\]

\begin{definition}[Intrinsic differentiated KS fibre]
\label{def:prequantum-intrinsic-differentiated-ks-fibre}
Form the categorical fibre in the spectral partition-Lie algebroid category
\[
\boxed{
\mathbb F_{\widehat\omega}^{\mathrm{KS}}
:=
\operatorname{fib}_{0}
\left(
\mathbb A_{\widehat\omega}^{\mathrm{PreQ}}
\xrightarrow{p_{\mathrm{KS}}}
\mathbb A_{\widehat\omega}^{\mathrm{Ham,\,lift}}
\right).
}
\]
Its anchor is canonically zero.
\end{definition}

Write
\[
U_{\mathrm{anch}}:
\operatorname{LieAlgd}^{\pi,E_\infty}_{Y/T}
\longrightarrow
(\operatorname{QC}^{\vee}_Y)_{/\,T^{\mathrm{pc}}_{Y/T}[1]}
\]
for the monadic forgetful functor to anchored pro-coherent coefficients, and write
\[
U_{\mathrm{src}}:
\operatorname{LieAlgd}^{\pi,E_\infty}_{Y/T}
\longrightarrow
\operatorname{QC}^{\vee}_Y
\]
for its composite with the forgetful functor from the anchored slice to the
source module.  The first functor creates the categorical fibre above.  The
second lands in a stable category and therefore gives the actual stable fibre
sequence
\[
\boxed{
U_{\mathrm{src}}\mathbb F_{\widehat\omega}^{\mathrm{KS}}
\longrightarrow
U_{\mathrm{src}}\mathbb A_{\widehat\omega}^{\mathrm{PreQ}}
\longrightarrow
U_{\mathrm{src}}\mathbb A_{\widehat\omega}^{\mathrm{Ham,\,lift}}.
}
\]
This distinction is used below whenever a connecting suspension or cofiber is
formed.

\begin{definition}[Prequantum-to-quantomorphism controller discrepancy]
\label{def:prequantum-quantomorphism-controller-discrepancy}
Define
\[
\boxed{
C_{\widehat\omega}^{\mathrm{PreQ/Quant}}
:=
\operatorname{cofib}
\left(U_{\mathrm{src}}a_{\mathrm{PQ}}\right),
}
\]
and
\[
\boxed{
\operatorname{PreQQuantEq}_{\widehat\omega}
:=
\operatorname{Inv}
\left(U_{\mathrm{src}}a_{\mathrm{PQ}}\right).
}
\]
A morphism in the anchored slice is an equivalence exactly when its source-module
map is an equivalence.  Hence the displayed inverse locus is also the inverse
locus of $a_{\mathrm{PQ}}$ in spectral partition-Lie algebroids.  This computes
the comparison with the independently differentiated quantomorphism controller;
no equality is inserted by notation.
\end{definition}

\begin{construction}[Differentiated flat-to-KS comparison]
\label{constr:prequantum-differentiated-flat-ks-comparison}
The composite of formal groupoids
\[
\operatorname{FlatGau}_{\widehat\omega}
\longrightarrow
\operatorname{HamPreQ}_{\widehat\omega}
\longrightarrow
\operatorname{Ham}^{\mathrm{lift}}_{\widehat\omega}
\]
is the zero/identity-base morphism.  Differentiation therefore gives a canonical
nullhomotopy of
\[
\mathbb A_{\widehat\omega}^{\mathrm{Flat}}
\longrightarrow
\mathbb A_{\widehat\omega}^{\mathrm{PreQ}}
\longrightarrow
\mathbb A_{\widehat\omega}^{\mathrm{Ham,\,lift}}.
\]
The categorical fibre universal property constructs
\[
\boxed{
\gamma_{\widehat\omega}^{\mathrm{KS,diff}}:
\mathbb A_{\widehat\omega}^{\mathrm{Flat}}
\longrightarrow
\mathbb F_{\widehat\omega}^{\mathrm{KS}}.
}
\]
\end{construction}

\begin{definition}[Differentiated KS discrepancy and inverse locus]
\label{def:prequantum-differentiated-ks-discrepancy}
Define the stable discrepancy in pro-coherent source modules by
\[
\boxed{
C_{\widehat\omega}^{\mathrm{KS,diff}}
:=
\operatorname{cofib}
\bigl(U_{\mathrm{src}}\gamma_{\widehat\omega}^{\mathrm{KS,diff}}\bigr)
}
\]
Restrict the arrow and walking-equivalence objects of
Construction~\ref{constr:prequantum-walking-equivalence-inverse-locus} to the
small coherent spectral-\'etale site over $Y$ and sheafify there.  Denote the
resulting pullback inverse locus by $\operatorname{Inv}_{\acute et}$.  Define
\[
\boxed{
\operatorname{KSDiffEq}_{\widehat\omega}
:=
\operatorname{Inv}_{\acute et}
\bigl(U_{\mathrm{src}}\gamma_{\widehat\omega}^{\mathrm{KS,diff}}\bigr).
}
\]
Because the forgetful functor from partition-Lie algebroids to their source
modules is conservative, the inverse locus is equivalently the locus on which
$\gamma_{\widehat\omega}^{\mathrm{KS,diff}}$ is an algebroid equivalence.
\end{definition}

\begin{theorem}[Differentiated KS comparison]
\label{thm:prequantum-differentiation-integrated-extension}
The constructions above give a canonical diagram
\[
\begin{tikzcd}[column sep=huge,row sep=large]
\mathbb A_{\widehat\omega}^{\mathrm{Flat}}
  \arrow[r] \arrow[d,"\gamma_{\widehat\omega}^{\mathrm{KS,diff}}"'] &
\mathbb A_{\widehat\omega}^{\mathrm{PreQ}}
  \arrow[r,"p_{\mathrm{KS}}"] &
\mathbb A_{\widehat\omega}^{\mathrm{Ham,\,lift}} \\
\mathbb F_{\widehat\omega}^{\mathrm{KS}}
  \arrow[ur] & &
\end{tikzcd}
\]
where the lower-left object is the intrinsic algebroid fibre.  Over
$\operatorname{KSDiffEq}_{\widehat\omega}$, differentiated actual flat gauge and
the intrinsic KS fibre are canonically equivalent.  Away from that locus their
difference is retained by $C_{\widehat\omega}^{\mathrm{KS,diff}}$.
\end{theorem}

\begin{proof}
Only functoriality of differentiation is used to construct the upper row and the
null composite from actual flat gauge.  The lower-left object is then formed
inside the algebroid category, whose monadic forgetful functor creates the
required finite limit.  The comparison is the categorical fibre factorization.
The cofiber is formed only after applying $U_{\mathrm{src}}$ to the stable
pro-coherent module category.  No claim that differentiation of the native
formal-groupoid fibre is itself the algebraic fibre occurs in the argument.
\end{proof}

\begin{proposition}[Spectral-\'etale transport]
\label{prop:prequantum-quantomorphism-controller-etale-descent}
The three differentiated controllers, the intrinsic KS fibre, the comparison
$\gamma_{\widehat\omega}^{\mathrm{KS,diff}}$, its discrepancy, and its inverse
locus are compatible with the coherent one- and two-variable spectral-\'etale
transition equivalences of Chapters~14--15 and descend on the global locally
coherent coefficient system.
\end{proposition}

\begin{proof}
Formalized spectral leaves, partition-Lie differentiation, the zero algebroid,
finite limits in the monadic algebroid category, and coherent pro-coherent mapping
sheaves all carry the Chapter~14--15 spectral-\'etale exchange equivalences.
Every displayed construction is functorial in those maps.
\end{proof}

\section{The differentiated KS obstruction tower}
\label{sec:prequantum-differentiated-kostant-souriau}

The differentiated extension need not split.  Its obstruction theory is built
from the connecting morphism in the stable category of source pro-coherent
modules and the complete algebroid-morphism latching tower.

\begin{construction}[First Kostant--Souriau connecting class]
\label{constr:prequantum-extension-class}
The stable source-module fibre sequence determines the canonical connecting
morphism
\[
\boxed{
\mathbf k_{\widehat\omega}^{(1)}:
U_{\mathrm{src}}\mathbb A_{\widehat\omega}^{\mathrm{Ham,\,lift}}
\longrightarrow
U_{\mathrm{src}}\mathbb F_{\widehat\omega}^{\mathrm{KS}}[1].
}
\]
We call it the first brave-new Kostant--Souriau obstruction class.
\end{construction}

\subsection{Full algebroid lift realization}
\label{subsec:prequantum-full-algebroid-lift-realization}

Let
\[
b\in\operatorname{LieAlgd}^{\pi,E_\infty}_{Y/T}
\]
and let
\[
\rho:b\longrightarrow
\mathbb A_{\widehat\omega}^{\mathrm{Ham,\,lift}}
\]
be an actual spectral partition-Lie algebroid morphism.

\begin{definition}[Differentiated prequantum lift space]
\label{def:prequantum-differentiated-lift-space}
Define
\[
\boxed{
\operatorname{Lift}^{\mathrm{PreQ}}_{\widehat\omega}(\rho)
:=
\operatorname{hofib}_{\rho}
\left(
\operatorname{Map}_{\operatorname{LieAlgd}}
(b,\mathbb A_{\widehat\omega}^{\mathrm{PreQ}})
\longrightarrow
\operatorname{Map}_{\operatorname{LieAlgd}}
(b,\mathbb A_{\widehat\omega}^{\mathrm{Ham,\,lift}})
\right).
}
\]
\end{definition}

A point is an actual spectral partition-Lie algebroid lift, including every
operation compatibility.  No splitting or admissible-Hamiltonian condition is
part of the input.

\begin{construction}[Finite-weight prequantum lift tower]
\label{constr:prequantum-finite-weight-lift-tower}
Apply the Chapter~15 morphism corolla--latching construction to an underlying
anchored candidate lift
\[
h:U_{\mathrm{anch}}b
\longrightarrow
U_{\mathrm{anch}}\mathbb A_{\widehat\omega}^{\mathrm{PreQ}}
\]
whose composite with $U_{\mathrm{anch}}p_{\mathrm{KS}}$ is coherently identified
with $U_{\mathrm{anch}}\rho$.  For $s\ge0$, let
$\operatorname{Lift}^{\mathrm{PreQ}}_{\le s}(\rho)$ be the space of coherent
morphism relations through weight $s$ whose image under $p_{\mathrm{KS}}$ is the
already fixed weight-$s$ structure of $\rho$.  At weight zero this is the
anchored lifting space.  For $s\ge1$, a point through weight $s-1$ has a
canonical latching obstruction
\[
\operatorname{ob}^{\mathrm{PreQ}}_s(h):
R_s^{\mathrm{mor}}(h)
\longrightarrow
U_{\mathrm{anch}}\mathbb A_{\widehat\omega}^{\mathrm{PreQ}},
\]
and the weight-$s$ extension space is the homotopy fibre, over the prescribed
weight-$s$ filler of $\rho$, of the corresponding target extension map.
\end{construction}

\begin{theorem}[Complete differentiated KS lift obstruction theory]
\label{thm:prequantum-complete-differentiated-ks-lift-tower}
There is a canonical equivalence
\[
\boxed{
\operatorname{Lift}^{\mathrm{PreQ}}_{\widehat\omega}(\rho)
\simeq
\lim_s\operatorname{Lift}^{\mathrm{PreQ}}_{\le s}(\rho).
}
\]
The source-module shadow of the weight-zero stage is the zero locus
\[
\boxed{
\operatorname{Lift}^{\mathrm{PreQ}}_{\le0}(\rho)
\simeq
Z\left(
\mathbf k_{\widehat\omega}^{(1)}
\circ U_{\mathrm{src}}\rho
\right).
}
\]
The lift represented by a point of this zero locus automatically has the correct
anchor: if $p_{\mathrm{KS}}h=\rho$, then anchor compatibility follows from the
fact that $p_{\mathrm{KS}}$ is already a morphism over the common anchor target.
For each $s\ge1$, extension from weight $s-1$ to weight $s$ is equivalent to a
specified nullhomotopy of $\operatorname{ob}^{\mathrm{PreQ}}_s$ together with its
corolla filler.  The universal morphism-defect syzygy supplies the corresponding
Bianchi cocycle at the next weight.
\end{theorem}

\begin{proof}
The inverse-limit formula is the Chapter~15 finite-weight reconstruction theorem
for algebroid morphisms, applied in the slice over $\rho$.  At weight zero,
forget the anchor and use the stable source-module fibre sequence; its universal
property identifies a source lift with a nullhomotopy of the pulled-back
connecting class.  Since the projection is a morphism of anchored objects, the
anchor equation for that source lift follows from $p_{\mathrm{KS}}h=\rho$ and
needs no separate filler.  Higher stages are precisely the Chapter~15 morphism
corolla--latching extension problems.  Exhaustiveness of the anchor-weight
filtration reconstructs the full mapping space.
\end{proof}

\section{Prequantum representations, observables, and Atiyah--Spencer connection}
\label{sec:prequantum-local-observables}

The stable flat Cartan coefficient and the independently differentiated actual-flat
groupoid are distinct constructions.  The stable coefficient is descended before
applying relative zeroth space; conjugation of the flat gauge group is recovered
as a comparison of its infinite-loop shadow rather than used to manufacture a
spectrum-level action.

\begin{construction}[Native stable flat-adjoint Cartan coefficient]
\label{constr:prequantum-conjugation-cartan-representation}
Over $\operatorname{CartCurvReal}_{\mathcal H}(\kappa,u)$, the stable fibre
$\mathscr F_\kappa[r-1]$ is already a $J_{\mathcal H}$-coalgebra because the
forgetful functor from jet-comonad coalgebras creates finite limits.  Restrict
this Cartan transport along the fixed-object groupoid morphism
\[
I^{\mathrm{Ham,raw}}_{\widehat\omega,\bullet}
\longrightarrow
\operatorname{HamRaw}_{\omega,\bullet}(\mathcal H)
\longrightarrow
\mathcal H_\bullet.
\]
Write
\[
q_I:Y\twoheadrightarrow Q_I
:=
\bigl|I^{\mathrm{Ham,raw}}_{\widehat\omega,\bullet}\bigr|
\]
for the effective quotient.  Effective Cartan descent constructs a stable
coefficient
\[
\boxed{
\operatorname{Ad}^{\mathrm{flat,raw}}_{\widehat\omega}
\in E_{Q_I}
}
\]
whose pullback along $q_I$ is literally $\mathscr F_\kappa[r-1]$ with the
restricted stable Cartan transport.

Let
\[
c_I:
Q_{\operatorname{Ham}^{\mathrm{lift}}_{\widehat\omega}}
\longrightarrow Q_I
\]
be the quotient map induced by the formalization counit
$\operatorname{Ham}^{\mathrm{lift}}_{\widehat\omega}\to
I^{\mathrm{Ham,raw}}_{\widehat\omega}$.  Define
\[
\boxed{
\operatorname{Ad}^{\mathrm{flat}}_{\widehat\omega}
:=
c_I^*\operatorname{Ad}^{\mathrm{flat,raw}}_{\widehat\omega}
\in E_{Q_{\operatorname{Ham}^{\mathrm{lift}}_{\widehat\omega}}}.
}
\]
\end{construction}

\begin{proposition}[Comparison with kernel conjugation]
\label{prop:prequantum-flat-adjoint-conjugation-comparison}
Apply relative zeroth space to the stable transport of
$\operatorname{Ad}^{\mathrm{flat,raw}}_{\widehat\omega}$.  Under the canonical
identification
\[
\Omega_Y^\infty\mathscr F_\kappa[r-1]
=G_{\mathrm{flat}}(\widehat\omega),
\]
the resulting group transport is canonically the conjugation action of the raw
Hamiltonian-prequantum extension on its flat kernel.
\end{proposition}

\begin{proof}
A raw prequantum arrow acts on the entire Cartan curvature diagram by transport.
Since $\mathscr F_\kappa$ is the fibre of the transported morphism $\kappa$, the
induced map on the fibre is the additive difference transport between two lifts.
Under relative zeroth space this is exactly the action obtained by conjugating a
flat modification by the chosen prequantum arrow.  If two prequantum lifts project
to the same arrow of the effective image, the kernel-bitorsor theorem identifies
their difference with a flat arrow.  The flat kernel is an infinite-loop abelian
group object, so conjugation by that difference is coherently trivial.  Hence the
action descends to $I^{\mathrm{Ham,raw}}$, agreeing with the already descended
stable Cartan transport.  No promotion of an arbitrary action on
$\Omega^\infty\mathscr F_\kappa$ to a spectrum action is used.
\end{proof}

\begin{definition}[Prequantum flat-adjoint tangent representation]
\label{constr:prequantum-tangent-representation}
Apply the Chapter~15 differentiation-of-Cartan-representations functor to define
\[
\boxed{
\mathbb K_{\widehat\omega}^{\mathrm{Ad,PreQ}}
\in
\operatorname{Rep}^{\pi,E_\infty}
\bigl(\mathbb A_{\widehat\omega}^{\mathrm{Ham,\,lift}}\bigr).
}
\]
This representation is kept distinct from the independently differentiated
actual-flat controller
$\mathbb A_{\widehat\omega}^{\mathrm{Flat}}$.  No equivalence between them is
asserted by notation.
\end{definition}

\begin{definition}[Brave-new prequantum observable controller]
\label{def:prequantum-observable-controller}
Define
\[
\boxed{
\operatorname{Obs}_{\mathrm{PreQ}}(Y,\widehat\omega)
:=
\mathbb A_{\widehat\omega}^{\mathrm{PreQ}}.
}
\]
It maps to the liftable Hamiltonian controller and to the independently
differentiated quantomorphism controller.  Over the formal inverse locus
$\operatorname{KSEq}^{\mathrm{for}}_{\widehat\omega}$ the liftable Hamiltonian
groupoid is the full formal Hamiltonian groupoid.
\end{definition}

\begin{construction}[Exact labelled observable operations]
\label{constr:prequantum-exact-labelled-observable-operations}
For a finite nonempty labelled set $I$ and tangent variables $\xi_i$ in
$\mathbb A_{\widehat\omega}^{\mathrm{PreQ}}$, let
$\mathcal C_I(\alpha_\bullet)$ be the exact labelled corolla object of Chapter~15.
Define
\[
\boxed{
\operatorname{Op}^{I}_{\widehat\omega}(\xi_\bullet):
\mathcal C_I(\alpha_\bullet)
\longrightarrow
U_{\mathrm{src}}\mathbb A_{\widehat\omega}^{\mathrm{PreQ}}
}
\]
by partition-Lie corolla evaluation.
\end{construction}

\begin{theorem}[Complete higher observable coherence]
\label{thm:prequantum-complete-higher-observable-coherence}
The labelled operations satisfy the complete substitution, symmetric relabelling,
anchor--Leibniz, and morphism corolla--latching hierarchy of the Chapter~15
partition-Lie monad.  Projection to the liftable Hamiltonian controller gives the
corresponding Hamiltonian operations.  The difference of two underlying
prequantum lifts with the same Hamiltonian projection factors canonically through
$U_{\mathrm{src}}\mathbb F_{\widehat\omega}^{\mathrm{KS}}$.  Over
$\operatorname{KSDiffEq}_{\widehat\omega}$ this factorization identifies, through
the constructed inverse of $\gamma^{\mathrm{KS,diff}}$, with one through
$U_{\mathrm{src}}\mathbb A_{\widehat\omega}^{\mathrm{Flat}}$.
\end{theorem}

\begin{proof}
The observable controller is an actual spectral partition-Lie algebroid, so the
operation identities are its monad-algebra structure.  Naturality of corolla
evaluation under the KS projection gives the Hamiltonian operations.  Equal
Hamiltonian projections have source-module difference with null composite under
$U_{\mathrm{src}}p_{\mathrm{KS}}$, so the stable fibre universal property
produces the $U_{\mathrm{src}}\mathbb F^{\mathrm{KS}}$ factorization.  The
actual-flat factorization uses the inverse only on $\operatorname{KSDiffEq}$.
\end{proof}

\subsection{Canonical coefficient connection}
\label{sec:prequantum-atiyah-spencer-connection}

Let $U^\pi_{\mathbb A_{\widehat\omega}^{\mathrm{Ham,\,lift}}}$ be the tangent
enveloping monad.  For a tangent representation write $U_{\mathrm{rep}}$ for its
underlying pro-coherent coefficient.  The representation
$\mathbb K_{\widehat\omega}^{\mathrm{Ad,PreQ}}$ carries the full tangent-monad
action.

\begin{definition}[Canonical prequantum coefficient connection]
\label{def:prequantum-canonical-coefficient-connection}
Define the first-weight truncation
\[
\boxed{
\nabla_{\widehat\omega}^{\mathrm{PreQ}}:
F_{\le1}
U^\pi_{\mathbb A_{\widehat\omega}^{\mathrm{Ham,\,lift}}}
\bigl(U_{\mathrm{rep}}\mathbb K_{\widehat\omega}^{\mathrm{Ad,PreQ}}\bigr)
\longrightarrow
U_{\mathrm{rep}}\mathbb K_{\widehat\omega}^{\mathrm{Ad,PreQ}}.
}
\]
Its weight-zero restriction is the identity.  It is the restriction of the full
constructed representation action, not additional connection data.
\end{definition}

\begin{theorem}[Prequantum curvature, Bianchi, and higher flatness]
\label{thm:prequantum-curvature-bianchi-hierarchy}
The weight-two connection obstruction carries the canonical filler supplied by
the weight-two component of the full tangent action.  Hence
\[
\boxed{F_{\nabla^{\mathrm{PreQ}}}\simeq0.}
\]
The universal monad-defect syzygy gives the Bianchi coherence, and every higher
weight supplies the corresponding higher flatness filler.
\end{theorem}

\begin{proof}
Chapter~15's corolla--latching extension theorem applies to the actual
representation action.  Since the displayed connection is its first-weight
truncation, the weight-$s$ action component is the canonical filler of the
weight-$s$ latching obstruction.  Weight two is curvature; the universal
monad-defect syzygy gives Bianchi at the next stage; iteration gives the complete
hierarchy.  No independent connection, flatness datum, or Bianchi datum is
retained.
\end{proof}

\section{The field-theoretic Hamiltonian transgression and its coherence problem}
\label{sec:prequantum-differential-lift-hamiltonian-pair}

Return to the field-theoretic situation.  For $q\ge1$ put
\[
P_\omega^q:=\operatorname{Lep}^q(L)
\]
and retain the complete closed presymplectic section
$\omega_L^{(q),\mathrm{cl}}$ and its prequantization family
$\operatorname{PreQ}_{\omega,L}^{q}$.
Base-change the formal symmetry groupoid along $P_\omega^q\to Q$ and suppress the
resulting structural pullbacks from the notation.

Chapter~\ref{ch:symmetry-noether} constructs two different pieces of Noether data.
First, over its Hamiltonian realization it constructs the current identity
\[
d_V^{\mathrm{mix}}J_{\mathcal G,L}^{(q)}
\simeq
\jmath_{\mathrm{tr}}
\bigl(\iota_{\mathcal G}\omega_L^{(q)}\bigr)
\]
in the horizontal-tail coefficient.  Second, over the compatible full
Euler-nullification object it constructs the complete-horizontal Noether
transgression
\[
U_{\mathcal G,L}^{(q)}:
\mathbf1\longrightarrow
\operatorname{Trans}^{1,q}_{\mathcal G,H}(A;M)[n]
\]
whose image under the transgression-fibre inclusion is the Noether current with
its selected complete nullhomotopy.  The generic Hamiltonian primitive of the
present chapter lives in this transgression fibre.  We therefore lift the
Chapter~17 Hamiltonian identity to that fibre rather than identifying the current
with the primitive by notation.

Let $\mathcal N_{\mathcal G}^{q}(L)$ denote the common Noether--Lepage parameter
object of Chapter~\ref{ch:symmetry-noether}.  Define the common parameter
\[
P_{\mathrm{HamTr}}^q
:=
\operatorname{EulNull}_{\mathcal G}^{q}(L)
\times_{\mathcal N_{\mathcal G}^{q}(L)}
\operatorname{HamReal}_{\mathcal G,L}^{q}.
\]
Over it both $U_{\mathcal G,L}^{(q)}$ and the Chapter~17 Hamiltonian homotopy are
present.

\subsection{Local adjoint of the complete Noether transgression}
\label{subsec:prequantum-local-noether-current}

Factor the stable dependent products in symmetry degree $p$ through
$\mathcal G_p$ as in the generic local mixed-coefficient construction.  The
complete Noether transgression adjoints to a local section
\[
U_{\mathcal G,L}^{(q),\mathrm{loc}}:
\mathbf1_{\mathcal G_1}
\longrightarrow
\underline{\operatorname{Trans}}^{1,0,q}_{\mathcal G,V,H}[n]
\]
whose image under the local horizontal-tail inclusion is the local adjoint of
$J_{\mathcal G,L}^{(q)}$.  Vertical variation of the complete transgression gives
\[
d_V^{\mathrm{tr}}U_{\mathcal G,L}^{(q),\mathrm{loc}}
\]
in the same local mixed transgression coefficient which contains
$\iota_{\mathcal G}\omega_L^{(q)}$.

Let
\[
\jmath_{\mathrm{tr}}:
\mathscr T_{\mathcal G}^{1,1,q}
\longrightarrow
\mathscr F_{\mathcal G}^{1,1,q}
\]
denote the local transgression-fibre inclusion.  The Chapter~17 Hamiltonian
realization and the defining comparison from $U^{(q)}$ to the current determine a
specified path
\[
\eta_{\mathrm{Ham}}:
\jmath_{\mathrm{tr}}
\bigl(d_V^{\mathrm{tr}}U_{\mathcal G,L}^{(q),\mathrm{loc}}\bigr)
\simeq
\jmath_{\mathrm{tr}}
\bigl(\iota_{\mathcal G}\omega_L^{(q)}\bigr)
\]
in $\mathscr F_{\mathcal G}^{1,1,q}$.

\begin{definition}[Transgressive Hamiltonian comparison realization]
\label{def:prequantum-hamiltonian-transgression-comparison-realization}
Define
\[
\boxed{
\operatorname{HamCompReal}_{\mathcal G,L}^{q}
}
\]
over $P_{\mathrm{HamTr}}^q$ to be the homotopy fibre, over
$\eta_{\mathrm{Ham}}$, of the map of path objects
\[
\operatorname{Path}_{\Omega^\infty\mathscr T_{\mathcal G}^{1,1,q}}
\left(
 d_V^{\mathrm{tr}}U_{\mathcal G,L}^{(q),\mathrm{loc}},
 \iota_{\mathcal G}\omega_L^{(q)}
\right)
\longrightarrow
\operatorname{Path}_{\Omega^\infty\mathscr F_{\mathcal G}^{1,1,q}}
\left(
 \jmath_{\mathrm{tr}}d_V^{\mathrm{tr}}U_{\mathcal G,L}^{(q),\mathrm{loc}},
 \jmath_{\mathrm{tr}}\iota_{\mathcal G}\omega_L^{(q)}
\right).
\]
A point is therefore an actual coherent lift
\[
\boxed{
d_V^{\mathrm{tr}}U_{\mathcal G,L}^{(q),\mathrm{loc}}
\simeq
\iota_{\mathcal G}\omega_L^{(q)}
}
\]
whose image is the already constructed Chapter~17 Hamiltonian path.  The
realization may be empty.
\end{definition}

Over $\operatorname{HamCompReal}_{\mathcal G,L}^{q}$ the local transgression and
its lifted Hamiltonian path define the degree-one arrow section
\[
\boxed{
j^{(q)}_{\mathcal G,L,1}:
\mathcal G_1
\longrightarrow
\operatorname{HamRaw}_{\omega_L^{(q)},1}(\mathcal G).
}
\]
Degeneracy normalization identifies its restriction to identity arrows with the
zero primitive, so it extends to a morphism of reflexive graphs
\[
j^{(q)}_{\mathcal G,L,\le1}:
\mathcal G_{\le1}
\longrightarrow
\operatorname{HamRaw}_{\omega_L^{(q)},\le1}(\mathcal G).
\]
Its horizontal-tail shadow is the Chapter~17 Noether current and Hamiltonian
identity.

\section{The native Hamiltonian-Noether coherence realization}
\label{sec:prequantum-hamiltonian-noether-coherence}

Use the fixed-object internal mapping object constructed in
Construction~\ref{constr:prequantum-internal-fixed-groupoid-mapping-object}.
For raw groupoids write it as
$\underline{\operatorname{Map}}^{\mathrm{raw}}_{/G}(A,B)$; restriction to
$\Delta_{\le1}$ is the corresponding mapping object of reflexive graphs.

\begin{definition}[Hamiltonian-Noether coherence realization]
\label{def:prequantum-hamiltonian-noether-coherence-realization}
On the variational branch put
\[
P_{\mathrm{mix}}^q
:=
\operatorname{VertClassReal}_V
\bigl(\omega_L^{(q),\mathrm{cl}}\bigr).
\]
Proposition~\ref{prop:prequantum-vertical-origin-descent-canonical-point} and the
construction of the local mixed transgression give
$\operatorname{HamCompReal}_{\mathcal G,L}^{q}$ its canonical map to
$P_{\mathrm{mix}}^q$.  The Hamiltonian multiplicativity realization is already
an object over that same base.  Put
\[
B_{\mathrm{HN}}^q
:=
\operatorname{HamCompReal}_{\mathcal G,L}^{q}
\times_{P_{\mathrm{mix}}^q}
\operatorname{HamMultReal}_{\omega_L^{(q)}}(\mathcal G).
\]
Over $B_{\mathrm{HN}}^q$ define
\[
\boxed{
\operatorname{HamNoethReal}_{\mathcal G,L}^{q}
:=
\operatorname{hofib}_{j^{(q)}_{\mathcal G,L,\le1}}
\left(
\underline{\operatorname{Map}}^{\mathrm{raw}}_{/\mathcal G}
(\mathcal G,\operatorname{HamRaw}_{\omega_L^{(q)}}(\mathcal G))
\longrightarrow
\underline{\operatorname{Map}}^{\mathrm{fix}}_{/\mathcal G_{\le1}}
(\mathcal G_{\le1},
 \operatorname{HamRaw}_{\omega_L^{(q)},\le1}(\mathcal G))
\right).
}
\]
A point is exactly a complete groupoid-functor structure extending the actual
transgressive Hamiltonian Noether arrow and its unit.
\end{definition}

\begin{construction}[Native Hamiltonian-Noether coherence tower]
\label{constr:prequantum-hamiltonian-noether-coherence-tower}
Filter the preceding mapping-space fibre by the simplicial skeleta of
$\mathcal G$.  Write $\operatorname{HamNoethReal}_{\le m}$ for extension through
simplicial degree $m$.

On $\mathcal G_2$, transport the three pulled-back complete transgression
degree-zero coefficients to the long-edge coefficient.  The underlying
primitive of the degree-two matching defect is
\[
\epsilon^{\mathrm{Ham}}_{012}
:=
U_{02}-U_{01}-\operatorname{tr}_{01}U_{12}
-K_{012}(\omega_L^{(q)}).
\]
Its target and Hamiltonian-equation compatibility are defined by the selected
point of $\operatorname{HamMultReal}$; no cancellation of a mixed horizontal
coface is used.  The degree-two extension space is the homotopy fibre of the
degree-two restriction map over the boundary simplex determined by
$j^{(q)}_{\mathcal G,L,\le1}$.  Thus its first equation is
$\epsilon^{\mathrm{Ham}}_{012}\simeq0$, now interpreted inside the actual
long-edge primitive fibre of the universal target groupoid.

For $m\ge3$, restrict to $\operatorname{sk}_{m-1}\mathcal G$ and form the Reedy
matching map in degree $m$; the stage-$m$ realization is the homotopy fibre over
the already specified boundary simplex.  All target coherences in these
matching maps are those of the universal groupoid over
$\operatorname{HamMultReal}$.
\end{construction}

\begin{theorem}[Complete native Hamiltonian-Noether coherence]
\label{thm:prequantum-field-hamiltonian-groupoid-realization}
There is a canonical equivalence
\[
\boxed{
\operatorname{HamNoethReal}_{\mathcal G,L}^{q}
\simeq
\lim_{m\ge1}\operatorname{HamNoethReal}_{\le m}.
}
\]
Over this realization the universal point is a fixed-object formal-groupoid
section
\[
\boxed{
h_{\mathcal G,L}^{q}:
\mathcal G_\bullet
\longrightarrow
\operatorname{Ham}_{\omega_L^{(q)},\bullet}(\mathcal G).
}
\]
Its degree-one Hamiltonian primitive is the complete Chapter~17 transgression
$U_{\mathcal G,L}^{(q)}$ with the selected lift in
$\operatorname{HamCompReal}$.  Its horizontal-tail shadow is the Chapter~17
Noether current.  At degree two the first new condition is
$\epsilon^{\mathrm{Ham}}_{012}\simeq0$.
\end{theorem}

\begin{proof}
A map of simplicial objects is the inverse limit of its restrictions to finite
skeleta, and extension across the next skeleton is the corresponding matching
homotopy fibre.  Degree zero and the unit part of degree one are already fixed.
In degree two, compatibility with the Hamiltonian multiplication of the generic
raw groupoid is precisely
\[
U_{02}\simeq U_{01}+\operatorname{tr}_{01}U_{12}+K_{012}(\omega).
\]
This is the degree-two matching fibre, not a consequence of the mixed prism
alone.  The selected point of $\operatorname{HamMultReal}$ supplies the target's
degree-three and higher simplicial coherences, so every subsequent extension is
the ordinary Reedy matching problem.  Taking the inverse limit constructs the
raw map; the Chapter~13 formalization coreflection converts it, from the already
formal source, into the displayed formal section.
\end{proof}

\section{Liftability of the Hamiltonian Noether section and the native Prequantum Noether realization}
\label{sec:prequantum-noether-realization}

Let
\[
P_{\mathrm{HN}}^q
:=
\operatorname{PreQ}_{\omega,L}^{q}
\times_{P_\omega^q}
\operatorname{HamNoethReal}_{\mathcal G,L}^{q}.
\]
Over it there is a tautological prequantization $\widehat\omega_L^{(q)}$ and a
formal Hamiltonian section $h_{\mathcal G,L}^q$.  Local raw arrowwise liftability,
formal factorization through the liftable image, and a coherent section of the KS
extension remain separate.

\begin{definition}[Coherent liftable-Hamiltonian realization]
\label{def:prequantum-noether-liftable-hamiltonian-realization}
Define
\[
\boxed{
\operatorname{NoethHamLiftReal}_{\mathcal G,L}^{q}
:=
\operatorname{hofib}_{h_{\mathcal G,L}^{q}}
\left(
\underline{\operatorname{Map}}^{\mathrm{fix}}_{/\mathcal G}
\bigl(\mathcal G,
      \operatorname{Ham}^{\mathrm{lift}}_{\widehat\omega_L^{(q)}}(\mathcal G)\bigr)
\longrightarrow
\underline{\operatorname{Map}}^{\mathrm{fix}}_{/\mathcal G}
\bigl(\mathcal G,
      \operatorname{Ham}_{\omega_L^{(q)}}(\mathcal G)\bigr)
\right).
}
\]
A point is a complete coherent factorization of the Hamiltonian Noether section
through the formal liftable Hamiltonian image.
\end{definition}

Over the formal inverse locus
$\operatorname{KSEq}^{\mathrm{for}}_{\widehat\omega_L^{(q)}}$, the preceding
factorization space is contractible over each Hamiltonian section.  This does not
assert raw liftability on
$\operatorname{KSEq}^{\mathrm{raw}}$ unless that separate inverse locus is also
selected.

\begin{definition}[Prequantum Noether-lift realization]
\label{def:prequantum-noether-lift-realization}
Over $\operatorname{NoethHamLiftReal}_{\mathcal G,L}^{q}$ let
$h_{\mathcal G,L}^{q,\mathrm{lift}}$ be the universal liftable-Hamiltonian
section and define
\[
\boxed{
\operatorname{NoethPreQReal}_{\mathcal G,L}^{q}
:=
\operatorname{hofib}_{h_{\mathcal G,L}^{q,\mathrm{lift}}}
\left(
\underline{\operatorname{Map}}^{\mathrm{fix}}_{/\mathcal G}
\bigl(\mathcal G,
      \operatorname{HamPreQ}_{\widehat\omega_L^{(q)}}(\mathcal G)\bigr)
\longrightarrow
\underline{\operatorname{Map}}^{\mathrm{fix}}_{/\mathcal G}
\bigl(\mathcal G,
      \operatorname{Ham}^{\mathrm{lift}}_{\widehat\omega_L^{(q)}}(\mathcal G)\bigr)
\right).
}
\]
\end{definition}

A point is an actual coherent splitting of the pulled-back formal KS extension
over the Hamiltonian Noether section.  No equivariant-prequantizability condition
is selected beforehand.

\begin{proposition}[Two-stage native liftability]
\label{prop:prequantum-noether-liftability-image}
The native lift problem factors canonically as
\[
\boxed{
\text{transgressive Hamiltonian coherence}
\longrightarrow
\text{coherent factorization through }\operatorname{Ham}^{\mathrm{lift}}
\longrightarrow
\text{prequantum lift}.
}
\]
Fibre-pasting identifies the direct fibre of
\[
\operatorname{Map}(\mathcal G,\operatorname{HamPreQ})
\longrightarrow
\operatorname{Map}(\mathcal G,\operatorname{Ham})
\]
over $h_{\mathcal G,L}^q$ with the total space of the two successive lift fibres.
Even on $\operatorname{KSEq}^{\mathrm{for}}$, the final prequantum lift fibre may
be empty.
\end{proposition}

\section{The Brave-New Prequantum Noether Theorem}
\label{sec:prequantum-noether-theorem}

Work over the common parameter object on which a point of
$\operatorname{NoethPreQReal}_{\mathcal G,L}^{q}$ has been selected.

\begin{theorem}[Brave-New Prequantum Noether Theorem]
\label{thm:prequantum-noether-theorem}
There is a canonical formal-groupoid morphism
\[
\boxed{
\widetilde h_{\mathcal G,L}^{q}:
\mathcal G_\bullet
\longrightarrow
\operatorname{HamPreQ}_{\widehat\omega_L^{(q)},\bullet}(\mathcal G).
}
\]
Its Hamiltonian projection is the complete Noether transgression
$U_{\mathcal G,L}^{(q)}$ with the selected Hamiltonian lift and its quantomorphism
projection is the corresponding prequantum automorphism.  The horizontal-tail
shadow of the Hamiltonian component is the Chapter~17 current
$J_{\mathcal G,L}^{(q)}$.

On the locally coherent differentiated theorem domain, spectral partition-Lie
differentiation gives
\[
\boxed{
\widetilde\rho_{\mathrm{Noeth}}:
\mathbb A_{\mathcal G}^{\pi,E_\infty}
\longrightarrow
\mathbb A_{\widehat\omega_L^{(q)}}^{\mathrm{PreQ}}
}
\]
fitting into
\[
\begin{tikzcd}[column sep=huge,row sep=large]
& \mathbb A_{\widehat\omega_L^{(q)}}^{\mathrm{PreQ}}
  \arrow[d,"p_{\mathrm{KS}}"] \\
\mathbb A_{\mathcal G}^{\pi,E_\infty}
  \arrow[r,"\rho_{\mathrm{Ham}}"']
  \arrow[ur,"\widetilde\rho_{\mathrm{Noeth}}"] &
\mathbb A_{\widehat\omega_L^{(q)}}^{\mathrm{Ham,\,lift}}.
\end{tikzcd}
\]
The complete Hamiltonian identity is
\[
\boxed{
d_V^{\mathrm{tr}}U_{\mathcal G,L}^{(q)}
\simeq
\iota_{\mathcal G}\omega_L^{(q)},
}
\]
and its horizontal-tail image is the Chapter~17 identity
\[
d_V^{\mathrm{mix}}J_{\mathcal G,L}^{(q)}
\simeq
\jmath_{\mathrm{tr}}
\bigl(\iota_{\mathcal G}\omega_L^{(q)}\bigr).
\]
No global splitting of the integrated or differentiated KS extension is assumed; the theorem is formed over the derived realization of lifts along the Hamiltonian Noether section.
\end{theorem}

\begin{proof}
The universal point of $\operatorname{NoethPreQReal}$ is the formal groupoid lift
$\widetilde h_{\mathcal G,L}^{q}$.  Its Hamiltonian projection is the section
constructed over $\operatorname{HamNoethReal}$, whose degree-one component is the
complete transgression $U^{(q)}$ by construction.  The defining point of
$\operatorname{HamCompReal}$ supplies its Hamiltonian equation, and the
transgression-fibre inclusion identifies its horizontal-tail shadow with the
Chapter~17 current.  Differentiation is functorial on fixed-object formal-groupoid
morphisms in the locally coherent theorem range, giving the displayed algebroid
triangle.
\end{proof}

\subsection{Universal differential-refinement symmetry and compatibility}
\label{subsec:prequantum-differential-symmetry-compatibility}

Symmetry must act on the differential refinement before an underlying
Lagrangian is frozen.  We therefore construct one universal refinement family
containing the underlying Lagrangian, its differential lift, the comparison path,
and the differential Lepage realization.

Put
\[
P_M:=\operatorname{LagCls}_Q^n(M),
\qquad
P_{\widehat M}:=\operatorname{LagCls}_Q^n(\widehat M),
\]
and retain
\[
\kappa_{\mathrm{Lag}}:P_{\widehat M}\longrightarrow P_M.
\]
Let $P_M^{\Delta^1}$ be the internal path object with endpoint maps
$\operatorname{ev}_0,\operatorname{ev}_1$.

\begin{definition}[Universal Lagrangian comparison object]
\label{def:prequantum-universal-lagrangian-comparison-object}
Define
\[
\boxed{
\operatorname{LagCmp}_{\kappa_M}
:=
P_{\widehat M}
\times_{P_M,\kappa_{\mathrm{Lag}},\operatorname{ev}_0}
P_M^{\Delta^1}.
}
\]
It carries the endpoint map
\[
L_{\mathrm{univ}}:=\operatorname{ev}_1:
\operatorname{LagCmp}_{\kappa_M}\longrightarrow P_M
\]
and the tautological tuple
\[
(\widehat L_{\mathrm{univ}},L_{\mathrm{univ}},
  \eta_{\mathrm{Lag}}),
\qquad
\eta_{\mathrm{Lag}}:
\kappa_{\mathrm{Lag}}(\widehat L_{\mathrm{univ}})
\simeq L_{\mathrm{univ}}.
\]
\end{definition}

Let
\[
\operatorname{Lep}_{\widehat M,\mathrm{univ}}^q
\longrightarrow P_{\widehat M}
\]
be the Chapter~12 universal relative Lepage family for the tautological
$\widehat M$-valued Lagrangian.

\begin{definition}[Universal differential variational refinement]
\label{def:prequantum-universal-differential-variational-refinement}
Define
\[
\boxed{
\operatorname{VarRef}_{\kappa_M}^q
:=
\operatorname{LagCmp}_{\kappa_M}
\times_{P_{\widehat M}}
\operatorname{Lep}_{\widehat M,\mathrm{univ}}^q.
}
\]
Its universal point is
\[
\mathbf r_{\partial,\mathrm{univ}}^q
=
(L_{\mathrm{univ}},\widehat L_{\mathrm{univ}},
 \eta_{\mathrm{Lag}},\widehat\ell_{\mathrm{univ}}).
\]
Functoriality of the Chapter~12 Lepage family in coefficient morphisms and in
paths of Lagrangians gives an underlying Lepage realization
$\ell_{\mathrm{univ}}$ for $L_{\mathrm{univ}}$.
\end{definition}

Apply the variational curvature and prequantization construction of
Sections~\ref{sec:prequantum-lagrangians}--\ref{subsec:prequantum-variational-curvature-morphism}
to the tautological Lagrangian over $P_M$.  Write
\[
\operatorname{PreQ}_{\omega,\mathrm{univ}}^q
\]
for the resulting universal presymplectic prequantization family over the
universal $M$-valued Lepage space.  Naturality gives
\[
\boxed{
\operatorname{VarRef}_{\kappa_M}^q
\longrightarrow
\operatorname{PreQ}_{\omega,\mathrm{univ}}^q.
}
\]
For a fixed Lagrangian section $L:Q\to P_M$, pullback along the endpoint
$L_{\mathrm{univ}}=L$ identifies the resulting fibre with
$\operatorname{PreQLep}_{\kappa_M}^q(L)$ and identifies the preceding map with
$\operatorname{PreQLep}_{\kappa_M}^q(L)\to\operatorname{PreQ}_{\omega,L}^q$.
Thus the fixed-$L$ construction is recovered rather than used as the definition
of symmetry.

\begin{construction}[Complete symmetry realization of the full differential refinement]
\label{constr:prequantum-differential-lagrangian-stabilizer}
Pull the field-symmetry groupoid $\mathcal G_\bullet$ to the universal refinement
base.  In symmetry degree $p$, pull every object and every arrow in the finite
homotopy-limit diagram defining $\operatorname{VarRef}_{\kappa_M}^q$ to the
$p$th symmetry relation and form the same finite limit there.  Functoriality in
the symmetry faces and degeneracies gives an actual cosimplicial
\emph{space-valued} refinement object in the native topos; denote it by
\[
\mathbf R_{\partial,\mathrm{var}}^{\bullet,q}.
\]
Its degree-zero object carries the universal point
$\mathbf r_{\partial,\mathrm{univ}}^q$.  For $r\ge0$ define the partial symmetry
realization by the homotopy fibre of restriction to degree zero,
\[
\operatorname{DiffVarSymReal}^{q,\le r}_{\mathcal G}
:=
\operatorname{hofib}_{\mathbf r_{\partial,\mathrm{univ}}^q}
\left(
\operatorname{Tot}_{\mathcal G,\le r}
\mathbf R_{\partial,\mathrm{var}}^{\bullet,q}
\longrightarrow
\mathbf R_{\partial,\mathrm{var}}^{0,q}
\right).
\]
Define the complete realization
\[
\boxed{
\operatorname{DiffVarSymReal}^{q}_{\mathcal G}
:=
\lim_r
\operatorname{DiffVarSymReal}^{q,\le r}_{\mathcal G}.
}
\]
A point simultaneously stabilizes
$L$, $\widehat L$, the path $\eta_{\mathrm{Lag}}$, the differential Lepage
realization $\widehat\ell$, and all higher symmetry coherences.  Successive
extension from the $(r-1)$st to the $r$th partial totalization is the explicit
matching/lift fibre over the already constructed boundary simplex.  On the
stable coefficient vertices this specializes to the Chapter~17
obstruction/nullification description.  No invariance of $L$ or of its
comparison path is assumed before the realization is formed.
\end{construction}

\begin{proposition}[Full differential refinement induces prequantum symmetry]
\label{prop:prequantum-differential-symmetry-to-quantomorphism}
Over $\operatorname{DiffVarSymReal}^{q}_{\mathcal G}$, naturality of
$\operatorname{VarRef}_{\kappa_M}^q\to
\operatorname{PreQ}_{\omega,\mathrm{univ}}^q$ sends the universal complete
symmetry to a canonical fixed-object formal-groupoid morphism
\[
\boxed{
\mathcal G_\bullet
\longrightarrow
\operatorname{Quant}_{\widehat\omega_{\partial,\mathrm{univ}}^{(q)},\bullet}
(\mathcal G).
}
\]
After pullback to a fixed $L$ this is a quantomorphism of the corresponding
prequantization in $\operatorname{PreQ}_{\omega,L}^q$.
\end{proposition}

\begin{proof}
Every object in the finite diagram defining
$\operatorname{VarRef}_{\kappa_M}^q$ is functorial in the Lagrangian, coefficient
morphism, comparison path, and Lepage realization.  The variational
prequantization map is a natural transformation of that entire diagram.
Applying symmetry partial totalization to a natural transformation and then
taking the inverse limit carries a complete symmetry of the refinement to a
complete symmetry of its closed differential presymplectic class, which is
precisely a quantomorphism.
\end{proof}

Independently, the native Prequantum Noether lift supplies a quantomorphism
section after pullback to a fixed Lagrangian and to the common refinement
parameter.

\begin{definition}[Differential-refinement/Prequantum-Noether compatibility realization]
\label{def:prequantum-differential-symmetry-compatibility-realization}
On the common fixed-$L$ parameter define
\[
\boxed{
\operatorname{DiffSymPreQCompat}_{\mathcal G,L}^{q}
:=
(\operatorname{DiffVarSymReal}^{q}_{\mathcal G})_L
\times_{
\underline{\operatorname{Map}}^{\mathrm{fix}}_{/\mathcal G}
(\mathcal G,
 \operatorname{Quant}_{\widehat\omega_L^{(q)}})}
\operatorname{NoethPreQReal}_{\mathcal G,L}^{q}.
}
\]
A point is a complete symmetry of the full differential variational refinement,
a complete Prequantum Noether lift, and a specified coherent identification of
their quantomorphism sections.
\end{definition}

\begin{proposition}[Differential refinement realizes the native lift on the compatibility locus]
\label{prop:prequantum-differential-symmetry-produces-native-lift}
Over $\operatorname{DiffSymPreQCompat}_{\mathcal G,L}^{q}$, the quantomorphism
induced by the universal differential-refinement symmetry is canonically the
quantomorphism component of the native Prequantum Noether lift.  Its curvature
is the curvature symmetry constructed from the complete Noether transgression.
\end{proposition}

\section{The complete differentiated Prequantum Noether obstruction tower}
\label{sec:prequantum-noether-obstruction-tower}

The differentiated lifting problem is formed only after the Hamiltonian Noether
section has coherently factored through
$\operatorname{Ham}^{\mathrm{lift}}$.  Over
$\operatorname{NoethHamLiftReal}_{\mathcal G,L}^{q}$ differentiate that section to
obtain
\[
\rho_{\mathrm{Ham}}:
\mathbb A_{\mathcal G}^{\pi,E_\infty}
\longrightarrow
\mathbb A_{\widehat\omega_L^{(q)}}^{\mathrm{Ham,\,lift}}.
\]

\begin{definition}[Differentiated Prequantum Noether lift space]
\label{def:prequantum-differentiated-noether-lift-space}
Define
\[
\boxed{
\operatorname{NoethLift}^{\mathrm{diff},q}_{\mathcal G,L}
:=
\operatorname{Lift}^{\mathrm{PreQ}}_{\widehat\omega_L^{(q)}}
(\rho_{\mathrm{Ham}}).
}
\]
\end{definition}

\begin{theorem}[First and higher Prequantum Noether obstructions]
\label{thm:prequantum-first-higher-noether-obstructions}
The differentiated lift space is the inverse limit of its finite-weight
morphism-corolla--latching tower.  Its weight-zero stage is the zero locus
\[
\boxed{
\operatorname{NoethLift}^{\mathrm{diff},q}_{\le0}
\simeq
Z\left(
\mathbf k_{\widehat\omega_L^{(q)}}^{(1)}
\circ U_{\mathrm{src}}\rho_{\mathrm{Ham}}
\right).
}
\]
For every $s\ge1$, a partial Prequantum Noether lift through weight $s-1$ has a
canonical obstruction
\[
\operatorname{ob}^{\mathrm{NoethPreQ}}_s:
R_s^{\mathrm{mor}}
\longrightarrow
U_{\mathrm{anch}}\mathbb A_{\widehat\omega_L^{(q)}}^{\mathrm{PreQ}},
\]
and extension through weight $s$ is equivalent to a specified nullhomotopy of
that obstruction together with the prescribed corolla filler over the Hamiltonian
morphism.  The universal morphism-defect syzygy supplies the next Bianchi
cocycle.  After pullback to
$\operatorname{NoethPreQReal}_{\mathcal G,L}^{q}$, the differentiated native
Prequantum Noether lift supplies canonical fillers at every stage.
\end{theorem}

\begin{proof}
Apply
Theorem~\ref{thm:prequantum-complete-differentiated-ks-lift-tower} with
$b=\mathbb A_{\mathcal G}^{\pi,E_\infty}$ and
$\rho=\rho_{\mathrm{Ham}}$.  The first stage is the stable lifting problem for
the underlying KS fibre sequence, hence the zero locus of the pulled-back
connecting class.  Higher stages are exactly the Chapter~15 morphism
corolla--latching extensions.  Over
$\operatorname{NoethPreQReal}_{\mathcal G,L}^{q}$,
Theorem~\ref{thm:prequantum-noether-theorem} gives an actual algebroid morphism
lifting $\rho_{\mathrm{Ham}}$, so its finite-weight restrictions provide the
required fillers.
\end{proof}

Thus the first connecting class is only the first differentiated obstruction.
The native groupoid-coherence tower of
Construction~\ref{constr:prequantum-hamiltonian-noether-coherence-tower} and the
partition-Lie morphism-latching tower above are distinct constructions.

\section{Higher Noether operations and prequantum observables}
\label{sec:prequantum-higher-noether-operations}

\begin{construction}[Exact labelled prequantum Noether class]
\label{constr:prequantum-exact-labelled-noether-class}
For a finite nonempty labelled set $I$ and a labelled tangent tuple in
$\mathbb A_{\mathcal G}^{\pi,E_\infty}$, let
\[
\mu^I_{\mathcal G,\xi}:
\mathcal C_I(\alpha_\bullet)
\longrightarrow
U_{\mathrm{src}}\mathbb A_{\mathcal G}^{\pi,E_\infty}
\]
be the Chapter~15 corolla evaluation.  Over
$\operatorname{NoethPreQReal}_{\mathcal G,L}^{q}$ define
\[
\boxed{
\operatorname{Tay}^{I,q}_{\mathrm{PreQ},L}
:=
U_{\mathrm{src}}\widetilde\rho_{\mathrm{Noeth}}
\circ
\mu^I_{\mathcal G,\xi}:
\mathcal C_I(\alpha_\bullet)
\longrightarrow
U_{\mathrm{src}}\mathbb A_{\widehat\omega_L^{(q)}}^{\mathrm{PreQ}}.
}
\]
\end{construction}

\begin{theorem}[Compatibility with the differentiated Noether hierarchy]
\label{thm:prequantum-compatibility-noether-hierarchy}
Projection of
$\operatorname{Tay}^{I,q}_{\mathrm{PreQ},L}$ to the liftable Hamiltonian controller
is the corresponding exact labelled Hamiltonian Noether class.  The difference of
two prequantum lifts with the same Hamiltonian projection factors canonically
through the intrinsic KS fibre.  Over
$\operatorname{KSDiffEq}_{\widehat\omega_L^{(q)}}$ the constructed equivalence
\[
\mathbb A_{\widehat\omega_L^{(q)}}^{\mathrm{Flat}}
\simeq
\mathbb F_{\widehat\omega_L^{(q)}}^{\mathrm{KS}}
\]
identifies this factorization with one through the differentiated actual-flat
controller.  The complete family satisfies the Chapter~15 morphism
corolla--latching hierarchy.
\end{theorem}

\begin{proof}
Apply exact labelled corolla evaluation to the commutative triangle of
Theorem~\ref{thm:prequantum-noether-theorem}.  Naturality gives the Hamiltonian
projection.  A difference with zero Hamiltonian projection has null composite
under $U_{\mathrm{src}}p_{\mathrm{KS}}$ and therefore factors through
$U_{\mathrm{src}}\mathbb F^{\mathrm{KS}}$ by the stable fibre universal property.  On
$\operatorname{KSDiffEq}$ use the constructed inverse of
$\gamma^{\mathrm{KS,diff}}$.  Full latching coherence is functorial in the actual
algebroid morphism $\widetilde\rho_{\mathrm{Noeth}}$.
\end{proof}

\section{Critical restriction over the universal prequantization parameter}
\label{sec:prequantum-critical-restriction}

First retain the universal Lepage parameter independently of any selected
prequantization:
\[
P_\omega^q:=\operatorname{Lep}^q(L),
\qquad
C_\omega^q:=C_L\times_QP_\omega^q,
\qquad
R_\omega^q:=(C_L//\mathcal G)\times_QP_\omega^q.
\]
Let
\[
\mathcal K_{\omega,\bullet}
:=
\mathcal K_{\mathcal G,L,P_\omega^q,\bullet}
\rightrightarrows C_\omega^q
\]
be the parameterized critical formal groupoid and
$q_\omega:C_\omega^q\twoheadrightarrow R_\omega^q$ its effective quotient.

The full prequantization parameter is
\[
\widehat P_\omega^q
:=
\operatorname{PreQ}_{\omega,L}^{q}.
\]
Base change the preceding critical diagram along
$\widehat P_\omega^q\to P_\omega^q$ and put
\[
C_{\widehat\omega}^q
:=
C_L\times_Q\widehat P_\omega^q,
\qquad
R_{\widehat\omega}^q
:=
(C_L//\mathcal G)\times_Q\widehat P_\omega^q,
\]
\[
\mathcal K_{\widehat\omega,\bullet}
:=
\mathcal K_{\omega,\bullet}\times_{P_\omega^q}\widehat P_\omega^q,
\qquad
q_{\mathrm{crit},\widehat\omega}:
C_{\widehat\omega}^q\twoheadrightarrow R_{\widehat\omega}^q.
\]
Thus no prequantization parameter is discarded when passing to the critical
object.

Pull the field-theoretic curvature diagram
\[
\mathscr D_{L,q}
\xrightarrow{\kappa_{\mathrm{tr}}^{(q)}}
\mathscr U_{L,q}
\xleftarrow{u_{L,q}}
\mathscr C_{L,q}
\]
first to $C_\omega^q$ and then, for the selected universal prequantization, to
$C_{\widehat\omega}^q$.

\begin{construction}[Raw critical prequantum restriction]
\label{constr:prequantum-raw-critical-restriction}
On $C_{\widehat\omega}^q$ write
\[
\mathscr D_{\mathrm{crit}}^{(q)},
\qquad
\mathscr U_{\mathrm{crit}}^{(q)},
\qquad
\mathscr C_{\mathrm{crit}}^{(q)}
\]
for the pulled-back coefficients and
\[
\kappa_{\mathrm{crit}}^{(q)}:
\mathscr D_{\mathrm{crit}}^{(q)}\to\mathscr U_{\mathrm{crit}}^{(q)},
\qquad
u_{\mathrm{crit}}^{(q)}:
\mathscr C_{\mathrm{crit}}^{(q)}\to\mathscr U_{\mathrm{crit}}^{(q)}
\]
for the pulled-back comparison maps.  Put
\[
\mathscr P_{\mathrm{crit}}^{(q)}
:=
\mathscr D_{\mathrm{crit}}^{(q)}
\times_{\mathscr U_{\mathrm{crit}}^{(q)}}
\mathscr C_{\mathrm{crit}}^{(q)}.
\]
The tautological prequantization restricts to a section
\[
\boxed{
\widehat\omega_{\mathrm{crit}}^{(q)}:
\mathbf1_{C_{\widehat\omega}^q}
\longrightarrow
\mathscr P_{\mathrm{crit}}^{(q)}.
}
\]
Its closed curvature is denoted
$\omega_{\mathrm{crit}}^{(q),\mathrm{cl}}$.
\end{construction}

\begin{proposition}[Curvature of critical restriction]
\label{prop:prequantum-curvature-critical-restriction}
The projection
$\mathscr P_{\mathrm{crit}}^{(q)}\to\mathscr C_{\mathrm{crit}}^{(q)}$
carries $\widehat\omega_{\mathrm{crit}}^{(q)}$ to
$\omega_{\mathrm{crit}}^{(q),\mathrm{cl}}$, whose underlying transgression is the
critical presymplectic transgression of
Chapter~\ref{ch:symmetry-noether}.
\end{proposition}

\begin{proof}
Everything is obtained by pullback from the defining prequantization square, and
the underlying presymplectic section commutes with critical restriction by
functoriality of the variational transgression.
\end{proof}

\section{Descent of the complete prequantum curvature diagram}
\label{sec:prequantum-curvature-diagram-descent}

The coefficient-level descent is constructed before selecting a prequantization.
The critical mixed-grid descent theorem of
Chapter~\ref{ch:symmetry-noether} is coefficient-functorial, so it applies to
$\widehat M$, to $M$, and to the coefficient morphism
$\kappa_M:\widehat M\to M$.

\begin{theorem}[Coefficient-level descent of the prequantum curvature diagram]
\label{thm:prequantum-curvature-diagram-descent}
There are canonically constructed reduced stable coefficients on $R_\omega^q$
\[
\mathscr D_{\mathrm{red},0}^{(q)},
\qquad
\mathscr U_{\mathrm{red},0}^{(q)},
\]
and a natural morphism
\[
\kappa_{\mathrm{red},0}^{(q)}:
\mathscr D_{\mathrm{red},0}^{(q)}
\longrightarrow
\mathscr U_{\mathrm{red},0}^{(q)}
\]
whose pullbacks are the symmetry-\v{C}ech totalizations of the corresponding
critical differential and underlying transgression coefficients and of their
curvature map.

Let
\[
B_\omega^q
:=
\operatorname{PreSymDesc}_{\omega,\mathcal G}^{q}(L)
\]
and put
\[
R_{\mathrm{curv}}^q
:=
(C_L//\mathcal G)\times_QB_\omega^q.
\]
After pullback from $R_\omega^q$ to $R_{\mathrm{curv}}^q$, there is in addition a
reduced complete closed coefficient
\[
\mathscr C_{\mathrm{red}}^{(q)}
\]
with underlying map
\[
u_{\mathrm{red}}^{(q)}:
\mathscr C_{\mathrm{red}}^{(q)}
\longrightarrow
\mathscr U_{\mathrm{red}}^{(q)},
\]
where
\[
\mathscr D_{\mathrm{red}}^{(q)}
:=
\mathscr D_{\mathrm{red},0}^{(q)}|_{R_{\mathrm{curv}}^q},
\qquad
\mathscr U_{\mathrm{red}}^{(q)}
:=
\mathscr U_{\mathrm{red},0}^{(q)}|_{R_{\mathrm{curv}}^q},
\]
and a reduced closed section
\[
\omega_{L,\mathrm{red}}^{(q),\mathrm{cl}}:
\mathbf1_{R_{\mathrm{curv}}^q}
\longrightarrow
\mathscr C_{\mathrm{red}}^{(q)}
\]
whose pullback is the complete closed critical presymplectic descent point.
\end{theorem}

\begin{proof}
For every stable coefficient pulled from $Q$, Chapter~\ref{ch:symmetry-noether}
identifies the symmetry direction of the critical mixed grid with the \v{C}ech
nerve of the corresponding reduced maximal infinitesimal relation and identifies
stable cochains on the reduced relation with the symmetry totalization of the
critical mixed cochains.  Apply that theorem to $\widehat M$ and $M$.
Normalization, horizontal partial and complete totalization, horizontal tails,
and transgression fibres commute with the resulting equivalence by Fubini for
limits.  Naturality in coefficient morphisms carries $\kappa_M$ to
$\kappa_{\mathrm{red},0}^{(q)}$.

For the complete closed coefficient, leave the vertical degree unnormalized before
totalizing the symmetry direction.  Fubini commutes the symmetry descent with the
complete vertical filtration.  The complete presymplectic descent point defining
$B_\omega^q$ is precisely the compatible point whose descent is
$\omega_{L,\mathrm{red}}^{(q),\mathrm{cl}}$.  Naturality supplies
$u_{\mathrm{red}}^{(q)}$.
\end{proof}

Thus on $R_{\mathrm{curv}}^q$ there is a universal reduced curvature diagram
\[
\boxed{
\mathscr D_{\mathrm{red}}^{(q)}
\xrightarrow{\kappa_{\mathrm{red}}^{(q)}}
\mathscr U_{\mathrm{red}}^{(q)}
\xleftarrow{u_{\mathrm{red}}^{(q)}}
\mathscr C_{\mathrm{red}}^{(q)},
}
\]
where $\kappa_{\mathrm{red}}^{(q)}$ is the pullback of
$\kappa_{\mathrm{red},0}^{(q)}$.  Put
\[
\mathscr P_{\mathrm{red}}^{(q)}
:=
\mathscr D_{\mathrm{red}}^{(q)}
\times_{\mathscr U_{\mathrm{red}}^{(q)}}
\mathscr C_{\mathrm{red}}^{(q)}.
\]

For the selected universal prequantization, form the common base
\[
P_{\mathrm{curv}}^q
:=
\widehat P_\omega^q
\times_{P_\omega^q}
B_\omega^q.
\]
All critical and reduced coefficients below are understood after pullback to this
base when a selected prequantization is present.

\section{Complete prequantum Cartan and \v{C}ech descent}
\label{sec:prequantum-complete-cech-descent}

We first record the Cartan realization of a stable coefficient--section pair and
then show that complete descent supplies it, as in the presymplectic construction
of Chapter~\ref{ch:symmetry-noether}.

\begin{definition}[Prequantum Cartan-lift realization]
\label{def:prequantum-cartan-lift-realization}
For a test $T\to P_{\mathrm{curv}}^q$, let
$J_{\mathcal K,T}$ be the stable jet comonad of the base-changed critical formal
groupoid.  Let $\mathsf{Pair}^{\mathrm{PreQ}}_{\mathrm{Cart}}(T)$ be the maximal
$\infty$-groupoid of pairs consisting of a
$J_{\mathcal K,T}$-coalgebra structure on the pulled-back coefficient
$\mathscr P_{\mathrm{crit}}^{(q)}$ and a section
$\mathbf1\to\mathscr P_{\mathrm{crit}}^{(q)}$ which is a coalgebra morphism.  Let
$\mathsf{Pair}^{\mathrm{PreQ}}_{\mathrm{und}}(T)$ be the maximal subgroupoid of
underlying coefficient--section pairs.  Define
\[
\boxed{
\operatorname{CartLift}^{\partial,q}_{\omega,\mathcal G}(L)
}
\]
to be the internal homotopy fibre of
$\mathsf{Pair}^{\mathrm{PreQ}}_{\mathrm{Cart}}\to
\mathsf{Pair}^{\mathrm{PreQ}}_{\mathrm{und}}$ over
$(\mathscr P_{\mathrm{crit}}^{(q)},
\widehat\omega_{\mathrm{crit}}^{(q)})$.
\end{definition}

\begin{construction}[Internal cosimplicial prequantum section object]
\label{constr:prequantum-cosimplicial-section-object}
For every $p\ge0$, let
\[
g_p:\mathcal K_{\widehat\omega,p}\longrightarrow P_{\mathrm{curv}}^q
\]
be the structural map and let
$\mathscr P_{\mathrm{crit},p}^{(q)}$ be the pulled-back prequantum coefficient on
$\mathcal K_{\widehat\omega,p}$.  Define the object of the slice over
$P_{\mathrm{curv}}^q$
\[
\boxed{
\mathscr S_{\mathrm{PreQ}}^{p,q}
:=
\Omega_{P_{\mathrm{curv}}^q}^\infty
\Pi_{g_p}\mathscr P_{\mathrm{crit},p}^{(q)}.
}
\]
Right Beck--Chevalley for the simplicial maps of
$\mathcal K_{\widehat\omega,\bullet}$ makes
$\mathscr S_{\mathrm{PreQ}}^{\bullet,q}$ a cosimplicial object in the single
slice $(\mathsf{XSp})_{/P_{\mathrm{curv}}^q}$.  The selected critical
prequantization gives a section
\[
P_{\mathrm{curv}}^q\longrightarrow
\mathscr S_{\mathrm{PreQ}}^{0,q}.
\]
Only after this internal object is constructed may one take global points; no
external global-section space is used as the base of an internal pullback.
\end{construction}

\begin{definition}[Complete prequantum descent realization]
\label{def:prequantum-complete-descent-realization}
Define the homotopy pullback in
$(\mathsf{XSp})_{/P_{\mathrm{curv}}^q}$
\[
\boxed{
\operatorname{PreQDesc}_{\omega,\mathcal G}^{q}(L)
:=
P_{\mathrm{curv}}^q
\times_{\mathscr S_{\mathrm{PreQ}}^{0,q}}
\operatorname{Tot}_{p}
\mathscr S_{\mathrm{PreQ}}^{p,q}.
}
\]
A point is precisely complete internal symmetry-\v{C}ech descent of the selected
critical closed differential prequantization, with all higher totalization data.
\end{definition}

\begin{theorem}[Complete descent supplies the prequantum Cartan pair]
\label{thm:prequantum-complete-descent-supplies-cartan-pair}
There is a canonical morphism
\[
\boxed{
\operatorname{PreQDesc}_{\omega,\mathcal G}^{q}(L)
\longrightarrow
\operatorname{CartLift}^{\partial,q}_{\omega,\mathcal G}(L).
}
\]
Curvature carries it to the Chapter~\ref{ch:symmetry-noether} Cartan lift of the
critical presymplectic pair.
\end{theorem}

\begin{proof}
Theorem~\ref{thm:prequantum-curvature-diagram-descent} constructs the reduced
coefficient $\mathscr P_{\mathrm{red}}^{(q)}$.  Effective descent identifies its
pullback with the totalized critical coefficient and equips that pullback with the
canonical jet-comonad coalgebra structure.  A point of
$\operatorname{PreQDesc}$ is a section in the same totalization, so full
faithfulness of effective descent makes it a coalgebra morphism from the tensor
unit.  This is a point of the Cartan-lift fibre.  Projection to the closed
coefficient is exactly the corresponding presymplectic construction.
\end{proof}

\begin{theorem}[Curvature of complete prequantum descent]
\label{thm:prequantum-curvature-complete-descent}
There is a canonical morphism
\[
\boxed{
\operatorname{PreQDesc}_{\omega,\mathcal G}^{q}(L)
\longrightarrow
\operatorname{PreSymDesc}_{\omega,\mathcal G}^{q}(L)
}
\]
which sends complete differential prequantum descent to the complete closed
presymplectic descent datum.
\end{theorem}

\begin{proof}
Projection
$\mathscr P_{\mathrm{crit}}^{(q)}\to\mathscr C_{\mathrm{crit}}^{(q)}$
is a morphism of cosimplicial stable section objects.  Totalization and the
homotopy pullback defining $\operatorname{PreQDesc}$ preserve it.  The image of
the selected degree-zero section is the critical closed presymplectic section, so
the resulting point is exactly the defining complete presymplectic descent datum.
\end{proof}

\section{The reduced prequantum field object}
\label{sec:prequantum-reduced-field-object}

Put
\[
B_{\widehat\omega}^q
:=
\operatorname{PreQDesc}_{\omega,\mathcal G}^{q}(L)
\]
and base-change the reduced critical quotient and curvature diagram to this object.

\begin{theorem}[Effective descent of the selected prequantum section]
\label{thm:prequantum-effective-differential-descent-pair}
The universal point of $B_{\widehat\omega}^q$ determines a unique reduced closed
differential section
\[
\boxed{
\widehat\omega_{L,\mathrm{red}}^{(q)}:
\mathbf1
\longrightarrow
\mathscr P_{\mathrm{red}}^{(q)}
}
\]
on
\[
(C_L//\mathcal G)\times_QB_{\widehat\omega}^q
\]
whose pullback is
$\widehat\omega_{\mathrm{crit}}^{(q)}$ with the selected complete descent datum.
\end{theorem}

\begin{proof}
The reduced coefficient is already constructed by
Theorem~\ref{thm:prequantum-curvature-diagram-descent}.  Effective descent for
parameterized spectra identifies its mapping spectrum from the tensor unit with
the totalization of the corresponding mapping spectra on the critical \v{C}ech
nerve.  The universal point of $\operatorname{PreQDesc}$ is exactly a point of
that totalization restricting to the selected critical section.  The inverse
descent equivalence reconstructs the unique reduced section.
\end{proof}

\begin{theorem}[Curvature of the reduced prequantization]
\label{thm:prequantum-curvature-reduced-prequantization}
The reduced closed differential section has curvature the reduced complete
presymplectic transgression:
\[
\boxed{
\operatorname{curv}
\bigl(\widehat\omega_{L,\mathrm{red}}^{(q)}\bigr)
\simeq
\omega_{L,\mathrm{red}}^{(q),\mathrm{cl}}.
}
\]
After applying the underlying-transgression morphism this gives
\[
\boxed{
 u_{\mathrm{red}}^{(q)}
\operatorname{curv}
\bigl(\widehat\omega_{L,\mathrm{red}}^{(q)}\bigr)
\simeq
\omega_{L,\mathrm{red}}^{(q)}.
}
\]
The equivalence retains the complete vertical closedness and symmetry-descent
coherence selected before quotienting.
\end{theorem}

\begin{proof}
Projection from the closed differential coefficient to the closed curvature
coefficient is a morphism of the effective descent diagrams.  Applying it to the
universal prequantum descent point gives the complete presymplectic descent point.
Theorem~\ref{thm:prequantum-curvature-diagram-descent} identifies its descended
section with $\omega_{L,\mathrm{red}}^{(q),\mathrm{cl}}$.  Full faithfulness of
effective descent on mapping spectra supplies the displayed equivalence.
\end{proof}

\section{Fixed-pair prequantum descent}
\label{sec:prequantum-fixed-pair-descent}

The preceding construction yields a precise fixed-pair descent statement.

\begin{theorem}[Fixed critical prequantization descent]
\label{thm:prequantum-fixed-critical-prequantization-descent}
Let $\widehat\omega_{\mathrm{crit}}$ be a selected critical prequantization over an
effective critical quotient $q:C\twoheadrightarrow R$, with descended closed
differential coefficient $\mathscr P_R$.  There is a canonical equivalence
\[
\boxed{
\left\{
\begin{array}{c}
\text{complete \v{C}ech descent data on}\;
\widehat\omega_{\mathrm{crit}}
\end{array}
\right\}
\simeq
\left\{
\begin{array}{c}
(\widehat\omega_R,\epsilon),\\
\widehat\omega_R:\mathbf1_R\to\mathscr P_R,\\
\epsilon:q^*\widehat\omega_R\simeq
\widehat\omega_{\mathrm{crit}}
\end{array}
\right\}.
}
\]
The pullback identification $\epsilon$ is part of the descended object and is not
forgotten.
\end{theorem}

\begin{proof}
Effective descent gives
\[
\Gamma_R^{\mathrm{st}}\mathscr P_R
\simeq
\operatorname{Tot}_{p}
\Gamma_{C_p}^{\mathrm{st}}\mathscr P_{C_p}.
\]
Taking relative zeroth spaces and then the homotopy fibre over the degree-zero
section $\widehat\omega_{\mathrm{crit}}$ identifies a point of the right-hand side
with a section on $R$ together with the specified equivalence between its pullback
and the selected critical section.  This is exactly the complete descent datum.
\end{proof}

\section{Prequantization commutes with reduction}
\label{sec:prequantum-commutes-reduction}

The theorem deserving this name concerns the entire prequantization moduli problem,
not a single already selected prequantization.

Over $B_\omega^q=\operatorname{PreSymDesc}_{\omega,\mathcal G}^{q}(L)$, put
\[
C_{\mathrm{curv}}^q:=C_L\times_QB_\omega^q
\]
and pull the universal critical differential and underlying coefficients from
$C_\omega^q$, together with the complete closed critical coefficient selected by
presymplectic descent, to $C_{\mathrm{curv}}^q$.  Denote the resulting curvature
diagram by
\[
\mathscr D_{\mathrm{crit,univ}}^{(q)}
\xrightarrow{\kappa_{\mathrm{crit,univ}}^{(q)}}
\mathscr U_{\mathrm{crit,univ}}^{(q)}
\xleftarrow{u_{\mathrm{crit,univ}}^{(q)}}
\mathscr C_{\mathrm{crit,univ}}^{(q)}
\]
and put
\[
\mathscr P_{\mathrm{crit,univ}}^{(q)}
:=
\mathscr D_{\mathrm{crit,univ}}^{(q)}
\times_{\mathscr U_{\mathrm{crit,univ}}^{(q)}}
\mathscr C_{\mathrm{crit,univ}}^{(q)}.
\]
The closed curvature section is the pulled-back
$\omega_{\mathrm{crit}}^{(q),\mathrm{cl}}$ supplied by complete presymplectic
descent.

Define the reduced prequantization moduli
\[
\boxed{
\operatorname{PreQ}_{\mathrm{red}}^q
:=
\operatorname{PreQ}_{\kappa_{\mathrm{red}}^{(q)},u_{\mathrm{red}}^{(q)}}
\bigl(\omega_{L,\mathrm{red}}^{(q),\mathrm{cl}}\bigr).
}
\]
Write
\[
q_{\mathrm{curv}}:
C_{\mathrm{curv}}^q\twoheadrightarrow R_{\mathrm{curv}}^q
\]
for the base-changed effective critical quotient.  On its \v{C}ech nerve let
$\operatorname{PreQ}_{\mathrm{crit},p}^q$ denote the prequantization fibre of the
pulled-back universal critical curvature diagram over the pulled-back closed
curvature section.

By effective descent there is an equivalence of slices
\[
(\mathsf{XSp})_{/R_{\mathrm{curv}}^q}
\simeq
\operatorname{Tot}_{[p]\in\Delta}
(\mathsf{XSp})_{/C_{\mathrm{curv},p}^q}.
\]
For a compatible object $Z_\bullet$ of the right-hand totalization, write
\[
\operatorname{DescTot}_{q_{\mathrm{curv}}}(Z_\bullet)
\]
for its image under the inverse effective-descent equivalence.  This notation
keeps distinct the categorical descent totalization from a limit formed inside a
single slice.

\begin{theorem}[Brave-new prequantization commutes with reduction]
\label{thm:prequantum-commutes-with-reduction}
There is a canonical equivalence in the slice over
$R_{\mathrm{curv}}^q$, natural over $B_\omega^q$,
\[
\boxed{
\operatorname{DescTot}_{q_{\mathrm{curv}}}
\bigl(\operatorname{PreQ}_{\mathrm{crit},\bullet}^q\bigr)
\simeq
\operatorname{PreQ}_{\mathrm{red}}^q.
}
\]
Under this equivalence, the subspace obtained by fixing a critical
prequantization $\widehat\omega_{\mathrm{crit}}$ is exactly the fixed-pair descent
space of
Theorem~\ref{thm:prequantum-fixed-critical-prequantization-descent}.
\end{theorem}

\begin{proof}
In every \v{C}ech degree $p$ the prequantization object is the homotopy pullback
\[
\Omega^\infty\mathscr P_{\mathrm{crit,univ},p}^{(q)}
\times_{\Omega^\infty\mathscr C_{\mathrm{crit,univ},p}^{(q)}}
\{\omega_{\mathrm{crit},p}^{(q),\mathrm{cl}}\}
\]
in the slice over $C_{\mathrm{curv},p}^q$.  The effective-descent equivalence is an
equivalence of $\infty$-categories and therefore preserves finite limits.
Consequently it carries this compatible family of pullback squares to the
pullback square
\[
\Omega^\infty\mathscr P_{\mathrm{red}}^{(q)}
\times_{\Omega^\infty\mathscr C_{\mathrm{red}}^{(q)}}
\{\omega_{L,\mathrm{red}}^{(q),\mathrm{cl}}\}.
\]
The latter is exactly
$\operatorname{PreQ}_{\mathrm{red}}^q$ by definition.  Taking the homotopy fibre
over a fixed degree-zero critical prequantization and retaining its pullback
identification gives the fixed-pair descent theorem.
\end{proof}

Construction~\ref{constr:prequantum-safronov-quantization-handoff} forms the
exact comparison-and-polarization input on which this reduced prequantization is
transported to shifted geometric quantization.  On the Hamiltonian $G$-scheme
domain of \cite[Proposition~4.13]{Safronov2020}, the resulting quantum reduction
is the invariant section object $\Gamma(B,\mathcal L)^G$.

\section{Reduced Prequantum Noether action}
\label{sec:prequantum-reduced-noether-automorphisms}

The critical formal symmetry groupoid has nontrivial formal isotropy.  Therefore
reduction does not canonically produce one unparameterized automorphism.  The
correct reduced object is the family of prequantum automorphisms parameterized
by the formal adjoint group, which Chapter~13 constructs as the contractibly
unique effective descent of formal isotropy.

\begin{definition}[Reduced prequantum Noether compatibility realization]
\label{def:prequantum-reduced-noether-compatibility-realization}
Write
\[
R_{\mathrm{desc}}:=\operatorname{PreQDesc}_{\omega,\mathcal G}^{q}(L),
\qquad
R_{\mathrm{Noeth}}:=\operatorname{NoethPreQReal}_{\mathcal G,L}^{q},
\]
and denote by
$R_{\mathrm{curr}}:=\operatorname{CurrCartDesc}_{\mathcal G,L}^{q}$ the
Chapter~\ref{ch:symmetry-noether} current Cartan-diagram realization which
descends the complete Noether transgression, its current shadow, and their
constructed nullhomotopies.

Pull $R_{\mathrm{Noeth}}$ and $R_{\mathrm{curr}}$ to
$P_{\mathrm{curv}}^q$ along their constructed critical-restriction maps; in
this definition the same symbols denote those base changes.  All objects below
therefore lie in the single slice $(\mathsf{XSp})_{/P_{\mathrm{curv}}^q}$.

Theorem~\ref{thm:prequantum-complete-descent-supplies-cartan-pair} gives
\[
d_{\mathrm{Cart}}:R_{\mathrm{desc}}
\longrightarrow
\operatorname{CartLift}^{\partial,q}_{\omega,\mathcal G}(L).
\]
The universal point of $R_{\mathrm{Noeth}}$ is a fixed-object formal-groupoid
morphism acting on the critical prequantum coefficient and its selected
section.  Restriction to the critical groupoid followed by the Chapter~13
formal-groupoid/jet-comonad equivalence gives the second actual map
\[
n_{\mathrm{Cart}}:R_{\mathrm{Noeth}}
\longrightarrow
\operatorname{CartLift}^{\partial,q}_{\omega,\mathcal G}(L).
\]

It remains to identify the complete transgression and current shadows with all
their homotopies.  Write
\[
j_{\mathrm{tr}}:A_{\mathrm{tr}}\to A_{\mathrm{tail}},
\qquad
d_V^{\mathrm{tr}}:A_{\mathrm{tr}}\to C,
\qquad
d_V^{\mathrm{mix}}:A_{\mathrm{tail}}\to C_{\mathrm{tail}},
\qquad
j_{\mathrm{cl}}:C\to C_{\mathrm{tail}}
\]
for the complete-to-tail and vertical maps constructed in the mixed
Noether diagram.  Their naturality supplies a specified square
$d_V^{\mathrm{mix}}j_{\mathrm{tr}}\simeq
j_{\mathrm{cl}}d_V^{\mathrm{tr}}$.  Define
$\operatorname{CurrNoethBase}_{\mathcal G,L}^q$ as the finite homotopy limit
classifying tuples
\[
(U,J,h_{\mathrm{tail}},h_{\mathrm{Ham}},h_{\mathrm{curr}},H)
\]
with
\[
h_{\mathrm{tail}}:j_{\mathrm{tr}}(U)\simeq J,
\qquad
h_{\mathrm{Ham}}:d_V^{\mathrm{tr}}U\simeq\iota_{\mathcal G}\omega,
\]
\[
h_{\mathrm{curr}}:
d_V^{\mathrm{mix}}J\simeq
j_{\mathrm{cl}}(\iota_{\mathcal G}\omega),
\]
and a $2$-path $H$ identifying the path obtained from
$d_V^{\mathrm{mix}}(h_{\mathrm{tail}})$ and $h_{\mathrm{curr}}$ with the image
$j_{\mathrm{cl}}(h_{\mathrm{Ham}})$ under the displayed naturality square.
Equivalently, this is the fixed-boundary mapping object from the walking
homotopy-commutative square with a chosen $2$-simplex into the mixed coefficient
diagram.  Projection of the universal Noether lift and projection of the
Chapter~17 current Cartan diagram give canonical maps
\[
n_{\mathrm{curr}}:R_{\mathrm{Noeth}}
\longrightarrow\operatorname{CurrNoethBase}_{\mathcal G,L}^q,
\qquad
c_{\mathrm{curr}}:R_{\mathrm{curr}}
\longrightarrow\operatorname{CurrNoethBase}_{\mathcal G,L}^q.
\]

Define the typed iterated pullback
\[
\boxed{
\operatorname{RedNoethPreQ}_{\mathcal G,L}^{q}
:=
R_{\mathrm{desc}}
\times_{\operatorname{CartLift}^{\partial,q}_{\omega,\mathcal G}(L)}
R_{\mathrm{Noeth}}
\times_{\operatorname{CurrNoethBase}_{\mathcal G,L}^q}
R_{\mathrm{curr}}.
}
\]
A $T$-point is therefore a point of each of the three named realizations together
with an identification of the two actual Cartan pairs and an identification of
the two actual complete-transgression/current squares, including their
nullhomotopies and comparison $2$-path.  Mapping $T$ into the displayed
pullback gives exactly this iterated pullback of mapping spaces.  No unnamed
overlap category or compatibility package is used.
\end{definition}

\begin{proposition}[Base change and descent of reduced Noether compatibility]
\label{prop:prequantum-reduced-noether-compatibility-locality}
$\operatorname{CurrNoethBase}_{\mathcal G,L}^q$ and
$\operatorname{RedNoethPreQ}_{\mathcal G,L}^q$ commute with native base change
and satisfy effective descent.
\end{proposition}

\begin{proof}
The current overlap is a finite limit of relative section objects, path objects,
and one fixed-boundary $2$-path object.  These commute with pullback in the
native locally Cartesian closed $\infty$-topos and descend along effective
epimorphisms.  The two maps to the Cartan-lift object are natural under base
change by effective descent and by the formal-groupoid/jet-comonad equivalence.
The defining iterated pullback therefore has the same properties.
\end{proof}

On the resulting reduced base write
\[
\boxed{
C_{\mathrm{red}}^q
:=
C_L\times_Q
\operatorname{RedNoethPreQ}_{\mathcal G,L}^{q},
}
\]
\[
R_{\mathrm{red}}^q
:=(C_L//\mathcal G)\times_Q
\operatorname{RedNoethPreQ}_{\mathcal G,L}^{q}
\]
and let
\[
q_{\mathrm{red}}:
C_{\mathrm{red}}^q\twoheadrightarrow R_{\mathrm{red}}^q
\]
be the base-changed effective critical quotient.  Let
\[
\mathcal K_{\mathrm{red},\bullet}
\rightrightarrows C_{\mathrm{red}}^q
\]
be the corresponding base change of the critical formal groupoid.

\begin{definition}[Reduced prequantum automorphism group object]
\label{def:prequantum-reduced-automorphism-group-object}
Define
\[
\boxed{
\operatorname{Aut}^{\partial}_{\mathrm{PreQ}}
\bigl(\widehat\omega_{L,\mathrm{red}}^{(q)}\bigr)
:=
\Omega_{\widehat\omega_{L,\mathrm{red}}^{(q)}}
\Omega_{R_{\mathrm{red}}^q}^\infty
\mathscr P_{\mathrm{red}}^{(q)}.
}
\]
This is the full based-loop automorphism group object of the selected reduced
closed differential section.  Put
\[
\mathscr F_{\kappa,\mathrm{red}}^{(q)}
:=
\operatorname{fib}
(\kappa_{\mathrm{red}}^{(q)}:
 \mathscr D_{\mathrm{red}}^{(q)}\to\mathscr U_{\mathrm{red}}^{(q)}).
\]
Its curvature-kernel subgroup is canonically the based loop of the reduced flat
fibre,
\[
\Omega_{R_{\mathrm{red}}^q}^\infty
\mathscr F_{\kappa,\mathrm{red}}^{(q)}[-1].
\]
\end{definition}

Retain from Chapter~13 the descended formal adjoint group
\[
\operatorname{Ad}^{\mathrm{for}}(\mathcal K_{\mathrm{red}})
\longrightarrow R_{\mathrm{red}}^q,
\]
characterized by
\[
q_{\mathrm{red}}^*
\operatorname{Ad}^{\mathrm{for}}(\mathcal K_{\mathrm{red}})
\simeq
I^{\mathrm{for}}_{\mathcal K_{\mathrm{red}}}.
\]
Its canonical \v{C}ech transport is formal groupoid conjugation.

\begin{theorem}[Reduced Prequantum Noether action]
\label{thm:prequantum-reduced-noether-theorem}
Over $\operatorname{RedNoethPreQ}_{\mathcal G,L}^{q}$, restriction of the
critical Prequantum Noether lift to formal isotropy descends uniquely to a
homomorphism of group objects
\[
\boxed{
\alpha_{\mathcal G,L,\mathrm{red}}^{(q)}:
\operatorname{Ad}^{\mathrm{for}}
(\mathcal K_{\mathrm{red}})
\longrightarrow
\operatorname{Aut}^{\partial}_{\mathrm{PreQ}}
\bigl(\widehat\omega_{L,\mathrm{red}}^{(q)}\bigr).
}
\]
Its Hamiltonian primitive component is the descended complete Noether
transgression; its horizontal-tail shadow is
$J_{\mathcal G,L,\mathrm{red}}^{(q)}$.  Pullback along $q_{\mathrm{red}}$ recovers
the formal-isotropy restriction of the selected critical Prequantum Noether
lift.
\end{theorem}

\begin{proof}
The universal point of $\operatorname{NoethPreQReal}$ is a fixed-object formal
groupoid morphism.  Restricting it to formal isotropy gives a homomorphism
\[
I^{\mathrm{for}}_{\mathcal K_{\mathrm{red}}}
\longrightarrow
\operatorname{Aut}^{\partial}_{\mathrm{PreQ}}
(\widehat\omega_{\mathrm{crit}}^{(q)}).
\]
Because it is induced by a groupoid functor, it intertwines conjugation:
\[
\widetilde h(ghg^{-1})
\simeq
\widetilde h(g)\widetilde h(h)\widetilde h(g)^{-1}.
\]
Chapter~13 identifies conjugation with the canonical effective-descent transport
of formal isotropy.  Complete prequantum descent identifies the pullback of the
reduced automorphism group object with the critical automorphism group object.
Thus the displayed isotropy homomorphism is a morphism of the two effective
\v{C}ech descent diagrams.  Effective descent for internal group objects is an
equivalence, so it descends through a contractible space to the stated
homomorphism on $\operatorname{Ad}^{\mathrm{for}}(\mathcal K_{\mathrm{red}})$.
The current Cartan-diagram realization identifies the Hamiltonian/transgression
projection of this descended map with the descended $U^{(q)}$ and its
horizontal-tail current $J^{(q)}$.
\end{proof}

\begin{definition}[Selected reduced symmetry realization]
\label{def:prequantum-selected-reduced-symmetry-realization}
Define the internal reduced-symmetry parameter object by
\[
\boxed{
\operatorname{RedSymReal}_{\mathcal G,L}^{q}
:=
\operatorname{Ad}^{\mathrm{for}}(\mathcal K_{\mathrm{red}})
\longrightarrow R_{\mathrm{red}}^q.
}
\]
After base change to its total space, the diagonal gives the tautological reduced
symmetry parameter $\sigma$.  Evaluation of
$\alpha_{\mathcal G,L,\mathrm{red}}^{(q)}$ at $\sigma$ gives an individual
reduced prequantum automorphism
\[
\boxed{
\alpha_{\mathcal G,L,\mathrm{red}}^{(q)}(\sigma):
\mathbf1
\longrightarrow
\operatorname{Aut}^{\partial}_{\mathrm{PreQ}}
\bigl(\widehat\omega_{L,\mathrm{red}}^{(q)}\bigr).
}
\]
Global reduced symmetries are the sections
$\operatorname{Map}_{R_{\mathrm{red}}^q}
(R_{\mathrm{red}}^q,\operatorname{RedSymReal}_{\mathcal G,L}^{q})$.
Thus an individual automorphism is evaluated only after a reduced symmetry
parameter has been constructed or selected.
\end{definition}

\section{The prequantum discrepancy system}
\label{sec:prequantum-discrepancy-system}

The constructions above expose typed failures of existence, comparison, or
coherence.  None is converted into an admissibility condition, and comparisons
between different discrepancy types are themselves represented below.

\begin{definition}[Prequantum discrepancies]
\label{def:prequantum-discrepancies}
The prequantum discrepancy system consists of the following separately
constructed fibres, zero loci, inverse loci, cofibres, and matching/latching
problems:
\begin{enumerate}
\item the prequantization connecting section
      $\mathfrak o_{\mathrm{PreQ}}(\omega)$;
\item the differential Lagrangian and differential Lepage refinement fibres
      $\operatorname{PreQLag}$ and $\operatorname{PreQLep}$;
\item the trace and fundamental-class shadow-lift fibres
      $\operatorname{TrLiftReal}_{\partial}$ and
      $\operatorname{FundLiftReal}_{\partial}$, the coherent-stage extension
      fibre $\operatorname{TrCohLiftReal}_{\partial}$, the universal underlying
      stage-system objects $\operatorname{StageSys}_Z^U$, the complete relative
      lift $\operatorname{TransStageCurvReal}_{\partial}$, and the fixed
      transgression-square path object $\operatorname{TransSqLiftReal}$;
\item the differential correspondence $2$-cell lift fibre
      $\operatorname{TwoCellReal}_{\partial}$ and its decorated refinement;
\item the functor-level BNV inverse locus $\operatorname{StokesEq}_p$, its small
      detection presentation $\operatorname{StokesTestEq}_p^\lambda$, and the
      evaluated cofiber $C^{\mathrm{Stokes}}_{p,\widehat E}$;
\item the Cartan-transport and vertical-origin realizations
      $\operatorname{CartCurvReal}$, $\operatorname{VertCurvReal}$, and
      $\operatorname{VertClassReal}$ when those structures are not already canonical;
\item the Hamiltonian primitive fibre, the long-edge obstruction
      $\bar\epsilon_2$, the full matching object
      $\operatorname{HamMultReal}$, and the sufficient functor-level inverse
      locus $\operatorname{HamEdgeEq}$;
\item the distinct raw/formal liftability
      inverse loci $\operatorname{KSEq}^{\mathrm{raw}}$ and
      $\operatorname{KSEq}^{\mathrm{for}}$, together with the raw-recovery fibre
      $\operatorname{KSRawLiftReal}$ over a selected formal inverse;
\item the prequantum-to-quantomorphism controller comparison
      $a_{\mathrm{PQ}}$, its source-module cofiber
      $C^{\mathrm{PreQ/Quant}}_{\widehat\omega}$, and its inverse locus
      $\operatorname{PreQQuantEq}_{\widehat\omega}$;
\item the differentiated comparison
      $\gamma^{\mathrm{KS,diff}}_{\widehat\omega}$, its source-module cofiber
      $C^{\mathrm{KS,diff}}_{\widehat\omega}$, and
      $\operatorname{KSDiffEq}_{\widehat\omega}$;
\item the first source-module KS connecting class
      $\mathbf k_{\widehat\omega}^{(1)}$ and the complete higher
      morphism-corolla--latching tower $\operatorname{ob}^{\mathrm{PreQ}}_s$;
\item the transgressive Hamiltonian comparison
      $\operatorname{HamCompReal}_{\mathcal G,L}^{q}$, the native Hamiltonian
      Noether coherence realization, and its skeletal matching tower;
\item the coherent liftable-Hamiltonian factorization
      $\operatorname{NoethHamLiftReal}$ and native prequantum lift fibre
      $\operatorname{NoethPreQReal}$;
\item the full differential-refinement/Prequantum-Noether compatibility realization
      $\operatorname{DiffSymPreQCompat}$;
\item the complete differentiated Prequantum Noether lift tower;
\item the internal complete descent fibre $\operatorname{PreQDesc}$ and the
      universal descent comparison
      $\operatorname{DescTot}_{q_{\mathrm{curv}}}
      (\operatorname{PreQ}_{\mathrm{crit},\bullet}^{q})
      \to\operatorname{PreQ}_{\mathrm{red}}^q$, whose inverse locus is
      canonically inhabited by the universal reduction theorem;
\item the realized de Rham exchange inverse loci $\operatorname{DREq}(c;F)$;
\item the indexed comparison frame, its exchange inverse loci, and
      $\operatorname{CmpPreQReal}_{\mathcal B,b}$;
\item the separately constructed represented nondegeneracy inverse locus.
\end{enumerate}
\end{definition}

\begin{construction}[Typed comparison realization for discrepancies]
\label{constr:prequantum-typed-discrepancy-comparison}
Let $D_i\to P_i$ and $D_j\to P_j$ be two discrepancy objects with displayed
derivation trees $d_i,d_j$.  Form the labelled derivation forest $F_{ij}$ by
identifying their genuinely shared leaves.  Let
$\operatorname{TgtCosp}_{ij}$ be the maximal subgroupoid of typed cospans
\[
s_i\longrightarrow s\longleftarrow s_j
\]
in the occurrence category of $F_{ij}$, where $s_i,s_j$ are the output sorts of
$D_i,D_j$ and both arrows are composites of actual comparison or base-transport
occurrences.  This is a small groupoid determined by the two derivation trees.

For $k\in\operatorname{TgtCosp}_{ij}$, take the finite limit of the two displayed
base-transport diagrams and the indexed comparison realization along the two
legs; denote it by $P_k$.  Over $P_k$ both discrepancies have constructed
transports
\[
\Phi_{k,i}D_i,\ \Phi_{k,j}D_j\in\mathcal C_k,
\]
where $\mathcal C_k$ is the actual target category at the apex sort $s$.  Define
the universal common-target realization
\[
\boxed{
\operatorname{AmbReal}_{ij}
:=
\int_{k\in\operatorname{TgtCosp}_{ij}}P_k.
}
\]
If there is no typed common-target cospan, this object is empty.  Over it define
\[
\boxed{
\operatorname{DiscCmpReal}_{ij}
:=
\int_{k\in\operatorname{TgtCosp}_{ij}}
\operatorname{Eqv}_{\mathcal C_k}
\left(\Phi_{k,i}D_i,\Phi_{k,j}D_j\right).
}
\]
For a constructed comparison arrow
$\chi_{ij}:\Phi_{k,i}D_i\to\Phi_{k,j}D_j$ along a specified cospan $k$, define
the sharper realization
\[
\boxed{
\operatorname{DiscCmpReal}_{ij}(\chi_{ij})
:=
\operatorname{Inv}(\chi_{ij}).
}
\]
A point is exactly a typed identification of the two discrepancies, including
the common-target cospan, coefficient-system transport, and all base-change
coherences.  Consequently
the unary universal property or inhabitation of one discrepancy supplies no
identification with another; every such identification is a point of the
displayed comparison realization.  For the pairs already related above---raw to
formal liftability, actual flat gauge to the intrinsic KS fibre, and differential
to underlying transgression---the constructed comparison arrows define the
corresponding canonical maps into $\operatorname{DiscCmpReal}_{ij}$.

The relative integrated isotropy of the KS projection is not a discrepancy: the
finite-pullback computation of
Theorem~\ref{prop:prequantum-universal-isotropy-exact-extension} identifies it
canonically with the flat kernel.
\end{construction}

\begin{proposition}[Locality of typed discrepancy comparison]
\label{prop:prequantum-typed-discrepancy-comparison-locality}
$\operatorname{AmbReal}_{ij}$ and $\operatorname{DiscCmpReal}_{ij}$ commute
with native base change and satisfy effective descent.
\end{proposition}

\begin{proof}
The cospan groupoid is a fixed small indexing groupoid.  For every cospan,
$P_k$ is a finite limit of base transports and indexed comparison realizations,
and the walking-equivalence object is a limit in the corresponding native
slice.  Pullback commutes with these limits and with the Grothendieck
construction over the constant small groupoid.  Effective descent reconstructs
each transported discrepancy and its equivalence data, hence reconstructs the
same Grothendieck construction.
\end{proof}

\section{Represented prequantum anomalies}
\label{sec:prequantum-anomalies}

\begin{definition}[Represented prequantum anomaly]
\label{def:prequantum-represented-anomaly}
Let $D$ be one of the native discrepancy objects constructed in this chapter,
let $b_D$ be its displayed derivation tree, and let
\[
P_D:=\operatorname{CmpPreQReal}_{\mathcal B,b_D}.
\]
Over $P_D$, apply in the target coefficient system the same constructors, in the
same order, to the transported leaves of $b_D$.  Structural recursion on $b_D$
produces a target discrepancy $D^{\mathrm{rep}}$ and the composite comparison
arrow
\[
\boxed{
\chi_D:\Phi(D)\longrightarrow D^{\mathrm{rep}}.
}
\]
Every exchange used in this recursion is one of the inverse loci included in
$P_D$, so $\chi_D$ is canonically an equivalence there.  Define
\[
\boxed{
\operatorname{Anom}^{\mathrm{rep}}(D):=D^{\mathrm{rep}},
\qquad
\operatorname{AnomCmpReal}(D):=\operatorname{Inv}(\chi_D).
}
\]
A represented prequantum anomaly is a section, morphism, or boundary point of
$\operatorname{Anom}^{\mathrm{rep}}(D)$ obtained by transporting the
corresponding native discrepancy datum.  In particular this construction applies
to $\mathfrak o_{\mathrm{PreQ}}$, the Hamiltonian lift fibre,
$C^{\mathrm{Stokes}}$, $C^{\mathrm{KS,diff}}$, the native
Hamiltonian--Noether matching obstructions, the differentiated Prequantum Noether
latching obstructions, and the complete descent comparison.  They respectively
measure failure of differential refinement, coherent prequantum lift, Stokes
compatibility, differentiated kernel recognition, native Hamiltonian--Noether
coherence, partition-Lie Noether liftability, and reduction descent.
\end{definition}

\begin{construction}[Typed anomaly-cancellation realization]
\label{constr:prequantum-typed-anomaly-cancellation}
Work over $P_D$ and write $\mathcal C_D$ for the target slice in which
$D^{\mathrm{rep}}$ lives.  Its cancellation object is determined by the actual
shape of the discrepancy:
\begin{enumerate}
\item for a pointed section $a:\mathbf1\to A$, put
      \[
      \operatorname{Cancel}(a)
      :=\operatorname{Path}_{\operatorname{Map}_{\mathcal C_D}
      (\mathbf1,A)}(a,0);
      \]
\item for a lifting problem $u:K\to Y$ against $p:X\to Y$, put
      \[
      \operatorname{Cancel}(u,p)
      :=\operatorname{fib}_{u}\!\left(
      \operatorname{Map}_{\mathcal C_D}(K,X)\to
      \operatorname{Map}_{\mathcal C_D}(K,Y)
      \right);
      \]
\item for a comparison arrow $v:X\to Y$ required to be invertible, put
      $\operatorname{Cancel}(v):=\operatorname{Inv}(v)$;
\item for a matching or latching boundary map
      $r:F\to\operatorname{Match}(\partial)$, put
      $\operatorname{Cancel}(r;\partial):=\operatorname{fib}_{\partial}(r)$;
\item for a stable cofiber discrepancy $C=\operatorname{cofib}(v)$, put
      \[
      \operatorname{Cancel}(C)
      :=\operatorname{Inv}(v)
      \simeq
      \operatorname{Inv}(C\to0);
      \]
\item for a selected descent problem with restriction map
      $r_{\mathrm{desc}}:\operatorname{Desc}(X^\bullet)\to
      \operatorname{DescBd}(X^\bullet)$ and prescribed boundary datum $\delta$,
      put
      \[
      \operatorname{Cancel}_{\mathrm{desc}}(\delta)
      :=\operatorname{fib}_{\delta}(r_{\mathrm{desc}}).
      \]
      For the prequantum reduction discrepancy this is the already constructed
      complete fibre $\operatorname{PreQDesc}$.  Only a theorem asserting
      equivalence of the entire descent functor uses
      $\operatorname{Inv}(X\to\operatorname{Tot}(X^\bullet))$.
\end{enumerate}
The represented cancellation locus of $D$ is the applicable object in this
list, pulled back along $\operatorname{AnomCmpReal}(D)$ when the comparison
inverse is to be retained as part of the point.
\end{construction}

\begin{proposition}[Cancellation has no additional axiom]
\label{prop:prequantum-anomaly-cancellation-universal-property}
A point of $\operatorname{Cancel}(D)$ is exactly a cancellation of the
represented discrepancy, with all required homotopies and higher coherences.
No anomaly-free condition or selected nullhomotopy is part of the input.
\end{proposition}

\begin{proof}
The six cases are the defining universal properties of, respectively, a path
object, a mapping-space fibre, the walking-equivalence inverse locus, a
matching/latching fibre, the stable equivalence
$\operatorname{cofib}(v)\simeq0$ if and only if $v$ is invertible, and the
mapping-space fibre defining complete descent.  Pullback along
$\operatorname{AnomCmpReal}(D)$ merely records the already constructed typed
identification $\Phi(D)\simeq D^{\mathrm{rep}}$.  Thus the point contains the
requested nullhomotopy, lift, inverse, filler, or descent inverse rather than a
predicate asserting that one exists.  Generalized differential cohomology is the
standard represented receptacle for the abelian cases \cite{Freed2000}.
\end{proof}

\section{Nondegenerate prequantum realization}
\label{sec:prequantum-nondegenerate-realization}

Prequantization is defined for complete closed presymplectic data.  Nondegeneracy
is imposed only through the exact represented $(1,1)$ contraction already
constructed in Chapter~\ref{ch:symmetry-noether}.

On the represented smooth scalar critical theorem domain retain
\[
\omega_{L,T}^{(q),\flat}:
T_{C_{L,T}/T}
\longrightarrow
\mathcal W_{\omega,T}^{(q)}
\]
and its reduced counterpart.

\begin{definition}[Nondegeneracy realization]
\label{def:prequantum-nondegeneracy-realization}
Define
\[
\operatorname{NonDeg}_{\omega}^{q}
:=
\operatorname{Inv}
\bigl(\omega_{L,T}^{(q),\flat}\bigr),
\qquad
\operatorname{NonDeg}_{\omega,\mathrm{red}}^{q}
:=
\operatorname{Inv}
\bigl(\omega_{L,\mathrm{red},T}^{(q),\flat}\bigr).
\]
\end{definition}

\begin{definition}[Nondegenerate prequantization]
\label{def:prequantum-nondegenerate-prequantization}
Let $P_{\mathrm{PreQ}}$ and $P_{\mathrm{flat}}$ be the displayed parameter
objects of the represented prequantization and of the exact contraction, and
let both map to the common represented curvature parameter $P_{\mathrm{curv}}$.
Form the explicit common base
\[
P_{\mathrm{common}}
:=
P_{\mathrm{PreQ}}\times_{P_{\mathrm{curv}}}P_{\mathrm{flat}}.
\]
After pulling both realizations to $P_{\mathrm{common}}$, define
\[
\boxed{
\operatorname{SympPreQ}_{\kappa,u}(\omega)
:=
P_{\mathrm{common}}
\times_{P_{\mathrm{PreQ}}}
\operatorname{PreQ}_{\kappa,u}(\omega)
\times_{P_{\mathrm{flat}}}
\operatorname{NonDeg}_{\omega}^{q},
}
\]
and analogously after reduction.
\end{definition}

\begin{proposition}[Prequantization precedes nondegeneracy]
\label{prop:prequantum-precedes-nondegeneracy}
The projection
$\operatorname{SympPreQ}_{\kappa,u}(\omega)\to
\operatorname{PreQ}_{\kappa,u}(\omega)$ is pullback along the separately
constructed nondegeneracy inverse locus.  The construction of prequantization uses
neither the represented tangent complex nor invertibility of the exact contraction.
\end{proposition}

This order matches higher prequantum geometry, where prequantization exists for
pre-$n$-plectic closed forms before nondegeneracy is imposed
\cite{FiorenzaRogersSchreiber2013,Schreiber2016}.

\section{Indexed comparison realization and transport of prequantum geometry}
\label{sec:prequantum-comparison-transport}

The native constructions use coefficient categories over many geometric bases:
relation degrees, mixed rectangles, Cartan quotients, critical objects, and
reduced quotients.  A single functor between two coefficient categories over one
base cannot compare this package.  We therefore realize an indexed comparison of
the entire coefficient systems and then represent, as inverse loci, precisely
the right-adjoint and complete-limit exchanges used by a branch.

For a displayed branch $b$, form its geometric occurrence category
$\mathcal B_b$ as follows.  Its objects are the occurrences of geometric bases
in the typed derivation tree of $b$; its generating arrows are the actual base
morphism occurrences in that tree; identities and composites are adjoined, and
two composites are identified precisely by the identity, composition, or higher
pasting coherence already displayed in the branch.  Shared labelled input
leaves give one object.  Since the derivation tree is small and well founded,
$\mathcal B_b$ is a canonical small $\infty$-category.  Write
$\mathcal B:=\mathcal B_b$.  Let
\[
\mathcal E,\mathcal D:
\mathcal B^{\mathrm{op}}
\longrightarrow
\operatorname{CAlg}(\operatorname{Pr}^{L}_{\mathrm{st}})
\]
be the source and target stable coefficient systems.

\begin{definition}[Indexed comparison frame]
\label{def:prequantum-indexed-comparison-frame}
Let $\Delta_{\mathsf{XSp}}:\mathcal S\to\mathsf{XSp}$ denote the constant-object
functor.  Define
\[
\boxed{
\operatorname{CmpFrame}_{\mathcal B}(\mathcal E,\mathcal D)
:=
\Delta_{\mathsf{XSp}}\operatorname{Map}_{
\operatorname{Fun}(\mathcal B^{\mathrm{op}},
\operatorname{CAlg}(\operatorname{Pr}^{L}_{\mathrm{st}}))}
(\mathcal E,\mathcal D).
}
\]
A point is an indexed exact strong symmetric monoidal comparison
$\Phi=(\Phi_Z)_{Z\in\mathcal B}$ together with the coherent pullback exchanges
\[
\chi_{g^*}:
 g_{\mathcal D}^*\Phi_Z
\simeq
\Phi_{Z'}g_{\mathcal E}^*
\]
for every $g:Z'\to Z$ in $\mathcal B$, including identity, composition, and
higher pasting coherences.
When the coefficient systems vary over a native parameter, the same mapping
object is formed affine-testwise and right Kan extended.  It is therefore an
object of the native slice, rather than an external space subsequently used as
the base of an internal pullback.
\end{definition}

\begin{construction}[Right-adjoint exchange and inverse locus]
\label{constr:prequantum-pi-exchange-realization}
For a morphism $g$ for which both indexed coefficient systems carry the
constructed adjunction $g^*\dashv\Pi_g$, the pullback
exchange has the canonical right-adjoint mate
\[
\boxed{
\chi_{\Pi,g}:
\Phi_Z\Pi_g^{\mathcal E}
\longrightarrow
\Pi_g^{\mathcal D}\Phi_{Z'}.
}
\]
Define
\[
\boxed{
\operatorname{PiEq}(g):=\operatorname{Inv}(\chi_{\Pi,g}).
}
\]
Thus compatibility with a dependent product is a constructed mate followed, when
needed, by an explicit inverse locus.
\end{construction}

\begin{construction}[Internal-Hom exchange and inverse locus]
\label{constr:prequantum-hom-exchange-realization}
For $A,B\in\mathcal E_Z$, strong monoidality of $\Phi_Z$ adjoints the evaluation
map to
\[
\boxed{
\chi_{\mathrm{Hom},A,B}:
\Phi_Z\underline{\operatorname{Hom}}_{\mathcal E}(A,B)
\longrightarrow
\underline{\operatorname{Hom}}_{\mathcal D}(\Phi_ZA,\Phi_ZB).
}
\]
Define
\[
\boxed{
\operatorname{HomEq}(A,B)
:=
\operatorname{Inv}(\chi_{\mathrm{Hom},A,B}).
}
\]
\end{construction}

\begin{construction}[Relative mapping- and section-spectrum exchange]
\label{constr:prequantum-map-section-exchange-realization}
For an affine test $t:h_U\to Z$ and $A,B\in\mathcal E_Z$, functoriality of the
exact comparison gives a morphism of spectra
\[
\chi_{\mathrm{Map},t;A,B}:
\operatorname{map}_{\mathcal E_U}(t^*A,t^*B)
\longrightarrow
\operatorname{map}_{\mathcal D_U}
(t^*\Phi_ZA,t^*\Phi_ZB).
\]
The pullback coherences of the indexed frame make these arrows natural in the
affine-point category.  Taking their affine right Kan extension gives an actual
arrow of relative mapping-spectrum objects, denoted
$\chi_{\mathrm{Map},A,B}$.  Define
\[
\boxed{
\operatorname{MapEq}(A,B)
:=
\operatorname{Inv}(\chi_{\mathrm{Map},A,B}).
}
\]
Since strong monoidality identifies $\Phi_Z\mathbf1_Z$ with $\mathbf1_Z$, put
\[
\chi_{\mathrm{Sec},A}:=\chi_{\mathrm{Map},\mathbf1_Z,A},
\qquad
\boxed{
\operatorname{SecEq}(A):=
\operatorname{Inv}(\chi_{\mathrm{Sec},A}).
}
\]
Applying $\Omega^\infty$ to the universal spectrum equivalences over these
inverse loci gives the required exchanges for section objects, mapping spaces,
path objects, and based-loop objects.
\end{construction}

\begin{proposition}[Mapping exchange is independent of exact monoidality]
\label{prop:prequantum-map-exchange-necessary}
The $\operatorname{MapEq}$ and $\operatorname{SecEq}$ factors are not implied
by exact strong symmetric monoidality or by $\operatorname{HomEq}$.
\end{proposition}

\begin{proof}
For a commutative ring spectrum $R$, extension of scalars
$\operatorname{Sp}\to\operatorname{Mod}_R$ is exact and strong symmetric
monoidal.  Its comparison on endomorphism spectra of the tensor unit is the unit
map $\mathbb S\to R$, which is not generally an equivalence.  Therefore neither
relative mapping spectra nor section spectra may be transported without their
own inverse loci.  The construction above records exactly the missing
equivalence, with its affine and pullback coherences.
\end{proof}

\begin{construction}[Complete-totalization exchange and inverse locus]
\label{constr:prequantum-totalization-exchange-realization}
For an actual cosimplicial diagram $C^\bullet$ in $\mathcal E_Z$, the universal
cone gives a canonical comparison
\[
\boxed{
\chi_{\mathrm{Tot},C}:
\Phi_Z(\operatorname{Tot}C^\bullet)
\longrightarrow
\operatorname{Tot}\Phi_Z(C^\bullet).
}
\]
Define
\[
\boxed{
\operatorname{TotEq}(C^\bullet)
:=
\operatorname{Inv}(\chi_{\mathrm{Tot},C}).
}
\]
This inverse locus is retained for every sequential or complete totalization
whose preservation is used.  Exactness of $\Phi$ alone is not substituted for
this construction.
\end{construction}

\begin{lemma}[Adjunction-mate comparison]
\label{lem:prequantum-adjunction-mate-comparison}
Suppose
\[
L_{\mathcal E}:\mathcal A_{\mathcal E}
\rightleftarrows\mathcal C_{\mathcal E}:R_{\mathcal E},
\qquad
L_{\mathcal D}:\mathcal A_{\mathcal D}
\rightleftarrows\mathcal C_{\mathcal D}:R_{\mathcal D}
\]
are adjunctions, $\Phi_{\mathcal A}$ and $\Phi_{\mathcal C}$ are the indexed
comparison functors, and
\[
\sigma:L_{\mathcal D}\Phi_{\mathcal A}
\Longrightarrow
\Phi_{\mathcal C}L_{\mathcal E}
\]
is the constructed left-adjoint exchange.  Its right mate is the natural
transformation
\[
\boxed{
\sigma^\sharp:
\Phi_{\mathcal A}R_{\mathcal E}
\longrightarrow
R_{\mathcal D}\Phi_{\mathcal C}
}
\]
whose component is the adjoint of
\[
L_{\mathcal D}\Phi_{\mathcal A}R_{\mathcal E}
\xrightarrow{\sigma}
\Phi_{\mathcal C}L_{\mathcal E}R_{\mathcal E}
\xrightarrow{\Phi_{\mathcal C}(\varepsilon_{\mathcal E})}
\Phi_{\mathcal C}.
\]
Mate formation preserves identity and pasting diagrams.  Hence the unit,
counit, monad, and comonad compatibility cells of $\sigma^\sharp$ are determined
by those of $\sigma$.
\end{lemma}

\begin{proof}
The displayed composite corresponds under
$L_{\mathcal D}\dashv R_{\mathcal D}$ to a unique natural transformation
$\sigma^\sharp$.  Applying the two adjunction equivalences to a pasted diagram
shows that the mate of a composite is the composite of the mates, with the usual
unit and counit whiskerings.  The triangle identities remove those whiskerings
for identities.  This proves naturality and all stated coherence assertions.
\end{proof}

\begin{construction}[Formal, Artinian, Lie, tangent, and descent exchanges]
\label{constr:prequantum-remaining-comparison-exchanges}
For unitwise formalization, apply
Lemma~\ref{lem:prequantum-adjunction-mate-comparison} to the inclusion of
unitwise formal objects and its right adjoint.  For every transported input $A$
this gives
\[
\chi_{\mathrm{Form},A}:
 \Phi(\operatorname{Form}_{\mathcal E}A)
 \longrightarrow
 \operatorname{Form}_{\mathcal D}(\Phi A).
\]
Restrict the indexed comparison and its pullback exchanges to the actual
Artinian test category used at $A$.  Applying the finite-limit exchanges
degreewise constructs
\[
\chi_{\mathrm{Art},A}:
 \Phi(\operatorname{Art}_{\mathcal E}A)
 \longrightarrow
 \operatorname{Art}_{\mathcal D}(\Phi A).
\]
On the locally coherent theorem domain, the source and target
partition-Lie differentiation equivalences identify maps of Artinian formal
moduli problems with maps of spectral partition-Lie algebroids.  The preceding
Artinian arrow therefore induces the unique arrow
\[
\chi_{\operatorname{Lie}^{\pi},A}:
 \Phi(\operatorname{Lie}^{\pi}_{\mathcal E}A)
 \longrightarrow
 \operatorname{Lie}^{\pi}_{\mathcal D}(\Phi A).
\]
For a transported action, differentiate its action map and apply
$\chi_{\operatorname{Lie}^{\pi}}$ to the controller and
$\chi_{\mathrm{Hom}}$ to its representation object.  The tangent
representation universal property then gives
\[
\chi_{\mathrm{Tan},A}:
 \Phi(\operatorname{Tan}_{\mathcal E}A)
 \longrightarrow
 \operatorname{Tan}_{\mathcal D}(\Phi A).
\]
Here $\Phi$ denotes the appropriate component of the indexed comparison; all
base-change whiskerings are those of
$\operatorname{CmpFrame}_{\mathcal B}$.  Define the objectwise inverse loci
\[
\boxed{
\begin{aligned}
\operatorname{FormEq}(A)&:=\operatorname{Inv}(\chi_{\mathrm{Form},A}),&
\operatorname{ArtEq}(A)&:=\operatorname{Inv}(\chi_{\mathrm{Art},A}),\\
\operatorname{LieEq}^{\pi}(A)&:=
 \operatorname{Inv}(\chi_{\operatorname{Lie}^{\pi},A}),&
\operatorname{TanEq}(A)&:=\operatorname{Inv}(\chi_{\mathrm{Tan},A}).
\end{aligned}
}
\]

For a critical or effective quotient with augmented simplicial descent object
$A\to A^\bullet$, functoriality of the augmentation and the totalization cone
constructs
\[
\begin{aligned}
\chi_{\mathrm{CritDesc},A}&:
 \Phi(\operatorname{Tot}_{\mathrm{crit}}A^\bullet)
 \longrightarrow
 \operatorname{Tot}_{\mathrm{crit}}\Phi(A^\bullet),\\
\chi_{\mathrm{EffDesc},A}&:
 \Phi(\operatorname{Tot}_{\mathrm{eff}}A^\bullet)
 \longrightarrow
 \operatorname{Tot}_{\mathrm{eff}}\Phi(A^\bullet).
\end{aligned}
\]
Set
\[
\boxed{
\operatorname{CritDescEq}(A):=
\operatorname{Inv}(\chi_{\mathrm{CritDesc},A}),
\qquad
\operatorname{EffDescEq}(A):=
\operatorname{Inv}(\chi_{\mathrm{EffDesc},A}).
}
\]

Finally, if $h:X\to Y$ is an occurrence to which effective image is applied,
write its effective image as the realization of its image groupoid
$\lvert\operatorname{Im}_\bullet(h)\rvert$.  The degreewise indexed comparison
and the universal realization cone give
\[
\chi_{\mathrm{EffIm},h}:
\Phi\lvert\operatorname{Im}_\bullet(h)\rvert
\longrightarrow
\lvert\operatorname{Im}_\bullet(\Phi h)\rvert,
\qquad
\boxed{
\operatorname{EffImEq}(h):=
\operatorname{Inv}(\chi_{\mathrm{EffIm},h}).
}
\]
When the endpoint comparisons are equivalences and $\Phi$ preserves the
realization in question, the universal colimit comparison supplies a canonical
point of $\operatorname{EffImEq}(h)$; otherwise its possible emptiness records
the exact preservation problem.
\end{construction}

\begin{construction}[Cartan, relation, and mixed exchanges]
\label{constr:prequantum-cartan-mixed-comparison-exchanges}
For a quotient map $q$, combine $\chi_{q^*}$ and $\chi_{\Pi,q}$ to obtain the jet
comonad comparison
\[
\boxed{
\chi_{J,q}:
\Phi J_q^{\mathcal E}
\longrightarrow
J_q^{\mathcal D}\Phi.
}
\]
The map $\chi_{J,q}$ sends a $J_q^{\mathcal E}$-coalgebra to a
$J_q^{\mathcal D}$-coalgebra and hence constructs
\[
\Phi_q^{\mathrm{Cart}}:
\operatorname{Coalg}_{J_q^{\mathcal E}}(\mathcal E)
\longrightarrow
\operatorname{Coalg}_{J_q^{\mathcal D}}(\mathcal D).
\]
The comonad counit and comultiplication squares are the mates of the identity and
composition squares for $q^*\dashv\Pi_q$; hence
$\Phi_q^{\mathrm{Cart}}$ is defined without an additional compatibility axiom.
For a selected source coalgebra $(A,a)$ and the independently constructed target
coalgebra $(A',a')$, the underlying comparison and its coalgebra-compatibility
square define
\[
\chi_{\mathrm{Cart},A}:
\Phi_q^{\mathrm{Cart}}(A,a)\longrightarrow(A',a')
\quad\text{in }\operatorname{Coalg}_{J_q^{\mathcal D}}(\mathcal D).
\]
When an equivalence with that selected target is required, define
\[
\boxed{
\operatorname{CartEq}(A,a;A',a')
:=
\operatorname{Inv}_{\operatorname{Coalg}_{J_q^{\mathcal D}}(\mathcal D)}
(\chi_{\mathrm{Cart},A}).
}
\]

Let $\mathcal S_{\mathrm{rel}}$ be the $\infty$-category of Cartesian sections
of the coefficient fibration over the relation simplex category, and let
$\mathcal S_{\mathrm{mix}}$ be the analogous section category over the mixed
bisimplex category.  The degreewise maps and all exchange squares assemble,
before any inverse is selected, to natural transformations
\[
\chi_{\mathrm{rel}}:
\Phi(C_{\mathrm{rel}}^{\mathcal E})
\Longrightarrow C_{\mathrm{rel}}^{\mathcal D}
\quad\text{in }\mathcal S_{\mathrm{rel}},
\]
\[
\chi_{\mathrm{mix}}:
\Phi(C_{\mathrm{mix}}^{\mathcal E})
\Longrightarrow C_{\mathrm{mix}}^{\mathcal D}
\quad\text{in }\mathcal S_{\mathrm{mix}}.
\]
Naturality of the geometric staircase gives a morphism
$\chi_{\mathrm{AW}}$ between the two staircase arrows in
$\operatorname{Fun}(\Delta^1,\mathcal S_{\mathrm{mix}})$.  Define
\[
\boxed{
\operatorname{RelEq}:=
\operatorname{Inv}_{\mathcal S_{\mathrm{rel}}}(\chi_{\mathrm{rel}}),
\qquad
\operatorname{MixEq}:=
\operatorname{Inv}_{\mathcal S_{\mathrm{mix}}}(\chi_{\mathrm{mix}}),
}
\]
and
\[
\boxed{
\operatorname{AWEq}:=
\operatorname{Inv}_{\operatorname{Fun}(\Delta^1,\mathcal S_{\mathrm{mix}})}
(\chi_{\mathrm{AW}}).
}
\]
The $\operatorname{TotEq}$ factors are then added separately for every complete
filtration/totalization used after these degreewise comparisons.  Thus each
relation, mixed, and Alexander--Whitney condition is the inverse locus of an
actual natural transformation.  In particular these are not independent
degreewise inverse choices: their simplicial, bisimplicial, and staircase
compatibilities are part of the ambient functor categories.
\end{construction}

\begin{definition}[Indexed curvature comparison realization]
\label{def:prequantum-morphism-curvature-diagrams}
Let distinguished curvature diagrams
\[
\mathbf C_{\mathcal E,Z}
=(\mathscr D_{\mathcal E,Z}\to\mathscr U_{\mathcal E,Z}
  \leftarrow\mathscr C_{\mathcal E,Z}),
\]
\[
\mathbf C_{\mathcal D,Z}
=(\mathscr D_{\mathcal D,Z}\to\mathscr U_{\mathcal D,Z}
  \leftarrow\mathscr C_{\mathcal D,Z})
\]
be given on the bases where a curvature comparison is required.  Over
$\operatorname{CmpFrame}_{\mathcal B}$ form the walking-equivalence spaces
between the images of the three source coefficients and the three target
coefficients and impose the two curvature-square homotopies.  Taking the
homotopy limit over the required bases defines
\[
\boxed{
\operatorname{CmpCurvReal}_{\mathcal B}
(\mathbf C_{\mathcal E},\mathbf C_{\mathcal D}).
}
\]
A point is an indexed comparison together with actual coherent equivalences of
the curvature diagrams.
\end{definition}

\begin{definition}[Full indexed prequantum comparison realization]
\label{def:prequantum-full-comparison-realization}
Fix a displayed construction branch $b$.  Its derivation tree $d_b$ is the
well-founded rooted typed tree whose leaves are the actual inputs of $b$ and whose
internal vertices are the successive displayed constructor occurrences used to
produce its output.  Thus a repeated use of the same constructor gives repeated
vertices, and a shared input is represented by the same labelled leaf.  Let
$\operatorname{Occ}_{\mathrm{ex}}(d_b)$ be the full occurrence poset on the
vertices labelled by one of
\[
\Pi,\quad
\underline{\operatorname{Hom}},\quad
\operatorname{Map},\quad
\operatorname{Sec},\quad
\operatorname{Tot},\quad
\operatorname{Cartan},\quad
\operatorname{Rel},\quad
\operatorname{Mix},\quad
\operatorname{AW},\quad
\operatorname{Form},\quad
\operatorname{Art},\quad
\operatorname{Lie}^{\pi},\quad
\operatorname{Tan},\quad
\operatorname{CritDesc},\quad
\operatorname{EffDesc},\quad
\operatorname{EffIm}.
\]
The order is generated by $o<o'$ when the output at $o$ is an input, possibly
through exchange-insensitive vertices, to the constructor at $o'$.  Write
\[
\mathcal K_b:=N(\operatorname{Occ}_{\mathrm{ex}}(d_b)).
\]
Because the displayed expression for $b$ is universe-small and well founded,
$\mathcal K_b$ is a small direct $\infty$-category.  The well-founded rank
\[
\rho:\operatorname{Ob}(\mathcal K_b)\longrightarrow\theta_b
\]
is defined recursively by
$\rho(o)=\sup_{o'<o}(\rho(o')+1)$, and
$\theta_b=\sup_o(\rho(o)+1)$.

Put
\[
P_0:=\operatorname{CmpCurvReal}_{\mathcal B}
(\mathbf C_{\mathcal E},\mathbf C_{\mathcal D}).
\]
Assume $P_\alpha$ has been constructed and let $o$ have rank $\alpha$.
All inputs of $o$ occur at smaller rank or are exchange-insensitive.  The former
have already been identified over $P_\alpha$, while the latter are transported
by functoriality and exactness.  The construction at $o$ therefore yields, over
$P_\alpha$, one of the actual arrows
\[
\chi_{\Pi},\ \chi_{\mathrm{Hom}},\ \chi_{\mathrm{Map}},\
\ \chi_{\mathrm{Sec}},\ \chi_{\mathrm{Tot}},\
\ \chi_{\mathrm{Cart}},\ \chi_{\mathrm{rel}},\ \chi_{\mathrm{mix}},\
\ \chi_{\mathrm{AW}},\ \chi_{\mathrm{Form}},\ \chi_{\mathrm{Art}},\
\ \chi_{\operatorname{Lie}^{\pi}},\ \chi_{\mathrm{Tan}},\
\ \chi_{\mathrm{CritDesc}},\ \chi_{\mathrm{EffDesc}},\
\ \chi_{\mathrm{EffIm}},
\]
with the occurrence-specific inputs suppressed.  Denote it by
$\chi_o:x_o\to y_o$ in its displayed $P_\alpha$-indexed ambient category.
Naturality of the exchanges and the mate-pasting lemma supply all squares and
higher cells along dependencies in $\mathcal K_b$; no such cells are extra data.

At a successor stage form the simultaneous relative inverse locus in the slice
over $P_\alpha$,
\[
P_{\alpha+1}
:=
\prod_{\substack{o\in\mathcal K_b\\\rho(o)=\alpha}}^{(\mathsf{XSp})_{/P_\alpha}}
\operatorname{Inv}_{P_\alpha}(\chi_o),
\]
and at a limit stage take the inverse limit of the preceding bases,
\[
P_\lambda:=\lim_{\alpha<\lambda}P_\alpha.
\]
Define
\[
\boxed{
\operatorname{CmpPreQReal}_{\mathcal B,b}
:=
P_{\theta_b}.
}
\]
It represents a point of
$\operatorname{CmpCurvReal}_{\mathcal B}$ together with a coherent inverse to
every exchange arrow at every actual occurrence in $d_b$.

For a specified small family $\mathfrak c$ of displayed branches, form the
labelled derivation forest
\[
F_{\mathfrak c}:=
\left(\coprod_{b\in\mathfrak c}d_b\right)\!\big/\!\sim,
\]
where equal labelled input leaves are identified and an output root is identified
with its occurrence as an input of a later named output.  Let
$\operatorname{Br}_{\mathfrak c}$ be the full subposet of named output roots.
For $r\leq r'$, restriction of inverse data from the subderivation below $r'$ to
the one below $r$ gives a map of comparison realizations over the common object
$P_0$.  Define
\[
\boxed{
\operatorname{CmpPreQReal}_{\mathcal B,\mathfrak c}
:=
\lim_{r\in\operatorname{Br}_{\mathfrak c}^{\mathrm{op}}}^{(\mathsf{XSp})_{/P_0}}
\operatorname{CmpPreQReal}_{\mathcal B,d_{\leq r}}.
}
\]
Its point consists of one comparison frame, one curvature comparison, and
compatible exchange inverses for all branches, including the identifications on
their genuinely shared subderivations.  Equivalently, apply the preceding rank
tower once to the exchange-sensitive occurrence poset of $F_{\mathfrak c}$;
Fubini for limits identifies the two constructions.
\end{definition}

\begin{proposition}[Rank independence of the comparison realization]
\label{prop:prequantum-comparison-rank-independence}
The object $\operatorname{CmpPreQReal}_{\mathcal B,b}$ is independent, up to a
canonical contractible choice, of the ordinal ranking or of any linear extension
used to compute the tower.  The same holds for
$\operatorname{CmpPreQReal}_{\mathcal B,\mathfrak c}$.
\end{proposition}

\begin{proof}
For every test object $T$, mapping $T$ into the terminal stage of the tower gives
the limit of the groupoid-valued diagram whose objects are: a map
$T\to\operatorname{CmpCurvReal}_{\mathcal B}$; for each occurrence $o$, an
inverse of the pulled-back arrow $\chi_o$; and for each dependency, the
naturality and mate-pasting coherence already present in the occurrence
diagram.  Any well-founded rank computes this same small limit by iterated
limits.  Fubini for limits in spaces gives a canonical equivalence between any
two such computations, and the space of comparison equivalences is contractible
because both represent the same functor of $T$.  The forest case is identical.
\end{proof}

\begin{proposition}[Locality of the indexed comparison realization]
\label{prop:prequantum-indexed-comparison-locality}
The parameterized form of
$\operatorname{CmpPreQReal}_{\mathcal B,b}$ commutes with native base change and
satisfies effective descent.  The same is true for a comparison forest.
\end{proposition}

\begin{proof}
The comparison frame, the mapping-spectrum and section-spectrum arrows, and all
other exchange arrows are formed affine-testwise from the indexed comparison
and its pullback coherences.  Right Kan extension makes them Cartesian under
native pullback.  Walking-equivalence and inverse loci are limits in the native
$\infty$-topos, and each successor or limit stage of the rank tower is a small
limit of such objects.  Pullback preserves these limits, while effective descent
reconstructs the frame, every exchange arrow, its inverse data, and their
dependency coherences.  Transfinite induction on the rank proves the assertion;
the forest case follows by Fubini for small limits.
\end{proof}

When $\mathcal B$ and the branch are understood, write simply
$\operatorname{CmpPreQReal}$.

\begin{theorem}[Prequantum comparison transport]
\label{thm:prequantum-comparison-transport}
Over $\operatorname{CmpPreQReal}_{\mathcal B,b}$ the universal indexed comparison
induces canonical equivalences of every construction in the corresponding
branch, including:
\begin{enumerate}
\item closed differential coefficients, flat fibres, prequantization moduli,
      obstruction zero loci, torsor actions, and higher automorphisms;
\item Cartan and vertical-origin realizations, quantomorphism and
      curvature-symmetry groupoids, and flat kernels;
\item global mixed coefficient systems, coherent and complete
      Alexander--Whitney morphisms, prism identities, long-edge obstruction and
      multiplicativity realizations, Hamiltonian groupoids over those
      realizations, raw and formal Hamiltonian-prequantum objects, effective
      liftable images, and their inverse loci;
\item native Hamiltonian-Noether, Prequantum Noether, and universal full
      differential-refinement symmetry realizations;
\item on the locally coherent domain, differentiated controllers, intrinsic KS
      fibres, latching towers, flat-adjoint tangent representations, observable
      operations, and coefficient connections;
\item after critical comparison, complete prequantum descent, reduced
      prequantization, the descended formal adjoint group, and the reduced
      Prequantum Noether action.
\end{enumerate}
A noninvertible indexed comparison gives the corresponding functorial maps but
no equivalence conclusion.
\end{theorem}

\begin{proof}
Proceed simultaneously by structural induction on the derivation tree $d_b$
and induction on the rank of $\mathcal K_b$.  At a successor rank, the
universal point of $P_{\alpha+1}$ selects simultaneous coherent inverses of the
canonical exchanges at that rank.  The dependency arrows of
$\operatorname{Occ}_{\mathrm{ex}}(d_b)$ and the mate-pasting coherences identify
their input data and naturality squares with the outputs already transported at
lower ranks.
At a limit stage the compatible transported objects assemble by the universal
property of $P_\lambda$.  This transports the whole diagram and all its pasting
coherences.

Finite stable limits use the objectwise exact comparison.  Dependent products use
$\operatorname{PiEq}$; internal Homs use $\operatorname{HomEq}$; relative
mapping spectra and section spectra use $\operatorname{MapEq}$ and
$\operatorname{SecEq}$; complete filtrations and symmetry/descent totalizations
use $\operatorname{TotEq}$;
Cartan coalgebras use the objectwise $\operatorname{CartEq}$; relation and mixed
constructions use $\operatorname{RelEq}$, $\operatorname{MixEq}$, and
$\operatorname{AWEq}$.  Formalization, Artinian restriction, partition-Lie
differentiation, and tangent formation use
$\operatorname{FormEq}$, $\operatorname{ArtEq}$,
$\operatorname{LieEq}^{\pi}$, and $\operatorname{TanEq}$.  Critical and
effective descent use $\operatorname{CritDescEq}$ and
$\operatorname{EffDescEq}$, while effective liftable images use
$\operatorname{EffImEq}$.  Mapping spaces, path and inverse objects,
prequantization fibres, torsors, automorphism loops, and groupoid mapping ends
are thereby transported through the explicit mapping- or section-spectrum
exchange rather than through exactness alone.  These are exactly the labels of
$\operatorname{Occ}_{\mathrm{ex}}(d_b)$; every remaining vertex is transported
by ordinary functoriality or finite exactness.  This exhausts the constructors of
$b$, proves transport of the complete branch term, and invokes no unstated
preservation theorem.
\end{proof}

\subsection{External recognition targets}
\label{subsec:prequantum-external-recognition-targets}

The comparison with an external theory is represented by the indexed object
$\operatorname{CmpPreQReal}_{\mathcal B,b}$ (or
$\operatorname{CmpPreQReal}$ when $\mathcal B$ and $b$ are understood).  A
target-side recognition theorem identifies
what a point of the target prequantization fibre means; a point of the comparison
realization then transports the native object to that recognized target.  No
bridge is postulated as a hypothesis on the native prequantum object.

In ordinary degree-two differential cohomology, the target prequantization fibre
is the groupoid of $U(1)$-principal bundles, equivalently complex line bundles,
with connection and prescribed curvature; its inhabitation is the ordinary
integrality condition and its automorphisms are flat gauge transformations
\cite{Kostant1970}.  Hence a point of $\operatorname{CmpPreQReal}$ with this
target identifies the native prequantization, obstruction, flat torsor, and
quantomorphism constructions with the ordinary ones.

For higher differential cohomology, the target fibres are higher circle
bundles/gerbes with connection and the infinitesimal higher quantomorphism
controller is the higher local-observable $L_\infty$-algebra on the theorem domain
of \cite{Rogers2012,FiorenzaRogersSchreiber2013,Schreiber2016}.  A point of the
mixed comparison realization transports the native integrated and differentiated
KS packages to those target objects.

Likewise the shifted closed-form, shifted-symplectic, boundary/Lagrangian, and
shifted-prequantization theories of
\cite{PantevToenVaquieVezzosi2013,Calaque2015,Safronov2020} provide target-side
recognition statements.  The present construction does not identify the native
complete infinitesimal coefficient with an exterior-power or HKR model by
notation.  Instead the corresponding recognition holds precisely over the
constructed comparison realization, whose possible emptiness records failure of
the bridge.

\section{Functoriality, base change, and locality}
\label{sec:prequantum-functoriality}

\begin{theorem}[Fixed-base prequantum functoriality]
\label{thm:prequantum-fixed-base-functoriality}
Every morphism of stable curvature diagrams induces the functorial maps of
Theorem~\ref{thm:prequantum-comparison-transport}.  These maps are compatible with
composition and identities.  A fixed-object formal-groupoid morphism compatible
with the curvature diagram carries quantomorphism, Hamiltonian, KS, and Prequantum
Noether objects covariantly; the induced maps on stable cochains have the
contravariance inherited from pullback along the relation morphism.
\end{theorem}

\begin{proof}
Represent the chosen output by its displayed derivation tree.  Define the map on
leaves by the given morphism of curvature diagrams and, for a formal-groupoid
leaf, by the given groupoid morphism.  Induct upward through the tree.  A finite
limit or stable fibre is mapped by the induced morphism of its defining diagram;
an internal mapping object is mapped by precomposition in its first variable and
postcomposition in its second; a dependent product is mapped by the mate of the
pullback square; and a totalization is mapped by the induced morphism of
cosimplicial diagrams.  Relation and mixed objects are mapped degreewise, so the
Alexander--Whitney staircase and prism homotopies map by naturality.

Effective image is functorial through the induced map of image groupoids and
geometric realizations.  Unitwise formalization and Artinian restriction use the
functorial units of their adjunctions; partition-Lie differentiation and tangent
formation use the functors constructed in Chapters~14--15.  Matching, latching,
zero, lift, path, and inverse loci are pullbacks of functorial arrow diagrams, so
they inherit maps.  Critical descent is mapped by the morphism of augmented
descent diagrams and its totalization cone.  At each vertex, identity and
composition follow from functoriality or from the mate-pasting identity of
Lemma~\ref{lem:prequantum-adjunction-mate-comparison}.  Induction reaches the
root and proves the assertion with the stated covariance and contravariance.
\end{proof}

\section{Arbitrary native base change}
\label{sec:prequantum-arbitrary-base-change}

\begin{theorem}[Arbitrary native base change]
\label{thm:prequantum-arbitrary-native-base-change}
Let $f:T'\to T$ be arbitrary in the native big spectral Nisnevich
$\infty$-topos.  It induces the exact strong symmetric monoidal pullback
\[
f^{*,\partial}:\mathcal R_T^\partial\longrightarrow\mathcal R_{T'}^\partial.
\]
The following constructions commute with that pullback by canonical coherent
base-change equivalences:
\begin{enumerate}
\item big differential mapping and section spectra;
\item stable curvature diagrams, closed differential class coefficients,
      prequantization fibres, flat coefficients, obstruction zero loci, and
      universal prequantization families;
\item differential Lagrangian and differential Lepage refinement objects and the
      complete variational prequantization map;
\item Cartan-transport and vertical-origin realizations, raw/formal
      quantomorphisms and curvature symmetries, and flat kernels;
\item mixed rectangular coefficients, the ordered Alexander--Whitney morphism,
      complete and normalized prism identities, the long-edge defect
      $\bar\epsilon_2$, $\operatorname{HamMultReal}$, and the universal raw and
      formal Hamiltonian groupoids over that realization;
\item raw/formal Hamiltonian-prequantum objects, raw liftable images, formal
      liftable groupoids, and both $\operatorname{KSEq}^{\mathrm{raw}}$ and
      $\operatorname{KSEq}^{\mathrm{for}}$;
\item $\operatorname{HamCompReal}$, native Hamiltonian-Noether matching,
      coherent liftability, Prequantum Noether, and full differential-refinement
      symmetry compatibility;
\item critical restriction, internal $\operatorname{PreQDesc}$, and stable
      coefficient--section descent along the base-changed critical quotient;
\item the comparison-realization constructions whenever both coefficient systems
      and the comparison problem are pulled back.
\end{enumerate}
No equivalence $\mathcal R_T^\partial\simeq\mathcal R_{T'}^\partial$ is asserted
for a general $f$; the statement is compatibility of the constructions with the
actual pullback functor.
\end{theorem}

\begin{proof}
Write $\mathscr X:=\mathsf{XSp}$.  Pullback between slices of the native
$\infty$-topos has both adjoints:
\[
\Sigma_f\dashv f^*\dashv\Pi_f.
\]
Here $\Sigma_f$ is composition with $f$, and $\Pi_f$ exists by local Cartesian
closedness.  Hence $f^*:\mathscr X_{/T}\to\mathscr X_{/T'}$ preserves all small
limits and all small colimits.  In the parameterized stable coefficient system,
the induced functor
\[
f_E^*:E_T\longrightarrow E_{T'}
\]
likewise occurs between its indexed adjoints
\[
f_{E,!}\dashv f_E^*\dashv f_{E,*}.
\]
It therefore preserves all small stable limits and colimits.  In particular,
for every simplicial object $D_\bullet$ and cosimplicial object $C^\bullet$ used
below,
\[
f_E^*|D_\bullet|
\simeq|f_E^*D_\bullet|,
\qquad
f_E^*\operatorname{Tot}(C^\bullet)
\simeq\operatorname{Tot}(f_E^*C^\bullet).
\]
The functor is exact and strong symmetric monoidal.  The closed and dependent
Beck--Chevalley equivalences constructed for the native coefficient system give
\[
f_E^*
\underline{\operatorname{Hom}}_{E_T}(D,E)
\simeq
\underline{\operatorname{Hom}}_{E_{T'}}
(f_E^*D,f_E^*E)
\]
and, for every Cartesian square
\[
\begin{tikzcd}
Y' \ar[r,"f_Y"] \ar[d,"g'"'] & Y \ar[d,"g"]\\
T' \ar[r,"f"'] & T,
\end{tikzcd}
\]
the equivalence
\[
f_E^*\Pi_g
\simeq
\Pi_{g'}f_Y^*.
\]
Thus closed internal Homs are transported by closed Beck--Chevalley, not merely
by strong monoidality.

Independently, $f^{*,\partial}$ on the big Thom-stable differential category is
restriction of Cartesian affine-point families along
$\mathcal C^{\mathrm{aff}}_{/T'}\to\mathcal C^{\mathrm{aff}}_{/T}$.  Its exact
strong symmetric monoidal structure is computed componentwise, as in
Proposition~\ref{prop:prequantum-big-differential-recovery-pullback}.
The big differential mapping spectrum is even more directly pointwise.  If
$t':h_U\to T'$, then
\[
\begin{aligned}
\bigl(f^*\operatorname{map}^{\partial}_T(\widehat E,\widehat F)\bigr)(t')
&=
\operatorname{map}_{\mathcal R_U^0}
(\widehat E_{ft'},\widehat F_{ft'})\\
&=
\operatorname{map}^{\partial}_{T'}
(f^{*,\partial}\widehat E,f^{*,\partial}\widehat F)(t').
\end{aligned}
\]
The comparison is therefore an equivalence before spectral Nisnevich descent
and hence after descent.  Strong monoidality identifies the pulled-back unit, so
the same calculation proves the assertion for section spectra.

Exactness of $f_E^*$ now transports stable fibres, cofibres, zero loci, curvature
pullbacks, flat coefficients, obstruction fibres, and prequantization fibres.
Complete filtrations, matching objects, path objects, walking-equivalence
objects, mapping ends, and descent totalizations are small limits; latching
objects, geometric realizations, and groupoid quotients are small colimits.
The preceding adjunction argument transports all of them.  Effective images
also commute with pullback: the pullback of an effective-epimorphism/monomorphism
factorization is again such a factorization, and uniqueness of image
factorizations identifies it with the image of the pulled-back map.

Chapters~10 and~13 identify the base-changed relative infinitesimal relation and
all its simplicial structure maps.  The Chapter~13 Cartesianity theorem then
identifies the base-changed unitwise formal relation.  Since unitwise
formalization is constructed from that relation by the displayed coreflection
limits, its base-change mate is an equivalence.  Applying $f_E^*$ degreewise
identifies the relation and mixed coefficient
diagrams.  Naturality of the ordered Alexander--Whitney transformation then
identifies its staircase, all prism homotopies, and the long-edge defect.
Cartan coalgebras are transported by the comonad mate together with its
counit/comultiplication coherences.  It follows that the raw and formal mapping
groupoids and every zero, path, lift, matching, multiplicativity, and inverse
locus in items~(4)--(7) are the pullbacks of their original defining diagrams.

For critical descent, base change carries the augmented critical relation to the
augmented relation of the base-changed quotient.  Degreewise base change, the
displayed dependent Beck--Chevalley equivalences, and preservation of
totalization identify $f^*\operatorname{PreQDesc}$ with the newly formed descent
fibre.  The same argument transports the stable coefficient and section descent
objects.  Finally, pulling back both indexed coefficient systems carries each
comparison arrow $\chi_o$ to the arrow for the base-changed occurrence.
Pullback preserves its walking-equivalence inverse locus and every small-limit
stage of the rank tower defining $\operatorname{CmpPreQReal}$.  Transfinite
induction on $d_b$ proves item~(9).

This argument concerns the intrinsic native stable branch.  It does not alter
the separate scope of the Chapter~10 realized de Rham exchange in
Proposition~\ref{prop:prequantum-thom-stable-realized-base-change-scope}.
\end{proof}

\begin{proposition}[Actual scope of Thom-stable realized de Rham base change]
\label{prop:prequantum-thom-stable-realized-base-change-scope}
For the Chapter~10 Thom-stable realizations of de Rham stages, normalized cochains,
and coherently closed cochains, arbitrary represented base change carries the
already constructed exchange morphisms
\[
c^{*,0}F_S\longrightarrow F_{S'}.
\]
These are equivalences on the represented spectral-\'etale finite-presentation
domain proved in Chapter~10.  The prequantum constructions do not replace these
exchange morphisms by arbitrary-base-change equivalences.
\end{proposition}

\begin{definition}[Realized de Rham exchange inverse locus]
\label{def:prequantum-realized-dr-exchange-inverse-locus}
For any one of the actually constructed Chapter~10 exchange morphisms
\[
\beta_{c,F}:c^{*,0}F_S\longrightarrow F_{S'},
\]
define
\[
\boxed{
\operatorname{DREq}(c;F)
:=
\operatorname{Inv}(\beta_{c,F}).
}
\]
A point is the contractibly coherent inverse data for the actual exchange map.
On the represented spectral-\'etale finite-presentation theorem domain,
Chapter~10 supplies the canonical point of this inverse locus.  Outside that
domain the inverse locus records the exact invertibility problem instead of a
guardrail assertion.
\end{definition}

\section{Formal-\'etale locality of the intrinsic stable branch}
\label{sec:prequantum-formal-etale-locality}

\begin{theorem}[Formal-\'etale locality]
\label{thm:prequantum-formal-etale-locality}
Let $f:Y'\to Y$ be formally \'etale over $T$.  For curvature diagrams and mixed
Hamiltonian coefficients constructed intrinsically from the relative de Rham
relations and parameterized stable coefficients, formal-\'etale locality of the
relative infinitesimal relation carries the closed coefficients, ordered
contractions, Hamiltonian primitive equations, the long-edge obstruction and
multiplicativity realization, and their prequantization pullbacks to the
corresponding restricted constructions by canonical equivalences.
Critical mixed calculus, complete stable presymplectic descent, and stable
prequantum reduction are transported after the corresponding critical base change.
\end{theorem}

\begin{proof}
Chapters~10, 12, and~\ref{ch:symmetry-noether} prove formal-\'etale Cartesianity
for the intrinsic relative de Rham relation and the stable constructions obtained
from it by limits.  The prequantum curvature pullback, Hamiltonian mixed rectangle,
and complete descent objects are further limits of those transported diagrams.
The statement concerns this intrinsic stable branch; the Thom-stable realized de
Rham exchange retains the separate scope of
Proposition~\ref{prop:prequantum-thom-stable-realized-base-change-scope}.
\end{proof}

\section{Finite spectral Nisnevich descent}
\label{sec:prequantum-finite-nisnevich-descent}

\begin{theorem}[Finite Nisnevich descent for prequantum objects]
\label{thm:prequantum-finite-nisnevich-descent}
For a finite spectral Nisnevich cover with \v{C}ech nerve $U_\bullet$, the
following internal objects are reconstructed by totalization from their
restrictions:
\begin{enumerate}
\item $\mathcal R^\partial$, its mapping spectra, and all stable curvature and
      prequantization coefficients;
\item $\operatorname{PreQLag}$, $\operatorname{PreQLep}$, the complete closed
      presymplectic coefficient, and variational prequantization;
\item trace and fundamental-class candidate maps and their shadow-lift fibres,
      coherent-stage diagram maps and extension fibres, fixed
      transgression-curvature square lifts, and differential correspondence
      $2$-cell lift realizations on represented correspondence charts;
\item $\operatorname{CartCurvReal}$, $\operatorname{VertCurvReal}$,
      $\operatorname{VertClassReal}$, raw/formal
      quantomorphism and curvature-symmetry objects, the Hamiltonian long-edge
      obstruction, $\operatorname{HamMultReal}$, the raw/formal Hamiltonian and
      Hamiltonian-prequantum groupoids over it, liftable-image, and raw/formal
      KS inverse-locus objects;
\item $\operatorname{HamCompReal}$, native Hamiltonian-Noether matching,
      liftable-Hamiltonian factorization, Prequantum Noether, full
      differential-refinement compatibility, and stable critical
      $\operatorname{PreQDesc}$ objects;
\item reduced stable prequantum coefficients and sections, the descended formal
      adjoint group, and the reduced Prequantum Noether action after the
      corresponding critical descent data have been glued;
\item $\operatorname{CmpPreQReal}_{\mathcal B,b}$, its exchange inverse loci,
      and $\operatorname{DREq}$ whenever their source and target diagrams are
      present on the cover.
\end{enumerate}
Global spaces of points are formed only after global sections of these descended
parameterized objects.
\end{theorem}

\begin{proof}
Chapter~9 gives finite Nisnevich descent of $\mathcal R^0$; mapping spectra in a
limit of stable categories are limits.  Parameterized stable coefficients,
relation-level coefficient systems, complete filtrations, internal Homs, path
objects, and walking-equivalence loci are computed by limits compatible with the
cover.  The mixed rectangles are finite limits of relation degrees.  Local
formalization and effective images reconstruct their global counterparts by the
already proved functoriality, and effective descent reconstructs the reduced
coefficient--section and automorphism diagrams.
\end{proof}

\section{Spectral-\'etale transport of the differentiated branch}
\label{sec:prequantum-spectral-etale-differentiated}

\begin{theorem}[Spectral-\'etale transport of differentiated prequantum geometry]
\label{thm:prequantum-spectral-etale-transport-differentiated}
On the locally coherent theorem domain, coherent one- and two-variable
spectral-\'etale transport carries:
\begin{enumerate}
\item $\mathbb A_{\widehat\omega}^{\mathrm{PreQ}}$,
      $\mathbb A_{\widehat\omega}^{\mathrm{Ham,\,lift}}$,
      $\mathbb A_{\widehat\omega}^{\mathrm{Flat}}$, and
      $\mathbb F_{\widehat\omega}^{\mathrm{KS}}$;
\item the comparison $\gamma^{\mathrm{KS,diff}}$, its cofiber, and its inverse
      locus;
\item the connecting class $\mathbf k_{\widehat\omega}^{(1)}$ and the complete
      prequantum morphism-latching tower;
\item the flat-adjoint prequantum tangent representation, exact labelled observable operations,
      coefficient connection, curvature filler, Bianchi coherence, and higher
      flatness hierarchy;
\item the differentiated Prequantum Noether lift, its finite-weight obstruction
      tower, and its higher labelled operation system.
\end{enumerate}
Consequently these affine constructions descend to the global coherent spectral
coefficient system.
\end{theorem}

\begin{proof}
Quotient formalization, formal leaves, spectral Artinian restriction, Schlessinger
reflection, partition-Lie differentiation, tangent representations, the tangent
enveloping monad, coherent mapping sheaves, and morphism corolla--latching all
carry the Chapter~14--15 spectral-\'etale exchange equivalences.  The remaining
objects are functorial finite limits, cofibers, inverse loci, and operation objects
inside those transported categories.
\end{proof}


\section{Audit of constructions and theorem scopes}
\label{sec:prequantum-audit-constructions-theorem-scopes}

The chapter has four constructional layers.  Potential input packages have been
replaced by mapping objects, walking-equivalence loci, Cartan and vertical-origin
realizations, homotopy fibres, effective images, zero loci, inverse loci, matching
towers, or morphism-corolla--latching towers whenever the preceding framework
permits.

\subsection{Unconditional native and stable constructions}
\label{subsec:prequantum-audit-native}

Once the Chapter~9 affine differential coefficient system and the Chapter~10
parameterized stable category are fixed, the following are constructions:
\begin{enumerate}
\item $\mathcal R_T^{\partial}$, arbitrary pullback, represented recovery,
      effective descent, and big differential mapping and section spectra;
\item $\operatorname{CurvDiagReal}$, the closed differential coefficient, flat
      fibre, connecting obstruction, prequantization zero locus, torsor, higher
      automorphisms, and universal family;
\item $\operatorname{DiffDenCurvReal}_x(M)$,
      $\operatorname{PreQLag}$, and the separate
      $\operatorname{PreQLep}$ lifting problem, the complete closed
      presymplectic coefficient, and variational prequantization;
\item $\operatorname{CartCurvReal}$, $\operatorname{VertCurvReal}$, and
      $\operatorname{VertClassReal}$, followed by
      raw/formal quantomorphism, curvature-symmetry, and flat-gauge groupoids;
\item the mixed rectangle, ordered Alexander--Whitney morphism, complete closed
      first prism, the degree-two long-edge obstruction,
      $\operatorname{HamMultReal}$, the sufficient locus
      $\operatorname{HamEdgeEq}$, and the universal raw and formal Hamiltonian
      groupoids over the multiplicativity realization;
\item $\operatorname{HamPreQRaw}$, the raw effective liftable image, the raw KS
      torsor extension, $\operatorname{HamPreQ}$, the formal homotopy-kernel
      sequence, the separate raw/formal KS inverse loci, and the raw-recovery
      fibre over a formal inverse;
\item exact integrated relative isotropy, with no extra isotropy-exactness datum;
\item $\operatorname{HamCompReal}$, the complete Hamiltonian-Noether matching
      tower, coherent formal liftability, the native Prequantum Noether lift,
      $\operatorname{VarRef}_{\kappa_M}^q$, its complete universal
      differential-refinement symmetry realization, and the resulting
      compatibility fibre;
\item critical restriction over the full prequantization parameter,
      coefficient-level curvature descent, internally parameterized
      $\operatorname{PreQDesc}$, the reduced closed differential section, fixed-pair
      descent, and the universal reduction theorem;
\item descent of formal isotropy to the formal adjoint group and the reduced
      Prequantum Noether action, with the reduced current recovered as its
      Hamiltonian primitive shadow;
\item the Cartan/current overlap limit $\operatorname{RedNoethPreQ}$,
      $\operatorname{CmpPreQReal}_{\mathcal B_b,b}$ indexed by
      the derivation-tree occurrence poset $\mathcal K_b$, including its
      $\operatorname{MapEq}$ and $\operatorname{SecEq}$ inverse loci, the typed
      discrepancy-comparison/anomaly realizations, $\operatorname{DREq}$, and the
      separately constructed represented nondegeneracy inverse locus.
\end{enumerate}
None of these constructions requires a line-bundle or gerbe presentation,
integral-lift axiom, ordinary connection, orientation, Hamiltonian admissibility,
prequantum liftability condition, reduction-compatibility condition, anomaly-free
condition, nondegeneracy condition, or BV complex as input.

\subsection{Thom-stable differential and exceptional constructions}
\label{subsec:prequantum-audit-differential-umkehr}

The following use the represented Thom-stable differential theorem range of
Chapters~9--10:
\begin{enumerate}
\item represented differential coefficients and their pullback and exceptional
      functors;
\item Chapter~12 differential density and action-value coefficients and the
      orientation-free action;
\item the pointed candidate mapping objects $\operatorname{TrReal}_{\partial}$
      and $\operatorname{FundReal}_{\partial}$ and their lift fibres over
      prescribed motivic shadows;
\item the explicitly indexed Chapter~10 realized stage diagram
      $\mathbf{DRStage}^{\partial}$, $\operatorname{TrCohReal}_{\partial}$,
      and the coherent extension fibre of a selected stage trace;
\item $\operatorname{SqReal}$, the differential transgression operator,
      the pointed $\operatorname{UTrReal}$ and
      $\operatorname{TransCurvReal}_{\partial}$, the universal
      $\operatorname{StageSys}_Z^U$, and the complete relative stage-square
      lift $\operatorname{TransStageCurvReal}_{\partial}$;
\item differential transgression, canonical trace/fundamental-class composition,
      and Fubini;
\item $\operatorname{TwoCellReal}_{\partial}$,
      $\operatorname{TwoCellStageReal}$, and the decorated traced
      correspondence $(\infty,2)$-category;
\item the BNV coefficientwise discrepancy, functor-level Stokes inverse locus,
      and its small $\lambda$-compact detection presentation.
\end{enumerate}
The exceptional adjunction constructs its counit but not a coefficient trace or
fundamental class.  The realized de Rham package retains the actual arbitrary
base-change exchange morphisms; their invertibility is represented by
$\operatorname{DREq}$ outside the theorem domain where Chapter~10 proves it.

\subsection{Formalized Artinian and differentiated constructions}
\label{subsec:prequantum-audit-differentiated}

On the locally coherent spectral coefficient sector the chapter constructs:
\begin{enumerate}
\item separate differentiation of prequantum, liftable-Hamiltonian,
      quantomorphism, and actual-flat formal groupoids;
\item the differentiated prequantum-to-quantomorphism comparison, its stable
      source-module cofiber, and its inverse locus;
\item the intrinsic categorical KS fibre in partition-Lie algebroids and the
      comparison from differentiated actual flat gauge;
\item its discrepancy and inverse locus after forgetting to stable source modules;
\item the source-module connecting KS class and the complete anchored
      morphism-corolla--latching lift tower;
\item the stable flat Cartan coefficient descended before applying
      $\Omega^\infty$, its comparison with kernel conjugation, and its tangent
      representation, kept distinct from actual-flat differentiation;
\item the observable controller, exact labelled operations, coefficient
      connection, curvature filler, Bianchi coherence, and higher flatness tower;
\item differentiation of an actual native Prequantum Noether lift and its complete
      finite-weight obstruction and higher-operation systems.
\end{enumerate}
No splitting of the KS extension, identification of actual flat gauge with the
intrinsic KS fibre, identification of the flat Cartan representation with actual
isotropy, or arbitrary varying-object partition-Lie pullback is inserted.

\subsection{Represented and external comparison constructions}
\label{subsec:prequantum-audit-comparison}

Ordinary line-bundle, higher-gerbe, $n$-plectic, shifted-symplectic,
shifted-prequantum, and dg/$L_\infty$ identifications are not guardrail clauses.
The bridge itself is the indexed realization
$\operatorname{CmpPreQReal}_{\mathcal B,b}$; its target-side
fibres are recognized by the corresponding external theorems
\cite{Kostant1970,Rogers2012,FiorenzaRogersSchreiber2013,
PantevToenVaquieVezzosi2013,Calaque2015,Safronov2020}.  Its possible emptiness
records failure of the actual indexed bridge.  On the Safronov branch,
Construction~\ref{constr:prequantum-safronov-quantization-handoff} further forms
the complete prequantum-Lagrangian-fibration input and the target quantum
functors; no first-order or associated-graded comparison is promoted to a full
bridge.

\subsection{No hidden input conditions}
\label{subsec:prequantum-audit-no-hidden-input}

None of the following is retained as a field of structure on a prequantum object,
formal symmetry, or field theory:
\begin{enumerate}
\item a ``prequantizable'' or ``integral'' closed-class condition;
\item a curvature or underlying-class morphism attached to arbitrary coefficients;
\item a differential refinement of an arbitrary Lepage realization;
\item a relation-level vertical coefficient origin or a vertical closed-class
      origin attached to an arbitrary closed coefficient/section;
\item closedness imposed only at unary cochain level;
\item an orientation, regulator, coefficient trace, coherent stage-trace package,
      or fundamental class;
\item a Cartan structure on an arbitrary curvature diagram;
\item a Hamiltonian/admissible symmetry subcategory;
\item a Hamiltonian multiplication, long-edge inverse, or primitive-exactness
      datum supplied in advance;
\item identification of raw and formal Hamiltonian liftability;
\item a splitting of the integrated or differentiated KS extension;
\item an isotropy-exactness condition;
\item an identification of actual flat gauge with the intrinsic differentiated
      fibre;
\item a stable representation manufactured from group-level conjugation;
\item an identification of the Chapter~17 current with the complete Hamiltonian
      transgression;
\item a coherent Hamiltonian-Noether groupoid action supplied in advance;
\item a fixed underlying Lagrangian silently assumed invariant before the full
      differential refinement symmetry problem is formed;
\item a compatible-reduction, a canonically selected reduced symmetry parameter,
      or anomaly-cancellation condition;
\item a nondegeneracy condition attached to prequantization;
\item a differential lift of every correspondence $2$-cell;
\item a line-bundle, gerbe, shifted-form, or dg/$L_\infty$ comparison supplied as
      a theorem-domain hypothesis rather than constructed through the indexed
      comparison frame and its exchange inverse loci.
\end{enumerate}
They are replaced by the explicit realization, fibre, inverse-locus, or
matching/latching constructions named above.  The lift fibres, zero loci,
inverse loci, and coherence-filler spaces may be empty.  The unrestricted stable
mapping realizations are instead canonically pointed by zero.  A point of a
genuine discrepancy object contains the required lift, nullhomotopy, inverse,
matching filler, latching filler, or descent coherence.

\subsection{Theorem-domain properties are not data}
\label{subsec:prequantum-audit-theorem-domain-properties}

Representability where quasi-coherent cotangent or ordinary tangent geometry is
used, the qcqs and smooth separated finite-presentation domains of differential
exceptional integration, coherence/local coherence for pro-coherent Koszul
duality, spectral-\'etaleness for the varying-object differentiated coefficient
system, rationality for dg rectification, and perfectness where ordinary duality
is invoked remain theorem domains inherited from Chapters~9--17.  They are
properties of objects on which proved constructions apply; none is attached as an
auxiliary field to a formal groupoid, prequantization, Hamiltonian symmetry, or
Noether lift.

\section{Boundary with quantization proper}
\label{sec:prequantum-boundary-quantization-proper}

\begin{definition}[Tagged native prequantum derivations]
\label{def:prequantum-output-closure}
Fix universes $\mathbb U\in\mathbb V$.  Let
$\Sigma_{\mathrm{PreQ}}^{\mathrm{occ}}$ be the occurrence signature of this
chapter.  Its sort symbols are the actual displayed occurrences of parameter
objects, parameterized stable and differential coefficient categories,
simplicial objects, internal groupoids, spectral partition-Lie algebroids,
representations, and slice realizations.  Its constant symbols are the displayed
standing inputs.  For every actual constructor occurrence in the chapter it has
one constructor symbol, with exactly that occurrence's input base, input sorts,
output base, output sort, and indexing shape.  The occurrences are of the
following kinds:
\begin{enumerate}
\item the finite limits and the particular small-limit shapes actually used,
      stable fibres and cofibres, internal mapping objects, tensor products,
      dependent products, and $\Omega^\infty$;
\item exact functor categories, coherent-cochain objects, complete
      totalizations, and affine-point descent;
\item effective images, unitwise formalization, partition-Lie differentiation,
      tangent-representation formation, and effective quotient descent;
\item walking-equivalence, path, lift, zero, inverse, matching, latching, and
      critical-descent realizations.
\end{enumerate}
Because the textual collection of displayed occurrences and each displayed
indexing shape are $\mathbb U$-small,
$\Sigma_{\mathrm{PreQ}}^{\mathrm{occ}}$ is $\mathbb U$-small.  In particular,
there is no single purported constructor whose arity ranges over an unbounded
collection of all small categories.  Put
$\Sigma_{\mathrm{PreQ}}:=\Sigma_{\mathrm{PreQ}}^{\mathrm{occ}}$.  The signature
adds neither an admissibility predicate nor a mathematical input.

A $\Sigma_{\mathrm{PreQ}}$-derivation is a universe-small well-founded rooted
typed tree.  Its leaves are labelled by standing constants; an internal vertex
labelled by $C$ has the arity prescribed for $C$, and its incoming subtrees
evaluate in the input sorts of $C$.  Evaluation is defined recursively: evaluate
the incoming trees in their displayed ambient categories and apply the actual
constructor $C$ there.  For an evaluated output $T$ in a specified displayed
ambient category define
\[
\boxed{
\operatorname{Der}_{\mathrm{PreQ}}(T)
}
\]
to be the maximal subgroupoid whose objects are pairs $(d,e)$ consisting of such
a derivation and an equivalence
$e:\operatorname{ev}(d)\simeq T$ in that same ambient category.  Its morphisms
are isomorphisms of labelled rooted trees together with the recursively induced
coherent equivalences at all vertices.  We use
\[
\boxed{
\mathsf{PreQData}:=\{(T,d,e)\mid(d,e)\in
\operatorname{Der}_{\mathrm{PreQ}}(T)\},
\qquad U(T,d,e):=T
}
\]
as notation for this sort-indexed family, not as an assertion that all displayed
ambient categories form one dependent term category.

For a derivation $d$, let $\operatorname{Occ}(d)$ be the poset of its internal
vertices, ordered by subtree dependence.  The exchange-sensitive full subposet
of $\operatorname{Occ}(d_b)$ is precisely $\mathcal K_b$ of
Definition~\ref{def:prequantum-full-comparison-realization}.
\end{definition}

\begin{proposition}[Provenance boundary of the chapter]
\label{prop:prequantum-output-type-boundary}
Every native construction $T$ preceding this section has a canonical point of
$\operatorname{Der}_{\mathrm{PreQ}}(T)$.  No such derivation contains a
constructor for
polarization, quantum-state formation, operator or factorization algebras,
measures, renormalized amplitudes, BV Laplacians, or quantum-master solutions.
This is a statement about constructional provenance, not about disjointness of
underlying categorical types: for example an underlying category equivalent to
$\operatorname{Mod}_k$ may occur on both sides of the boundary with different
tags.
\end{proposition}

\begin{proof}
Proceed by induction in the textual dependency order of the displayed
constructions.  The standing geometric data, big differential coefficients, and
stable curvature span are labelled leaves.  If the inputs of a displayed
construction already have derivation trees, graft their roots below a new vertex
labelled by its actual constructor.  Well-foundedness is preserved because every
edge points from an earlier construction to a later one, and the recursive
evaluation is the displayed output by definition.  The identity equivalence of
that output supplies the required point of
$\operatorname{Der}_{\mathrm{PreQ}}(T)$.  Functoriality of the constructor and
its universal-property coherences give the morphisms in the maximal subgroupoid.

The constructors for polarization, quantum-state formation, operator or
factorization algebras, measures, renormalized amplitudes, BV Laplacians, and
quantum-master solutions do not occur in $\Sigma_{\mathrm{PreQ}}$.  By the
definition of a labelled derivation tree, none of its vertices can carry a label
outside the signature.  This proves the provenance statement.  Since
$\mathsf{PreQData}$ is only a sort-indexed family and $U$ forgets the tree, the
proof makes no claim that equivalent underlying objects in different displayed
ambient categories have identical provenance.
\end{proof}

\begin{construction}[Shifted-geometric-quantization handoff]
\label{constr:prequantum-safronov-quantization-handoff}
Let $\operatorname{PreLagFib}^{n}_{\mathrm{Saf}}$ denote the target-side
realization of prequantum $n$-shifted Lagrangian fibrations of
\cite[Definition~2.13]{Safronov2020}: its point consists of the morphism
$\pi:X\to B$, the $n$-gerbe on $B$, the extended connective structure, and the
inverse data for the induced Lagrangian comparison.  Let
$\operatorname{PreQ}^{n}_{\mathrm{Saf}}$ be its underlying target
prequantization realization.  On the indexed native-to-Safronov comparison
branch define
\[
\boxed{
\operatorname{QuantIn}_{\mathrm{Saf}}^{n}
:=
\operatorname{CmpPreQReal}_{\mathrm{Saf}}
\times_{\operatorname{PreQ}^{n}_{\mathrm{Saf}}}
\operatorname{PreLagFib}^{n}_{\mathrm{Saf}}.
}
\]
Thus a point is a native prequantum object, an actual indexed bridge to the
target theory, and the additional Lagrangian-fibration/polarization data required
there.  The indexed comparison in this fibre is the prequantum comparison in
symmetric-monoidal coefficient systems; the quantum output below is a separate
module-valued handoff.

The quantum target is now separated by shift.  On the $n=1$ branch put
\[
A_B:=\operatorname{QCoh}(B),
\qquad
\mathsf{Lin}(B)
:=
\operatorname{Mod}_{A_B}
(\operatorname{Pr}^{L}_{k,\mathrm{st}}).
\]
Following \cite[Definition~2.3]{Safronov2020}, let
\[
\operatorname{Gerb}^{1}(B)
:=
\operatorname{Map}(B,\mathrm B^{2}\mathbf G_m).
\]
For $\mathcal G\in\operatorname{Gerb}^{1}(B)$ define
\[
\mathcal C_{\mathcal G}
:=
\operatorname{QCoh}^{\mathcal G}(B).
\]
Tensoring an untwisted quasi-coherent sheaf with a $\mathcal G$-twisted one is
compatible with the descent transition functors and therefore gives a
colimit-preserving action
\[
A_B\otimes\mathcal C_{\mathcal G}
\longrightarrow
\mathcal C_{\mathcal G}.
\]
Thus $\mathcal C_{\mathcal G}$ is canonically an object of $\mathsf{Lin}(B)$.
This module-valued assignment is the universally defined target; define its
tagged graph by
\[
\boxed{
\operatorname{TwCat}(B)
:=
\int_{\mathcal G\in\operatorname{Gerb}^{1}(B)}
\{(\mathcal G,\mathcal C_{\mathcal G})\}
\longrightarrow\mathsf{Lin}(B).
}
\]

The tensor product of twisted sheaves gives natural $A_B$-linear morphisms
\[
m_{\mathcal G,\mathcal H}:
\mathcal C_{\mathcal G}\otimes_{A_B}\mathcal C_{\mathcal H}
\longrightarrow
\mathcal C_{\mathcal G+\mathcal H}.
\]
For the inverse gerbe $-\mathcal G$, compose $m_{\mathcal G,-\mathcal G}$ with
the canonical trivialization
$\mathcal G+(-\mathcal G)\simeq0$ and the identification
$\mathcal C_0\simeq A_B$ to obtain
\[
\mu_{\mathcal G}:
\mathcal C_{\mathcal G}\otimes_{A_B}\mathcal C_{-\mathcal G}
\longrightarrow A_B.
\]
These arrows assemble into a classifying morphism
\[
\mu_B:
\operatorname{Gerb}^{1}(B)
\longrightarrow
\operatorname{Mor}(\mathsf{Lin}(B)).
\]
Form its walking-equivalence inverse locus
\[
\boxed{
\operatorname{Gerb}^{1}_{\mathrm{Pic}}(B)
:=
\operatorname{Gerb}^{1}(B)
\mathop{\times}_{\operatorname{Mor}(\mathsf{Lin}(B)),\,\mu_B}
\operatorname{Eqv}(\mathsf{Lin}(B)).
}
\]
A point is a gerbe together with coherent inverse data for
$\mu_{\mathcal G}$.  Over this locus,
\[
\mathcal C_{\mathcal G}\otimes_{A_B}\mathcal C_{-\mathcal G}
\simeq A_B,
\]
and symmetry gives the reverse tensor equivalence.  Hence
$\mathcal C_{\mathcal G}$ is invertible in $\mathsf{Lin}(B)$ with inverse
$\mathcal C_{-\mathcal G}$.  The construction therefore induces
\[
\tau_B^{\mathrm{Pic}}:
\operatorname{Gerb}^{1}_{\mathrm{Pic}}(B)
\longrightarrow
\operatorname{Pic}(\mathsf{Lin}(B))
\]
and the genuinely Picard-valued tagged graph
\[
\boxed{
\operatorname{TwMod}(B)
:=
\int_{(\mathcal G,\epsilon)\in
\operatorname{Gerb}^{1}_{\mathrm{Pic}}(B)}
\{(\mathcal G,\epsilon,\mathcal C_{\mathcal G})\}
\longrightarrow
\operatorname{Pic}(\mathsf{Lin}(B)).
}
\]

The point of $\operatorname{PreLagFib}^{1}_{\mathrm{Saf}}$ supplies the gerbe on
$B$ together with the separately retained extended connective structure on its
pullback to $X$ \cite[Definition~2.13]{Safronov2020}.  Let
$\mathsf{Base}_{\mathrm{Saf}}:=\mathcal B_{b_{\mathrm{Saf}}}$ be the geometric
occurrence category of the chosen Safronov comparison branch and put
\[
\operatorname{TwCat}_{\mathrm{Saf}}
:=
\int_{B\in\mathsf{Base}_{\mathrm{Saf}}}\operatorname{TwCat}(B),
\qquad
\operatorname{TwMod}_{\mathrm{Saf}}
:=
\int_{B\in\mathsf{Base}_{\mathrm{Saf}}}\operatorname{TwMod}(B).
\]
The target-side construction of
\cite[Definition~3.1]{Safronov2020} gives the universally defined handoff
\[
\boxed{
\operatorname{GQ}_{\mathrm{Saf}}^1:
\operatorname{QuantIn}_{\mathrm{Saf}}^1
\longrightarrow\operatorname{TwCat}_{\mathrm{Saf}},
\qquad
(X\to B,\mathcal G)\longmapsto
(B,\mathcal G,\mathcal C_{\mathcal G}).
}
\]
Let
\[
\operatorname{Gerb}^{1}_{\mathrm{Saf}}
:=
\int_{B\in\mathsf{Base}_{\mathrm{Saf}}}\operatorname{Gerb}^{1}(B),
\qquad
\operatorname{Gerb}^{1,\mathrm{Pic}}_{\mathrm{Saf}}
:=
\int_{B\in\mathsf{Base}_{\mathrm{Saf}}}
\operatorname{Gerb}^{1}_{\mathrm{Pic}}(B),
\]
and let
$\operatorname{QuantIn}_{\mathrm{Saf}}^1\to
\operatorname{Gerb}^{1}_{\mathrm{Saf}}$ forget to $(B,\mathcal G)$.  Define the
Picard-recognized quantum input by
\[
\boxed{
\operatorname{QuantIn}_{\mathrm{Saf}}^{1,\mathrm{Pic}}
:=
\operatorname{QuantIn}_{\mathrm{Saf}}^1
\times_{\operatorname{Gerb}^{1}_{\mathrm{Saf}}}
\operatorname{Gerb}^{1,\mathrm{Pic}}_{\mathrm{Saf}}.
}
\]
Restriction of the preceding handoff constructs the stronger target exactly on
that realization:
\[
\boxed{
\operatorname{GQ}_{\mathrm{Saf}}^{1,\mathrm{Pic}}:
\operatorname{QuantIn}_{\mathrm{Saf}}^{1,\mathrm{Pic}}
\longrightarrow\operatorname{TwMod}_{\mathrm{Saf}}.
}
\]
The separate $0$-shifted construction is
\[
\operatorname{GQ}_{\mathrm{Saf}}^0:
\operatorname{QuantIn}_{\mathrm{Saf}}^0
\longrightarrow\operatorname{Mod}_k,
\qquad
(X\to B,\mathcal L)\longmapsto\Gamma(B,\mathcal L).
\]
This is the construction of
\cite[Definition~3.7]{Safronov2020}; it does not factor through
$\operatorname{TwCat}_{\mathrm{Saf}}$ or
$\operatorname{TwMod}_{\mathrm{Saf}}$.
On the Hamiltonian $G$-scheme branch with the constructed equivariant
prequantization and the stated qcqs target, transport of the reduction diagram
followed by \cite[Proposition~4.13]{Safronov2020} gives
\[
\boxed{
\operatorname{GQ}_{\mathrm{Saf}}^0
(X/\!/G\longrightarrow B/G)
\simeq
\Gamma(B,\mathcal L)^G.
}
\]
This is an explicit enlargement of both the prequantum input and the codomain on
the chosen quantum branch.  On the $n=1$ branch the universally constructed
codomain is $\operatorname{TwCat}_{\mathrm{Saf}}$; the genuinely Picard-valued
codomain $\operatorname{TwMod}_{\mathrm{Saf}}$ is constructed over the explicit
inverse locus $\operatorname{Gerb}^{1,\mathrm{Pic}}_{\mathrm{Saf}}$.  Forgetting
the module and twist tags gives the underlying presentable category, but that
forgetful step is not used as a commutative-monoidal comparison.  The handoff is
not an omitted field of the native prequantum object.  No quantum-output functor
is asserted here for $n>1$: \cite[Remark~3.14]{Safronov2020} describes the
expected higher-categorical target but does not construct it.
\end{construction}

\section{The Brave-New Prequantum Field Theory theorem}
\label{sec:prequantum-final-package}

\begin{theorem}[Brave-New Prequantum Field Theory]
\label{thm:prequantum-final-package}
Retain
\[
x:X\to S,
\qquad
Q=Q_{X/S},
\qquad
a:\overline A\to Q,
\]
the standing Cartan field object $(A,\rho_A)$ whose descent is $a$, the stable
variational coefficient $M\in E_Q$, a degree-$n$ Lagrangian
\[
L:\mathbf1_Q\to\operatorname{Sec}^0(A;M)[n],
\]
and a formal field-symmetry groupoid
\[
\mathcal G_\bullet\in\operatorname{FormLieGpd}(\overline A/Q).
\]
Also retain, by their cited constructions in Chapters~12--\ref{ch:symmetry-noether},
the relative Lepage objects, universal presymplectic transgressions, critical
object, canonically formalized critical symmetry groupoid, and its effective
quotient.  Then the constructions of this chapter provide the following outputs.
Every assertion below is relative to the complete realization object displayed
for the data it uses.  When several outputs are combined, their base is the
iterated homotopy pullback over their shared parameter objects.  In particular,
the standing tuple $(x,A,M,L,\mathcal G)$ does not silently supply a differential
coefficient, a coefficient morphism, a curvature span, an underlying stage
system, a fundamental class, a comparison frame, or a selected prequantization.
The differential Lepage and presymplectic-prequantization constructions are
defined for every stage $q$ in the Chapter~12 Lepage range.  For every
field-theoretic Hamiltonian, Noether, and ensuing critical-reduction output using
$U_{\mathcal G,L}^{(q)}$, $J_{\mathcal G,L}^{(q)}$,
$\operatorname{HamCompReal}_{\mathcal G,L}^{q}$, or a construction derived from
them, fix $q\ge1$.

\medskip
\noindent\emph{Big differential coefficients.}
For every native parameter object $T$,
\[
\mathcal R_T^{\partial}
=
\lim_{(h_U\to T)^{op}}\mathcal R_U^0
\]
is a stable presentable symmetric monoidal category with arbitrary pullback,
represented recovery, effective descent, and big differential mapping and section
spectra.

\medskip
\noindent\emph{Stable curvature and prequantization.}
A stable curvature span
\[
\mathscr D\xrightarrow{\kappa}\mathscr U\xleftarrow{u}\mathscr C
\]
is a point of $\operatorname{CurvDiagReal}$.  It constructs
\[
\mathscr D^{\mathrm{cl}}_{\kappa,u}
=
\mathscr D\times_{\mathscr U}\mathscr C
\]
and the stable fibre sequence
\[
\mathscr F_\kappa
\longrightarrow
\mathscr D^{\mathrm{cl}}_{\kappa,u}
\longrightarrow
\mathscr C
\xrightarrow{\partial_{\kappa,u}}
\mathscr F_\kappa[1].
\]
For every degree-$r$ complete closed class $\omega$,
\[
\boxed{
\operatorname{PreQ}_{\kappa,u}(\omega)
\simeq
Z\bigl(\partial_{\kappa,u}[r]\omega\bigr).
}
\]
If inhabited, this fibre is a torsor for
$\Omega^\infty\mathscr F_\kappa[r]$, with $j$th higher automorphisms
$\Omega^\infty\mathscr F_\kappa[r-j]$.

\medskip
\noindent\emph{Differential variational refinement.}
Over $\operatorname{DiffDenCurvReal}_x(M)$ the universal pair
$(M^\partial_{x,D},\kappa_M)$ is constructed.  The internal pullback
$\operatorname{PreQLag}_{\kappa_M}(L)$ represents differential Lagrangian lifts.
The separate object $\operatorname{PreQLep}_{\kappa_M}^{q}(L)$ represents lifts
of actual relative Lepage realizations.  Before fixing $L$, the universal object
$\operatorname{VarRef}_{\kappa_M}^{q}$ retains
$(L,\widehat L,\eta_{\mathrm{Lag}},\widehat\ell)$ as one refinement family.  The complete vertical filtration
constructs
$\operatorname{ClTrans}^{2,q}_H(A;M;n)$ and
$\omega_L^{(q),\mathrm{cl}}$.  Over $\operatorname{PreQLep}$,
\[
\kappa_{\mathrm{tr}}^{(q)}\omega_{\widehat L}^{(q)}
\simeq
u_{L,q}\omega_L^{(q),\mathrm{cl}},
\]
so there is a canonical map
\[
\boxed{
\operatorname{PreQLep}_{\kappa_M}^{q}(L)
\longrightarrow
\operatorname{PreQ}_{\omega,L}^{q}.
}
\]
The zero-comparison Lepage point gives the canonical special case for every
differential Lagrangian refinement, while
$\operatorname{PreQLep}_{\kappa_M}^{q}(L)$ is the complete derived lifting object
for an arbitrary underlying Lepage realization.

\medskip
\noindent\emph{Differential action and traced correspondences.}
The Chapter~12 action remains orientation-free as
\[
x_{e!}\mathbf1_X^\partial\longrightarrow D[n].
\]
$\operatorname{FundReal}_{\partial}$ and
$\operatorname{TrReal}_{\partial}$ are pointed mapping realizations of
fundamental-class and coefficient-trace candidates.  Their fibres
$\operatorname{FundLiftReal}_{\partial}$ and
$\operatorname{TrLiftReal}_{\partial}$ over prescribed motivic shadows are the
actual existence problems.  The explicitly constructed Chapter~10 diagram
$\mathbf{DRStage}^{\partial}$ and
$\operatorname{TrCohReal}_{\partial}$ represent diagram maps, while
$\operatorname{TrCohLiftReal}_{\partial}$ represents coherent extension of a
selected stage trace.
$\operatorname{SqReal}$ and $\operatorname{TransCurvReal}_{\partial}$ compare
the actual differential transgression operator with the independently constructed
underlying stable transgression rather than introducing an untyped exceptional
pullback.  The objects $\operatorname{StageSys}_X^U$ and
$\operatorname{StageSys}_B^U$ construct the underlying stage diagrams and their
output equivalences; $\operatorname{TransStageCurvReal}_{\partial}$ is the
relative lift of the three terminal maps and their curvature path.  Refined
correspondence morphisms are the fixed-square fibres
$\operatorname{TwoCellStageReal}$, rather than a pullback action of a raw
decorated $2$-cell.  Trace/fundamental-class pasting gives coherent
composition and Fubini transgression.  Differential lifts of geometric
correspondence $2$-cells are represented by
$\operatorname{TwoCellReal}_{\partial}$; the trace and fundamental-class
compatibilities are retained as the two opposite-variance path factors of
$\operatorname{TwoCellDecReal}_{\partial}$.  These data produce the decorated
$\operatorname{Corr}^{\partial}_{\mathrm{tr,fund}}$ as an actual
$(\infty,2)$-category.  BNV/Stokes compatibility is controlled separately by
$\operatorname{StokesEq}_p$, its small detection presentation
$\operatorname{StokesTestEq}_p^\lambda$, and
$C^{\mathrm{Stokes}}_{p,\widehat E}$.

\medskip
\noindent\emph{Cartan transport, vertical origin, and Hamiltonian geometry.}
$\operatorname{CartCurvReal}_{\mathcal H}$ represents coherent Cartan transport
of an arbitrary curvature span.  $\operatorname{VertCurvReal}$ represents the
relation-level coefficient origin of a generic complete closed coefficient, while
$\operatorname{VertClassReal}$ represents the separate lift of its selected closed
section to that origin.  The variational closed coefficient and its universal
presymplectic section have canonical points of these two realizations.
The ordered staircase family is the ordered Alexander--Whitney morphism.  Its
complete closed degree-two identity is
\[
\boxed{
\Delta^{\mathrm{cl}}_{\mathcal H}\omega
\simeq
-\,d^{\mathrm{mix}}_{V,\mathrm{cl}}
\bigl(\iota_{\mathcal H}\omega\bigr),
}
\]
with normalized shadows
\[
d_{\mathcal H}\bar\omega
\simeq-d_V^{\mathrm{mix}}\iota_{\mathcal H}\omega,
\qquad
 d_{\mathcal H}\iota_{\mathcal H}\omega
\simeq-d_V^{\mathrm{mix}}K_{\mathcal H}(\omega),
\qquad
 d_{\mathcal H}K_{\mathcal H}(\omega)\simeq0.
\]
In degree two the candidate primitive
\[
\mu_K
=
J_{01}+\operatorname{tr}_{01}J_{12}+K_{012},
\]
has a defect $\epsilon_2$ satisfying
$\beta_{02}(\epsilon_2)\simeq0$.  Its lift
$\bar\epsilon_2\in\operatorname{fib}(\beta_{02})$ defines the dependent-product
zero locus $\operatorname{HamMultReal}_{\omega,2}$.  The fixed-object Reedy
tower, groupoid core locus, and curvature mapping fibre define
$\operatorname{HamMultReal}_{\omega}$.  Over that full realization the
universal raw Hamiltonian groupoid has multiplication $\mu_K$ and a canonical
morphism to $\operatorname{CurvSymRaw}_{\omega}$; formalization gives the formal
Hamiltonian groupoid.  The stronger functor-level inverse locus
$\operatorname{HamEdgeEq}_{\omega}$ maps canonically to this realization.

\medskip
\noindent\emph{Integrated Kostant--Souriau geometry.}
Over the same Hamiltonian multiplicativity realization, the raw
Hamiltonian-prequantum pullback has the effective extension
\[
\boxed{
B_YG_{\mathrm{flat}}
\longrightarrow
\operatorname{HamPreQRaw}_{\widehat\omega}
\twoheadrightarrow
I^{\mathrm{Ham,raw}}_{\widehat\omega}.
}
\]
Formalization gives
\[
\boxed{
\operatorname{FlatGau}_{\widehat\omega}
\longrightarrow
\operatorname{HamPreQ}_{\widehat\omega}
\longrightarrow
\operatorname{Ham}^{\mathrm{lift}}_{\widehat\omega}
}
\]
as an unconditional formal homotopy-kernel sequence.  Raw and formal recovery of
the full Hamiltonian target are governed by the distinct inverse loci
$\operatorname{KSEq}^{\mathrm{raw}}$ and
$\operatorname{KSEq}^{\mathrm{for}}$.  The fibre
$\operatorname{KSRawLiftReal}$ over a selected formal inverse is the complete
raw-recovery problem.  Relative integrated isotropy is exactly the flat kernel.

\medskip
\noindent\emph{Differentiated prequantum geometry.}
Separate partition-Lie differentiation constructs
$\mathbb A^{\mathrm{PreQ}}_{\widehat\omega}$,
$\mathbb A^{\mathrm{Ham,\,lift}}_{\widehat\omega}$, and
$\mathbb A^{\mathrm{Flat}}_{\widehat\omega}$.  Projection to the independently
differentiated quantomorphism controller gives $a_{\mathrm{PQ}}$, with discrepancy
$C^{\mathrm{PreQ/Quant}}_{\widehat\omega}$ and inverse locus
$\operatorname{PreQQuantEq}_{\widehat\omega}$.  The intrinsic algebroid fibre
$\mathbb F^{\mathrm{KS}}_{\widehat\omega}$ receives
$\gamma^{\mathrm{KS,diff}}_{\widehat\omega}$ from differentiated actual flat
gauge.  Its cofiber is formed after applying $U_{\mathrm{src}}$.  The stable
source-module fibre sequence gives
\[
\boxed{
\mathbf k_{\widehat\omega}^{(1)}:
U_{\mathrm{src}}\mathbb A_{\widehat\omega}^{\mathrm{Ham,\,lift}}
\longrightarrow
U_{\mathrm{src}}\mathbb F_{\widehat\omega}^{\mathrm{KS}}[1].
}
\]
The complete Chapter~15 morphism-corolla--latching tower reconstructs every
algebroid lift; anchor compatibility at weight zero follows from the actual
anchored KS projection and is not an extra datum.

\medskip
\noindent\emph{Stable flat representation and observables.}
The stable fibre $\mathscr F_\kappa[r-1]$ already carries Cartan transport.  Its
restriction to the raw liftable image descends as
$\operatorname{Ad}^{\mathrm{flat,raw}}_{\widehat\omega}$ before applying
$\Omega^\infty$.  Its infinite-loop transport is then proved to coincide with
flat-kernel conjugation.  Differentiation gives
$\mathbb K^{\mathrm{Ad,PreQ}}_{\widehat\omega}$, kept separate from actual-flat
differentiation.  The observable controller is
$\operatorname{Obs}_{\mathrm{PreQ}}(Y,\widehat\omega)
=\mathbb A_{\widehat\omega}^{\mathrm{PreQ}}$, with exact labelled operations and
the canonical first-weight coefficient connection; higher tangent-action weights
supply curvature nullification, Bianchi, and all higher flatness fillers.

\medskip
\noindent\emph{Prequantum Noether theory.}
Chapter~\ref{ch:symmetry-noether}'s complete transgression
$U_{\mathcal G,L}^{(q)}$ and its current shadow $J_{\mathcal G,L}^{(q)}$ are kept
distinct.  $\operatorname{HamCompReal}_{\mathcal G,L}^{q}$ represents an actual
lift
\[
\boxed{
d_V^{\mathrm{tr}}U_{\mathcal G,L}^{(q)}
\simeq
\iota_{\mathcal G}\omega_L^{(q)}
}
\]
of the Chapter~17 Hamiltonian path.  The native matching tower then constructs
$\operatorname{HamNoethReal}$; its degree-two defect is
\[
U_{02}-U_{01}-\operatorname{tr}_{01}U_{12}-K_{012}.
\]
$\operatorname{NoethHamLiftReal}$ represents formal factorization through the
liftable Hamiltonian image and $\operatorname{NoethPreQReal}$ represents the
actual coherent prequantum lift.  Over the latter,
\[
\widetilde\rho_{\mathrm{Noeth}}:
\mathbb A_{\mathcal G}^{\pi,E_\infty}
\longrightarrow
\mathbb A_{\widehat\omega_L^{(q)}}^{\mathrm{PreQ}}
\]
is the differentiated Prequantum Noether morphism.  Differential symmetry is formed universally before $L$ is frozen: it stabilizes
$(L,\widehat L,\eta_{\mathrm{Lag}},\widehat\ell)$ through
$\operatorname{DiffVarSymReal}^{q}_{\mathcal G}$, and its compatibility with the
independent Noether lift is $\operatorname{DiffSymPreQCompat}$.

\medskip
\noindent\emph{Critical descent and reduction.}
The full prequantization parameter is retained under critical restriction.
Critical mixed-grid descent constructs the reduced differential, underlying, and
complete closed curvature coefficients.  The cosimplicial prequantum section
object is formed internally as
\[
\mathscr S_{\mathrm{PreQ}}^{p,q}
=
\Omega_{P_{\mathrm{curv}}^q}^\infty
\Pi_{g_p}\mathscr P_{\mathrm{crit},p}^{(q)},
\]
and $\operatorname{PreQDesc}$ is its internal totalization fibre.  It descends
the selected section to
$\widehat\omega_{L,\mathrm{red}}^{(q)}$ with
\[
\boxed{
\operatorname{curv}
\bigl(\widehat\omega_{L,\mathrm{red}}^{(q)}\bigr)
\simeq
\omega_{L,\mathrm{red}}^{(q),\mathrm{cl}}.
}
\]
Universally,
\[
\boxed{
\operatorname{DescTot}_{q_{\mathrm{curv}}}
\bigl(\operatorname{PreQ}_{\mathrm{crit},\bullet}^{q}\bigr)
\simeq
\operatorname{PreQ}_{\mathrm{red}}^q.
}
\]
Critical formal isotropy descends as
$\operatorname{Ad}^{\mathrm{for}}(\mathcal K_{\mathrm{red}})$, and the reduced
Prequantum Noether output is
\[
\boxed{
\alpha_{\mathcal G,L,\mathrm{red}}^{(q)}:
\operatorname{Ad}^{\mathrm{for}}(\mathcal K_{\mathrm{red}})
\longrightarrow
\operatorname{Aut}^{\partial}_{\mathrm{PreQ}}
(\widehat\omega_{L,\mathrm{red}}^{(q)}).
}
\]
Its complete Hamiltonian primitive shadow is the descended $U^{(q)}$, its
horizontal-tail shadow is the reduced current, and evaluation at a point of
$\operatorname{RedSymReal}_{\mathcal G,L}^{q}$ gives an individual reduced
prequantum automorphism.
The compatibility base $\operatorname{RedNoethPreQ}_{\mathcal G,L}^q$ is the
typed iterated pullback over the actual prequantum Cartan-lift realization and
the explicit finite limit $\operatorname{CurrNoethBase}_{\mathcal G,L}^q$.

\medskip
\noindent\emph{Comparison and locality.}
$\operatorname{CmpPreQReal}_{\mathcal B_b,b}$ represents an indexed comparison
over the geometric occurrence category of the branch, together with the
explicit $\operatorname{HomEq}$, $\operatorname{MapEq}$,
$\operatorname{SecEq}$, $\operatorname{PiEq}$, $\operatorname{TotEq}$,
Cartan, relation, mixed,
formalization, Artinian, partition-Lie, tangent, descent, and effective-image
inverse loci indexed by the derivation-tree occurrence poset $\mathcal K_b$; over it the corresponding
prequantum branch transports.  The intrinsic stable branch
has arbitrary native pullback and effective descent.  The Chapter~10 realized de
Rham branch retains its actual exchange maps, with invertibility represented by
$\operatorname{DREq}(c;F)$ and canonically realized on the proved
spectral-\'etale finite-presentation domain.  The differentiated branch retains
the separately scoped coherent spectral-\'etale transport of Chapters~14--15.
\end{theorem}

\begin{proof}
Every item is one of the constructions established in the preceding sections:
affine-point limits of the Chapter~9 differential categories; parameterized stable
mapping objects, fibres, complete filtrations and dependent products; the
Chapter~10 coherent stage diagram; exceptional Beck--Chevalley and trace pasting;
trace and fundamental-class shadow fibres; walking-equivalence and fixed
commutative-square lift realizations; universal underlying stage systems and
fixed-boundary stage-square morphisms; Cartan Eilenberg--Moore
fibres; relation-level coefficient limits; ordered Alexander--Whitney maps;
endpoint-fixed cotensors; effective images and unitwise formalization; raw/formal
groupoid mapping ends and fibres; source-module stable fibre sequences;
prequantum-to-quantomorphism comparison; partition-Lie
differentiation and the Chapter~15 monad/morphism-latching calculus; and effective
critical descent.  The reduction compatibility object and the full indexed
comparison, including the mapping- and section-spectrum exchanges, are the
explicit small limits constructed above.  The typed
discrepancy-comparison realizations represent every cross-identification.  Hence
no additional axiom, admissibility category, or hidden package is invoked.
\end{proof}

\subsection{Handoff to brave-new quantization}
\label{subsec:prequantum-handoff-quantization}

The differential and infinitesimal geometry required for prequantization is now
closed under the operations constructed in the preceding chapters.  Its output
includes the differential lift problem, differential Lepage refinement, flat
torsor, Cartan and vertical-origin realizations, quantomorphism and Hamiltonian
groupoids, raw/formal KS geometry, differentiated discrepancy and lift towers,
observable controller, native and differentiated Prequantum Noether lifts,
traced-correspondence transgression, comparison realization, and reduced critical
prequantization.

Proposition~\ref{prop:prequantum-output-type-boundary} proves that every native
construction carries a prequantum provenance tag and uses no quantization
constructor; it does not claim disjoint underlying types.  On the
shifted-geometric branch,
Construction~\ref{constr:prequantum-safronov-quantization-handoff} supplies the
actual enlarged input $\operatorname{QuantIn}_{\mathrm{Saf}}^n$.  Its constructed
quantum outputs are $\operatorname{TwCat}_{\mathrm{Saf}}$ on the unrestricted
$n=1$ branch,
$\operatorname{TwMod}_{\mathrm{Saf}}$ on the realized Picard inverse locus, and
$\operatorname{Mod}_k$ on the $n=0$ branch.  Its Hamiltonian reduction formula
is the cited $0$-shifted quantization-commutes-with-reduction theorem.  No
higher-shifted quantum-output functor is asserted.
